\input{preamble}

% OK, start here.
%
\begin{document}

\title{Cohomologie sur les sites}


\maketitle

\phantomsection
\label{section-phantom}

\tableofcontents

\section{Introduction}
\label{section-introduction}

\noindent
Dans ce document, nous développons quelques thèmes relatifs à la cohomologie
des faisceaux. Nous étudions le cas des faisceaux sur des sites,
bien que nous reprenions souvent simplement les raisonnements,
les constructions et les démonstrations du cas
topologique.
Les références de base sont \cite{SGA4}, \cite{Godement} et \cite{Iversen}.




\section{Cohomologie des faisceaux}
\label{section-cohomology-sheaves}

\noindent
Soit $\mathcal{C}$ un site, voir
Sites, définition \ref{sites-definition-site}.
Soit $\mathcal{F}$ un faisceau abélien sur $\mathcal{C}$.
Nous savons que la catégorie des faisceaux abéliens sur $\mathcal{C}$
possède suffisamment d'injectifs, voir
Injectifs, théorème \ref{injectives-theorem-sheaves-injectives}.
Nous pouvons donc choisir une résolution injective
$\mathcal{F}[0] \to \mathcal{I}^\bullet$.
Pour tout objet $U$ du site $\mathcal{C}$, nous définissons
\begin{equation}
\label{equation-cohomology-object-site}
H^i(U, \mathcal{F}) = H^i(\Gamma(U, \mathcal{I}^\bullet))
\end{equation}
comme le {\it $i$-ième groupe de cohomologie du faisceau abélien
$\mathcal{F}$ sur l'objet $U$}. Autrement dit, ce sont les
foncteurs dérivés à droite du foncteur $\mathcal{F} \mapsto \mathcal{F}(U)$.
La famille de foncteurs $H^i(U, -)$ forme un $\delta$-foncteur universel
$\textit{Ab}(\mathcal{C}) \to \textit{Ab}$.

\medskip\noindent
Il arrive parfois que
le site $\mathcal{C}$ ne possède pas d'objet final. Dans ce
cas, nous définissons les {\it sections globales} d'un préfaisceau
d'ensembles $\mathcal{F}$ sur $\mathcal{C}$ comme l'ensemble
\begin{equation}
\label{equation-global-sections}
\Gamma(\mathcal{C}, \mathcal{F}) =
\Mor_{\textit{PSh}(\mathcal{C})}(e, \mathcal{F})
\end{equation}
où $e$ est un objet final de la catégorie des préfaisceaux sur $\mathcal{C}$.
Dans ce cas, étant donné un faisceau abélien $\mathcal{F}$ sur $\mathcal{C}$,
nous définissons le {\it $i$-ième groupe de cohomologie de $\mathcal{F}$ sur $\mathcal{C}$}
comme suit :
\begin{equation}
\label{equation-cohomology}
H^i(\mathcal{C}, \mathcal{F}) = H^i(\Gamma(\mathcal{C}, \mathcal{I}^\bullet))
\end{equation}
autrement dit, il s'agit du $i$-ième foncteur dérivé à droite du
foncteur des sections globales.
La famille de foncteurs $H^i(\mathcal{C}, -)$ forme un $\delta$-foncteur universel
$\textit{Ab}(\mathcal{C}) \to \textit{Ab}$.

\medskip\noindent
Soit $f : \Sh(\mathcal{C}) \to \Sh(\mathcal{D})$ un morphisme de topos, voir
Sites, définition \ref{sites-definition-topos}.
Avec $\mathcal{F}[0] \to \mathcal{I}^\bullet$ comme ci-dessus,
nous définissons
\begin{equation}
\label{equation-higher-direct-image}
R^if_*\mathcal{F} = H^i(f_*\mathcal{I}^\bullet)
\end{equation}
comme la {\it $i$-ième image directe supérieure de $\mathcal{F}$}.
Ce sont les foncteurs dérivés à droite de $f_*$.
La famille de foncteurs $R^if_*$ forme un $\delta$-foncteur universel
de type $\textit{Ab}(\mathcal{C}) \to \textit{Ab}(\mathcal{D})$.

\medskip\noindent
Soit $(\mathcal{C}, \mathcal{O})$ un site annelé, voir
Modules sur les sites, définition \ref{sites-modules-definition-ringed-site}.
Soit $\mathcal{F}$ un $\mathcal{O}$-module.
Nous savons que la catégorie des $\mathcal{O}$-modules
possède suffisamment d'injectifs, voir
Injectifs, théorème \ref{injectives-theorem-sheaves-modules-injectives}.
Nous pouvons donc choisir une résolution injective
$\mathcal{F}[0] \to \mathcal{I}^\bullet$.
Pour tout objet $U$ du site $\mathcal{C}$, nous définissons
\begin{equation}
\label{equation-cohomology-object-site-modules}
H^i(U, \mathcal{F}) = H^i(\Gamma(U, \mathcal{I}^\bullet))
\end{equation}
comme le {\it $i$-ième groupe de cohomologie de $\mathcal{F}$ sur $U$}.
La famille de foncteurs $H^i(U, -)$ forme un $\delta$-foncteur universel
$\textit{Mod}(\mathcal{O}) \to \text{Mod}_{\mathcal{O}(U)}$. De même,
\begin{equation}
\label{equation-cohomology-modules}
H^i(\mathcal{C}, \mathcal{F}) = H^i(\Gamma(\mathcal{C}, \mathcal{I}^\bullet))
\end{equation}
est le {\it $i$-ième groupe de cohomologie de $\mathcal{F}$ sur $\mathcal{C}$}.
La famille de foncteurs $H^i(\mathcal{C}, -)$ forme un
$\delta$-foncteur universel
$\textit{Mod}(\mathcal{O}) \to \text{Mod}_{\Gamma(\mathcal{C}, \mathcal{O})}$.

\medskip\noindent
Soit $f : (\Sh(\mathcal{C}), \mathcal{O}) \to (\Sh(\mathcal{D}), \mathcal{O}')$
un morphisme de topos annelés, voir
Modules sur les sites, définition \ref{sites-modules-definition-ringed-topos}.
Avec $\mathcal{F}[0] \to \mathcal{I}^\bullet$ comme ci-dessus,
nous définissons
\begin{equation}
\label{equation-higher-direct-image-modules}
R^if_*\mathcal{F} = H^i(f_*\mathcal{I}^\bullet)
\end{equation}
comme la {\it $i$-ième image directe supérieure de $\mathcal{F}$}.
Ce sont les foncteurs dérivés à droite de $f_*$.
La famille de foncteurs $R^if_*$ forme un $\delta$-foncteur universel
de type $\textit{Mod}(\mathcal{O}) \to \textit{Mod}(\mathcal{O}')$.








\section{Foncteurs dérivés}
\label{section-derived-functors}

\noindent
Nous exposons brièvement une approche des foncteurs dérivés à droite au moyen
de foncteurs de résolution. Supposons donc que $(\mathcal{C}, \mathcal{O})$ soit un site annelé.
Dans ce chapitre, nous écrirons
$$
K(\mathcal{O}) = K(\textit{Mod}(\mathcal{O}))
\quad
\text{and}
\quad
D(\mathcal{O}) = D(\textit{Mod}(\mathcal{O}))
$$
et de même pour les versions bornées des catégories triangulées
introduites dans
Catégories dérivées, définition \ref{derived-definition-complexes-notation} et
définition \ref{derived-definition-unbounded-derived-category}.
D'après
Catégories dérivées, remarque \ref{derived-remark-big-abelian-category},
il existe un foncteur de résolution
$$
j = j_{(\mathcal{C}, \mathcal{O})} :
K^{+}(\textit{Mod}(\mathcal{O}))
\longrightarrow
K^{+}(\mathcal{I})
$$
où $\mathcal{I}$ est la sous-catégorie additive strictement pleine de
$\textit{Mod}(\mathcal{O})$ constituée des $\mathcal{O}$-modules injectifs.
Pour tout foncteur exact à gauche $F : \textit{Mod}(\mathcal{O}) \to \mathcal{B}$
à valeurs dans une catégorie abélienne $\mathcal{B}$, nous noterons $RF$ le
foncteur dérivé à droite de
Catégories dérivées, section \ref{derived-section-right-derived-functor},
construit à l'aide du foncteur de résolution $j$ que nous venons de décrire :
\begin{equation}
\label{equation-RF}
RF = F \circ j' : D^{+}(\mathcal{O}) \longrightarrow D^{+}(\mathcal{B})
\end{equation}
voir
Catégories dérivées, lemme \ref{derived-lemma-right-derived-functor},
pour les notations. Remarquons que nous pouvons considérer $RF$ comme défini sur
$\textit{Mod}(\mathcal{O})$, $\text{Comp}^{+}(\textit{Mod}(\mathcal{O}))$, ou
$K^{+}(\mathcal{O})$, selon la situation. D'après
Catégories dérivées, définition \ref{derived-definition-higher-derived-functors},
nous obtenons le $i$-ième foncteur dérivé à droite
\begin{equation}
\label{equation-RFi}
R^iF = H^i \circ RF : \textit{Mod}(\mathcal{O}) \longrightarrow \mathcal{B}
\end{equation}
de sorte que $R^0F = F$ et que $\{R^iF, \delta\}_{i \geq 0}$ est un
$\delta$-foncteur universel, voir
Catégories dérivées, lemme \ref{derived-lemma-higher-derived-functors}.

\medskip\noindent
Voici deux cas particuliers de cette construction. Pour un anneau $R$, nous écrivons
$K(R) = K(\text{Mod}_R)$ et $D(R) = D(\text{Mod}_R)$, et de même pour les
versions bornées. Pour tout objet $U$ de $\mathcal{C}$, nous avons un foncteur exact à gauche
$
\Gamma(U, -) :
\textit{Mod}(\mathcal{O})
\longrightarrow
\text{Mod}_{\mathcal{O}(U)}
$
qui donne lieu à
$$
R\Gamma(U, -) :
D^{+}(\mathcal{O})
\longrightarrow
D^{+}(\mathcal{O}(U))
$$
d'après la discussion ci-dessus. Remarquons que $H^i(U, -) = R^i\Gamma(U, -)$
est compatible avec (\ref{equation-cohomology-object-site-modules}) ci-dessus.
De même, nous avons
$$
R\Gamma(\mathcal{C}, -) :
D^{+}(\mathcal{O})
\longrightarrow
D^{+}(\Gamma(\mathcal{C}, \mathcal{O}))
$$
compatible avec (\ref{equation-cohomology-modules}). Si
$f : (\Sh(\mathcal{C}), \mathcal{O}) \to (\Sh(\mathcal{D}), \mathcal{O}')$
est un morphisme de topos annelés, nous obtenons alors un foncteur exact à gauche
$f_* : \textit{Mod}(\mathcal{O}) \to \textit{Mod}(\mathcal{O}')$
qui donne lieu à l'{\it image directe dérivée}
$$
Rf_* : D^{+}(\mathcal{O}) \to D^+(\mathcal{O}')
$$
Le $i$-ième faisceau de cohomologie de $Rf_*\mathcal{F}^\bullet$ est noté
$R^if_*\mathcal{F}^\bullet$ et appelé la $i$-ième {\it image directe supérieure},
conformément à (\ref{equation-higher-direct-image-modules}).
Les foncteurs affichés ci-dessus sont des foncteurs exacts
de catégories dérivées.







\section{Première cohomologie et torseurs}
\label{section-h1-torsors}

\begin{definition}
\label{definition-torsor}
Soit $\mathcal{C}$ un site.
Soit $\mathcal{G}$ un faisceau de groupes (éventuellement non commutatifs)
sur $\mathcal{C}$.
Un {\it pseudo-torseur}, ou plus précisément un
{\it pseudo-$\mathcal{G}$-torseur}, est un faisceau
d'ensembles $\mathcal{F}$ sur $\mathcal{C}$ muni d'une action
$\mathcal{G} \times \mathcal{F} \to \mathcal{F}$ telle que
\begin{enumerate}
\item dès que $\mathcal{F}(U)$ est non vide, l'action
$\mathcal{G}(U) \times \mathcal{F}(U) \to \mathcal{F}(U)$
est simplement transitive.
\end{enumerate}
Un {\it morphisme de pseudo-$\mathcal{G}$-torseurs}
$\mathcal{F} \to \mathcal{F}'$
est un morphisme de faisceaux d'ensembles compatible aux
actions de $\mathcal{G}$.
Un {\it torseur}, ou plus précisément un
{\it $\mathcal{G}$-torseur}, est un pseudo-$\mathcal{G}$-torseur qui vérifie
en outre
\begin{enumerate}
\item[(2)] pour tout $U \in \Ob(\mathcal{C})$, il existe
un recouvrement $\{U_i \to U\}_{i \in I}$ de $U$
tel que $\mathcal{F}(U_i)$ soit non vide pour tout $i \in I$.
\end{enumerate}
Un {\it morphisme de $\mathcal{G}$-torseurs} est un morphisme de
pseudo-$\mathcal{G}$-torseurs.
Le {\it $\mathcal{G}$-torseur trivial}
est le faisceau $\mathcal{G}$ muni de l'action à gauche
évidente de $\mathcal{G}$.
\end{definition}

\noindent
Il est clair qu'un morphisme de torseurs est automatiquement un isomorphisme.

\begin{lemma}
\label{lemma-trivial-torsor}
Soit $\mathcal{C}$ un site.
Soit $\mathcal{G}$ un faisceau de groupes (éventuellement non commutatifs)
sur $\mathcal{C}$.
Un $\mathcal{G}$-torseur $\mathcal{F}$ est trivial si et seulement si
$\Gamma(\mathcal{C}, \mathcal{F}) \not = \emptyset$.
\end{lemma}

\begin{proof}
Omis.
\end{proof}

\begin{lemma}
\label{lemma-torsors-h1}
Soit $\mathcal{C}$ un site.
Soit $\mathcal{H}$ un faisceau abélien sur $\mathcal{C}$.
Il existe une bijection canonique entre l'ensemble des classes
d'isomorphisme de $\mathcal{H}$-torseurs et $H^1(\mathcal{C}, \mathcal{H})$.
\end{lemma}

\begin{proof}
Soit $\mathcal{F}$ un $\mathcal{H}$-torseur.
Considérons le faisceau abélien libre $\mathbf{Z}[\mathcal{F}]$
sur $\mathcal{F}$. C'est la faisceautisation de la règle
qui associe à $U \in \Ob(\mathcal{C})$ l'ensemble des sommes formelles
finies $\sum n_i[s_i]$, où $n_i \in \mathbf{Z}$
et $s_i \in \mathcal{F}(U)$. Il existe un morphisme naturel
$$
\sigma : \mathbf{Z}[\mathcal{F}] \longrightarrow \underline{\mathbf{Z}}
$$
qui, à une section locale $\sum n_i[s_i]$, associe $\sum n_i$.
Le noyau de $\sigma$ est engendré par les sections de la forme
$[s] - [s']$. Il existe un morphisme canonique
$a : \Ker(\sigma) \to \mathcal{H}$
qui envoie $[s] - [s'] \mapsto h$, où $h$ est la section locale de
$\mathcal{H}$ telle que $h \cdot s' = s$. Considérons le diagramme cocartésien
$$
\xymatrix{
0 \ar[r] &
\Ker(\sigma) \ar[r] \ar[d]^a &
\mathbf{Z}[\mathcal{F}] \ar[r] \ar[d] &
\underline{\mathbf{Z}} \ar[r] \ar[d] &
0 \\
0 \ar[r] &
\mathcal{H} \ar[r] &
\mathcal{E} \ar[r] &
\underline{\mathbf{Z}} \ar[r] &
0
}
$$
Ici, $\mathcal{E}$ est l'extension obtenue par somme amalgamée.
À partir de la suite exacte longue de cohomologie associée à la
suite exacte courte inférieure, nous obtenons un élément
$\xi = \xi_\mathcal{F} \in H^1(\mathcal{C}, \mathcal{H})$
en appliquant le morphisme de connexion à
$1 \in H^0(\mathcal{C}, \underline{\mathbf{Z}})$.

\medskip\noindent
Réciproquement, étant donné $\xi \in H^1(\mathcal{C}, \mathcal{H})$, nous
pouvons associer à $\xi$ un torseur comme suit. Choisissons un plongement
$\mathcal{H} \to \mathcal{I}$
de $\mathcal{H}$ dans un faisceau abélien injectif $\mathcal{I}$. Posons
$\mathcal{Q} = \mathcal{I}/\mathcal{H}$, de sorte que nous avons une suite
exacte courte
$$
\xymatrix{
0 \ar[r] &
\mathcal{H} \ar[r] &
\mathcal{I} \ar[r] &
\mathcal{Q} \ar[r] &
0
}
$$
L'élément $\xi$ est l'image d'une section globale
$q \in H^0(\mathcal{C}, \mathcal{Q})$
car $H^1(\mathcal{C}, \mathcal{I}) = 0$ (voir
Catégories dérivées, lemme \ref{derived-lemma-higher-derived-functors}).
Soit $\mathcal{F} \subset \mathcal{I}$ le sous-faisceau (d'ensembles) des sections
qui s'envoient sur $q$ dans le faisceau $\mathcal{Q}$. Il est facile de vérifier que
$\mathcal{F}$ est un $\mathcal{H}$-torseur.

\medskip\noindent
Nous omettons la vérification que les deux constructions données
ci-dessus sont réciproques l'une de l'autre.
\end{proof}






\section{Première cohomologie et extensions}
\label{section-h1-extensions}

\begin{lemma}
\label{lemma-h1-extensions}
Soit $(\mathcal{C}, \mathcal{O})$ un site annelé.
Soit $\mathcal{F}$ un faisceau de $\mathcal{O}$-modules sur $\mathcal{C}$.
Il existe une bijection canonique
$$
\Ext^1_{\textit{Mod}(\mathcal{O})}(\mathcal{O}, \mathcal{F})
\longrightarrow
H^1(\mathcal{C}, \mathcal{F})
$$
qui associe à l'extension
$$
0 \to \mathcal{F} \to \mathcal{E} \to \mathcal{O} \to 0
$$
l'image de $1 \in \Gamma(\mathcal{C}, \mathcal{O})$ dans
$H^1(\mathcal{C}, \mathcal{F})$.
\end{lemma}

\begin{proof}
Construisons l'inverse de l'application donnée dans le lemme.
Soit $\xi \in H^1(\mathcal{C}, \mathcal{F})$.
Choisissons une injection $\mathcal{F} \subset \mathcal{I}$, avec
$\mathcal{I}$ injectif dans $\textit{Mod}(\mathcal{O})$.
Posons $\mathcal{Q} = \mathcal{I}/\mathcal{F}$.
La suite exacte longue de cohomologie montre que
$\xi$ est l'image d'une section
$\tilde \xi \in \Gamma(\mathcal{C}, \mathcal{Q}) =
\Hom_\mathcal{O}(\mathcal{O}, \mathcal{Q})$.
Il suffit alors de former le produit fibré
$$
\xymatrix{
0 \ar[r] &
\mathcal{F} \ar[r] \ar@{=}[d] &
\mathcal{E} \ar[r] \ar[d] &
\mathcal{O} \ar[r] \ar[d]^{\tilde \xi} &
0 \\
0 \ar[r] &
\mathcal{F} \ar[r] &
\mathcal{I} \ar[r] &
\mathcal{Q} \ar[r] &
0
}
$$
voir Homologie, section \ref{homology-section-extensions}.
\end{proof}

\noindent
Le lemme suivant sera remplacé par le lemme plus général
\ref{lemma-cohomology-modules-abelian-agree}.

\begin{lemma}
\label{lemma-h1-mod-ab-agree}
Soit $(\mathcal{C}, \mathcal{O})$ un site annelé.
Soit $\mathcal{F}$ un faisceau de $\mathcal{O}$-modules sur $\mathcal{C}$.
Notons $\mathcal{F}_{ab}$ le faisceau de groupes abéliens
sous-jacent. Il existe alors un isomorphisme fonctoriel
$$
H^1(\mathcal{C}, \mathcal{F}_{ab})
=
H^1(\mathcal{C}, \mathcal{F})
$$
où le membre de gauche est la cohomologie calculée dans
$\textit{Ab}(\mathcal{C})$ et le membre de droite
est la cohomologie calculée dans $\textit{Mod}(\mathcal{O})$.
\end{lemma}

\begin{proof}
Notons $\underline{\mathbf{Z}}$ le faisceau constant
$\mathbf{Z}$. Comme
$\textit{Ab}(\mathcal{C}) = \textit{Mod}(\underline{\mathbf{Z}})$
nous pouvons appliquer deux fois le
lemme \ref{lemma-h1-extensions},
et il suffit donc de montrer
$$
\Ext^1_{\textit{Mod}(\mathcal{O})}(\mathcal{O}, \mathcal{F})
=
\Ext^1_{\textit{Mod}(\underline{\mathbf{Z}})}(
\underline{\mathbf{Z}}, \mathcal{F}_{ab}).
$$
Supposons que $0 \to \mathcal{F} \to \mathcal{E} \to \mathcal{O} \to 0$
soit une extension dans $\textit{Mod}(\mathcal{O})$. Nous pouvons alors utiliser
le morphisme évident de faisceaux abéliens
$1 : \underline{\mathbf{Z}} \to \mathcal{O}$
et former le produit fibré pour obtenir une extension $\mathcal{E}_{ab}$, comme suit :
$$
\xymatrix{
0 \ar[r] &
\mathcal{F}_{ab} \ar[r] \ar@{=}[d] &
\mathcal{E}_{ab} \ar[r] \ar[d] &
\underline{\mathbf{Z}} \ar[r] \ar[d]^{1} &
0 \\
0 \ar[r] &
\mathcal{F} \ar[r] &
\mathcal{E} \ar[r] &
\mathcal{O} \ar[r] &
0
}
$$
La réciproque est un peu plus plaisante. Supposons que
$0 \to \mathcal{F}_{ab} \to \mathcal{E}_{ab} \to \underline{\mathbf{Z}} \to 0$
soit une extension dans $\textit{Mod}(\underline{\mathbf{Z}})$.
Puisque $\underline{\mathbf{Z}}$ est un $\underline{\mathbf{Z}}$-module plat,
la suite
$$
0 \to \mathcal{F}_{ab} \otimes_{\underline{\mathbf{Z}}} \mathcal{O}
\to \mathcal{E}_{ab} \otimes_{\underline{\mathbf{Z}}} \mathcal{O}
\to \underline{\mathbf{Z}} \otimes_{\underline{\mathbf{Z}}} \mathcal{O}
\to 0
$$
est exacte, voir
Modules sur les sites, lemme \ref{sites-modules-lemma-flat-tor-zero}.
Bien entendu,
$\underline{\mathbf{Z}} \otimes_{\underline{\mathbf{Z}}} \mathcal{O}
= \mathcal{O}$.
Nous pouvons donc former la somme amalgamée par le morphisme de multiplication ($\mathcal{O}$-linéaire)
$\mu : \mathcal{F} \otimes_{\underline{\mathbf{Z}}} \mathcal{O}
\to \mathcal{F}$ afin d'obtenir une extension de $\mathcal{O}$ par
$\mathcal{F}$, comme suit :
$$
\xymatrix{
0 \ar[r] &
\mathcal{F}_{ab} \otimes_{\underline{\mathbf{Z}}} \mathcal{O}
\ar[r] \ar[d]^\mu &
\mathcal{E}_{ab} \otimes_{\underline{\mathbf{Z}}} \mathcal{O}
\ar[r] \ar[d] &
\mathcal{O} \ar[r] \ar@{=}[d] &
0 \\
0 \ar[r] &
\mathcal{F} \ar[r] &
\mathcal{E} \ar[r] &
\mathcal{O} \ar[r] &
0
}
$$
qui est l'extension cherchée. Nous omettons de vérifier que ces
constructions sont réciproques l'une de l'autre.
\end{proof}





\section{Première cohomologie et faisceaux inversibles}
\label{section-invertible-sheaves}

\noindent
Le groupe de Picard d'un site annelé est défini dans
Modules sur les sites, section \ref{sites-modules-section-invertible}.

\begin{lemma}
\label{lemma-h1-invertible}
Soit $(\mathcal{C}, \mathcal{O})$ un site localement annelé.
Il existe un isomorphisme canonique de groupes abéliens
$$
H^1(\mathcal{C}, \mathcal{O}^*) = \Pic(\mathcal{O}).
$$

\end{lemma}

\begin{proof}
Soit $\mathcal{L}$ un $\mathcal{O}$-module inversible.
Considérons le préfaisceau $\mathcal{L}^*$ défini par la règle
$$
U \longmapsto \{s \in \mathcal{L}(U)
\text{ such that } \mathcal{O}_U \xrightarrow{s \cdot -} \mathcal{L}_U
\text{ is an isomorphism}\}
$$
Ce préfaisceau vérifie la condition de faisceau. De plus, si
$f \in \mathcal{O}^*(U)$ et $s \in \mathcal{L}^*(U)$, alors il est clair que
$fs \in \mathcal{L}^*(U)$. De même, si $s, s' \in \mathcal{L}^*(U)$,
il existe un unique $f \in \mathcal{O}^*(U)$ tel que
$fs = s'$. En outre, le faisceau $\mathcal{L}^*$ possède localement des sections
d'après Modules sur les sites, lemme
\ref{sites-modules-lemma-invertible-is-locally-free-rank-1}.
Autrement dit, nous
voyons que $\mathcal{L}^*$ est un $\mathcal{O}^*$-torseur. Nous obtenons ainsi
une application
$$
\begin{matrix}
\text{set of invertible sheaves on }(\mathcal{C}, \mathcal{O}) \\
\text{ up to isomorphism}
\end{matrix}
\longrightarrow
\begin{matrix}
\text{set of }\mathcal{O}^*\text{-torsors} \\
\text{ up to isomorphism}
\end{matrix}
$$
Nous omettons de vérifier qu'il s'agit d'un morphisme de groupes abéliens.
D'après le
lemme \ref{lemma-torsors-h1},
le membre de droite est canoniquement
en bijection avec $H^1(\mathcal{C}, \mathcal{O}^*)$.
Il reste donc à montrer que cette application est injective et surjective.

\medskip\noindent
Injectivité. Si le torseur $\mathcal{L}^*$ est trivial, cela signifie, d'après le
lemme \ref{lemma-trivial-torsor},
que $\mathcal{L}^*$ possède une section globale.
Cela signifie donc précisément que $\mathcal{L} \cong \mathcal{O}$ est
l'élément neutre de $\Pic(\mathcal{O})$.

\medskip\noindent
Surjectivité. Soit $\mathcal{F}$ un $\mathcal{O}^*$-torseur.
Considérons le préfaisceau d'ensembles
$$
\mathcal{L}_1 : U \longmapsto
(\mathcal{F}(U) \times \mathcal{O}(U))/\mathcal{O}^*(U)
$$
où l'action de $f \in \mathcal{O}^*(U)$ sur
$(s, g)$ est donnée par $(fs, f^{-1}g)$. Alors $\mathcal{L}_1$ est un préfaisceau
de $\mathcal{O}$-modules si l'on pose
$(s, g) + (s', g') = (s, g + (s'/s)g')$, où $s'/s$ est la section locale
$f$ de $\mathcal{O}^*$ telle que $fs = s'$, et
$h(s, g) = (s, hg)$ pour toute section locale $h$ de $\mathcal{O}$.
Nous omettons de vérifier que la faisceautisation
$\mathcal{L} = \mathcal{L}_1^\#$ est un $\mathcal{O}$-module inversible
dont le $\mathcal{O}^*$-torseur associé $\mathcal{L}^*$ est isomorphe
à $\mathcal{F}$.
\end{proof}









\section{Localité de la cohomologie}
\label{section-locality}

\noindent
Le lemme suivant affirme qu'il n'y a aucune ambiguïté dans la définition de la cohomologie
d'un faisceau $\mathcal{F}$ sur un objet du site.

\begin{lemma}
\label{lemma-cohomology-of-open}
Soit $(\mathcal{C}, \mathcal{O})$ un site annelé.
Soit $U$ un objet de $\mathcal{C}$.
\begin{enumerate}
\item Si $\mathcal{I}$ est un $\mathcal{O}$-module injectif,
alors $\mathcal{I}|_U$ est un $\mathcal{O}_U$-module injectif.
\item Pour tout faisceau de $\mathcal{O}$-modules $\mathcal{F}$, nous avons
$H^p(U, \mathcal{F}) = H^p(\mathcal{C}/U, \mathcal{F}|_U)$.
\end{enumerate}
\end{lemma}

\begin{proof}
Rappelons que le foncteur $j_U^{-1}$ de restriction à $U$ est adjoint à droite
du foncteur $j_{U!}$ de prolongement par $0$, voir
Modules sur les sites, section
\ref{sites-modules-section-localize}.
De plus, $j_{U!}$ est exact. Ainsi, (1) découle de
Homologie, lemme \ref{homology-lemma-adjoint-preserve-injectives}.

\medskip\noindent
Par définition, $H^p(U, \mathcal{F}) = H^p(\mathcal{I}^\bullet(U))$,
où $\mathcal{F} \to \mathcal{I}^\bullet$ est une résolution injective
dans $\textit{Mod}(\mathcal{O})$.
D'après ce qui précède, $\mathcal{F}|_U \to \mathcal{I}^\bullet|_U$
est une résolution injective dans $\textit{Mod}(\mathcal{O}_U)$.
Par conséquent, $H^p(U, \mathcal{F}|_U)$ est égal à
$H^p(\mathcal{I}^\bullet|_U(U))$.
Bien entendu, $\mathcal{F}(U) = \mathcal{F}|_U(U)$ pour
tout faisceau $\mathcal{F}$ sur $\mathcal{C}$.
D'où l'égalité de (2).
\end{proof}

\noindent
Le lemme suivant servira à étudier ce qui se passe lorsque l'on change
d'univers partiel, ou à comparer la cohomologie des petits et grands sites
\'etales.

\begin{lemma}
\label{lemma-cohomology-bigger-site}
Soient $\mathcal{C}$ et $\mathcal{D}$ des sites.
Soit $u : \mathcal{C} \to \mathcal{D}$ un foncteur.
Supposons que $u$ vérifie les hypothèses de
Sites, lemme \ref{sites-lemma-bigger-site}.
Soit $g : \Sh(\mathcal{C}) \to \Sh(\mathcal{D})$
le morphisme de topos associé.
Pour tout faisceau abélien $\mathcal{F}$ sur $\mathcal{D}$, nous avons des
isomorphismes
$$
R\Gamma(\mathcal{C}, g^{-1}\mathcal{F}) = R\Gamma(\mathcal{D}, \mathcal{F}),
$$
en particulier
$H^p(\mathcal{C}, g^{-1}\mathcal{F}) = H^p(\mathcal{D}, \mathcal{F})$
et, pour tout $U \in \Ob(\mathcal{C})$, nous avons des isomorphismes
$$
R\Gamma(U, g^{-1}\mathcal{F}) = R\Gamma(u(U), \mathcal{F}),
$$
en particulier
$H^p(U, g^{-1}\mathcal{F}) = H^p(u(U), \mathcal{F})$. Tous ces
isomorphismes sont fonctoriels en $\mathcal{F}$.
\end{lemma}

\begin{proof}
Puisqu'il est clair que
$\Gamma(\mathcal{C}, g^{-1}\mathcal{F}) = \Gamma(\mathcal{D}, \mathcal{F})$
d'après l'hypothèse (e), il suffit de montrer que $g^{-1}$ envoie les faisceaux
abéliens injectifs sur des faisceaux abéliens injectifs. Comme d'habitude, nous utilisons
Homologie, lemme \ref{homology-lemma-adjoint-preserve-injectives},
pour le voir. L'adjoint à gauche de $g^{-1}$ est $g_! = f^{-1}$, avec les
notations de
Sites, lemme \ref{sites-lemma-bigger-site},
et ce foncteur est exact. Le lemme s'applique donc bien.
\end{proof}

\noindent
Soit $(\mathcal{C}, \mathcal{O})$ un site annelé.
Soit $\mathcal{F}$ un faisceau de $\mathcal{O}$-modules.
Soit $\varphi : U \to V$ un morphisme de $\mathcal{C}$.
Il existe alors un {\it morphisme de restriction} canonique
\begin{equation}
\label{equation-restriction-mapping}
H^n(V, \mathcal{F})
\longrightarrow
H^n(U, \mathcal{F}), \quad
\xi \longmapsto \xi|_U
\end{equation}
fonctoriel en $\mathcal{F}$. En effet, choisissons une résolution
injective quelconque $\mathcal{F} \to \mathcal{I}^\bullet$. Les morphismes de
restriction des faisceaux $\mathcal{I}^p$ donnent un morphisme de complexes
$$
\Gamma(V, \mathcal{I}^\bullet)
\longrightarrow
\Gamma(U, \mathcal{I}^\bullet)
$$
Le membre de gauche est un complexe représentant $R\Gamma(V, \mathcal{F})$
et celui de droite un complexe représentant $R\Gamma(U, \mathcal{F})$.
Nous obtenons le morphisme sur les groupes de cohomologie en appliquant le foncteur $H^n$.
Comme indiqué, nous noterons ce morphisme $\xi \mapsto \xi|_U$.
Ainsi, la règle $U \mapsto H^n(U, \mathcal{F})$ est un préfaisceau de
$\mathcal{O}$-modules. Ce préfaisceau est habituellement noté
$\underline{H}^n(\mathcal{F})$. Nous donnerons une autre interprétation
de ce préfaisceau dans le lemme \ref{lemma-include}.

\medskip\noindent
Le lemme suivant affirme qu'il est possible d'annuler les classes de cohomologie
supérieure en passant à un recouvrement.

\begin{lemma}
\label{lemma-kill-cohomology-class-on-covering}
Soit $(\mathcal{C}, \mathcal{O})$ un site annelé.
Soit $\mathcal{F}$ un faisceau de $\mathcal{O}$-modules.
Soit $U$ un objet de $\mathcal{C}$.
Soient $n > 0$ et $\xi \in H^n(U, \mathcal{F})$.
Il existe alors un recouvrement $\{U_i \to U\}$ de $\mathcal{C}$
tel que $\xi|_{U_i} = 0$ pour tout $i \in I$.
\end{lemma}

\begin{proof}
Soit $\mathcal{F} \to \mathcal{I}^\bullet$ une résolution injective.
Alors
$$
H^n(U, \mathcal{F}) =
\frac{\Ker(\mathcal{I}^n(U) \to \mathcal{I}^{n + 1}(U))}
{\Im(\mathcal{I}^{n - 1}(U) \to \mathcal{I}^n(U))}.
$$
Choisissons un élément $\tilde \xi \in \mathcal{I}^n(U)$ représentant la
classe de cohomologie dans la présentation ci-dessus. Puisque $\mathcal{I}^\bullet$
est une résolution injective de $\mathcal{F}$ et que $n > 0$, nous voyons que
le complexe $\mathcal{I}^\bullet$ est exact en degré $n$. Ainsi,
$\Im(\mathcal{I}^{n - 1} \to \mathcal{I}^n) =
\Ker(\mathcal{I}^n \to \mathcal{I}^{n + 1})$ comme faisceaux.
Puisque $\tilde \xi$ est une section sur $U$ du faisceau noyau,
nous en déduisons qu'il existe un recouvrement $\{U_i \to U\}$ du site
tel que $\tilde \xi|_{U_i}$ soit l'image par $d$ d'une section
$\xi_i \in \mathcal{I}^{n - 1}(U_i)$. Par définition, la
restriction $\xi|_{U_i}$ correspond à la classe de
$\tilde \xi|_{U_i}$, ce qui permet de conclure.
\end{proof}

\begin{lemma}
\label{lemma-higher-direct-images}
Soit $f : (\mathcal{C}, \mathcal{O}_\mathcal{C}) \to
(\mathcal{D}, \mathcal{O}_\mathcal{D})$ un morphisme de sites annelés
correspondant au foncteur continu $u : \mathcal{D} \to \mathcal{C}$.
Pour tout $\mathcal{F} \in \Ob(\textit{Mod}(\mathcal{O}_\mathcal{C}))$,
le faisceau $R^if_*\mathcal{F}$ est le faisceau associé au
préfaisceau
$$
V \longmapsto H^i(u(V), \mathcal{F})
$$
\end{lemma}

\begin{proof}
Soit $\mathcal{F} \to \mathcal{I}^\bullet$ une résolution injective.
Alors $R^if_*\mathcal{F}$ est par définition le $i$-ième faisceau de cohomologie
du complexe
$$
f_*\mathcal{I}^0 \to f_*\mathcal{I}^1 \to f_*\mathcal{I}^2 \to \ldots
$$
Par définition de la structure de catégorie abélienne sur les
$\mathcal{O}_\mathcal{D}$-modules,
ce faisceau de cohomologie est le faisceau associé au préfaisceau
$$
V
\longmapsto
\frac{\Ker(f_*\mathcal{I}^i(V) \to f_*\mathcal{I}^{i + 1}(V))}
{\Im(f_*\mathcal{I}^{i - 1}(V) \to f_*\mathcal{I}^i(V))}
$$
et ceci est manifestement égal à
$$
\frac{\Ker(\mathcal{I}^i(u(V)) \to \mathcal{I}^{i + 1}(u(V)))}
{\Im(\mathcal{I}^{i - 1}(u(V)) \to \mathcal{I}^i(u(V)))}
$$
qui est égal à $H^i(u(V), \mathcal{F})$,
ce qui conclut.
\end{proof}






\section{Le complexe de {\v C}ech et la cohomologie de {\v C}ech}
\label{section-cech}

\noindent
Soit $\mathcal{C}$ une catégorie. Soit $\mathcal{U} = \{U_i \to U\}_{i \in I}$
une famille de morphismes de but fixé, voir
Sites, définition \ref{sites-definition-family-morphisms-fixed-target}.
Supposons que tous les produits fibrés $U_{i_0} \times_U \ldots \times_U U_{i_p}$
existent dans $\mathcal{C}$. Soit $\mathcal{F}$ un préfaisceau abélien sur
$\mathcal{C}$. Posons
$$
\check{\mathcal{C}}^p(\mathcal{U}, \mathcal{F})
=
\prod\nolimits_{(i_0, \ldots, i_p) \in I^{p + 1}}
\mathcal{F}(U_{i_0} \times_U \ldots \times_U U_{i_p}).
$$
C'est un groupe abélien. Pour
$s \in \check{\mathcal{C}}^p(\mathcal{U}, \mathcal{F})$, nous notons
$s_{i_0\ldots i_p}$ sa composante dans le facteur
$\mathcal{F}(U_{i_0} \times_U \ldots \times_U U_{i_p})$.
Nous définissons
$$
d : \check{\mathcal{C}}^p(\mathcal{U}, \mathcal{F})
\longrightarrow
\check{\mathcal{C}}^{p + 1}(\mathcal{U}, \mathcal{F})
$$
par la formule
\begin{equation}
\label{equation-d-cech}
d(s)_{i_0\ldots i_{p + 1}} =
\sum\nolimits_{j = 0}^{p + 1}
(-1)^j s_{i_0\ldots \hat i_j \ldots i_{p + 1}}
|_{U_{i_0} \times_U \ldots \times_U U_{i_{p + 1}}}
\end{equation}
où la restriction est induite par la projection
$$
U_{i_0} \times_U \ldots \times_U U_{i_{p + 1}} \longrightarrow
U_{i_0} \times_U \ldots \times_U \widehat{U_{i_j}} \times_U
\ldots \times_U U_{i_{p + 1}}.
$$
On vérifie immédiatement que $d \circ d = 0$. Autrement dit,
$\check{\mathcal{C}}^\bullet(\mathcal{U}, \mathcal{F})$ est un complexe.

\begin{definition}
\label{definition-cech-complex}
Soit $\mathcal{C}$ une catégorie. Soit $\mathcal{U} = \{U_i \to U\}_{i \in I}$
une famille de morphismes de but fixé telle que tous les produits fibrés
$U_{i_0} \times_U \ldots \times_U U_{i_p}$ existent dans $\mathcal{C}$.
Soit $\mathcal{F}$ un préfaisceau abélien sur $\mathcal{C}$.
Le complexe $\check{\mathcal{C}}^\bullet(\mathcal{U}, \mathcal{F})$
est le {\it complexe de {\v C}ech} associé à $\mathcal{F}$ et à la
famille $\mathcal{U}$. Ses groupes de cohomologie
$H^i(\check{\mathcal{C}}^\bullet(\mathcal{U}, \mathcal{F}))$ sont
appelés les {\it groupes de cohomologie de {\v C}ech} de $\mathcal{F}$ relativement
à $\mathcal{U}$. Ils sont notés $\check H^i(\mathcal{U}, \mathcal{F})$.
\end{definition}

\noindent
Observons que tout recouvrement $\{U_i \to U\}$ d'un site $\mathcal{C}$
est une famille de morphismes de but fixé à laquelle la définition s'applique.

\begin{lemma}
\label{lemma-cech-h0}
Soit $\mathcal{C}$ un site.
Soit $\mathcal{F}$ un préfaisceau abélien sur $\mathcal{C}$.
Les assertions suivantes sont équivalentes :
\begin{enumerate}
\item $\mathcal{F}$ est un faisceau abélien sur $\mathcal{C}$ ;
\item pour tout recouvrement $\mathcal{U} = \{U_i \to U\}_{i \in I}$
du site $\mathcal{C}$, l'application naturelle
$$
\mathcal{F}(U) \to \check{H}^0(\mathcal{U}, \mathcal{F})
$$
(voir Sites, section \ref{sites-section-sheafification}) est bijective.
\end{enumerate}
\end{lemma}

\begin{proof}
Cela résulte du fait que la condition de faisceau exige précisément que
$\mathcal{F}(U) \to \check{H}^0(\mathcal{U}, \mathcal{F})$
soit bijective pour tout recouvrement de $\mathcal{C}$.
\end{proof}

\noindent
Soit $\mathcal{C}$ une catégorie. Soit $\mathcal{U} = \{U_i \to U\}_{i\in I}$
une famille de morphismes de $\mathcal{C}$ de but fixé telle que
tous les produits fibrés $U_{i_0} \times_U \ldots \times_U U_{i_p}$ existent dans
$\mathcal{C}$. Soit $\mathcal{V} = \{V_j \to V\}_{j\in J}$ une autre telle famille.
Les données $f : U \to V$, $\alpha : I \to J$ et $f_i : U_i \to V_{\alpha(i)}$
forment un morphisme de familles de morphismes de but fixé, voir
Sites, section \ref{sites-section-refinements}.
Dans ce cas, nous obtenons un morphisme de complexes de {\v C}ech
\begin{equation}
\label{equation-map-cech-complexes}
\varphi : \check{\mathcal{C}}^\bullet(\mathcal{V}, \mathcal{F})
\longrightarrow
\check{\mathcal{C}}^\bullet(\mathcal{U}, \mathcal{F})
\end{equation}
qui, en degré $p$, est donné par
$$
\varphi(s)_{i_0 \ldots i_p} =
(f_{i_0} \times \ldots \times f_{i_p})^*s_{\alpha(i_0) \ldots \alpha(i_p)}
$$


\section{La cohomologie de {\v C}ech comme foncteur sur les préfaisceaux}
\label{section-cech-functor}

\noindent
Attention : dans cette section, nous travaillons exclusivement avec des préfaisceaux abéliens
sur une catégorie. Les résultats sont complètement faux dans le
cadre des faisceaux et des catégories de faisceaux !

\medskip\noindent
Soit $\mathcal{C}$ une catégorie. Soit $\mathcal{U} = \{U_i \to U\}_{i \in I}$
une famille de morphismes de but fixé telle que tous les produits fibrés
$U_{i_0} \times_U \ldots \times_U U_{i_p}$ existent dans $\mathcal{C}$.
Soit $\mathcal{F}$ un préfaisceau abélien sur $\mathcal{C}$.
La construction
$$
\mathcal{F} \longmapsto \check{\mathcal{C}}^\bullet(\mathcal{U}, \mathcal{F})
$$
est fonctorielle en $\mathcal{F}$. En fait, c'est un foncteur
\begin{equation}
\label{equation-cech-functor}
\check{\mathcal{C}}^\bullet(\mathcal{U}, -) :
\textit{PAb}(\mathcal{C})
\longrightarrow
\text{Comp}^{+}(\textit{Ab})
\end{equation}
voir
Catégories dérivées, définition \ref{derived-definition-complexes-notation},
pour les notations. Rappelons que la catégorie des complexes bornés inférieurement
dans une catégorie abélienne est une catégorie abélienne, voir
Homologie, lemme \ref{homology-lemma-cat-cochain-abelian}.

\begin{lemma}
\label{lemma-cech-exact-presheaves}
Le foncteur donné par l'équation (\ref{equation-cech-functor})
est exact (voir Homologie, lemme \ref{homology-lemma-exact-functor}).
\end{lemma}

\begin{proof}
Pour tout objet $W$ de $\mathcal{C}$, le foncteur
$\mathcal{F} \mapsto \mathcal{F}(W)$ est un foncteur additif exact
de $\textit{PAb}(\mathcal{C})$ dans $\textit{Ab}$.
Les foncteurs donnant les termes $\check{\mathcal{C}}^p(\mathcal{U}, \mathcal{F})$
du complexe sont des produits de ces foncteurs exacts et sont donc exacts.
De plus, une suite de complexes est exacte si et seulement si la suite
des termes en chaque degré est exacte. Le lemme en découle.
\end{proof}

\begin{lemma}
\label{lemma-cech-cohomology-delta-functor-presheaves}
Soit $\mathcal{C}$ une catégorie.
Soit $\mathcal{U} = \{U_i \to U\}_{i \in I}$ une famille de morphismes
de but fixé telle que tous les produits fibrés
$U_{i_0} \times_U \ldots \times_U U_{i_p}$ existent dans $\mathcal{C}$.
Les foncteurs $\mathcal{F} \mapsto \check{H}^n(\mathcal{U}, \mathcal{F})$
forment un $\delta$-foncteur de la catégorie abélienne $\textit{PAb}(\mathcal{C})$
vers la catégorie des $\mathbf{Z}$-modules (voir
Homologie, définition \ref{homology-definition-cohomological-delta-functor}).
\end{lemma}

\begin{proof}
D'après le
lemme \ref{lemma-cech-exact-presheaves},
une suite exacte courte de préfaisceaux abéliens
$0 \to \mathcal{F}_1 \to \mathcal{F}_2 \to \mathcal{F}_3 \to 0$
est transformée en une suite exacte courte de complexes de
$\mathbf{Z}$-modules. Nous pouvons donc utiliser
Homologie, lemme \ref{homology-lemma-long-exact-sequence-cochain},
pour obtenir les morphismes de connexion
$\delta_{\mathcal{F}_1 \to \mathcal{F}_2 \to \mathcal{F}_3} :
\check{H}^n(\mathcal{U}, \mathcal{F}_3) \to
\check{H}^{n + 1}(\mathcal{U}, \mathcal{F}_1)$
et une suite exacte longue correspondante. Nous omettons de vérifier
que ces morphismes sont compatibles avec les morphismes entre suites exactes
courtes de préfaisceaux.
\end{proof}

\begin{lemma}
\label{lemma-cech-map-into}
Soit $\mathcal{C}$ une catégorie. Soit $\mathcal{U} = \{U_i \to U\}_{i \in I}$
une famille de morphismes de but fixé telle que tous les produits fibrés
$U_{i_0} \times_U \ldots \times_U U_{i_p}$ existent dans $\mathcal{C}$.
Considérons le complexe de chaînes $\mathbf{Z}_{\mathcal{U}, \bullet}$
de préfaisceaux abéliens
$$
\ldots
\to
\bigoplus_{i_0i_1i_2} \mathbf{Z}_{U_{i_0} \times_U U_{i_1} \times_U U_{i_2}}
\to
\bigoplus_{i_0i_1} \mathbf{Z}_{U_{i_0} \times_U U_{i_1}}
\to
\bigoplus_{i_0} \mathbf{Z}_{U_{i_0}}
\to 0 \to \ldots
$$
où le dernier terme non nul est placé en degré $0$
et où le morphisme
$$
\mathbf{Z}_{U_{i_0} \times_U \ldots \times_U U_{i_{p + 1}}}
\longrightarrow
\mathbf{Z}_{U_{i_0} \times_U
\ldots \widehat{U_{i_j}} \ldots \times_U U_{i_{p + 1}}}
$$
est égal au morphisme canonique multiplié par $(-1)^j$.
Il existe alors un isomorphisme
$$
\Hom_{\textit{PAb}(\mathcal{C})}(\mathbf{Z}_{\mathcal{U}, \bullet}, \mathcal{F})
=
\check{\mathcal{C}}^\bullet(\mathcal{U}, \mathcal{F})
$$
fonctoriel en $\mathcal{F} \in \Ob(\textit{PAb}(\mathcal{C}))$.
\end{lemma}

\begin{proof}
C'est une tautologie qui repose sur le fait que
\begin{align*}
\Hom_{\textit{PAb}(\mathcal{C})}(
\bigoplus_{i_0 \ldots i_p}
\mathbf{Z}_{U_{i_0} \times_U \ldots \times_U U_{i_p}},
\mathcal{F})
& =
\prod_{i_0 \ldots i_p}
\Hom_{\textit{PAb}(\mathcal{C})}(
\mathbf{Z}_{U_{i_0} \times_U \ldots \times_U U_{i_p}},
\mathcal{F}) \\
& =
\prod_{i_0 \ldots i_p}
\mathcal{F}(U_{i_0} \times_U \ldots \times_U U_{i_p})
\end{align*}
voir Modules sur les sites, lemme \ref{sites-modules-lemma-obvious-adjointness}.
\end{proof}

\begin{lemma}
\label{lemma-homology-complex}
Soit $\mathcal{C}$ une catégorie. Soit
$\mathcal{U} = \{f_i : U_i \to U\}_{i \in I}$ une famille de morphismes
de but fixé telle que tous les produits fibrés
$U_{i_0} \times_U \ldots \times_U U_{i_p}$ existent dans $\mathcal{C}$.
Le complexe de chaînes de préfaisceaux $\mathbf{Z}_{\mathcal{U}, \bullet}$
du lemme \ref{lemma-cech-map-into} ci-dessus est exact en degrés
strictement positifs, c'est-à-dire que les préfaisceaux d'homologie
$H_i(\mathbf{Z}_{\mathcal{U}, \bullet})$ sont nuls pour $i > 0$.
\end{lemma}

\begin{proof}
Soit $V$ un objet de $\mathcal{C}$. Nous devons montrer que le complexe de chaînes
de groupes abéliens $\mathbf{Z}_{\mathcal{U}, \bullet}(V)$ est exact en
degrés $> 0$. Il s'agit du complexe
$$
\xymatrix{
\ldots \ar[d] \\
\bigoplus_{i_0i_1i_2}
\mathbf{Z}[
\Mor_\mathcal{C}(V, U_{i_0} \times_U U_{i_1} \times_U U_{i_2})
]
\ar[d] \\
\bigoplus_{i_0i_1}
\mathbf{Z}[
\Mor_\mathcal{C}(V, U_{i_0} \times_U U_{i_1})
]
\ar[d] \\
\bigoplus_{i_0}
\mathbf{Z}[
\Mor_\mathcal{C}(V, U_{i_0})
] \ar[d] \\
0
}
$$
Pour tout morphisme $\varphi : V \to U$, notons
$\Mor_\varphi(V, U_i) = \{\varphi_i : V \to U_i \mid
f_i \circ \varphi_i = \varphi\}$. Nous utiliserons une notation analogue
pour $\Mor_\varphi(V, U_{i_0} \times_U \ldots \times_U U_{i_p})$.
Remarquons que la composition avec les différentes projections entre les
produits fibrés $U_{i_0} \times_U \ldots \times_U U_{i_p}$ préserve
ces ensembles de morphismes. Nous voyons donc que le complexe ci-dessus
est le même que le complexe
$$
\xymatrix{
\ldots \ar[d] \\
\bigoplus_\varphi
\bigoplus_{i_0i_1i_2}
\mathbf{Z}[
\Mor_\varphi(V, U_{i_0} \times_U U_{i_1} \times_U U_{i_2})
]
\ar[d] \\
\bigoplus_\varphi
\bigoplus_{i_0i_1}
\mathbf{Z}[
\Mor_\varphi(V, U_{i_0} \times_U U_{i_1})
]
\ar[d] \\
\bigoplus_\varphi
\bigoplus_{i_0}
\mathbf{Z}[
\Mor_\varphi(V, U_{i_0})
] \ar[d] \\
0
}
$$
Ensuite, remarquons que nous avons
$$
\Mor_\varphi(V, U_{i_0} \times_U \ldots \times_U U_{i_p})
=
\Mor_\varphi(V, U_{i_0}) \times \ldots
\times \Mor_\varphi(V, U_{i_p})
$$
En utilisant ceci et le fait que $\mathbf{Z}[A] \oplus \mathbf{Z}[B] =
\mathbf{Z}[A \amalg B]$, nous voyons que le complexe devient
$$
\xymatrix{
\ldots \ar[d] \\
\bigoplus_\varphi
\mathbf{Z}\left[
\coprod_{i_0i_1i_2}
\Mor_\varphi(V, U_{i_0}) \times \Mor_\varphi(V, U_{i_1}) \times
\Mor_\varphi(V, U_{i_2})
\right]
\ar[d] \\
\bigoplus_\varphi
\mathbf{Z}\left[
\coprod_{i_0i_1}
\Mor_\varphi(V, U_{i_0}) \times \Mor_\varphi(V, U_{i_1})
\right]
\ar[d] \\
\bigoplus_\varphi
\mathbf{Z}\left[
\coprod_{i_0}
\Mor_\varphi(V, U_{i_0})
\right] \ar[d] \\
0
}
$$
Enfin, en posant $S_\varphi = \coprod_{i \in I} \Mor_\varphi(V, U_i)$,
nous obtenons
$$
\bigoplus\nolimits_\varphi \left(\ldots \to
\mathbf{Z}[S_\varphi \times S_\varphi \times S_\varphi] \to
\mathbf{Z}[S_\varphi \times S_\varphi] \to
\mathbf{Z}[S_\varphi] \to 0 \to \ldots
\right)
$$
Nous avons ainsi simplifié notre tâche. Il suffit en effet de montrer que,
pour tout ensemble non vide $S$, le complexe (augmenté) de groupes abéliens libres
$$
\ldots \to
\mathbf{Z}[S \times S \times S] \to
\mathbf{Z}[S \times S] \to
\mathbf{Z}[S] \xrightarrow{\Sigma} \mathbf{Z} \to 0 \to \ldots
$$
est exact en tout degré. Pour le voir, fixons un élément $s \in S$ et
utilisons l'homotopie
$$
n_{(s_0, \ldots, s_p)} \longmapsto n_{(s, s_0, \ldots, s_p)}
$$
avec les notations évidentes.
\end{proof}

\begin{lemma}
\label{lemma-complex-tensored-still-exact}
\begin{slogan}
Le complexe de préfaisceaux de {\v C}ech à coefficients entiers est une résolution plate du
préfaisceau constant des entiers.
\end{slogan}
Soit $\mathcal{C}$ une catégorie. Soit
$\mathcal{U} = \{f_i : U_i \to U\}_{i \in I}$ une famille de morphismes
de but fixé telle que tous les produits fibrés
$U_{i_0} \times_U \ldots \times_U U_{i_p}$ existent dans $\mathcal{C}$.
Soit $\mathcal{O}$ un préfaisceau d'anneaux sur $\mathcal{C}$.
Le complexe de chaînes
$$
\mathbf{Z}_{\mathcal{U}, \bullet}
\otimes_{p, \mathbf{Z}}
\mathcal{O}
$$
est exact en degrés strictement positifs. Ici, $\mathbf{Z}_{\mathcal{U}, \bullet}$
est le complexe de chaînes du lemme \ref{lemma-cech-map-into}, et
le produit tensoriel est pris sur le préfaisceau constant d'anneaux
de valeur $\mathbf{Z}$.
\end{lemma}

\begin{proof}
Soit $V$ un objet de $\mathcal{C}$.
Dans la démonstration du lemme \ref{lemma-homology-complex}, nous avons vu que
$\mathbf{Z}_{\mathcal{U}, \bullet}(V)$ est isomorphe, comme complexe,
à une somme directe de complexes homotopes à $\mathbf{Z}$
placé en degré zéro. Par conséquent,
$\mathbf{Z}_{\mathcal{U}, \bullet}(V) \otimes_\mathbf{Z} \mathcal{O}(V)$
est également isomorphe, comme complexe, à une somme directe de complexes homotopes
à $\mathcal{O}(V)$ placé en degré zéro.
On peut aussi utiliser
Modules sur les sites, lemme \ref{sites-modules-lemma-flat-resolution-of-flat},
qui s'applique puisque les préfaisceaux $\mathbf{Z}_{\mathcal{U}, i}$ sont plats,
et la démonstration du lemme \ref{lemma-homology-complex} montre que
$H_0(\mathbf{Z}_{\mathcal{U}, \bullet})$ est lui aussi un préfaisceau plat.
\end{proof}

\begin{lemma}
\label{lemma-cech-cohomology-derived-presheaves}
Soit $\mathcal{C}$ une catégorie. Soit
$\mathcal{U} = \{f_i : U_i \to U\}_{i \in I}$ une famille de morphismes
de but fixé telle que tous les produits fibrés
$U_{i_0} \times_U \ldots \times_U U_{i_p}$ existent dans $\mathcal{C}$.
Les foncteurs de cohomologie de {\v C}ech $\check{H}^p(\mathcal{U}, -)$
sont canoniquement isomorphes, en tant que $\delta$-foncteur, aux
foncteurs dérivés à droite du foncteur
$$
\check{H}^0(\mathcal{U}, -) :
\textit{PAb}(\mathcal{C})
\longrightarrow
\textit{Ab}.
$$
De plus, il existe un quasi-isomorphisme fonctoriel
$$
\check{\mathcal{C}}^\bullet(\mathcal{U}, \mathcal{F})
\longrightarrow
R\check{H}^0(\mathcal{U}, \mathcal{F})
$$
où le membre de droite désigne le foncteur dérivé
$$
R\check{H}^0(\mathcal{U}, -) :
D^{+}(\textit{PAb}(\mathcal{C}))
\longrightarrow
D^{+}(\mathbf{Z})
$$
du foncteur exact à gauche $\check{H}^0(\mathcal{U}, -)$.
\end{lemma}

\begin{proof}
Remarquons que la catégorie des préfaisceaux abéliens possède suffisamment d'injectifs, voir
Injectifs, proposition \ref{injectives-proposition-presheaves-injectives}.
Remarquons que $\check{H}^0(\mathcal{U}, -)$ est un foncteur exact à gauche
de la catégorie des préfaisceaux abéliens
vers la catégorie des $\mathbf{Z}$-modules.
Le foncteur dérivé et le foncteur dérivé à droite existent donc, voir
Catégories dérivées, section \ref{derived-section-right-derived-functor}.

\medskip\noindent
Soit $\mathcal{I}$ un préfaisceau abélien injectif.
Dans ce cas, le foncteur
$\Hom_{\textit{PAb}(\mathcal{C})}(-, \mathcal{I})$
est exact sur $\textit{PAb}(\mathcal{C})$. D'après le
lemme \ref{lemma-cech-map-into}, nous avons
$$
\Hom_{\textit{PAb}(\mathcal{C})}(
\mathbf{Z}_{\mathcal{U}, \bullet}, \mathcal{I})
=
\check{\mathcal{C}}^\bullet(\mathcal{U}, \mathcal{I}).
$$
D'après le lemme \ref{lemma-homology-complex}, le complexe
$\mathbf{Z}_{\mathcal{U}, \bullet}$ est exact en degrés strictement positifs.
L'exactitude du foncteur Hom à valeurs dans $\mathcal{I}$ mentionnée ci-dessus montre donc
que $\check{H}^i(\mathcal{U}, \mathcal{I}) = 0$ pour tout
$i > 0$. Ainsi, le $\delta$-foncteur $(\check{H}^n, \delta)$
(voir le lemme \ref{lemma-cech-cohomology-delta-functor-presheaves})
vérifie les hypothèses de
Homologie, lemme \ref{homology-lemma-efface-implies-universal},
et est donc un $\delta$-foncteur universel.

\medskip\noindent
D'après
Catégories dérivées, lemme \ref{derived-lemma-higher-derived-functors},
la famille $R^i\check{H}^0(\mathcal{U}, -)$
forme elle aussi un $\delta$-foncteur universel. Par unicité des
$\delta$-foncteurs universels, voir
Homologie, lemme \ref{homology-lemma-uniqueness-universal-delta-functor},
nous concluons que
$R^i\check{H}^0(\mathcal{U}, -) = \check{H}^i(\mathcal{U}, -)$.
Cela suffit pour la plupart des applications,
et nous suggérons au lecteur de passer la suite de la démonstration.

\medskip\noindent
Soit $\mathcal{F}$ un préfaisceau abélien quelconque sur $\mathcal{C}$.
Choisissons une résolution injective $\mathcal{F} \to \mathcal{I}^\bullet$
dans la catégorie $\textit{PAb}(\mathcal{C})$.
Considérons le complexe double
$\check{\mathcal{C}}^\bullet(\mathcal{U}, \mathcal{I}^\bullet)$
de termes $\check{\mathcal{C}}^p(\mathcal{U}, \mathcal{I}^q)$.
Considérons ensuite le complexe total
$\text{Tot}(\check{\mathcal{C}}^\bullet(\mathcal{U}, \mathcal{I}^\bullet))$
associé à ce complexe double, voir
Homologie, section \ref{homology-section-double-complexes}.
Il existe un morphisme de complexes
$$
\check{\mathcal{C}}^\bullet(\mathcal{U}, \mathcal{F})
\longrightarrow
\text{Tot}(\check{\mathcal{C}}^\bullet(\mathcal{U}, \mathcal{I}^\bullet))
$$
provenant des morphismes
$\check{\mathcal{C}}^p(\mathcal{U}, \mathcal{F})
\to \check{\mathcal{C}}^p(\mathcal{U}, \mathcal{I}^0)$
et il existe un morphisme de complexes
$$
\check{H}^0(\mathcal{U}, \mathcal{I}^\bullet)
\longrightarrow
\text{Tot}(\check{\mathcal{C}}^\bullet(\mathcal{U}, \mathcal{I}^\bullet))
$$
provenant des morphismes
$\check{H}^0(\mathcal{U}, \mathcal{I}^q) \to
\check{\mathcal{C}}^0(\mathcal{U}, \mathcal{I}^q)$.
Ces deux morphismes sont des quasi-isomorphismes par application de
Homologie, lemme \ref{homology-lemma-double-complex-gives-resolution}.
En effet, les colonnes du complexe double sont exactes en degrés strictement positifs
parce que le complexe de {\v C}ech est exact en tant que foncteur
(lemme \ref{lemma-cech-exact-presheaves}),
et les lignes du complexe double sont exactes en degrés strictement positifs
puisque, comme nous venons de le voir, les groupes de cohomologie supérieure de {\v C}ech des
préfaisceaux injectifs $\mathcal{I}^q$ sont nuls.
Puisque les quasi-isomorphismes deviennent inversibles
dans $D^{+}(\mathbf{Z})$, ceci donne le dernier morphisme affiché
du lemme. Nous omettons de vérifier que ce morphisme est
fonctoriel.
\end{proof}





\section{Cohomologie de {\v C}ech et cohomologie}
\label{section-cech-cohomology-cohomology}

\noindent
La relation entre cohomologie et cohomologie de {\v C}ech vient de ce que la
cohomologie de {\v C}ech d'un faisceau abélien injectif est nulle. Pour le voir,
on remarque qu'un faisceau abélien injectif est un préfaisceau abélien injectif,
puis on applique les résultats de cohomologie de {\v C}ech établis dans la
section précédente.

\begin{lemma}
\label{lemma-injective-abelian-sheaf-injective-presheaf}
Soit $\mathcal{C}$ un site. Tout faisceau abélien injectif est également
injectif en tant qu'objet de la catégorie $\textit{PAb}(\mathcal{C})$.
\end{lemma}

\begin{proof}
Appliquons le lemme \ref{homology-lemma-adjoint-preserve-injectives}
d'Homologie aux catégories $\mathcal{A} = \textit{Ab}(\mathcal{C})$ et
$\mathcal{B} = \textit{PAb}(\mathcal{C})$, au foncteur d'inclusion et à la
faisceautisation. (Voir Modules sur les sites, section
\ref{sites-modules-section-abelian-sheaves}, pour constater que toutes les
hypothèses du lemme sont satisfaites.)
\end{proof}

\begin{lemma}
\label{lemma-injective-trivial-cech}
Soit $\mathcal{C}$ un site.
Soit $\mathcal{U} = \{U_i \to U\}_{i \in I}$ un recouvrement de
$\mathcal{C}$. Soit $\mathcal{I}$ un faisceau abélien injectif, c'est-à-dire
un objet injectif de $\textit{Ab}(\mathcal{C})$. Alors
$$
\check{H}^p(\mathcal{U}, \mathcal{I}) =
\left\{
\begin{matrix}
\mathcal{I}(U) & \text{if} & p = 0 \\
0 & \text{if} & p > 0
\end{matrix}
\right.
$$
\end{lemma}

\begin{proof}
D'après le lemme \ref{lemma-injective-abelian-sheaf-injective-presheaf},
$\mathcal{I}$ est un objet injectif de $\textit{PAb}(\mathcal{C})$.
Nous pouvons donc appliquer le
lemme \ref{lemma-cech-cohomology-derived-presheaves} (ou sa preuve) pour
obtenir l'annulation des groupes de cohomologie de {\v C}ech supérieurs.
Pour le degré zéro, voir le lemme \ref{lemma-cech-h0}.
\end{proof}

\begin{lemma}
\label{lemma-cech-cohomology}
Soit $\mathcal{C}$ un site.
Soit $\mathcal{U} = \{U_i \to U\}_{i \in I}$ un recouvrement de
$\mathcal{C}$. Il existe une transformation
$$
\check{\mathcal{C}}^\bullet(\mathcal{U}, -)
\longrightarrow
R\Gamma(U, -)
$$
de foncteurs $\textit{Ab}(\mathcal{C}) \to D^{+}(\mathbf{Z})$.
En particulier, elle fournit une transformation de foncteurs
$\check{H}^p(\mathcal{U}, \mathcal{F}) \to H^p(U, \mathcal{F})$, lorsque
$\mathcal{F}$ parcourt $\textit{Ab}(\mathcal{C})$.
\end{lemma}

\begin{proof}
Soit $\mathcal{F}$ un faisceau abélien. Choisissons une résolution injective
$\mathcal{F} \to \mathcal{I}^\bullet$. Considérons le bicomplexe
$\check{\mathcal{C}}^\bullet(\mathcal{U}, \mathcal{I}^\bullet)$
dont les termes sont $\check{\mathcal{C}}^p(\mathcal{U}, \mathcal{I}^q)$.
Considérons ensuite le complexe total associé
$\text{Tot}(\check{\mathcal{C}}^\bullet(\mathcal{U}, \mathcal{I}^\bullet))$,
voir Homologie, définition
\ref{homology-definition-associated-simple-complex}. Il existe un morphisme
de complexes
$$
\alpha :
\Gamma(U, \mathcal{I}^\bullet)
\longrightarrow
\text{Tot}(\check{\mathcal{C}}^\bullet(\mathcal{U}, \mathcal{I}^\bullet))
$$
provenant des morphismes
$\mathcal{I}^q(U) \to \check{H}^0(\mathcal{U}, \mathcal{I}^q)$
et un morphisme de complexes
$$
\beta :
\check{\mathcal{C}}^\bullet(\mathcal{U}, \mathcal{F})
\longrightarrow
\text{Tot}(\check{\mathcal{C}}^\bullet(\mathcal{U}, \mathcal{I}^\bullet))
$$
provenant du morphisme $\mathcal{F} \to \mathcal{I}^0$.
Le lemme \ref{homology-lemma-double-complex-gives-resolution} d'Homologie
montre que $\alpha$ est un quasi-isomorphisme.
En effet, le lemme \ref{lemma-injective-trivial-cech} implique que la
$q$-ième ligne du bicomplexe
$\check{\mathcal{C}}^\bullet(\mathcal{U}, \mathcal{I}^\bullet)$ est une
résolution de $\Gamma(U, \mathcal{I}^q)$.
Ainsi, $\alpha$ devient inversible dans $D^{+}(\mathbf{Z})$, et la
transformation du lemme est la composée de $\beta$ avec l'inverse de
$\alpha$. Nous omettons la vérification de la fonctorialité.
\end{proof}

\begin{lemma}
\label{lemma-cech-h1}
Soit $\mathcal{C}$ un site. Soit $\mathcal{G}$ un faisceau abélien sur
$\mathcal{C}$. Soit $\mathcal{U} = \{U_i \to U\}_{i \in I}$ un recouvrement
de $\mathcal{C}$. Le morphisme
$$
\check{H}^1(\mathcal{U}, \mathcal{G})
\longrightarrow
H^1(U, \mathcal{G})
$$
est injectif et, par la bijection du lemme \ref{lemma-torsors-h1}, identifie
$\check{H}^1(\mathcal{U}, \mathcal{G})$ à l'ensemble des classes
d'isomorphisme de $\mathcal{G}|_U$-torseurs dont les restrictions à chaque
$U_i$ sont des torseurs triviaux.
\end{lemma}

\begin{proof}
Pour le voir, construisons un morphisme inverse. Soit $\mathcal{F}$ un
$\mathcal{G}|_U$-torseur sur $\mathcal{C}/U$ dont la restriction à
$\mathcal{C}/U_i$ est triviale. D'après le lemme
\ref{lemma-trivial-torsor}, cela signifie qu'il existe une section
$s_i \in \mathcal{F}(U_i)$. Sur $U_{i_0} \times_U U_{i_1}$, il existe une
unique section $s_{i_0i_1}$ de $\mathcal{G}$ telle que
$s_{i_0i_1} \cdot s_{i_0}|_{U_{i_0} \times_U U_{i_1}} =
s_{i_1}|_{U_{i_0} \times_U U_{i_1}}$. Un calcul immédiat montre que
$s_{i_0i_1}$ est un cocycle de {\v C}ech et que sa classe est bien définie
(c'est-à-dire indépendante du choix des sections $s_i$).
L'inverse envoie la classe d'isomorphisme de $\mathcal{F}$ sur la classe de
cohomologie du cocycle $(s_{i_0i_1})$. Nous omettons de vérifier que ce
morphisme est bien un inverse.
\end{proof}

\begin{lemma}
\label{lemma-include}
Soit $\mathcal{C}$ un site.
Considérons le foncteur
$i : \textit{Ab}(\mathcal{C}) \to \textit{PAb}(\mathcal{C})$.
C'est un foncteur exact à gauche dont les foncteurs dérivés à droite sont
donnés par
$$
R^pi(\mathcal{F}) = \underline{H}^p(\mathcal{F}) :
U \longmapsto H^p(U, \mathcal{F})
$$
voir la discussion de la section \ref{section-locality}.
\end{lemma}

\begin{proof}
Il est clair que $i$ est exact à gauche.
Choisissons une résolution injective $\mathcal{F} \to \mathcal{I}^\bullet$.
Par définition, $R^pi$ est le {\it préfaisceau} de cohomologie de degré $p$
du complexe $\mathcal{I}^\bullet$. Autrement dit, les sections de
$R^pi(\mathcal{F})$ sur un objet $U$ de $\mathcal{C}$ sont données par
$$
\frac{\Ker(\mathcal{I}^n(U) \to \mathcal{I}^{n + 1}(U))}
{\Im(\mathcal{I}^{n - 1}(U) \to \mathcal{I}^n(U))}.
$$
ce qui est la définition de $H^p(U, \mathcal{F})$.
\end{proof}

\begin{lemma}
\label{lemma-cech-spectral-sequence}
Soit $\mathcal{C}$ un site. Soit $\mathcal{U} = \{U_i \to U\}_{i \in I}$
un recouvrement de $\mathcal{C}$. Pour tout faisceau abélien $\mathcal{F}$,
il existe une suite spectrale $(E_r, d_r)_{r \geq 0}$ telle que
$$
E_2^{p, q} = \check{H}^p(\mathcal{U}, \underline{H}^q(\mathcal{F}))
$$
qui converge vers $H^{p + q}(U, \mathcal{F})$.
Cette suite spectrale est fonctorielle en $\mathcal{F}$.
\end{lemma}

\begin{proof}
C'est la suite spectrale de Grothendieck (voir Catégories dérivées,
lemme \ref{derived-lemma-grothendieck-spectral-sequence}) associée aux
foncteurs
$$
i :  \textit{Ab}(\mathcal{C}) \to \textit{PAb}(\mathcal{C})
\quad\text{and}\quad
\check{H}^0(\mathcal{U}, - ) : \textit{PAb}(\mathcal{C})
\to \textit{Ab}.
$$
En effet, $\check{H}^0(\mathcal{U}, i(\mathcal{F})) = \mathcal{F}(U)$ d'après
le lemme \ref{lemma-cech-h0}. De plus, $i(\mathcal{I})$ est acyclique pour la
cohomologie de {\v C}ech d'après le
lemme \ref{lemma-injective-trivial-cech}. Enfin,
$\check{H}^p(\mathcal{U}, -) = R^p\check{H}^0(\mathcal{U}, -)$ comme foncteurs
sur $\textit{PAb}(\mathcal{C})$, d'après le
lemme \ref{lemma-cech-cohomology-derived-presheaves}.
La réunion de ces observations prouve le lemme.
\end{proof}

\begin{lemma}
\label{lemma-cech-spectral-sequence-application}
Soit $\mathcal{C}$ un site.
Soit $\mathcal{U} = \{U_i \to U\}_{i \in I}$ un recouvrement.
Soit $\mathcal{F} \in \Ob(\textit{Ab}(\mathcal{C}))$.
Supposons que $H^i(U_{i_0} \times_U \ldots \times_U U_{i_p}, \mathcal{F}) = 0$
pour tous $i > 0$, $p \geq 0$ et $i_0, \ldots, i_p \in I$.
Alors $\check{H}^p(\mathcal{U}, \mathcal{F}) = H^p(U, \mathcal{F})$.
\end{lemma}

\begin{proof}
Nous utilisons la suite spectrale du
lemme \ref{lemma-cech-spectral-sequence}.
Les hypothèses signifient que $E_2^{p, q} = 0$ pour tout $(p, q)$ tel que
$q \not = 0$. La suite spectrale dégénère donc en $E_2$, d'où le résultat.
\end{proof}

\begin{lemma}
\label{lemma-ses-cech-h1}
Soit $\mathcal{C}$ un site.
Soit
$$
0 \to \mathcal{F} \to \mathcal{G} \to \mathcal{H} \to 0
$$
une suite exacte courte de faisceaux abéliens sur $\mathcal{C}$.
Soit $U$ un objet de $\mathcal{C}$. S'il existe un système cofinal de
recouvrements $\mathcal{U}$ de $U$ tels que
$\check{H}^1(\mathcal{U}, \mathcal{F}) = 0$,
alors le morphisme $\mathcal{G}(U) \to \mathcal{H}(U)$ est surjectif.
\end{lemma}

\begin{proof}
Prenons un élément $s \in \mathcal{H}(U)$. Choisissons un recouvrement
$\mathcal{U} = \{U_i \to U\}_{i \in I}$ tel que
(a) $\check{H}^1(\mathcal{U}, \mathcal{F}) = 0$ et (b) $s|_{U_i}$ soit
l'image d'une section $s_i \in \mathcal{G}(U_i)$.
Comme nous pouvons certainement trouver un recouvrement vérifiant (b), les
hypothèses du lemme permettent d'en trouver un qui vérifie à la fois (a) et
(b). Considérons les sections
$$
s_{i_0i_1} =
s_{i_1}|_{U_{i_0} \times_U U_{i_1}} - s_{i_0}|_{U_{i_0} \times_U U_{i_1}}.
$$
Puisque $s_i$ relève $s$, nous avons
$s_{i_0i_1} \in \mathcal{F}(U_{i_0} \times_U U_{i_1})$.
L'annulation de $\check{H}^1(\mathcal{U}, \mathcal{F})$ fournit des sections
$t_i \in \mathcal{F}(U_i)$ telles que
$$
s_{i_0i_1} =
t_{i_1}|_{U_{i_0} \times_U U_{i_1}} - t_{i_0}|_{U_{i_0} \times_U U_{i_1}}.
$$
Les sections $s_i - t_i$ satisfont alors manifestement à la condition de
faisceau et se recollent en une section de $\mathcal{G}$ sur $U$ dont l'image
est $s$. Ceci conclut la preuve.
\end{proof}

\begin{lemma}
\label{lemma-cech-vanish-collection}
(Variante de Cohomologie, lemme \ref{cohomology-lemma-cech-vanish}.)
Soit $\mathcal{C}$ un site. Soit $\text{Cov}_\mathcal{C}$ l'ensemble des
recouvrements de $\mathcal{C}$ (voir Sites, définition
\ref{sites-definition-site}). Soient $\mathcal{B} \subset \Ob(\mathcal{C})$
et $\text{Cov} \subset \text{Cov}_\mathcal{C}$ des sous-ensembles. Soit
$\mathcal{F}$ un faisceau abélien sur $\mathcal{C}$. Supposons que
\begin{enumerate}
\item pour tout $\mathcal{U} \in \text{Cov}$,
$\mathcal{U} = \{U_i \to U\}_{i \in I}$, on ait
$U, U_i \in \mathcal{B}$ et tout
$U_{i_0} \times_U \ldots \times_U U_{i_p} \in \mathcal{B}$.
\item pour tout $U \in \mathcal{B}$, les recouvrements de $U$ appartenant à
$\text{Cov}$ forment un système cofinal de recouvrements de $U$ ;
\item pour tout $\mathcal{U} \in \text{Cov}$, on ait
$\check{H}^p(\mathcal{U}, \mathcal{F}) = 0$ pour tout $p > 0$.
\end{enumerate}
Alors $H^p(U, \mathcal{F}) = 0$ pour tout $p > 0$ et tout
$U \in \mathcal{B}$.
\end{lemma}

\begin{proof}
Soient $\mathcal{F}$ et $\text{Cov}$ comme dans le lemme. Nous exprimerons
cette situation en disant que « la cohomologie de {\v C}ech supérieure de
$\mathcal{F}$ s'annule pour tout $\mathcal{U} \in \text{Cov}$ ».
Choisissons un plongement $\mathcal{F} \to \mathcal{I}$ dans un faisceau
abélien injectif. D'après le lemme \ref{lemma-injective-trivial-cech}, la
cohomologie de {\v C}ech supérieure de $\mathcal{I}$ s'annule pour tout
$\mathcal{U} \in \text{Cov}$. Posons $\mathcal{Q} = \mathcal{I}/\mathcal{F}$,
de sorte que nous ayons une suite exacte courte
$$
0 \to \mathcal{F} \to \mathcal{I} \to \mathcal{Q} \to 0.
$$
D'après le lemme \ref{lemma-ses-cech-h1} et notre hypothèse (2), cette suite
donne naissance à une suite exacte
$$
0 \to \mathcal{F}(U) \to \mathcal{I}(U) \to \mathcal{Q}(U) \to 0.
$$
pour tout $U \in \mathcal{B}$. Par conséquent, pour tout
$\mathcal{U} \in \text{Cov}$, nous obtenons une suite exacte courte de
complexes de {\v C}ech
$$
0 \to
\check{\mathcal{C}}^\bullet(\mathcal{U}, \mathcal{F}) \to
\check{\mathcal{C}}^\bullet(\mathcal{U}, \mathcal{I}) \to
\check{\mathcal{C}}^\bullet(\mathcal{U}, \mathcal{Q}) \to 0
$$
puisque, d'après l'hypothèse (1), chaque terme du complexe de {\v C}ech est
formé d'un produit de valeurs sur des éléments de $\mathcal{B}$.
En particulier, pour tout recouvrement $\mathcal{U} \in \text{Cov}$, nous
avons une suite exacte longue de groupes de cohomologie de {\v C}ech.
Il s'ensuit que $\mathcal{Q}$ est lui aussi un faisceau abélien dont la
cohomologie de {\v C}ech supérieure s'annule pour tout
$\mathcal{U} \in \text{Cov}$.

\medskip\noindent
Considérons ensuite la suite exacte longue de cohomologie
$$
\xymatrix{
0 \ar[r] &
H^0(U, \mathcal{F}) \ar[r] &
H^0(U, \mathcal{I}) \ar[r] &
H^0(U, \mathcal{Q}) \ar[lld] \\
&
H^1(U, \mathcal{F}) \ar[r] &
H^1(U, \mathcal{I}) \ar[r] &
H^1(U, \mathcal{Q}) \ar[lld] \\
&
\ldots & \ldots & \ldots \\
}
$$
pour tout $U \in \mathcal{B}$. Puisque $\mathcal{I}$ est injectif, nous avons
$H^n(U, \mathcal{I}) = 0$ pour $n > 0$ (voir Catégories dérivées,
lemme \ref{derived-lemma-higher-derived-functors}).
D'après ce qui précède, $H^0(U, \mathcal{I}) \to H^0(U, \mathcal{Q})$ est
surjectif, et donc $H^1(U, \mathcal{F}) = 0$.
Comme $\mathcal{F}$ était un faisceau abélien arbitraire dont la cohomologie de
{\v C}ech supérieure s'annule pour tout $\mathcal{U} \in \text{Cov}$, nous
concluons aussi que $H^1(U, \mathcal{Q}) = 0$, puisque $\mathcal{Q}$ est un
autre faisceau de ce type (voir ci-dessus). La suite exacte longue implique à
son tour que $H^2(U, \mathcal{F}) = 0$, et ainsi de suite.
\end{proof}














\section{Deuxième cohomologie et gerbes}
\label{section-gerbes}

\noindent
Soit $p : \mathcal{S} \to \mathcal{C}$ une gerbe sur un site dont tous les
groupes d'automorphismes sont commutatifs. Dans cette situation, les premier
et deuxième groupes de cohomologie du faisceau des automorphismes (Champs,
lemme \ref{stacks-lemma-gerbe-abelian-auts}) régissent l'existence d'objets.

\medskip\noindent
Le lemme suivant deviendra obsolète lorsqu'une étude plus complète de cette
relation sera ajoutée.

\begin{lemma}
\label{lemma-existence}
Soit $\mathcal{C}$ un site. Soit $p : \mathcal{S} \to \mathcal{C}$ une gerbe
sur un site dont les faisceaux d'automorphismes sont abéliens.
Soit $\mathcal{G}$ le faisceau de groupes abéliens construit dans Champs,
lemme \ref{stacks-lemma-gerbe-abelian-auts}.
Soit $U$ un objet de $\mathcal{C}$ tel que
\begin{enumerate}
\item il existe un système cofinal de recouvrements $\{U_i \to U\}$ de $U$
dans $\mathcal{C}$ tel que $H^1(U_i, \mathcal{G}) = 0$ et
$H^1(U_i \times_U U_j, \mathcal{G}) = 0$
pour tous $i, j$ ; et
\item $H^2(U, \mathcal{G}) = 0$.
\end{enumerate}
Alors il existe un objet de $\mathcal{S}$ au-dessus de $U$.
\end{lemma}

\begin{proof}
D'après Champs, définition \ref{stacks-definition-gerbe}, il existe un
recouvrement $\mathcal{U} = \{U_i \to U\}$ et des $x_i$ dans $\mathcal{S}$
au-dessus de $U_i$. Posons $U_{ij} = U_i \times_U U_j$. En raffinant le
recouvrement conformément à (1), nous pouvons supposer que
$H^1(U_i, \mathcal{G}) = 0$ et $H^1(U_{ij}, \mathcal{G}) = 0$.
Considérons le faisceau
$$
\mathcal{F}_{ij} =
\mathit{Isom}(x_i|_{U_{ij}}, x_j|_{U_{ij}})
$$
sur $\mathcal{C}/U_{ij}$. Puisque
$\mathcal{G}|_{U_{ij}} = \mathit{Aut}(x_i|_{U_{ij}})$, il existe une action
$$
\mathcal{G}|_{U_{ij}} \times \mathcal{F}_{ij} \to \mathcal{F}_{ij}
$$
par précomposition. Il est clair que $\mathcal{F}_{ij}$ est un
pseudo-$\mathcal{G}|_{U_{ij}}$-torseur, et même un torseur, car deux objets
quelconques d'une gerbe sont localement isomorphes.
Par notre choix du recouvrement et le lemme \ref{lemma-torsors-h1}, ces
torseurs sont triviaux (et admettent donc des sections globales d'après le
lemme \ref{lemma-trivial-torsor}). Autrement dit, nous pouvons choisir des
isomorphismes
$$
\varphi_{ij} : x_i|_{U_{ij}} \longrightarrow x_j|_{U_{ij}}
$$
Pour obtenir un objet $x$ au-dessus de $U$, nous allons modifier notre choix
des $\varphi_{ij}$ afin d'obtenir une donnée de descente (nécessairement
effective puisque $p : \mathcal{S} \to \mathcal{C}$ est un champ).
L'obstruction à ce que ces données forment une donnée de descente est le
possible défaut de la condition de cocycle. Posons
$U_{ijk} = U_i \times_U U_j \times_U U_k$. Nous pouvons considérer
$$
g_{ijk} = \varphi_{ik}^{-1}|_{U_{ijk}} \circ \varphi_{jk}|_{U_{ijk}} \circ
\varphi_{ij}|_{U_{ijk}}
$$
qui est un automorphisme de $x_i$ au-dessus de $U_{ijk}$. Nous considérons
donc $g_{ijk}$ comme une section de $\mathcal{G}$ sur $U_{ijk}$.
Un calcul, que nous omettons, montre que $(g_{i_0i_1i_2})$ est un $2$-cocycle
du complexe de {\v C}ech ${\check C}^\bullet(\mathcal{U}, \mathcal{G})$ de
$\mathcal{G}$ relativement au recouvrement $\mathcal{U}$.
D'après la suite spectrale du lemme \ref{lemma-cech-spectral-sequence} et
puisque $H^1(U_i, \mathcal{G}) = 0$ pour tout $i$, le morphisme
${\check H}^2(\mathcal{U}, \mathcal{G}) \to H^2(U, \mathcal{G})$ est injectif.
Par conséquent, $(g_{i_0i_1i_2})$ est un cobord, puisque nous supposons que
$H^2(U, \mathcal{G}) = 0$. Il existe donc des sections
$g_{ij} \in \mathcal{G}(U_{ij})$ telles que
$g_{ik}^{-1}|_{U_{ijk}} g_{jk}|_{U_{ijk}}  g_{ij}|_{U_{ijk}} = g_{ijk}$
pour tous $i, j, k$.
Après avoir remplacé $\varphi_{ij}$ par $\varphi_{ij}g_{ij}^{-1}$, les
$\varphi_{ij}$ définissent une donnée de descente sur les objets $x_i$
au-dessus de $U_i$, ce qui achève la preuve.
\end{proof}












\section{Cohomologie des modules}
\label{section-cohomology-modules}

\noindent
Tout ce qui a été dit pour la cohomologie des faisceaux abéliens vaut pour la
cohomologie des modules, puisque les deux coïncident.

\begin{lemma}
\label{lemma-injective-module-injective-presheaf}
Soit $(\mathcal{C}, \mathcal{O})$ un site annelé.
Tout faisceau injectif de modules est également injectif en tant qu'objet de
la catégorie $\textit{PMod}(\mathcal{O})$.
\end{lemma}

\begin{proof}
Appliquons le lemme \ref{homology-lemma-adjoint-preserve-injectives}
d'Homologie aux catégories $\mathcal{A} = \textit{Mod}(\mathcal{O})$ et
$\mathcal{B} = \textit{PMod}(\mathcal{O})$, au foncteur d'inclusion et à la
faisceautisation. (Voir Modules sur les sites, section
\ref{sites-modules-section-sheafification-presheaves-modules}, pour constater
que toutes les hypothèses du lemme sont satisfaites.)
\end{proof}

\begin{lemma}
\label{lemma-include-modules}
Soit $(\mathcal{C}, \mathcal{O})$ un site annelé.
Considérons le foncteur
$i : \textit{Mod}(\mathcal{C}) \to \textit{PMod}(\mathcal{C})$.
C'est un foncteur exact à gauche dont les foncteurs dérivés à droite sont
donnés par
$$
R^pi(\mathcal{F}) = \underline{H}^p(\mathcal{F}) :
U \longmapsto H^p(U, \mathcal{F})
$$
voir la discussion de la section \ref{section-locality}.
\end{lemma}

\begin{proof}
Il est clair que $i$ est exact à gauche.
Choisissons une résolution injective $\mathcal{F} \to \mathcal{I}^\bullet$
dans $\textit{Mod}(\mathcal{O})$.
Par définition, $R^pi$ est le {\it préfaisceau} de cohomologie de degré $p$
du complexe $\mathcal{I}^\bullet$. Autrement dit, les sections de
$R^pi(\mathcal{F})$ sur un objet $U$ de $\mathcal{C}$ sont données par
$$
\frac{\Ker(\mathcal{I}^n(U) \to \mathcal{I}^{n + 1}(U))}
{\Im(\mathcal{I}^{n - 1}(U) \to \mathcal{I}^n(U))}.
$$
ce qui est la définition de $H^p(U, \mathcal{F})$.
\end{proof}

\begin{lemma}
\label{lemma-injective-module-trivial-cech}
Soit $(\mathcal{C}, \mathcal{O})$ un site annelé.
Soit $\mathcal{U} = \{U_i \to U\}_{i \in I}$ un recouvrement de
$\mathcal{C}$. Soit $\mathcal{I}$ un $\mathcal{O}$-module injectif,
c'est-à-dire un objet injectif de $\textit{Mod}(\mathcal{O})$. Alors
$$
\check{H}^p(\mathcal{U}, \mathcal{I}) =
\left\{
\begin{matrix}
\mathcal{I}(U) & \text{if} & p = 0 \\
0 & \text{if} & p > 0
\end{matrix}
\right.
$$
\end{lemma}

\begin{proof}
Le lemme \ref{lemma-cech-map-into} donne la première égalité de la chaîne
d'égalités suivante :
\begin{align*}
\check{\mathcal{C}}^\bullet(\mathcal{U}, \mathcal{I})
& =
\Mor_{\textit{PAb}(\mathcal{C})}(
\mathbf{Z}_{\mathcal{U}, \bullet}, \mathcal{I}) \\
& =
\Mor_{\textit{PMod}(\mathbf{Z})}(
\mathbf{Z}_{\mathcal{U}, \bullet}, \mathcal{I}) \\
& =
\Mor_{\textit{PMod}(\mathcal{O})}(
\mathbf{Z}_{\mathcal{U}, \bullet} \otimes_{p, \mathbf{Z}} \mathcal{O},
\mathcal{I})
\end{align*}
La troisième égalité découle de Modules sur les sites,
lemme \ref{sites-modules-lemma-adjointness-tensor-restrict-presheaves}.
D'après le lemme \ref{lemma-injective-module-injective-presheaf},
$\mathcal{I}$ est un objet injectif de $\textit{PMod}(\mathcal{O})$.
Par conséquent, $\Hom_{\textit{PMod}(\mathcal{O})}(-, \mathcal{I})$ est un
foncteur exact. Le lemme \ref{lemma-complex-tensored-still-exact} donne alors
l'annulation des groupes de cohomologie de {\v C}ech supérieurs.
Pour le degré zéro, voir le lemme \ref{lemma-cech-h0}.
\end{proof}

\begin{lemma}
\label{lemma-cohomology-modules-abelian-agree}
Soit $\mathcal{C}$ un site.
Soit $\mathcal{O}$ un faisceau d'anneaux sur $\mathcal{C}$.
Soit $\mathcal{F}$ un $\mathcal{O}$-module, et notons $\mathcal{F}_{ab}$ le
faisceau de groupes abéliens sous-jacent. Alors
$$
H^i(\mathcal{C}, \mathcal{F}_{ab})
=
H^i(\mathcal{C}, \mathcal{F})
$$
et, pour tout objet $U$ de $\mathcal{C}$, nous avons également
$$
H^i(U, \mathcal{F}_{ab})
=
H^i(U, \mathcal{F}).
$$
Ici, le membre de gauche est la cohomologie calculée dans
$\textit{Ab}(\mathcal{C})$, et le membre de droite celle calculée dans
$\textit{Mod}(\mathcal{O})$.
\end{lemma}

\begin{proof}
D'après Catégories dérivées,
lemme \ref{derived-lemma-higher-derived-functors}, le $\delta$-foncteur
$(\mathcal{F} \mapsto H^p(U, \mathcal{F}))_{p \geq 0}$ est universel. Le
foncteur
$\textit{Mod}(\mathcal{O}) \to \textit{Ab}(\mathcal{C})$,
$\mathcal{F} \mapsto \mathcal{F}_{ab}$ est exact. Ainsi,
$(\mathcal{F} \mapsto H^p(U, \mathcal{F}_{ab}))_{p \geq 0}$
est également un $\delta$-foncteur. Supposons que nous montrions que
$(\mathcal{F} \mapsto H^p(U, \mathcal{F}_{ab}))_{p \geq 0}$
est lui aussi universel. La deuxième assertion du lemme découlera alors de
l'unicité des $\delta$-foncteurs universels ; voir Homologie,
lemme \ref{homology-lemma-uniqueness-universal-delta-functor}.
Puisque $\textit{Mod}(\mathcal{O})$ possède assez d'injectifs, il suffit de
montrer que $H^i(U, \mathcal{I}_{ab}) = 0$ pour tout objet injectif
$\mathcal{I}$ de $\textit{Mod}(\mathcal{O})$ ; voir Homologie,
lemme \ref{homology-lemma-efface-implies-universal}.

\medskip\noindent
Soit $\mathcal{I}$ un objet injectif de $\textit{Mod}(\mathcal{O})$.
Appliquons le lemme \ref{lemma-cech-vanish-collection} avec
$\mathcal{F} = \mathcal{I}$, $\mathcal{B} = \mathcal{C}$ et
$\text{Cov} = \text{Cov}_\mathcal{C}$.
L'hypothèse (3) de ce lemme est satisfaite d'après le
lemme \ref{lemma-injective-module-trivial-cech}.
Ainsi, $H^i(U, \mathcal{I}_{ab}) = 0$ pour tout objet $U$ de $\mathcal{C}$.

\medskip\noindent
Si $\mathcal{C}$ possède un objet final, cela implique également la première
égalité. Sinon, d'après Sites, lemme \ref{sites-lemma-topos-good-site}, le
topos annelé $(\Sh(\mathcal{C}), \mathcal{O})$ est équivalent à un topos annelé
dont le site sous-jacent possède un objet final. Le lemme en découle.
\end{proof}

\begin{lemma}
\label{lemma-cohomology-products}
Soit $\mathcal{C}$ un site. Soit $I$ un ensemble. Pour $i \in I$, soit
$\mathcal{F}_i$ un faisceau abélien sur $\mathcal{C}$. Soit
$U \in \Ob(\mathcal{C})$. Le morphisme canonique
$$
H^p(U, \prod\nolimits_{i \in I} \mathcal{F}_i)
\longrightarrow
\prod\nolimits_{i \in I} H^p(U, \mathcal{F}_i)
$$
est un isomorphisme pour $p = 0$ et est injectif pour $p = 1$.
\end{lemma}

\begin{proof}
L'assertion pour $p = 0$ est vraie parce que le produit des faisceaux est égal
au produit des préfaisceaux sous-jacents ; voir Sites,
lemme \ref{sites-lemma-limit-sheaf}.
Traitons $p = 1$. Posons $\mathcal{F} = \prod \mathcal{F}_i$.
Soit $\xi \in H^1(U, \mathcal{F})$ dont l'image dans
$\prod H^1(U, \mathcal{F}_i)$ est nulle. Par localité de la cohomologie, voir
le lemme \ref{lemma-kill-cohomology-class-on-covering}, il existe un
recouvrement $\mathcal{U} = \{U_j \to U\}$ tel que $\xi|_{U_j} = 0$ pour tout
$j$. D'après le lemme \ref{lemma-cech-h1}, cela signifie que $\xi$ provient
d'un élément
$\check \xi \in \check H^1(\mathcal{U}, \mathcal{F})$.
Puisque les morphismes
$\check H^1(\mathcal{U}, \mathcal{F}_i) \to H^1(U, \mathcal{F}_i)$
sont injectifs pour tout $i$ (d'après le lemme \ref{lemma-cech-h1}) et puisque
l'image de $\xi$ dans $\prod H^1(U, \mathcal{F}_i)$ est nulle, nous voyons
que l'image
$\check \xi_i = 0$ dans $\check H^1(\mathcal{U}, \mathcal{F}_i)$.
Toutefois, comme $\mathcal{F} = \prod \mathcal{F}_i$, le complexe
$\check{\mathcal{C}}^\bullet(\mathcal{U}, \mathcal{F})$ est le produit des
complexes $\check{\mathcal{C}}^\bullet(\mathcal{U}, \mathcal{F}_i)$.
D'après Homologie,
lemme \ref{homology-lemma-product-abelian-groups-exact}, nous concluons que
$\check \xi = 0$, comme voulu.
\end{proof}

\begin{lemma}
\label{lemma-restriction-along-monomorphism-surjective}
Soit $(\mathcal{C}, \mathcal{O})$ un site annelé. Soit $a : U' \to U$ un
monomorphisme dans $\mathcal{C}$. Alors, pour tout $\mathcal{O}$-module
injectif $\mathcal{I}$, le morphisme de restriction
$\mathcal{I}(U) \to \mathcal{I}(U')$ est surjectif.
\end{lemma}

\begin{proof}
Soient $j : \mathcal{C}/U \to \mathcal{C}$ et
$j' : \mathcal{C}/U' \to \mathcal{C}$ les morphismes de localisation
(Modules sur les sites, section \ref{sites-modules-section-localize}).
Comme $j_!$ est adjoint à gauche de la restriction, nous avons, pour tout
faisceau $\mathcal{F}$ de $\mathcal{O}$-modules,
$$
\Hom_\mathcal{O}(j_!\mathcal{O}_U, \mathcal{F})
=
\Hom_{\mathcal{O}_U}(\mathcal{O}_U, \mathcal{F}|_U)
=
\mathcal{F}(U)
$$
De même, le faisceau $j'_!\mathcal{O}_{U'}$ représente le foncteur
$\mathcal{F} \mapsto \mathcal{F}(U')$.
Nous décrivons ci-dessous un morphisme canonique de $\mathcal{O}$-modules
$$
j'_!\mathcal{O}_{U'} \longrightarrow j_!\mathcal{O}_U
$$
qui correspond, par le lemme de Yoneda, au morphisme de restriction
$\mathcal{F}(U) \to \mathcal{F}(U')$ (Catégories,
lemme \ref{categories-lemma-yoneda}).
Il suffit de prouver que le morphisme de modules affiché est injectif ; voir
Homologie, lemme \ref{homology-lemma-characterize-injectives}.

\medskip\noindent
Pour construire notre morphisme, il suffit de construire un morphisme entre
les préfaisceaux qui associent à un objet $V$ de $\mathcal{C}$ le
$\mathcal{O}(V)$-module
$$
\bigoplus\nolimits_{\varphi' \in \Mor_\mathcal{C}(V, U')} \mathcal{O}(V)
\quad\text{and}\quad
\bigoplus\nolimits_{\varphi \in \Mor_\mathcal{C}(V, U)} \mathcal{O}(V)
$$
voir Modules sur les sites,
lemme \ref{sites-modules-lemma-extension-by-zero}.
Nous prenons le morphisme qui envoie le facteur correspondant à $\varphi'$
sur celui correspondant à $\varphi = a \circ \varphi'$ par l'identité de
$\mathcal{O}(V)$. Comme $a$ est un monomorphisme, ce morphisme est injectif.
La faisceautisation étant exacte, le résultat en découle.
\end{proof}






\section{Faisceaux totalement acycliques}
\label{section-limp}

\noindent
Soit $(\mathcal{C}, \mathcal{O})$ un site annelé.
Soit $K$ un préfaisceau d'ensembles sur $\mathcal{C}$ (nous employons
intentionnellement ici une capitale romaine pour le distinguer des faisceaux
abéliens). Étant donné un faisceau de $\mathcal{O}$-modules $\mathcal{F}$,
posons
$$
\mathcal{F}(K) =
\Mor_{\textit{PSh}(\mathcal{C})}(K, \mathcal{F}) =
\Mor_{\Sh(\mathcal{C})}(K^\#, \mathcal{F})
$$
Le foncteur $\mathcal{F} \mapsto \mathcal{F}(K)$ est un foncteur exact à
gauche $\textit{Mod}(\mathcal{O}) \to \textit{Ab}$ ; il possède donc des
foncteurs dérivés à droite. Nous les noterons $H^p(K, \mathcal{F})$, de sorte
que $H^0(K, \mathcal{F}) = \mathcal{F}(K)$.

\medskip\noindent
Voici quelques observations :
\begin{enumerate}
\item Puisque $\mathcal{F}(K) = \mathcal{F}(K^\#)$, nous avons
$H^p(K, \mathcal{F}) = H^p(K^\#, \mathcal{F})$.
Autoriser $K$ à être un préfaisceau dans la définition ci-dessus n'est qu'une
commodité de notation.
\item Supposons que $K = h_U$ ou $K = h_U^\#$ pour un objet $U$ de
$\mathcal{C}$. Alors $H^p(K, \mathcal{F}) = H^p(U, \mathcal{F})$, car
$\Mor_{\Sh(\mathcal{C})}(h_U^\#, \mathcal{F}) = \mathcal{F}(U)$ ; voir
Sites, section \ref{sites-section-representable-sheaves}.
\item Si $\mathcal{O} = \mathbf{Z}$ (le faisceau constant), alors les groupes
de cohomologie sont des foncteurs
$H^p(K, - ) : \textit{Ab}(\mathcal{C}) \to \textit{Ab}$
car $\textit{Mod}(\mathcal{O}) = \textit{Ab}(\mathcal{C})$ dans ce cas.
\end{enumerate}
Nous pouvons reformuler dans ce langage certains résultats déjà démontrés.

\begin{lemma}
\label{lemma-compute-cohomology-on-sheaf-sets}
Soit $(\mathcal{C}, \mathcal{O})$ un site annelé.
Soit $K$ un préfaisceau d'ensembles sur $\mathcal{C}$.
Soit $\mathcal{F}$ un $\mathcal{O}$-module, et notons $\mathcal{F}_{ab}$ le
faisceau de groupes abéliens sous-jacent. Alors
$H^p(K, \mathcal{F}) = H^p(K, \mathcal{F}_{ab})$.
\end{lemma}

\begin{proof}
Nous pouvons remplacer $K$ par son faisceautisé et supposer que $K$ est un
faisceau. Remarquons que $H^p(K, \mathcal{F})$ et
$H^p(K, \mathcal{F}_{ab})$ ne dépendent que du topos, et non du site
sous-jacent. D'après Sites, lemme \ref{sites-lemma-topos-good-site}, nous
pouvons donc remplacer $\mathcal{C}$ par un site « plus grand » tel que
$K = h_U$ pour un objet $U$ de $\mathcal{C}$. Dans ce cas, le résultat découle
du lemme \ref{lemma-cohomology-modules-abelian-agree}.
\end{proof}

\begin{lemma}
\label{lemma-cech-to-cohomology-sheaf-sets}
Soit $\mathcal{C}$ un site. Soit $K' \to K$ un morphisme de préfaisceaux
d'ensembles sur $\mathcal{C}$ dont le faisceautisé est surjectif. Posons
$K'_p = K' \times_K \ldots \times_K K'$ ($p + 1$ facteurs).
Pour tout faisceau abélien $\mathcal{F}$, il existe une suite spectrale
vérifiant $E_1^{p, q} = H^q(K'_p, \mathcal{F})$ et convergeant vers
$H^{p + q}(K, \mathcal{F})$.
\end{lemma}

\begin{proof}
Puisque la faisceautisation est exacte, nous voyons que
$(K_p')^\#$ est égal à
$(K')^\# \times_{K^\#} \ldots \times_{K^\#} (K')^\#$
($p + 1$ facteurs). Nous pouvons donc remplacer $K$ et $K'$ par leurs
faisceautisés et supposer que $K' \to K$ est un morphisme surjectif de
faisceaux. Après avoir remplacé $\mathcal{C}$ par un site « plus grand » comme
dans Sites, lemme \ref{sites-lemma-topos-good-site}, nous pouvons supposer que
$K, K'$ sont des objets de $\mathcal{C}$ et que
$\mathcal{U} = \{K' \to K\}$ est un recouvrement. Nous disposons alors de la
suite spectrale de la cohomologie de {\v C}ech vers la cohomologie du
lemme \ref{lemma-cech-spectral-sequence}, dont la page $E_1$ est celle indiquée
dans l'énoncé.
\end{proof}

\begin{lemma}
\label{lemma-cohomology-on-sheaf-sets}
Soit $\mathcal{C}$ un site. Soit $K$ un faisceau d'ensembles sur
$\mathcal{C}$. Considérons le morphisme de topos
$j : \Sh(\mathcal{C}/K) \to \Sh(\mathcal{C})$ ; voir
Sites, lemme \ref{sites-lemma-localize-topos-site}.
Alors $j^{-1}$ préserve les injectifs et
$H^p(K, \mathcal{F}) = H^p(\mathcal{C}/K, j^{-1}\mathcal{F})$
pour tout faisceau abélien $\mathcal{F}$ sur $\mathcal{C}$.
\end{lemma}

\begin{proof}
D'après Sites, lemmes \ref{sites-lemma-localize-topos} et
\ref{sites-lemma-localize-topos-site}, le morphisme de topos $j$ est équivalent
à une localisation. Le résultat découle donc du
lemme \ref{lemma-cohomology-of-open}.
\end{proof}

\noindent
Compte tenu du lemme \ref{lemma-compute-cohomology-on-sheaf-sets}, nous voyons
que la définition suivante est également la « bonne » pour les faisceaux de
modules sur les sites annelés.

\begin{definition}
\label{definition-limp}
Soit $\mathcal{C}$ un site.
On dit qu'un faisceau abélien $\mathcal{F}$ est
{\it totalement acyclique}\footnote{Bien que cette terminologie soit employée
dans \cite[Vbis, Proposition 1.3.10]{SGA4}, elle n'est probablement pas
standard. Dans \cite[V, Definition 4.1]{SGA4}, cette propriété est appelée
« flasque », mais nous ne pouvons pas employer ce terme car il entrerait en
conflit avec notre définition des faisceaux flasques sur les espaces
topologiques. Veuillez écrire à
\href{mailto:stacks.project@gmail.com}{stacks.project@gmail.com}
si vous avez une meilleure suggestion.} si, pour tout faisceau d'ensembles
$K$, nous avons $H^p(K, \mathcal{F}) = 0$ pour tout $p \geq 1$.
\end{definition}

\noindent
Il est clair que la propriété d'être totalement acyclique est intrinsèque,
c'est-à-dire préservée par les équivalences de topos.
La cohomologie supérieure d'un faisceau totalement acyclique s'annule sur tous
les objets du site, mais, en général, la condition d'acyclicité totale est
strictement plus forte. Voici une caractérisation parfois utile des faisceaux
totalement acycliques.

\begin{lemma}
\label{lemma-characterize-limp}
Soit $\mathcal{C}$ un site. Soit $\mathcal{F}$ un faisceau abélien. Si
\begin{enumerate}
\item $H^p(U, \mathcal{F}) = 0$ pour $p > 0$ et
$U \in \Ob(\mathcal{C})$ ; et
\item pour toute surjection $K' \to K$ de faisceaux d'ensembles, le complexe
de {\v C}ech augmenté
$$
0 \to H^0(K, \mathcal{F}) \to H^0(K', \mathcal{F}) \to
H^0(K' \times_K K', \mathcal{F}) \to \ldots
$$
soit exact,
\end{enumerate}
alors $\mathcal{F}$ est totalement acyclique (et la réciproque est vraie).
\end{lemma}

\begin{proof}
D'après l'hypothèse (1), nous avons
$H^p(h_U^\#, \mathcal{F}) = 0$ pour tout $p > 0$ et tout objet $U$ de
$\mathcal{C}$. Remarquons que si $K = \coprod K_i$ est un coproduit de
faisceaux d'ensembles sur $\mathcal{C}$, alors
$H^p(K, \mathcal{F}) = \prod H^p(K_i, \mathcal{F})$.
Pour tout faisceau d'ensembles $K$, il existe une surjection
$$
K' = \coprod h_{U_i}^\# \longrightarrow K
$$
voir Sites, lemme \ref{sites-lemma-sheaf-coequalizer-representable}.
Nous en concluons que : (*) pour tout faisceau d'ensembles $K$, il existe une
surjection $K' \to K$ de faisceaux d'ensembles telle que
$H^p(K', \mathcal{F}) = 0$ pour $p > 0$. Nous affirmons que (*) et la
condition (2) impliquent que $\mathcal{F}$ est totalement acyclique.
Remarquons que les conditions (*) et (2) ne dépendent que de $\mathcal{F}$ en
tant qu'objet du topos $\Sh(\mathcal{C})$, et non du site sous-jacent.
(Nous n'utiliserons plus la propriété (1) dans la suite de la preuve.)

\medskip\noindent
Nous allons montrer par récurrence sur $n \geq 0$ que (*) et (2) impliquent
l'hypothèse de récurrence suivante $IH_n$ :
$H^p(K, \mathcal{F}) = 0$ pour tout $0 < p \leq n$ et tout faisceau
d'ensembles $K$. Remarquons que $IH_0$ est vraie. Supposons $IH_n$.
Choisissons un faisceau d'ensembles $K$, puis une surjection $K' \to K$ telle
que $H^p(K', \mathcal{F}) = 0$ pour tout $p > 0$. Nous disposons d'une suite
spectrale vérifiant
$$
E_1^{p, q} = H^q(K'_p, \mathcal{F})
$$
et convergeant vers $H^{p + q}(K, \mathcal{F})$ ; voir le
lemme \ref{lemma-cech-to-cohomology-sheaf-sets}.
D'après $IH_n$, $E_1^{p, q} = 0$ pour $0 < q \leq n$, et l'hypothèse (2)
donne $E_2^{p, 0} = 0$ pour $p > 0$. Enfin, $E_1^{0, q} = 0$ pour $q > 0$,
car $H^q(K', \mathcal{F}) = 0$ par choix de $K'$. Nous concluons donc que
$H^{n + 1}(K, \mathcal{F}) = 0$, puisque tous les termes $E_2^{p, q}$ tels
que $p + q = n + 1$ sont nuls.
\end{proof}







\section{La suite spectrale de Leray}
\label{section-leray}

\noindent
Le lemme suivant est la clé de la preuve de l'existence de la suite spectrale
de Leray.

\begin{lemma}
\label{lemma-direct-image-injective-sheaf}
Soit $f : (\Sh(\mathcal{C}), \mathcal{O}_\mathcal{C}) \to
(\Sh(\mathcal{D}), \mathcal{O}_\mathcal{D})$ un morphisme de topos annelés.
Alors, pour tout objet injectif $\mathcal{I}$ de
$\textit{Mod}(\mathcal{O}_\mathcal{C})$, l'image directe
$f_*\mathcal{I}$ est totalement acyclique.
\end{lemma}

\begin{proof}
Soit $K$ un faisceau d'ensembles sur $\mathcal{D}$.
D'après Modules sur les sites, lemme
\ref{sites-modules-lemma-morphism-ringed-topoi-comes-from-morphism-ringed-sites}
nous pouvons remplacer $\mathcal{C}$ et $\mathcal{D}$ par des sites « plus
grands » de sorte que $f$ provienne d'un morphisme de sites annelés induit par
un foncteur continu $u : \mathcal{D} \to \mathcal{C}$ et que $K = h_V$ pour
un objet $V$ de $\mathcal{D}$.

\medskip\noindent
Il nous suffit donc de montrer que $H^q(V, f_*\mathcal{I})$ est nul pour
$q > 0$ et tout objet $V$ de $\mathcal{D}$ lorsque $f$ est donné par un
morphisme de sites annelés. Soit $\mathcal{V} = \{V_j \to V\}$ un
recouvrement quelconque de $\mathcal{D}$. Comme $u$ est continu,
$\mathcal{U} = \{u(V_j) \to u(V)\}$ est un recouvrement de $\mathcal{C}$.
Nous avons alors une égalité de complexes de {\v C}ech
$$
\check{\mathcal{C}}^\bullet(\mathcal{V}, f_*\mathcal{I})
=
\check{\mathcal{C}}^\bullet(\mathcal{U}, \mathcal{I})
$$
par définition de $f_*$. D'après le
lemme \ref{lemma-injective-module-trivial-cech}, la cohomologie de ce complexe
est nulle en degrés positifs. Le résultat découle alors du
lemme \ref{lemma-cech-vanish-collection}.
\end{proof}

\noindent
Pour les morphismes plats, le foncteur $f_*$ préserve les modules injectifs.
En particulier, le foncteur
$f_* : \textit{Ab}(\mathcal{C}) \to \textit{Ab}(\mathcal{D})$ transforme
toujours les faisceaux abéliens injectifs en faisceaux abéliens injectifs.

\begin{lemma}
\label{lemma-pushforward-injective-flat}
Soit $f : (\Sh(\mathcal{C}), \mathcal{O}_\mathcal{C}) \to
(\Sh(\mathcal{D}), \mathcal{O}_\mathcal{D})$ un morphisme de topos annelés.
Si $f$ est plat, alors $f_*\mathcal{I}$ est un
$\mathcal{O}_\mathcal{D}$-module injectif pour tout
$\mathcal{O}_\mathcal{C}$-module injectif $\mathcal{I}$.
\end{lemma}

\begin{proof}
Dans ce cas, le foncteur $f^*$ est exact ; voir Modules sur les sites,
lemme \ref{sites-modules-lemma-flat-pullback-exact}.
Le résultat découle donc d'Homologie,
lemme \ref{homology-lemma-adjoint-preserve-injectives}.
\end{proof}

\begin{lemma}
\label{lemma-limp-acyclic}
Soit $(\Sh(\mathcal{C}), \mathcal{O}_\mathcal{C})$ un topos annelé.
Un faisceau totalement acyclique est acyclique à droite pour les foncteurs
suivants :
\begin{enumerate}
\item le foncteur $H^0(U, -)$ pour tout objet $U$ de $\mathcal{C}$ ;
\item le foncteur $\mathcal{F} \mapsto \mathcal{F}(K)$ pour tout préfaisceau
d'ensembles $K$ ;
\item le foncteur de sections globales $\Gamma(\mathcal{C}, -)$ ;
\item le foncteur $f_*$ pour tout morphisme
$f : (\Sh(\mathcal{C}), \mathcal{O}_\mathcal{C}) \to
(\Sh(\mathcal{D}), \mathcal{O}_\mathcal{D})$ de topos annelés.
\end{enumerate}
\end{lemma}

\begin{proof}
La partie (2) est la définition d'un faisceau totalement acyclique.
La partie (1) découle de (2), comme indiqué dans la discussion qui suit la
définition des faisceaux totalement acycliques.
La partie (3) est le cas particulier de (2) où $K = e$ est l'objet final de
$\Sh(\mathcal{C})$.

\medskip\noindent
Pour prouver (4), nous pouvons supposer, d'après Modules sur les sites, lemme
\ref{sites-modules-lemma-morphism-ringed-topoi-comes-from-morphism-ringed-sites}
que $f$ est donné par un morphisme de sites. Dans ce cas, la description des
images directes supérieures du lemme \ref{lemma-higher-direct-images} montre
que les $R^if_*$, $i > 0$, d'un faisceau totalement acyclique sont nuls.
\end{proof}

\begin{remark}
\label{remark-before-Leray}
Les résultats précédents montrent que Catégories dérivées,
lemme \ref{derived-lemma-compose-derived-functors}, s'applique dans plusieurs
situations. Par exemple, pour tout morphisme
$f : (\Sh(\mathcal{C}), \mathcal{O}_\mathcal{C}) \to
(\Sh(\mathcal{D}), \mathcal{O}_\mathcal{D})$ de topos annelés, nous avons
$$
R\Gamma(\mathcal{D}, Rf_*\mathcal{F}) = R\Gamma(\mathcal{C}, \mathcal{F})
$$
pour tout faisceau de $\mathcal{O}_\mathcal{C}$-modules $\mathcal{F}$.
En effet, pour un $\mathcal{O}_\mathcal{C}$-module injectif $\mathcal{I}$, le
$\mathcal{O}_\mathcal{D}$-module $f_*\mathcal{I}$ est totalement acyclique
d'après le lemme \ref{lemma-direct-image-injective-sheaf}, et un faisceau
totalement acyclique est acyclique pour $\Gamma(\mathcal{D}, -)$ d'après le
lemme \ref{lemma-limp-acyclic}.
\end{remark}

\begin{lemma}[Suite spectrale de Leray]
\label{lemma-Leray}
Soit $f : (\Sh(\mathcal{C}), \mathcal{O}_\mathcal{C}) \to
(\Sh(\mathcal{D}), \mathcal{O}_\mathcal{D})$ un morphisme de topos annelés.
Soit $\mathcal{F}^\bullet$ un complexe borné inférieurement de
$\mathcal{O}_\mathcal{C}$-modules. Il existe une suite spectrale
$$
E_2^{p, q} = H^p(\mathcal{D}, R^qf_*(\mathcal{F}^\bullet))
$$
qui converge vers $H^{p + q}(\mathcal{C}, \mathcal{F}^\bullet)$.
\end{lemma}

\begin{proof}
Il s'agit simplement de la suite spectrale de Grothendieck de Catégories
dérivées, lemme \ref{derived-lemma-grothendieck-spectral-sequence}, provenant
de la composition de foncteurs
$\Gamma(\mathcal{C}, -) = \Gamma(\mathcal{D}, -) \circ f_*$.
Pour vérifier que les hypothèses de Catégories dérivées,
lemme \ref{derived-lemma-grothendieck-spectral-sequence}, sont satisfaites,
voir les lemmes \ref{lemma-direct-image-injective-sheaf} et
\ref{lemma-limp-acyclic}.
\end{proof}

\begin{lemma}
\label{lemma-apply-Leray}
Soit $f : (\Sh(\mathcal{C}), \mathcal{O}_\mathcal{C}) \to
(\Sh(\mathcal{D}), \mathcal{O}_\mathcal{D})$ un morphisme de topos annelés.
Soit $\mathcal{F}$ un $\mathcal{O}_\mathcal{C}$-module.
\begin{enumerate}
\item Si $R^qf_*\mathcal{F} = 0$ pour $q > 0$, alors
$H^p(\mathcal{C}, \mathcal{F}) = H^p(\mathcal{D}, f_*\mathcal{F})$ pour tout
$p$.
\item Si $H^p(\mathcal{D}, R^qf_*\mathcal{F}) = 0$ pour tout $q$ et tout
$p > 0$, alors
$H^q(\mathcal{C}, \mathcal{F}) = H^0(\mathcal{D}, R^qf_*\mathcal{F})$ pour
tout $q$.
\end{enumerate}
\end{lemma}

\begin{proof}
Ce sont deux conditions simples qui forcent la suite spectrale de Leray à
dégénérer. On peut aussi démontrer ces faits directement, sans utiliser la
suite spectrale ; c'est un bon exercice de cohomologie des faisceaux.
\end{proof}

\begin{lemma}[Suite spectrale de Leray relative]
\label{lemma-relative-Leray}
Soient
$f : (\Sh(\mathcal{C}), \mathcal{O}_\mathcal{C}) \to
(\Sh(\mathcal{D}), \mathcal{O}_\mathcal{D})$
et
$g : (\Sh(\mathcal{D}), \mathcal{O}_\mathcal{D}) \to
(\Sh(\mathcal{E}), \mathcal{O}_\mathcal{E})$
soient des morphismes de topos annelés.
Soit $\mathcal{F}$ un $\mathcal{O}_\mathcal{C}$-module.
Il existe une suite spectrale telle que
$$
E_2^{p, q} = R^pg_*(R^qf_*\mathcal{F})
$$
qui converge vers $R^{p + q}(g \circ f)_*\mathcal{F}$.
Cette suite spectrale est fonctorielle en $\mathcal{F}$, et il en existe une
version pour les complexes bornés inférieurement de
$\mathcal{O}_\mathcal{C}$-modules.
\end{lemma}

\begin{proof}
C'est une suite spectrale de Grothendieck pour la composition de foncteurs ;
voir Catégories dérivées,
lemme \ref{derived-lemma-grothendieck-spectral-sequence}, ainsi que les lemmes
\ref{lemma-direct-image-injective-sheaf} et \ref{lemma-limp-acyclic}.
\end{proof}







\section{Le morphisme de changement de base}
\label{section-base-change-map}

\noindent
Dans cette section, nous construisons le morphisme de changement de base dans
certains cas ; le cas général est traité dans la remarque
\ref{remark-base-change}. Afin d'éviter ici le recours à l'image inverse
dérivée, nous nous restreignons au cas d'un changement de base par un morphisme
plat de sites annelés. Avant d'énoncer le résultat, examinons l'image inverse
plate sur la catégorie dérivée. Supposons que
$g : (\Sh(\mathcal{C}), \mathcal{O}_\mathcal{C})
\to (\Sh(\mathcal{D}), \mathcal{O}_\mathcal{D})$
soit un morphisme plat de topos annelés. D'après Modules sur les sites,
lemme \ref{sites-modules-lemma-flat-pullback-exact}, le foncteur
$g^* : \textit{Mod}(\mathcal{O}_\mathcal{D}) \to
\textit{Mod}(\mathcal{O}_\mathcal{C})$ est exact.
Il admet donc un foncteur dérivé
$$
g^* : D(\mathcal{O}_\mathcal{D}) \to D(\mathcal{O}_\mathcal{C})
$$
qui se calcule simplement en prenant l'image inverse d'un représentant d'un
objet donné de $D(\mathcal{O}_\mathcal{D})$ ; voir Catégories dérivées,
lemme \ref{derived-lemma-right-derived-exact-functor}.
Il préserve les sous-catégories bornées (supérieurement ou inférieurement).
Ainsi, comme annoncé, nous notons ce foncteur $g^*$ plutôt que $Lg^*$.

\begin{lemma}
\label{lemma-base-change-map-flat-case}
Soit
$$
\xymatrix{
(\Sh(\mathcal{C}'), \mathcal{O}_{\mathcal{C}'})
\ar[r]_{g'} \ar[d]_{f'} &
(\Sh(\mathcal{C}), \mathcal{O}_\mathcal{C}) \ar[d]^f \\
(\Sh(\mathcal{D}'), \mathcal{O}_{\mathcal{D}'})
\ar[r]^g &
(\Sh(\mathcal{D}), \mathcal{O}_\mathcal{D})
}
$$
un diagramme commutatif de topos annelés.
Soit $\mathcal{F}^\bullet$ un complexe borné inférieurement de
$\mathcal{O}_\mathcal{C}$-modules.
Supposons que $g$ et $g'$ soient plats.
Il existe alors un morphisme canonique de changement de base
$$
g^*Rf_*\mathcal{F}^\bullet
\longrightarrow
R(f')_*(g')^*\mathcal{F}^\bullet
$$
dans $D^{+}(\mathcal{O}_{\mathcal{D}'})$.
\end{lemma}

\begin{proof}
Choisissons des résolutions injectives $\mathcal{F}^\bullet \to \mathcal{I}^\bullet$
et $(g')^*\mathcal{F}^\bullet \to \mathcal{J}^\bullet$.
D'après le lemme \ref{lemma-pushforward-injective-flat},
$(g')_*\mathcal{J}^\bullet$ est un complexe d'injectifs qui représente
$R(g')_*(g')^*\mathcal{F}^\bullet$. Par Catégories dérivées, lemmes
\ref{derived-lemma-morphisms-lift}
et \ref{derived-lemma-morphisms-equal-up-to-homotopy},
le morphisme $\beta$ du diagramme
$$
\xymatrix{
(g')_*(g')^*\mathcal{F}^\bullet \ar[r] &
(g')_*\mathcal{J}^\bullet \\
\mathcal{F}^\bullet \ar[u]^{adjunction} \ar[r] &
\mathcal{I}^\bullet \ar[u]_\beta
}
$$
existe et est unique à homotopie près.
En prenant l'image directe vers $\mathcal{D}$, nous obtenons
$$
f_*\beta :
f_*\mathcal{I}^\bullet
\longrightarrow
f_*(g')_*\mathcal{J}^\bullet
=
g_*(f')_*\mathcal{J}^\bullet
$$
Par adjonction entre $g^*$ et $g_*$, nous obtenons un morphisme de complexes
$g^*f_*\mathcal{I}^\bullet \to (f')_*\mathcal{J}^\bullet$.
Ce morphisme est unique à homotopie près, puisque le seul choix effectué au
cours de la construction est celui du morphisme $\beta$ et que tout a été fait
au niveau des complexes.
\end{proof}








\section{Cohomologie et colimites}
\label{section-limits}

\noindent
Soit $(\mathcal{C}, \mathcal{O})$ un site annelé.
Soit $\mathcal{I} \to \textit{Mod}(\mathcal{O})$, $i \mapsto \mathcal{F}_i$,
un diagramme indexé par la catégorie $\mathcal{I}$ ; voir
Catégories, section \ref{categories-section-limits}.
Pour tout $i$, il existe un morphisme canonique
$\mathcal{F}_i \to \colim_i \mathcal{F}_i$ qui induit
un morphisme en cohomologie. Nous obtenons donc un morphisme canonique
$$
\colim_i H^p(U, \mathcal{F}_i)
\longrightarrow
H^p(U, \colim_i \mathcal{F}_i)
$$
pour tout $p \geq 0$ et tout objet $U$ de $\mathcal{C}$.
En général, ces morphismes ne sont pas des isomorphismes, même pour $p = 0$.

\medskip\noindent
Le lemme suivant est l'analogue, pour la cohomologie, du
lemme \ref{sites-lemma-directed-colimits-sections} de Sites.

\begin{lemma}
\label{lemma-colim-works-over-collection}
Soit $\mathcal{C}$ un site. Soit $\text{Cov}_\mathcal{C}$ l'ensemble
des recouvrements de $\mathcal{C}$ (voir
Sites, définition \ref{sites-definition-site}). Soient
$\mathcal{B} \subset \Ob(\mathcal{C})$ et
$\text{Cov} \subset \text{Cov}_\mathcal{C}$
des sous-ensembles. Supposons que
\begin{enumerate}
\item pour tout $\mathcal{U} \in \text{Cov}$, on ait
$\mathcal{U} = \{U_i \to U\}_{i \in I}$ avec $I$ fini,
$U, U_i \in \mathcal{B}$ et que tout
$U_{i_0} \times_U \ldots \times_U U_{i_p} \in \mathcal{B}$.
\item pour tout $U \in \mathcal{B}$, les recouvrements de $U$
appartenant à $\text{Cov}$ forment un système cofinal de recouvrements de $U$.
\end{enumerate}
Alors le morphisme
$$
\colim_i H^p(U, \mathcal{F}_i)
\longrightarrow
H^p(U, \colim_i \mathcal{F}_i)
$$
est un isomorphisme pour tout $p \geq 0$, tout $U \in \mathcal{B}$ et
tout diagramme filtrant $\mathcal{I} \to \textit{Ab}(\mathcal{C})$.
\end{lemma}

\begin{proof}
Nous raisonnons par récurrence sur $p$.
Remarquons qu'en (1), nous exigeons que les recouvrements
$\mathcal{U} \in \text{Cov}$ soient finis, de sorte que tous les éléments de
$\mathcal{B}$ soient quasi-compacts. Ainsi, (2) et (1) impliquent que tout
$U \in \mathcal{B}$ satisfait à l'hypothèse (4) du lemme
\ref{sites-lemma-directed-colimits-sections} de Sites.
Le résultat est donc vrai pour $p = 0$. Supposons-le vrai pour $p$ et
démontrons-le pour $p + 1$.

\medskip\noindent
Choisissons un diagramme filtrant
$\mathcal{F} : \mathcal{I} \to \textit{Ab}(\mathcal{C})$,
$i \mapsto \mathcal{F}_i$.
Puisque $\textit{Ab}(\mathcal{C})$ possède des plongements injectifs
fonctoriels, voir Injectifs, théorème
\ref{injectives-theorem-sheaves-injectives}, nous pouvons trouver un
morphisme de diagrammes filtrants
$\mathcal{F} \to \mathcal{I}$
tel que chaque $\mathcal{F}_i \to \mathcal{I}_i$ soit un monomorphisme de
faisceaux abéliens dans un faisceau abélien injectif. Notons $\mathcal{Q}_i$
son conoyau ; nous avons ainsi des suites exactes courtes
$$
0 \to
\mathcal{F}_i \to
\mathcal{I}_i \to
\mathcal{Q}_i \to 0.
$$
Comme les colimites de faisceaux sont les faisceautisés des colimites au
niveau des préfaisceaux, comme la faisceautisation est exacte et comme les
colimites filtrantes de groupes abéliens sont exactes
(voir Algèbre, lemme \ref{algebra-lemma-directed-colimit-exact}),
la suite
$$
0 \to
\colim_i \mathcal{F}_i \to
\colim_i \mathcal{I}_i \to
\colim_i \mathcal{Q}_i \to 0.
$$
est elle aussi exacte courte. Nous affirmons que
$H^q(U, \colim_i \mathcal{I}_i) = 0$ pour tout $U \in \mathcal{B}$
pour tout $q \geq 1$. Admettons provisoirement cette affirmation et
considérons le diagramme
$$
\xymatrix{
\colim_i H^p(U, \mathcal{I}_i) \ar[d] \ar[r] &
\colim_i H^p(U, \mathcal{Q}_i) \ar[d] \ar[r] &
\colim_i H^{p + 1}(U, \mathcal{F}_i) \ar[d] \ar[r] &
0 \ar[d] \\
H^p(U, \colim_i \mathcal{I}_i) \ar[r] &
H^p(U, \colim_i \mathcal{Q}_i) \ar[r] &
H^{p + 1}(U, \colim_i \mathcal{F}_i) \ar[r] &
0
}
$$
Le zéro situé en bas à droite provient de l'affirmation, et celui situé en
haut à droite du fait que les faisceaux $\mathcal{I}_i$ sont injectifs.
La ligne du haut est exacte par application du
lemme \ref{algebra-lemma-directed-colimit-exact} d'Algèbre.
Le lemme du serpent donne alors le résultat pour $p + 1$.

\medskip\noindent
Il reste à démontrer l'affirmation. Nous utiliserons le
lemme \ref{lemma-cech-vanish-collection}.
D'après le cas $p = 0$, pour $\mathcal{U} \in \text{Cov}$, nous avons
$$
\check{\mathcal{C}}^\bullet(\mathcal{U}, \colim_i \mathcal{I}_i)
=
\colim_i \check{\mathcal{C}}^\bullet(\mathcal{U}, \mathcal{I}_i)
$$
car tous les $U_{j_0} \times_U \ldots \times_U U_{j_p}$
appartiennent à $\mathcal{B}$. D'après le
lemme \ref{lemma-injective-trivial-cech}, chacun des complexes de la colimite
de complexes de {\v C}ech est acyclique en degrés $\geq 1$. Par le
lemme \ref{algebra-lemma-directed-colimit-exact} d'Algèbre, le complexe de
{\v C}ech
$\check{\mathcal{C}}^\bullet(\mathcal{U}, \colim_i \mathcal{I}_i)$
est donc lui aussi acyclique en degrés $\geq 1$. Autrement dit,
$\check{H}^p(\mathcal{U}, \colim_i \mathcal{I}_i) = 0$
pour tout $p \geq 1$. Les hypothèses du
lemme \ref{lemma-cech-vanish-collection} sont donc satisfaites, d'où
l'affirmation.
\end{proof}

\begin{lemma}
\label{lemma-colim-global}
Soit $\mathcal{C}$ un site. Soit $S \subset \Ob(\Sh(\mathcal{C}))$
un sous-ensemble. Notons $*$ l'objet final de $\Sh(\mathcal{C})$. Supposons que
\begin{enumerate}
\item pour un certain $K \in S$, le morphisme $K \to *$ soit surjectif ;
\item pour tout morphisme surjectif de faisceaux $\mathcal{F} \to K$ avec
$K \in S$, il existe $K' \in S$ et un morphisme $K' \to \mathcal{F}$ tels
que la composée $K' \to K$ soit surjective ;
\item pour tous $K, K' \in S$, il existe une surjection
$K'' \to K \times K'$ avec $K'' \in S$ ;
\item pour tous $a, b : K \to K'$ avec $K, K' \in S$, il existe une
surjection $K'' \to \text{Equalizer}(a, b)$ avec $K'' \in S$ ; et
\item tout $K \in S$ soit quasi-compact
(Sites, définition \ref{sites-definition-quasi-compact-topos}).
\end{enumerate}
Alors, pour tout $p \geq 0$, le morphisme
$$
\colim_\lambda H^p(\mathcal{C}, \mathcal{F}_\lambda)
\longrightarrow
H^p(\mathcal{C}, \colim_\lambda \mathcal{F}_\lambda)
$$
est un isomorphisme pour tout diagramme filtrant
$\Lambda \to \textit{Ab}(\mathcal{C})$, $\lambda \mapsto \mathcal{F}_\lambda$.
\end{lemma}

\begin{proof}
Nous raisonnons par récurrence sur $p$. Le cas initial $p = 0$ découle de la
partie (4) du lemme \ref{sites-lemma-directed-colimits-global-sections} de
Sites. Vérifions les hypothèses, bien que nous invitions le lecteur à omettre
cette partie. Supposons que $\mathcal{F} \to *$ soit surjectif.
Choisissons $K \in S$ et une surjection $K \to *$ comme en (1). Alors
$\mathcal{F} \times K \to K$ est surjectif. Comme en (2), choisissons
$K' \to \mathcal{F} \times K$ avec $K' \in S$ et tel que $K' \to K$
soit surjectif.
Il existe alors un morphisme $K' \to \mathcal{F}$, et $K' \to *$ est
surjectif. L'hypothèse (4)(a) du lemme
\ref{sites-lemma-directed-colimits-global-sections} de Sites est donc
satisfaite. D'après le lemme \ref{sites-lemma-image-of-quasi-compact} de Sites
et les hypothèses (3) et (5), $K \times K$ est quasi-compact pour tout
$K \in S$. L'hypothèse (4)(b) du lemme
\ref{sites-lemma-directed-colimits-global-sections} de Sites est donc
satisfaite. Ceci achève la preuve du cas initial.

\medskip\noindent
Étape de récurrence. Supposons le résultat vrai pour $H^p$, lorsque
$p \leq p_0$, et pour tout topos $\Sh(\mathcal{C})$ pour lequel il existe un
ensemble $S \subset \Ob(\Sh(\mathcal{C}))$ satisfaisant à (1) -- (5).
En raisonnant exactement comme dans la preuve du
lemme \ref{lemma-colim-works-over-collection}, il suffit de montrer que, pour
toute colimite filtrante $\mathcal{I} = \colim \mathcal{I}_\lambda$ de
faisceaux abéliens injectifs $\mathcal{I}_\lambda$, on a
$H^{p_0 + 1}(\mathcal{C}, \mathcal{I}) = 0$.
Comme en (1), choisissons une surjection $K \to *$ avec $K \in S$.
Notons $K^n$ le produit de $n$ copies de $K$.
Considérons la suite spectrale
$$
E_1^{p, q} = H^q(K^{p + 1}, \mathcal{I}) \Rightarrow
H^{p + q}(*, \mathcal{I}) = H^{p + q}(\mathcal{C}, \mathcal{I})
$$
du lemme \ref{lemma-cech-to-cohomology-sheaf-sets}. Rappelons que
$H^q(K^{p + 1}, \mathcal{F}) = H^q(\mathcal{C}/K^{p + 1}, j^{-1}\mathcal{F})$,
pour tout faisceau abélien sur $\mathcal{C}$ ; voir le
lemme \ref{lemma-cohomology-on-sheaf-sets}. Nous avons
$j^{-1}\mathcal{I} = \colim j^{-1}\mathcal{I}_\lambda$
car $j^{-1}$ commute aux colimites. D'après le
lemme \ref{lemma-cohomology-of-open}, les restrictions
$j^{-1}\mathcal{I}_\lambda$ sont des faisceaux abéliens injectifs sur
$\mathcal{C}/K^{p + 1}$. Nous montrerons ci-dessous que l'hypothèse de
récurrence s'applique à $\mathcal{C}/K^{p + 1}$ ; par conséquent,
$H^q(K^{p + 1}, \mathcal{I}) = \colim H^q(K^{p + 1}, \mathcal{I}_\lambda) = 0$
pour $q < p_0 + 1$ (l'annulation provenant de l'injectivité de
$\mathcal{I}_\lambda$). Il s'ensuit que
$$
H^{p_0 + 1}(\mathcal{C}, \mathcal{I}) =
H^{p_0 + 1}\left(\ldots \to
H^0(K^{p_0}, \mathcal{I}) \to
H^0(K^{p_0 + 1}, \mathcal{I}) \to
H^0(K^{p_0 + 2}, \mathcal{I}) \to \ldots\right)
$$
En utilisant de nouveau l'hypothèse de récurrence, le complexe figurant au
membre de droite est la colimite, sur $\Lambda$, des complexes
$$
\ldots \to
H^0(K^{p_0}, \mathcal{I}_\lambda) \to
H^0(K^{p_0 + 1}, \mathcal{I}_\lambda) \to
H^0(K^{p_0 + 2}, \mathcal{I}_\lambda) \to \ldots
$$
Ces complexes sont exacts puisque $\mathcal{I}_\lambda$ est un faisceau
abélien injectif (cela découle par exemple de la suite spectrale). Comme les
colimites filtrantes sont exactes dans la catégorie des groupes abéliens, nous
obtenons l'annulation voulue.

\medskip\noindent
Il reste à montrer que l'hypothèse de récurrence s'applique au site
$\mathcal{C}/K^n$ pour tout $n \geq 1$. Rappelons que
$\Sh(\mathcal{C}/K^n) = \Sh(\mathcal{C})/K^n$ ; voir Sites,
lemme \ref{sites-lemma-localize-topos-site}.
Nous pouvons donc travailler dans $\Sh(\mathcal{C})/K^n$.
Notons $S_n \subset \Ob(\Sh(\mathcal{C}/K^n))$ l'ensemble des objets de la
forme $K' \to K^n$. Vérifions successivement chaque propriété :
\begin{enumerate}
\item Par (3) et récurrence, il existe une surjection $K' \to K^n$ avec
$K' \in S$. Alors $(K' \to K^n) \to (K^n \to K^n)$ est une surjection dans
$\Sh(\mathcal{C})/K^n$, et $K^n \to K^n$ est l'objet final de
$\Sh(\mathcal{C})/K^n$. Ainsi, (1) est vérifiée pour $S_n$.
\item La propriété (2) pour $S_n$ découle immédiatement de (2) pour $S$.
\item Soient $a : K_1 \to K^n$ et $b : K_2 \to K^n$ dans $S_n$.
Alors $(K_1 \to K^n) \times (K_2 \to K^n)$ est l'objet
$K_1 \times_{K^n} K_2 \to K^n$ de $\Sh(\mathcal{C})/K^n$.
Le sous-faisceau $K_1 \times_{K^n} K_2 \subset K_1 \times K_2$ est
l'égalisateur de $a \circ \text{pr}_1$ et $b \circ \text{pr}_2$.
Écrivons $a = (a_1, \ldots, a_n)$ et $b = (b_1, \ldots, b_n)$.
Choisissons une surjection $K_3 \to K_1 \times K_2$ avec $K_3 \in S$ ;
c'est possible par l'hypothèse (3) pour $\mathcal{C}$.
Choisissons une surjection
$$
K_4
\longrightarrow
\text{Equalizer}(K_3 \to K_1 \times K_2 \xrightarrow{a_1, b_1} K)
$$
avec $K_4 \in S$.
C'est possible par l'hypothèse (4) pour $\mathcal{C}$.
Choisissons une surjection
$$
K_5
\longrightarrow
\text{Equalizer}(K_4 \to K_1 \times K_2 \xrightarrow{a_2, b_2} K)
$$
avec $K_5 \in S$.
C'est encore possible. Poursuivons ainsi jusqu'à obtenir une surjection
$$
K_{3 + n}
\longrightarrow
\text{Equalizer}(K_{2 + n} \to K_1 \times K_2 \xrightarrow{a_n, b_n} K)
$$
avec $K_{3 + n} \in S$. Par construction,
$K_{3 + n} \to K_1 \times_{K^n} K_2$
est surjectif. Ainsi, $(K_{3 + n} \to K^n)$ appartient à $S_n$ et se
surjecte sur le produit $(K_1 \to K^n) \times (K_2 \to K^n)$.
La propriété (3) est donc vérifiée pour $S_n$.
\item La propriété (4) pour $S_n$ découle immédiatement de la propriété
(4) pour $S$.
\item La propriété (5) pour $S_n$ découle de la propriété (5) pour $S$.
En effet, un objet $\mathcal{F} \to K^n$ de $\Sh(\mathcal{C})/K^n$
correspond à un objet quasi-compact de $\Sh(\mathcal{C}/K^n)$ si et seulement
si $\mathcal{F}$ est un objet quasi-compact de $\Sh(\mathcal{C})$.
\end{enumerate}
Ceci achève la preuve du lemme.
\end{proof}

\begin{remark}
\label{remark-colim-global}
Soit $\mathcal{C}$ un site. Soit $\mathcal{B} \subset \Ob(\mathcal{C})$
un sous-ensemble. Soit $S \subset \Ob(\Sh(\mathcal{C}))$
l'ensemble des faisceaux $K$ de la forme
$$
K = \coprod\nolimits_{i = 1, \ldots, n} h_{U_i}^\#
$$
où $U_1, \ldots, U_n \in \mathcal{B}$. On peut alors se demander quand cet
ensemble satisfait aux hypothèses du lemme \ref{lemma-colim-global}. Une
condition suffisante est que
\begin{enumerate}
\item pour certains $n \geq 0$ et $U_1, \ldots, U_n \in \mathcal{B}$,
le morphisme $\coprod_{i = 1, \ldots, n} h_{U_i}^\# \to *$ soit surjectif ;
\item tout recouvrement de $U \in \mathcal{B}$ admette un raffinement par un
recouvrement de la forme $\{U_i \to U\}_{i = 1, \ldots, n}$ avec
$U_i \in \mathcal{B}$ ;
\item étant donnés $U, U' \in \mathcal{B}$, il existe $n \geq 0$,
$U_1, \ldots, U_n \in \mathcal{B}$ et des morphismes $U_i \to U$ et
$U_i \to U'$ tels que $\coprod_{i = 1, \ldots, n} h_{U_i}^\# \to
h_U^\# \times h_{U'}^\#$ soit surjectif ;
\item étant donnés des morphismes $a, b : U \to U'$ dans $\mathcal{C}$ avec
$U, U' \in \mathcal{B}$, il existe $U_1, \ldots, U_n \in \mathcal{B}$ et
des morphismes $U_i \to U$ qui égalisent $a, b$, tels que
$\coprod_{i = 1, \ldots, n} h_{U_i}^\# \to
\text{Equalizer}(h_a^\#, h_b^\# : h_U^\# \to h_{U'}^\#)$
soit surjectif.
\end{enumerate}
Nous omettons la vérification détaillée ; signalons seulement que la partie
(2) ci-dessus assure que tout élément de $\mathcal{B}$ est quasi-compact et
donc, d'après le lemme \ref{sites-lemma-coproduct-quasi-compact} de Sites,
que tout $K \in S$ est également quasi-compact.
\end{remark}

\begin{lemma}
\label{lemma-higher-direct-image-colimit}
Soit $f : \mathcal{C} \to \mathcal{D}$ un morphisme de sites correspondant
au foncteur continu $u : \mathcal{D} \to \mathcal{C}$.
Soit $(\mathcal{F}_i, \varphi_{ii'})$ un système de faisceaux abéliens sur
$\mathcal{C}$. Posons $\mathcal{F} = \colim \mathcal{F}_i$.
Soit $p \geq 0$ un entier. Notons $\mathcal{B}$ l'ensemble des
$V \in \Ob(\mathcal{D})$ tels que
$H^p(u(V), \mathcal{F}) = \colim H^p(u(V), \mathcal{F}_i)$.
Si tout objet de $\mathcal{D}$ possède un recouvrement par des éléments de
$\mathcal{B}$, alors $R^pf_*\mathcal{F} = \colim R^pf_*\mathcal{F}_i$.
\end{lemma}

\begin{proof}
Rappelons que $R^pf_*\mathcal{F}$ est le faisceautisé du préfaisceau
$\mathcal{G}$ qui envoie $V$ sur $H^p(u(V), \mathcal{F})$ ; voir le
lemme \ref{lemma-higher-direct-images}. De même,
$R^pf_*\mathcal{F}_i$ est le faisceautisé du préfaisceau
$\mathcal{G}_i$ qui envoie $V$ sur $H^p(u(V), \mathcal{F}_i)$.
Rappelons que la faisceautisation est adjointe à gauche de l'inclusion des
faisceaux dans les préfaisceaux ; voir Sites, section
\ref{sites-section-sheafification}. Elle commute donc aux colimites ; voir
Catégories, lemme \ref{categories-lemma-adjoint-exact}.
Il suffit par conséquent de montrer que le morphisme de préfaisceaux (la
colimite étant prise dans la catégorie des préfaisceaux)
$$
\colim \mathcal{G}_i \longrightarrow \mathcal{G}
$$
induit un isomorphisme sur les faisceautisés. Cela découle du
lemme \ref{sites-lemma-isomorphism-sheafifications} de Sites et de notre
hypothèse sur $\mathcal{B}$.
\end{proof}

\begin{lemma}
\label{lemma-colim-sites-injective}
Soit $\mathcal{I}$ une catégorie d'indices cofiltrante, et soit
$(\mathcal{C}_i, f_a)$ un système projectif de sites indexé par $\mathcal{I}$,
comme dans Sites, situation \ref{sites-situation-inverse-limit-sites}.
Posons $\mathcal{C} = \colim \mathcal{C}_i$ comme dans Sites, lemmes
\ref{sites-lemma-colimit-sites} et
\ref{sites-lemma-compute-pullback-to-limit}.
Supposons de plus donnés
\begin{enumerate}
\item un faisceau abélien $\mathcal{F}_i$ sur $\mathcal{C}_i$ pour tout
$i \in \Ob(\mathcal{I})$ ;
\item pour $a : j \to i$, un morphisme
$\varphi_a : f_a^{-1}\mathcal{F}_i \to \mathcal{F}_j$
de faisceaux abéliens sur $\mathcal{C}_j$,
\end{enumerate}
de sorte que $\varphi_c = \varphi_b \circ f_b^{-1}\varphi_a$ dès que
$c = a \circ b$. Il existe alors un morphisme de systèmes
$(\mathcal{F}_i, \varphi_a) \to (\mathcal{G}_i, \psi_a)$
tel que $\mathcal{F}_i \to \mathcal{G}_i$ soit injectif et que
$\mathcal{G}_i$ soit un faisceau abélien injectif.
\end{lemma}

\begin{proof}
Pour chaque $i$, choisissons une injection $\mathcal{F}_i \to \mathcal{A}_i$,
où $\mathcal{A}_i$ est un faisceau abélien injectif sur $\mathcal{C}_i$.
Nous pouvons alors considérer la famille de morphismes
$$
\gamma_i :
\mathcal{F}_i
\longrightarrow
\prod\nolimits_{b : k \to i} f_{b, *}\mathcal{A}_k = \mathcal{G}_i
$$
dont les composantes sont les morphismes adjoints aux morphismes
$f_b^{-1}\mathcal{F}_i \to \mathcal{F}_k \to \mathcal{A}_k$.
Pour $a : j \to i$ dans $\mathcal{I}$, il existe un morphisme canonique
$$
\psi_a : f_a^{-1}\mathcal{G}_i \to \mathcal{G}_j
$$
dont les composantes sont les morphismes canoniques
$f_b^{-1}f_{a \circ b, *}\mathcal{A}_k \to f_{b, *}\mathcal{A}_k$
pour $b : k \to j$. Nous obtenons ainsi une injection
$(\gamma_i) : (\mathcal{F}_i, \varphi_a) \to (\mathcal{G}_i, \psi_a)$
de systèmes de faisceaux abéliens. Remarquons que $\mathcal{G}_i$ est un
faisceau injectif de groupes abéliens sur $\mathcal{C}_i$ ; voir le
lemme \ref{lemma-pushforward-injective-flat} et Homologie,
lemme \ref{homology-lemma-product-injectives}.
Ceci achève la construction.
\end{proof}

\begin{lemma}
\label{lemma-colimit}
Dans la situation du lemme \ref{lemma-colim-sites-injective}, posons
$\mathcal{F} = \colim f_i^{-1}\mathcal{F}_i$.
Soient $i \in \Ob(\mathcal{I})$ et $X_i \in \text{Ob}(\mathcal{C}_i)$. Alors
$$
\colim_{a : j \to i} H^p(u_a(X_i), \mathcal{F}_j) =
H^p(u_i(X_i), \mathcal{F})
$$
pour tout $p \geq 0$.
\end{lemma}

\begin{proof}
Le cas $p = 0$ est le lemme \ref{sites-lemma-colimit} de Sites.

\medskip\noindent
Choisissons $(\mathcal{F}_i, \varphi_a) \to (\mathcal{G}_i, \psi_a)$
comme dans le lemme \ref{lemma-colim-sites-injective}.
En raisonnant exactement comme dans la preuve du
lemme \ref{lemma-colim-works-over-collection}, il suffit de montrer que
$H^p(X, \colim f_i^{-1}\mathcal{G}_i) = 0$ pour $p > 0$.

\medskip\noindent
Posons $\mathcal{G} = \colim f_i^{-1}\mathcal{G}_i$.
Pour montrer l'annulation de la cohomologie de $\mathcal{G}$ sur tout objet de
$\mathcal{C}$, nous montrons que la cohomologie de {\v C}ech de $\mathcal{G}$
est nulle pour tout recouvrement $\mathcal{U}$ de $\mathcal{C}$
(lemme \ref{lemma-cech-vanish-collection}).
Le recouvrement $\mathcal{U}$ provient d'un recouvrement
$\mathcal{U}_i$ de $\mathcal{C}_i$ pour un certain $i$. Nous avons
$$
\check{\mathcal{C}}^\bullet(\mathcal{U}, \mathcal{G}) =
\colim_{a : j \to i}
\check{\mathcal{C}}^\bullet(u_a(\mathcal{U}_i), \mathcal{G}_j)
$$
d'après le cas $p = 0$. Le membre de droite est acyclique en degrés positifs,
car il est une colimite filtrante de complexes acycliques, d'après le
lemme \ref{lemma-injective-trivial-cech}. Voir Algèbre,
lemme \ref{algebra-lemma-directed-colimit-exact}.
\end{proof}















\section{Résolutions plates}
\label{section-flat}

\noindent
Dans cette section, nous reprenons les arguments de Cohomologie,
section \ref{cohomology-section-flat}, dans le cadre des sites annelés et des
topos annelés.

\begin{lemma}
\label{lemma-derived-tor-exact}
Soit $(\mathcal{C}, \mathcal{O})$ un site annelé.
Soit $\mathcal{G}^\bullet$ un complexe de $\mathcal{O}$-modules.
Les foncteurs
$$
K(\textit{Mod}(\mathcal{O}))
\longrightarrow
K(\textit{Mod}(\mathcal{O})),
\quad
\mathcal{F}^\bullet \longmapsto
\text{Tot}(\mathcal{G}^\bullet \otimes_\mathcal{O} \mathcal{F}^\bullet)
$$
et
$$
K(\textit{Mod}(\mathcal{O}))
\longrightarrow
K(\textit{Mod}(\mathcal{O})),
\quad
\mathcal{F}^\bullet \longmapsto
\text{Tot}(\mathcal{F}^\bullet \otimes_\mathcal{O} \mathcal{G}^\bullet)
$$
sont des foncteurs exacts de catégories triangulées.
\end{lemma}

\begin{proof}
Cela résulte de Catégories dérivées, remarque
\ref{derived-remark-double-complex-as-tensor-product-of}.
\end{proof}

\begin{definition}
\label{definition-K-flat}
Soit $(\mathcal{C}, \mathcal{O})$ un site annelé.
Un complexe $\mathcal{K}^\bullet$ de $\mathcal{O}$-modules est dit
{\it K-plat} si, pour tout complexe acyclique $\mathcal{F}^\bullet$ de
$\mathcal{O}$-modules, le complexe
$$
\text{Tot}(\mathcal{F}^\bullet \otimes_\mathcal{O} \mathcal{K}^\bullet)
$$
est acyclique.
\end{definition}

\begin{lemma}
\label{lemma-K-flat-quasi-isomorphism}
Soit $(\mathcal{C}, \mathcal{O})$ un site annelé.
Soit $\mathcal{K}^\bullet$ un complexe K-plat.
Alors le foncteur
$$
K(\textit{Mod}(\mathcal{O}))
\longrightarrow
K(\textit{Mod}(\mathcal{O})), \quad
\mathcal{F}^\bullet
\longmapsto
\text{Tot}(\mathcal{F}^\bullet \otimes_\mathcal{O} \mathcal{K}^\bullet)
$$
transforme les quasi-isomorphismes en quasi-isomorphismes.
\end{lemma}

\begin{proof}
Cela résulte du lemme \ref{lemma-derived-tor-exact} et du fait que les
quasi-isomorphismes se caractérisent par l'acyclicité de leurs cônes.
\end{proof}

\begin{lemma}
\label{lemma-restriction-K-flat}
Soit $(\mathcal{C}, \mathcal{O})$ un site annelé.
Soit $U$ un objet de $\mathcal{C}$.
Si $\mathcal{K}^\bullet$ est un complexe K-plat de $\mathcal{O}$-modules,
alors $\mathcal{K}^\bullet|_U$ est un complexe K-plat de
$\mathcal{O}_U$-modules.
\end{lemma}

\begin{proof}
Soit $\mathcal{G}^\bullet$ un complexe exact de $\mathcal{O}_U$-modules.
Puisque $j_{U!}$ est exact (Modules sur les sites,
lemme \ref{sites-modules-lemma-extension-by-zero-exact}) et que
$\mathcal{K}^\bullet$ est un complexe K-plat de $\mathcal{O}$-modules, le
complexe
$$
j_{U!}(\text{Tot}(\mathcal{G}^\bullet \otimes_{\mathcal{O}_U}
\mathcal{K}^\bullet|_U)) =
\text{Tot}(j_{U!}\mathcal{G}^\bullet \otimes_\mathcal{O} \mathcal{K}^\bullet)
$$
est exact. Ici, l'égalité découle de Modules sur les sites,
lemme \ref{sites-modules-lemma-j-shriek-and-tensor}, et du fait que $j_{U!}$
commute aux sommes directes (en tant qu'adjoint à gauche).
Nous concluons puisque $j_{U!}$ reflète l'exactitude, d'après Modules sur les
sites, lemme \ref{sites-modules-lemma-j-shriek-reflects-exactness}.
\end{proof}

\begin{lemma}
\label{lemma-tensor-product-K-flat}
Soit $(\mathcal{C}, \mathcal{O})$ un site annelé.
Si $\mathcal{K}^\bullet$ et $\mathcal{L}^\bullet$ sont des complexes K-plats
de $\mathcal{O}$-modules, alors
$\text{Tot}(\mathcal{K}^\bullet \otimes_\mathcal{O} \mathcal{L}^\bullet)$
est un complexe K-plat de $\mathcal{O}$-modules.
\end{lemma}

\begin{proof}
Cela résulte de l'isomorphisme
$$
\text{Tot}(\mathcal{M}^\bullet \otimes_\mathcal{O}
\text{Tot}(\mathcal{K}^\bullet \otimes_\mathcal{O} \mathcal{L}^\bullet))
=
\text{Tot}(\text{Tot}(\mathcal{M}^\bullet \otimes_\mathcal{O}
\mathcal{K}^\bullet) \otimes_\mathcal{O} \mathcal{L}^\bullet)
$$
et de la définition.
\end{proof}

\begin{lemma}
\label{lemma-K-flat-two-out-of-three}
Soit $(\mathcal{C}, \mathcal{O})$ un site annelé.
Soit $(\mathcal{K}_1^\bullet, \mathcal{K}_2^\bullet, \mathcal{K}_3^\bullet)$
un triangle distingué dans $K(\textit{Mod}(\mathcal{O}))$.
Si deux des trois complexes $\mathcal{K}_i^\bullet$ sont K-plats, le
troisième l'est aussi.
\end{lemma}

\begin{proof}
Cela résulte du lemme \ref{lemma-derived-tor-exact} et du fait que, dans un
triangle distingué de $K(\textit{Mod}(\mathcal{O}))$, si deux des trois objets
sont acycliques, le troisième l'est aussi.
\end{proof}

\begin{lemma}
\label{lemma-K-flat-two-out-of-three-ses}
Soit $(\mathcal{C}, \mathcal{O})$ un site annelé. Soit
$0 \to \mathcal{K}_1^\bullet \to \mathcal{K}_2^\bullet \to
\mathcal{K}_3^\bullet \to 0$ une suite exacte courte de complexes telle que
les termes de $\mathcal{K}_3^\bullet$ soient des $\mathcal{O}$-modules plats.
Si deux des trois complexes $\mathcal{K}_i^\bullet$ sont K-plats, le troisième
l'est aussi.
\end{lemma}

\begin{proof}
D'après Modules sur les sites,
lemme \ref{sites-modules-lemma-flat-tor-zero}, pour tout complexe
$\mathcal{L}^\bullet$, nous obtenons une suite exacte courte
$$
0 \to
\text{Tot}(\mathcal{L}^\bullet \otimes_\mathcal{O} \mathcal{K}_1^\bullet) \to
\text{Tot}(\mathcal{L}^\bullet \otimes_\mathcal{O} \mathcal{K}_2^\bullet) \to
\text{Tot}(\mathcal{L}^\bullet \otimes_\mathcal{O} \mathcal{K}_3^\bullet) \to 0
$$
de complexes. Le lemme résulte donc de la suite exacte longue des faisceaux de
cohomologie et de la définition des complexes K-plats.
\end{proof}

\begin{lemma}
\label{lemma-bounded-flat-K-flat}
Soit $(\mathcal{C}, \mathcal{O})$ un site annelé. Tout complexe borné
supérieurement de $\mathcal{O}$-modules plats est K-plat.
\end{lemma}

\begin{proof}
Soit $\mathcal{K}^\bullet$ un complexe borné supérieurement de
$\mathcal{O}$-modules plats. Soit $\mathcal{L}^\bullet$ un complexe acyclique
de $\mathcal{O}$-modules. Remarquons que
$\mathcal{L}^\bullet = \colim_m \tau_{\leq m}\mathcal{L}^\bullet$
où les colimites sont prises terme à terme. Ainsi,
$$
\text{Tot}(\mathcal{K}^\bullet \otimes_\mathcal{O} \mathcal{L}^\bullet)
=
\colim_m \text{Tot}(
\mathcal{K}^\bullet \otimes_\mathcal{O} \tau_{\leq m}\mathcal{L}^\bullet)
$$
terme à terme. Pour prouver que le complexe de gauche est acyclique, il suffit
donc de montrer que chacun des complexes de droite l'est. Comme
$\tau_{\leq m}\mathcal{L}^\bullet$ est acyclique, nous sommes ramenés au cas
où $\mathcal{L}^\bullet$ est borné supérieurement.
Dans ce cas, la suite spectrale du
lemme \ref{homology-lemma-first-quadrant-ss} d'Homologie vérifie
$$
{}'E_1^{p, q} = H^p(\mathcal{L}^\bullet \otimes_\mathcal{O} \mathcal{K}^q)
$$
qui est nul puisque $\mathcal{K}^q$ est plat et $\mathcal{L}^\bullet$
acyclique. Ceci conclut la preuve.
\end{proof}

\begin{lemma}
\label{lemma-colimit-K-flat}
Soit $(\mathcal{C}, \mathcal{O})$ un site annelé.
Soit $\mathcal{K}_1^\bullet \to \mathcal{K}_2^\bullet \to \ldots$
un système de complexes K-plats.
Alors $\colim_i \mathcal{K}_i^\bullet$ est K-plat.
\end{lemma}

\begin{proof}
Puisque nous prenons les colimites terme à terme, il est clair que
$$
\colim_i \text{Tot}(
\mathcal{F}^\bullet \otimes_\mathcal{O} \mathcal{K}_i^\bullet)
=
\text{Tot}(\mathcal{F}^\bullet \otimes_\mathcal{O}
\colim_i \mathcal{K}_i^\bullet)
$$
Le lemme découle donc de l'exactitude des colimites filtrantes.
\end{proof}

\begin{lemma}
\label{lemma-resolution-by-direct-sums-extensions-by-zero}
Soit $(\mathcal{C}, \mathcal{O})$ un site annelé.
Pour tout complexe $\mathcal{G}^\bullet$ de $\mathcal{O}$-modules, il existe
un diagramme commutatif de complexes de $\mathcal{O}$-modules
$$
\xymatrix{
\mathcal{K}_1^\bullet \ar[d] \ar[r] &
\mathcal{K}_2^\bullet \ar[d] \ar[r] & \ldots \\
\tau_{\leq 1}\mathcal{G}^\bullet \ar[r] &
\tau_{\leq 2}\mathcal{G}^\bullet \ar[r] & \ldots
}
$$
ayant les propriétés suivantes : (1) les flèches verticales sont des
quasi-isomorphismes et sont surjectives terme à terme ;
(2) chaque $\mathcal{K}_n^\bullet$ est un complexe borné supérieurement dont
les termes sont des sommes directes de $\mathcal{O}$-modules de la forme
$j_{U!}\mathcal{O}_U$ ; et (3) les morphismes
$\mathcal{K}_n^\bullet \to \mathcal{K}_{n + 1}^\bullet$ sont des injections
scindées terme à terme dont les conoyaux sont des sommes directes de
$\mathcal{O}$-modules de la forme $j_{U!}\mathcal{O}_U$. De plus, le morphisme
$\colim \mathcal{K}_n^\bullet \to \mathcal{G}^\bullet$ est un quasi-isomorphisme.
\end{lemma}

\begin{proof}
L'existence du diagramme et les propriétés (1), (2), (3) découlent
immédiatement de Modules sur les sites,
lemme \ref{sites-modules-lemma-module-quotient-flat}, et de Catégories
dérivées, lemme \ref{derived-lemma-special-direct-system}.
Le morphisme induit
$\colim \mathcal{K}_n^\bullet \to \mathcal{G}^\bullet$
est un quasi-isomorphisme parce que les colimites filtrantes sont exactes.
\end{proof}

\begin{lemma}
\label{lemma-K-flat-resolution}
Soit $(\mathcal{C}, \mathcal{O})$ un site annelé. Pour tout complexe
$\mathcal{G}^\bullet$, il existe un complexe $K$-plat $\mathcal{K}^\bullet$
dont les termes sont des $\mathcal{O}$-modules plats, ainsi qu'un
quasi-isomorphisme $\mathcal{K}^\bullet \to \mathcal{G}^\bullet$ surjectif
terme à terme.
\end{lemma}

\begin{proof}
Choisissons un diagramme comme dans le
lemme \ref{lemma-resolution-by-direct-sums-extensions-by-zero}.
Chaque complexe $\mathcal{K}_n^\bullet$ est un complexe borné supérieurement
de modules plats ; voir Modules sur les sites,
lemme \ref{sites-modules-lemma-j-shriek-flat}.
Ainsi, $\mathcal{K}_n^\bullet$ est K-plat d'après le
lemme \ref{lemma-bounded-flat-K-flat}.
Par conséquent, $\colim \mathcal{K}_n^\bullet$ est K-plat d'après le
lemme \ref{lemma-colimit-K-flat}.
Le morphisme induit
$\colim \mathcal{K}_n^\bullet \to \mathcal{G}^\bullet$
est, par construction, un quasi-isomorphisme surjectif terme à terme.
La propriété (3) du
lemme \ref{lemma-resolution-by-direct-sums-extensions-by-zero} montre que
$\colim \mathcal{K}_n^m$ est une somme directe de modules plats, donc est plat,
ce qui prouve la dernière assertion.
\end{proof}

\begin{lemma}
\label{lemma-derived-tor-quasi-isomorphism-other-side}
Soit $(\mathcal{C}, \mathcal{O})$ un site annelé. Soit
$\alpha : \mathcal{P}^\bullet \to \mathcal{Q}^\bullet$ un quasi-isomorphisme
de complexes K-plats de $\mathcal{O}$-modules.
Pour tout complexe $\mathcal{F}^\bullet$ de $\mathcal{O}$-modules, le
morphisme induit
$$
\text{Tot}(\text{id}_{\mathcal{F}^\bullet} \otimes \alpha) :
\text{Tot}(\mathcal{F}^\bullet \otimes_\mathcal{O} \mathcal{P}^\bullet)
\longrightarrow
\text{Tot}(\mathcal{F}^\bullet \otimes_\mathcal{O} \mathcal{Q}^\bullet)
$$
est un quasi-isomorphisme.
\end{lemma}

\begin{proof}
Choisissons un quasi-isomorphisme $\mathcal{K}^\bullet \to \mathcal{F}^\bullet$,
où $\mathcal{K}^\bullet$ est un complexe K-plat ; voir le
lemme \ref{lemma-K-flat-resolution}.
Considérons le diagramme commutatif
$$
\xymatrix{
\text{Tot}(\mathcal{K}^\bullet
\otimes_\mathcal{O} \mathcal{P}^\bullet) \ar[r] \ar[d] &
\text{Tot}(\mathcal{K}^\bullet
\otimes_\mathcal{O} \mathcal{Q}^\bullet) \ar[d] \\
\text{Tot}(\mathcal{F}^\bullet
\otimes_\mathcal{O} \mathcal{P}^\bullet) \ar[r] &
\text{Tot}(\mathcal{F}^\bullet
\otimes_\mathcal{O} \mathcal{Q}^\bullet)
}
$$
Le résultat découle du lemme \ref{lemma-K-flat-quasi-isomorphism}, puisque les
flèches verticales et la flèche horizontale supérieure sont des
quasi-isomorphismes.
\end{proof}

\noindent
Soit $(\mathcal{C}, \mathcal{O})$ un site annelé.
Soit $\mathcal{F}^\bullet$ un objet de $D(\mathcal{O})$.
Choisissons une résolution K-plate
$\mathcal{K}^\bullet \to \mathcal{F}^\bullet$ ; voir le
lemme \ref{lemma-K-flat-resolution}.
D'après le lemme \ref{lemma-derived-tor-exact}, nous obtenons un foncteur exact
de catégories triangulées
$$
K(\mathcal{O})
\longrightarrow
K(\mathcal{O}),
\quad
\mathcal{G}^\bullet
\longmapsto
\text{Tot}(\mathcal{G}^\bullet \otimes_\mathcal{O} \mathcal{K}^\bullet)
$$
D'après le lemme \ref{lemma-K-flat-quasi-isomorphism}, ce foncteur induit un
foncteur $D(\mathcal{O}) \to D(\mathcal{O})$, puisque $D(\mathcal{O})$ est la
localisation de $K(\mathcal{O})$ relativement aux quasi-isomorphismes.
Comme la catégorie des résolutions $K$-plates de $\mathcal{F}^\bullet$ est
cofiltrante et compte tenu du
lemme \ref{lemma-derived-tor-quasi-isomorphism-other-side}, le foncteur obtenu,
à isomorphisme près, ne dépend pas du choix de $\mathcal{K}^\bullet$.

\begin{definition}
\label{definition-derived-tor}
Soit $(\mathcal{C}, \mathcal{O})$ un site annelé.
Soit $\mathcal{F}^\bullet$ un objet de $D(\mathcal{O})$.
Le {\it produit tensoriel dérivé}
$$
- \otimes_\mathcal{O}^{\mathbf{L}} \mathcal{F}^\bullet :
D(\mathcal{O})
\longrightarrow
D(\mathcal{O})
$$
est le foncteur exact de catégories triangulées décrit ci-dessus.
\end{definition}

\noindent
Il ressort de nos constructions explicites qu'il existe un isomorphisme
canonique
$$
\mathcal{F}^\bullet \otimes_\mathcal{O}^{\mathbf{L}} \mathcal{G}^\bullet
\cong
\mathcal{G}^\bullet \otimes_\mathcal{O}^{\mathbf{L}} \mathcal{F}^\bullet
$$
pour $\mathcal{G}^\bullet$ et $\mathcal{F}^\bullet$ dans $D(\mathcal{O})$.
Nous utilisons ici les règles de signes données dans Compléments sur l'algèbre,
section \ref{more-algebra-section-sign-rules}. Ainsi, lorsque nous écrivons
$\mathcal{F}^\bullet \otimes_\mathcal{O}^{\mathbf{L}} \mathcal{G}^\bullet$
nous ne préciserons généralement pas par rapport à laquelle des deux variables
nous définissons le produit tensoriel dérivé.

\begin{definition}
\label{definition-tor}
Soit $(\mathcal{C}, \mathcal{O})$ un site annelé.
Soient $\mathcal{F}$ et $\mathcal{G}$ des $\mathcal{O}$-modules.
Les {\it Tor} de $\mathcal{F}$ et $\mathcal{G}$ sont définis par la formule
$$
\text{Tor}_p^\mathcal{O}(\mathcal{F}, \mathcal{G}) =
H^{-p}(\mathcal{F} \otimes_\mathcal{O}^\mathbf{L} \mathcal{G})
$$
où le produit tensoriel dérivé est celui défini ci-dessus.
\end{definition}

\noindent
Cette définition implique que, pour toute suite exacte courte de
$\mathcal{O}$-modules
$0 \to \mathcal{F}_1 \to \mathcal{F}_2 \to \mathcal{F}_3 \to 0$
nous avons une suite exacte longue de cohomologie
$$
\xymatrix{
\mathcal{F}_1 \otimes_\mathcal{O} \mathcal{G} \ar[r] &
\mathcal{F}_2 \otimes_\mathcal{O} \mathcal{G} \ar[r] &
\mathcal{F}_3 \otimes_\mathcal{O} \mathcal{G} \ar[r] & 0 \\
\text{Tor}_1^\mathcal{O}(\mathcal{F}_1, \mathcal{G}) \ar[r] &
\text{Tor}_1^\mathcal{O}(\mathcal{F}_2, \mathcal{G}) \ar[r] &
\text{Tor}_1^\mathcal{O}(\mathcal{F}_3, \mathcal{G}) \ar[ull]
}
$$
pour tout $\mathcal{O}$-module $\mathcal{G}$. On l'appelle la suite exacte
longue de $\text{Tor}$ associée à la situation.

\begin{lemma}
\label{lemma-flat-tor-zero}
Soit $(\mathcal{C}, \mathcal{O})$ un site annelé.
Soit $\mathcal{F}$ un $\mathcal{O}$-module.
Les conditions suivantes sont équivalentes :
\begin{enumerate}
\item $\mathcal{F}$ est un $\mathcal{O}$-module plat ;
\item $\text{Tor}_1^\mathcal{O}(\mathcal{F}, \mathcal{G}) = 0$
pour tout $\mathcal{O}$-module $\mathcal{G}$.
\end{enumerate}
\end{lemma}

\begin{proof}
Si $\mathcal{F}$ est plat, alors $\mathcal{F} \otimes_\mathcal{O} -$ est un
foncteur exact et ses foncteurs dérivés s'annulent. Réciproquement, supposons
(2). Si $\mathcal{G} \to \mathcal{H}$ est injectif, de conoyau $\mathcal{Q}$,
la suite exacte longue de $\text{Tor}$ montre que le noyau de
$\mathcal{F} \otimes_\mathcal{O} \mathcal{G} \to
\mathcal{F} \otimes_\mathcal{O} \mathcal{H}$
est un quotient de
$\text{Tor}_1^\mathcal{O}(\mathcal{F}, \mathcal{Q})$
qui est nul par hypothèse. Par conséquent, $\mathcal{F}$ est plat.
\end{proof}

\begin{lemma}
\label{lemma-K-flat-flat-acyclic}
Soit $(\mathcal{C}, \mathcal{O})$ un site annelé. Soit $\mathcal{K}^\bullet$
un complexe K-plat, acyclique et à termes plats. Alors
$\mathcal{F} = \Ker(\mathcal{K}^n \to \mathcal{K}^{n + 1})$
est un $\mathcal{O}$-module plat.
\end{lemma}

\begin{proof}
Observons que
$$
\ldots \to \mathcal{K}^{n - 2} \to \mathcal{K}^{n - 1} \to
\mathcal{F} \to 0
$$
est une résolution plate de notre module $\mathcal{F}$. Comme un complexe
borné supérieurement de modules plats est K-plat
(lemme \ref{lemma-bounded-flat-K-flat}), nous pouvons utiliser cette résolution
pour calculer $\text{Tor}_i(\mathcal{F}, \mathcal{G})$ pour tout
$\mathcal{O}$-module $\mathcal{G}$. D'une part,
$\mathcal{K}^\bullet \otimes_\mathcal{O}^\mathbf{L} \mathcal{G}$ est nul dans
$D(\mathcal{O})$ parce que $\mathcal{K}^\bullet$ est acyclique et, d'autre
part, il est représenté par
$\mathcal{K}^\bullet \otimes_\mathcal{O} \mathcal{G}$.
Nous voyons donc que
$$
\mathcal{K}^{n - 3} \otimes_\mathcal{O} \mathcal{G} \to
\mathcal{K}^{n - 2} \otimes_\mathcal{O} \mathcal{G} \to
\mathcal{K}^{n - 1} \otimes_\mathcal{O} \mathcal{G}
$$
est exacte. Ainsi, $\text{Tor}_1(\mathcal{F}, \mathcal{G}) = 0$, et nous
concluons par le lemme \ref{lemma-flat-tor-zero}.
\end{proof}

\begin{lemma}
\label{lemma-factor-through-K-flat}
Soit $(\mathcal{C}, \mathcal{O})$ un site annelé.
Soit $a : \mathcal{K}^\bullet \to \mathcal{L}^\bullet$ un morphisme de
complexes de $\mathcal{O}$-modules. Si $\mathcal{K}^\bullet$ est K-plat, il
existe un complexe $\mathcal{N}^\bullet$ et des morphismes de complexes
$b : \mathcal{K}^\bullet \to \mathcal{N}^\bullet$
et $c : \mathcal{N}^\bullet \to \mathcal{L}^\bullet$ tels que
\begin{enumerate}
\item $\mathcal{N}^\bullet$ est K-plat ;
\item $c$ est un quasi-isomorphisme ;
\item $a$ est homotope à $c \circ b$.
\end{enumerate}
Si les termes de $\mathcal{K}^\bullet$ sont plats, nous pouvons choisir
$\mathcal{N}^\bullet$, $b$ et $c$ de sorte qu'il en soit de même pour
$\mathcal{N}^\bullet$.
\end{lemma}

\begin{proof}
Nous utiliserons le fait que la catégorie d'homotopie
$K(\textit{Mod}(\mathcal{O}))$ est une catégorie triangulée ; voir Catégories
dérivées, proposition
\ref{derived-proposition-homotopy-category-triangulated}.
Choisissons un triangle distingué
$\mathcal{K}^\bullet \to \mathcal{L}^\bullet \to
\mathcal{C}^\bullet \to \mathcal{K}^\bullet[1]$.
Choisissons un quasi-isomorphisme
$\mathcal{M}^\bullet \to \mathcal{C}^\bullet$, où $\mathcal{M}^\bullet$ est
K-plat à termes plats ; voir le lemme \ref{lemma-K-flat-resolution}.
Par les axiomes des catégories triangulées, nous pouvons inscrire la composée
$\mathcal{M}^\bullet \to \mathcal{C}^\bullet \to \mathcal{K}^\bullet[1]$
dans un triangle distingué
$\mathcal{K}^\bullet \to \mathcal{N}^\bullet \to
\mathcal{M}^\bullet \to \mathcal{K}^\bullet[1]$.
D'après le lemme \ref{lemma-K-flat-two-out-of-three}, $\mathcal{N}^\bullet$ est
K-plat. En utilisant de nouveau les axiomes des catégories triangulées, nous
pouvons choisir un morphisme $\mathcal{N}^\bullet \to \mathcal{L}^\bullet$
qui s'insère dans le morphisme suivant de triangles distingués
$$
\xymatrix{
\mathcal{K}^\bullet \ar[r] \ar[d] &
\mathcal{N}^\bullet \ar[r] \ar[d] &
\mathcal{M}^\bullet \ar[r] \ar[d] &
\mathcal{K}^\bullet[1] \ar[d] \\
\mathcal{K}^\bullet \ar[r] &
\mathcal{L}^\bullet \ar[r] &
\mathcal{C}^\bullet \ar[r] &
\mathcal{K}^\bullet[1]
}
$$
Puisque deux des trois flèches sont des quasi-isomorphismes, la troisième
$\mathcal{N}^\bullet \to \mathcal{L}^\bullet$ l'est aussi, d'après les suites
exactes longues de cohomologie associées à ces triangles distingués (on peut
aussi considérer l'image de ce diagramme dans $D(\mathcal{O})$ et utiliser
Catégories dérivées, lemme \ref{derived-lemma-third-isomorphism-triangle}).
Ceci achève la preuve de (1), (2) et (3).
Pour prouver la dernière assertion, nous pouvons choisir
$\mathcal{N}^\bullet$ de sorte que
$\mathcal{N}^n \cong \mathcal{M}^n \oplus \mathcal{K}^n$ ; voir
Catégories dérivées, lemme
\ref{derived-lemma-improve-distinguished-triangle-homotopy}.
Nous obtenons ainsi la planéité voulue si les termes de
$\mathcal{K}^\bullet$ sont plats.
\end{proof}
\section{Image inverse dérivée}
\label{section-derived-pullback}

\noindent
Soit
$f : (\Sh(\mathcal{C}), \mathcal{O}) \to
(\Sh(\mathcal{C}'), \mathcal{O}')$
un morphisme de topos annelés. Nous pouvons utiliser des résolutions
K-plates pour définir un foncteur d'image inverse dérivée
$$
Lf^* : D(\mathcal{O}') \to D(\mathcal{O})
$$

\begin{lemma}
\label{lemma-pullback-K-flat}
Soit $f : (\Sh(\mathcal{C}'), \mathcal{O}') \to (\Sh(\mathcal{C}), \mathcal{O})$
un morphisme de topos annelés. Soit $\mathcal{K}^\bullet$ un complexe K-plat
de $\mathcal{O}$-modules dont les termes sont des $\mathcal{O}$-modules plats.
Alors $f^*\mathcal{K}^\bullet$ est un complexe K-plat de
$\mathcal{O}'$-modules dont les termes sont des $\mathcal{O}'$-modules plats.
\end{lemma}

\begin{proof}
Les termes $f^*\mathcal{K}^n$ sont des $\mathcal{O}'$-modules plats d'après
Modules sur les sites, Lemme \ref{sites-modules-lemma-pullback-flat}.
Choisissons un diagramme
$$
\xymatrix{
\mathcal{K}_1^\bullet \ar[d] \ar[r] &
\mathcal{K}_2^\bullet \ar[d] \ar[r] & \ldots \\
\tau_{\leq 1}\mathcal{K}^\bullet \ar[r] &
\tau_{\leq 2}\mathcal{K}^\bullet \ar[r] & \ldots
}
$$
comme dans le Lemme \ref{lemma-resolution-by-direct-sums-extensions-by-zero}.
Nous utiliserons sans autre mention toutes les propriétés énoncées dans ce
lemme. Chaque $\mathcal{K}_n^\bullet$ est un complexe borné supérieurement
de modules plats, voir
Modules sur les sites, Lemme \ref{sites-modules-lemma-j-shriek-flat}.
Considérons la suite exacte courte de complexes
$$
0 \to \mathcal{M}^\bullet \to
\colim \mathcal{K}_n^\bullet \to
\mathcal{K}^\bullet \to 0
$$
qui définit $\mathcal{M}^\bullet$. D'après les Lemmes
\ref{lemma-bounded-flat-K-flat} et \ref{lemma-colimit-K-flat}, le complexe
$\colim \mathcal{K}_n^\bullet$ est K-plat et, d'après
Modules sur les sites, Lemme \ref{sites-modules-lemma-colimits-flat},
ses termes sont plats. D'après Modules sur les sites, Lemme
\ref{sites-modules-lemma-flat-ses}, les termes de $\mathcal{M}^\bullet$
sont plats ; d'après le Lemme \ref{lemma-K-flat-two-out-of-three-ses},
$\mathcal{M}^\bullet$ est K-plat ; enfin, d'après la suite exacte longue
de cohomologie, $\mathcal{M}^\bullet$ est acyclique (car la seconde flèche
est un quasi-isomorphisme). L'image inverse
$f^*(\colim \mathcal{K}_n^\bullet) = \colim f^*\mathcal{K}_n^\bullet$
est une colimite de complexes bornés supérieurement de
$\mathcal{O}'$-modules plats, donc elle est K-plate (d'après les mêmes
lemmes que ci-dessus). L'image inverse de notre suite exacte courte
$$
0 \to f^*\mathcal{M}^\bullet \to
f^*(\colim \mathcal{K}_n^\bullet) \to
f^*\mathcal{K}^\bullet \to 0
$$
est une suite exacte courte de complexes d'après
Modules sur les sites, Lemme \ref{sites-modules-lemma-pullback-ses}.
Ainsi, d'après le Lemme \ref{lemma-K-flat-two-out-of-three-ses}, il suffit
de montrer que $f^*\mathcal{M}^\bullet$ est K-plat. Cela nous ramène au cas
traité dans le paragraphe suivant.

\medskip\noindent
Supposons que $\mathcal{K}^\bullet$ soit acyclique et K-plat, et que ses
termes soient plats. Le Lemme \ref{lemma-K-flat-flat-acyclic} assure alors
que tous les termes de $\tau_{\leq n}\mathcal{K}^\bullet$ sont des
$\mathcal{O}$-modules plats. Choisissons un diagramme comme ci-dessus et
utilisons toutes les propriétés démontrées plus haut pour ce diagramme.
Notons $\mathcal{M}_n^\bullet$ le noyau du morphisme de complexes
$\mathcal{K}_n^\bullet \to \tau_{\leq n}\mathcal{K}^\bullet$
de sorte que nous disposons de suites exactes courtes de complexes
$$
0 \to \mathcal{M}_n^\bullet \to \mathcal{K}_n^\bullet \to
\tau_{\leq n}\mathcal{K}^\bullet \to 0
$$
D'après Modules sur les sites, Lemme \ref{sites-modules-lemma-flat-ses},
les termes du complexe $\mathcal{M}_n^\bullet$ sont plats. Ainsi, dans ce
cas, $\mathcal{M} = \colim \mathcal{M}_n^\bullet$ est une colimite filtrante
de complexes bornés supérieurement de modules plats. Par conséquent,
$f^*\mathcal{M}^\bullet$ est K-plat (par le même argument que ci-dessus),
ce qui conclut.
\end{proof}

\begin{lemma}
\label{lemma-derived-base-change}
Soit $f : (\Sh(\mathcal{C}), \mathcal{O}) \to (\Sh(\mathcal{C}'), \mathcal{O}')$
un morphisme de topos annelés. Il existe un foncteur exact
$$
Lf^* : D(\mathcal{O}') \longrightarrow D(\mathcal{O})
$$
entre catégories triangulées tel que
$Lf^*\mathcal{K}^\bullet = f^*\mathcal{K}^\bullet$ pour tout complexe
K-plat $\mathcal{K}^\bullet$ à termes plats et, en particulier, pour tout
complexe borné supérieurement de $\mathcal{O}'$-modules plats.
\end{lemma}

\begin{proof}
Pour le voir, utilisons la théorie générale développée dans
Catégories dérivées, Section \ref{derived-section-derived-functors}.
Posons $\mathcal{D} = K(\mathcal{O}')$ et $\mathcal{D}' = D(\mathcal{O})$.
Notons $F : \mathcal{D} \to \mathcal{D}'$ le foncteur exact de catégories
triangulées défini par la règle
$F(\mathcal{G}^\bullet) = f^*\mathcal{G}^\bullet$.
Soit $S$ l'ensemble des quasi-isomorphismes de
$\mathcal{D} = K(\mathcal{O}')$. Nous obtenons ainsi une situation comme
celle de Catégories dérivées, Situation
\ref{derived-situation-derived-functor}, si bien que Catégories dérivées,
Définition \ref{derived-definition-right-derived-functor-defined},
s'applique. Nous affirmons que $LF$ est défini partout. Cela résulte de
Catégories dérivées, Lemme \ref{derived-lemma-find-existence-computes},
en prenant pour $\mathcal{P} \subset \Ob(\mathcal{D})$ la collection des
complexes K-plats $\mathcal{K}^\bullet$ à termes plats. En effet, (1)
résulte du Lemme \ref{lemma-K-flat-resolution} et, pour établir (2), nous
devons montrer que, pour tout quasi-isomorphisme
$\mathcal{K}_1^\bullet  \to \mathcal{K}_2^\bullet$ entre objets de
$\mathcal{P}$, le morphisme
$f^*\mathcal{K}_1^\bullet  \to f^*\mathcal{K}_2^\bullet$ est un
quasi-isomorphisme. Écrivons ce morphisme sous la forme
$$
f^{-1}\mathcal{K}_1^\bullet \otimes_{f^{-1}\mathcal{O}'} \mathcal{O}
\longrightarrow
f^{-1}\mathcal{K}_2^\bullet \otimes_{f^{-1}\mathcal{O}'} \mathcal{O}
$$
Le foncteur $f^{-1}$ est exact ; le morphisme
$f^{-1}\mathcal{K}_1^\bullet  \to f^{-1}\mathcal{K}_2^\bullet$ est donc un
quasi-isomorphisme. Les complexes
 $f^{-1}\mathcal{K}_1^\bullet$ et $f^{-1}\mathcal{K}_2^\bullet$
de $f^{-1}\mathcal{O}'$-modules sont K-plats d'après le Lemme
\ref{lemma-pullback-K-flat}, car nous pouvons considérer le morphisme de
topos annelés
$(\Sh(\mathcal{C}), f^{-1}\mathcal{O}') \to
(\Sh(\mathcal{C}'), \mathcal{O}')$. Le Lemme
\ref{lemma-derived-tor-quasi-isomorphism-other-side} assure donc que le
morphisme affiché est un quasi-isomorphisme. Nous obtenons ainsi un
foncteur dérivé
$$
LF :
D(\mathcal{O}') = S^{-1}\mathcal{D}
\longrightarrow
\mathcal{D}' = D(\mathcal{O})
$$
voir Catégories dérivées, Équation (\ref{derived-equation-everywhere}).
Enfin, Catégories dérivées, Lemme
\ref{derived-lemma-find-existence-computes}, assure également que
$LF(\mathcal{K}^\bullet) = F(\mathcal{K}^\bullet) = f^*\mathcal{K}^\bullet$
lorsque $\mathcal{K}^\bullet$ appartient à $\mathcal{P}$.
La démonstration s'achève en observant que les complexes bornés
supérieurement de modules plats appartiennent à $\mathcal{P}$ d'après le
Lemme \ref{lemma-bounded-flat-K-flat}.
\end{proof}

\begin{lemma}
\label{lemma-derived-pullback-composition}
Considérons des morphismes de topos annelés
$f : (\Sh(\mathcal{C}), \mathcal{O}_\mathcal{C}) \to
(\Sh(\mathcal{D}), \mathcal{O}_\mathcal{D})$
et
$g : (\Sh(\mathcal{D}), \mathcal{O}_\mathcal{D}) \to
(\Sh(\mathcal{E}), \mathcal{O}_\mathcal{E})$.
Alors $Lf^* \circ Lg^* = L(g \circ f)^*$ comme foncteurs
$D(\mathcal{O}_\mathcal{E}) \to D(\mathcal{O}_\mathcal{C})$.
\end{lemma}

\begin{proof}
Soit $E$ un objet de $D(\mathcal{O}_\mathcal{E})$. Nous pouvons représenter
$E$ par un complexe K-plat $\mathcal{K}^\bullet$ à termes plats, voir le
Lemme \ref{lemma-K-flat-resolution}. Par construction, $Lg^*E$ se calcule
par $g^*\mathcal{K}^\bullet$, voir le Lemme
\ref{lemma-derived-base-change}. D'après le Lemme
\ref{lemma-pullback-K-flat}, le complexe $g^*\mathcal{K}^\bullet$ est
K-plat à termes plats. Ainsi $Lf^*Lg^*E$ est représenté par
$f^*g^*\mathcal{K}^\bullet$. Comme $L(g \circ f)^*E$ est lui aussi
représenté par
$(g \circ f)^*\mathcal{K}^\bullet = f^*g^*\mathcal{K}^\bullet$
nous concluons.
\end{proof}

\begin{lemma}
\label{lemma-pullback-tensor-product}
Soit $f : (\Sh(\mathcal{C}), \mathcal{O}) \to (\Sh(\mathcal{D}), \mathcal{O}')$
un morphisme de topos annelés. Il existe un isomorphisme bifonctoriel
canonique
$$
Lf^*(
\mathcal{F}^\bullet \otimes_{\mathcal{O}'}^{\mathbf{L}} \mathcal{G}^\bullet
) =
Lf^*\mathcal{F}^\bullet 
\otimes_{\mathcal{O}}^{\mathbf{L}}
Lf^*\mathcal{G}^\bullet 
$$
pour $\mathcal{F}^\bullet, \mathcal{G}^\bullet \in \Ob(D(\mathcal{O}'))$.
\end{lemma}

\begin{proof}
Grâce à notre construction de l'image inverse dérivée dans le Lemme
\ref{lemma-derived-base-change} et à l'existence de résolutions donnée par
le Lemme \ref{lemma-K-flat-resolution}, nous pouvons remplacer
$\mathcal{F}^\bullet$ et $\mathcal{G}^\bullet$ par des complexes de
$\mathcal{O}'$-modules qui sont K-plats et dont les termes sont plats.
Dans ce cas,
$\mathcal{F}^\bullet \otimes_{\mathcal{O}'}^{\mathbf{L}} \mathcal{G}^\bullet$
est simplement le complexe total associé au bicomplexe
$\mathcal{F}^\bullet \otimes_{\mathcal{O}'} \mathcal{G}^\bullet$.
Le complexe
$\text{Tot}(\mathcal{F}^\bullet \otimes_{\mathcal{O}'} \mathcal{G}^\bullet)$
est K-plat à termes plats d'après le Lemme
\ref{lemma-tensor-product-K-flat} et Modules sur les sites, Lemme
\ref{sites-modules-lemma-tensor-flats}. Ainsi, l'isomorphisme du lemme
provient de l'isomorphisme
$$
\text{Tot}(f^*\mathcal{F}^\bullet \otimes_{\mathcal{O}}
f^*\mathcal{G}^\bullet)
\longrightarrow
f^*\text{Tot}(\mathcal{F}^\bullet \otimes_{\mathcal{O}'} \mathcal{G}^\bullet)
$$
dont les composantes sont les isomorphismes
$f^*\mathcal{F}^p \otimes_{\mathcal{O}} f^*\mathcal{G}^q \to
f^*(\mathcal{F}^p \otimes_{\mathcal{O}'} \mathcal{G}^q)$ fournis par
Modules sur les sites, Lemme
\ref{sites-modules-lemma-tensor-product-pullback}.
\end{proof}

\begin{lemma}
\label{lemma-variant-derived-pullback}
Soit $f : (\Sh(\mathcal{C}), \mathcal{O}) \to (\Sh(\mathcal{C}'), \mathcal{O}')$
un morphisme de topos annelés. Il existe un isomorphisme bifonctoriel
canonique
$$
\mathcal{F}^\bullet
\otimes_\mathcal{O}^{\mathbf{L}}
Lf^*\mathcal{G}^\bullet
=
\mathcal{F}^\bullet 
\otimes_{f^{-1}\mathcal{O}'}^{\mathbf{L}}
f^{-1}\mathcal{G}^\bullet 
$$
pour $\mathcal{F}^\bullet$ dans $D(\mathcal{O})$ et
$\mathcal{G}^\bullet$ dans $D(\mathcal{O}')$.
\end{lemma}

\begin{proof}
Soient $\mathcal{F}$ un $\mathcal{O}$-module et $\mathcal{G}$ un
$\mathcal{O}'$-module. Alors
$\mathcal{F} \otimes_{\mathcal{O}} f^*\mathcal{G} =
\mathcal{F} \otimes_{f^{-1}\mathcal{O}'} f^{-1}\mathcal{G}$
car
$f^*\mathcal{G} =
\mathcal{O} \otimes_{f^{-1}\mathcal{O}'} f^{-1}\mathcal{G}$.
Le lemme résulte de cette égalité et des définitions.
\end{proof}











\begin{lemma}
\label{lemma-check-K-flat-stalks}
Soit $(\mathcal{C}, \mathcal{O})$ un site annelé.
Soit $\mathcal{K}^\bullet$ un complexe de $\mathcal{O}$-modules.
\begin{enumerate}
\item Si $\mathcal{K}^\bullet$ est K-plat, alors, pour tout point $p$ du
site $\mathcal{C}$, le complexe de $\mathcal{O}_p$-modules
$\mathcal{K}_p^\bullet$ est K-plat au sens de
Compléments d'algèbre, Définition \ref{more-algebra-definition-K-flat}.
\item Si $\mathcal{C}$ possède suffisamment de points, la réciproque est
vraie.
\end{enumerate}
\end{lemma}

\begin{proof}
Démonstration de (2). Supposons que $\mathcal{C}$ possède suffisamment de
points et que $\mathcal{K}_p^\bullet$ soit K-plat pour tout point $p$ de
$\mathcal{C}$. Alors $\mathcal{K}^\bullet$ est K-plat, car $\otimes$ et les
sommes directes commutent à la prise des fibres, et l'exactitude peut se
vérifier sur les fibres ; voir Modules sur les sites, Lemme
\ref{sites-modules-lemma-check-exactness-stalks}.

\medskip\noindent
Démonstration de (1). Supposons $\mathcal{K}^\bullet$ K-plat. Choisissons
un quasi-isomorphisme $a : \mathcal{L}^\bullet \to \mathcal{K}^\bullet$ tel
que $\mathcal{L}^\bullet$ soit K-plat à termes plats, voir le Lemme
\ref{lemma-K-flat-resolution}. Toute image inverse de
$\mathcal{L}^\bullet$ est K-plate, voir le Lemme
\ref{lemma-pullback-K-flat}. En particulier, la fibre
$\mathcal{L}_p^\bullet$ est un complexe K-plat de
$\mathcal{O}_p$-modules. Le cône $C(a)$ de $a$ est donc un complexe
acyclique K-plat de $\mathcal{O}$-modules (Lemme
\ref{lemma-K-flat-two-out-of-three}), et il suffit de montrer que la fibre
de $C(a)$ est K-plate (d'après Compléments d'algèbre, Lemme
\ref{more-algebra-lemma-K-flat-two-out-of-three}). Nous pouvons donc
supposer que $\mathcal{K}^\bullet$ est K-plat et acyclique.

\medskip\noindent
Supposons $\mathcal{K}^\bullet$ acyclique et K-plat. Avant de poursuivre,
remplaçons le site $\mathcal{C}$ par un autre comme dans Sites, Lemme
\ref{sites-lemma-topos-good-site}, afin que $\mathcal{C}$ admette toutes
les limites finies. Il en résulte que la catégorie des voisinages de $p$
est filtrante (Sites, Lemme \ref{sites-lemma-neighbourhoods-directed}) et
que la colimite définissant la fibre d'un faisceau est filtrante. Soit $M$
un $\mathcal{O}_p$-module de présentation finie. Il suffit de montrer que
$\mathcal{K}^\bullet \otimes_{\mathcal{O}_p} M$ est acyclique ; voir
Compléments d'algèbre, Lemme
\ref{more-algebra-lemma-universally-acyclic-K-flat}. Comme
$\mathcal{O}_p$ est la colimite filtrante des $\mathcal{O}(U)$, où $U$
parcourt les voisinages de $p$, nous pouvons trouver un voisinage $(U, x)$
de $p$ et un $\mathcal{O}(U)$-module de présentation finie $M'$ dont le
changement de base à $\mathcal{O}_p$ est $M$ ; voir Algèbre, Lemme
\ref{algebra-lemma-colimit-category-fp-modules}. D'après le Lemme
\ref{lemma-restriction-K-flat}, nous pouvons remplacer
$\mathcal{C}, \mathcal{O}, \mathcal{K}^\bullet$ par
$\mathcal{C}/U, \mathcal{O}_U, \mathcal{K}^\bullet|_U$. Nous en déduisons
que nous pouvons supposer qu'il existe un $\mathcal{O}$-module
$\mathcal{F}$ tel que $M \cong \mathcal{F}_p$. Puisque
$\mathcal{K}^\bullet$ est K-plat et acyclique, le complexe
$\mathcal{K}^\bullet \otimes_\mathcal{O} \mathcal{F}$ est acyclique (car
il calcule, par définition, le produit tensoriel dérivé). La prise des
fibres étant un foncteur exact, il en résulte que
$\mathcal{K}^\bullet \otimes_{\mathcal{O}_p} M$ est acyclique, comme voulu.
\end{proof}

\begin{lemma}
\label{lemma-pullback-K-flat-points}
Soit $f : (\Sh(\mathcal{C}), \mathcal{O}) \to (\Sh(\mathcal{C}'), \mathcal{O}')$
un morphisme de topos annelés. Si $\mathcal{C}$ possède suffisamment de
points, l'image inverse d'un complexe K-plat de $\mathcal{O}'$-modules est
un complexe K-plat de $\mathcal{O}$-modules.
\end{lemma}

\begin{proof}
Cela résulte du Lemme \ref{lemma-check-K-flat-stalks}, de Modules sur les
sites, Lemme \ref{sites-modules-lemma-pullback-stalk}, et de Compléments
d'algèbre, Lemme \ref{more-algebra-lemma-base-change-K-flat}.
\end{proof}

\begin{lemma}
\label{lemma-tensor-pull-compatibility}
Soit $f : (\Sh(\mathcal{C}), \mathcal{O}_\mathcal{C}) \to
(\Sh(\mathcal{D}), \mathcal{O}_\mathcal{D})$ un morphisme de topos annelés.
Soient $\mathcal{K}^\bullet$ et $\mathcal{M}^\bullet$ des complexes de
$\mathcal{O}_\mathcal{D}$-modules. Le diagramme
$$
\xymatrix{
Lf^*(\mathcal{K}^\bullet
\otimes_{\mathcal{O}_\mathcal{D}}^\mathbf{L}
\mathcal{M}^\bullet) \ar[r] \ar[d] &
Lf^*\text{Tot}(\mathcal{K}^\bullet
\otimes_{\mathcal{O}_\mathcal{D}}
\mathcal{M}^\bullet) \ar[d] \\
Lf^*\mathcal{K}^\bullet \otimes_{\mathcal{O}_\mathcal{C}}^\mathbf{L}
Lf^*\mathcal{M}^\bullet \ar[d] &
f^*\text{Tot}(\mathcal{K}^\bullet
\otimes_{\mathcal{O}_\mathcal{D}}
\mathcal{M}^\bullet) \ar[d] \\
f^*\mathcal{K}^\bullet \otimes_{\mathcal{O}_\mathcal{C}}^\mathbf{L}
f^*\mathcal{M}^\bullet \ar[r] &
\text{Tot}(f^*\mathcal{K}^\bullet \otimes_{\mathcal{O}_\mathcal{C}}
f^*\mathcal{M}^\bullet)
}
$$
est commutatif.
\end{lemma}

\begin{proof}
Nous utiliserons l'existence de résolutions K-plates à termes plats (Lemme
\ref{lemma-K-flat-resolution}), le fait que l'image inverse dérivée se
calcule à l'aide de tels complexes (Lemme
\ref{lemma-derived-base-change}), ainsi que la préservation de ces
propriétés par image inverse (Lemme \ref{lemma-pullback-K-flat}). Si nous
choisissons de telles résolutions
$\mathcal{P}^\bullet \to \mathcal{K}^\bullet$ et
$\mathcal{Q}^\bullet \to \mathcal{M}^\bullet$, nous voyons alors que
$$
\xymatrix{
Lf^*\text{Tot}(\mathcal{P}^\bullet
\otimes_{\mathcal{O}_\mathcal{D}}
\mathcal{Q}^\bullet) \ar[r] \ar[d] &
Lf^*\text{Tot}(\mathcal{K}^\bullet
\otimes_{\mathcal{O}_\mathcal{D}}
\mathcal{M}^\bullet) \ar[d] \\
f^*\text{Tot}(\mathcal{P}^\bullet
\otimes_{\mathcal{O}_\mathcal{D}}
\mathcal{Q}^\bullet) \ar[d] \ar[r] &
f^*\text{Tot}(\mathcal{K}^\bullet
\otimes_{\mathcal{O}_\mathcal{D}}
\mathcal{M}^\bullet) \ar[d] \\
\text{Tot}(f^*\mathcal{P}^\bullet \otimes_{\mathcal{O}_\mathcal{C}}
f^*\mathcal{Q}^\bullet) \ar[r] &
\text{Tot}(f^*\mathcal{K}^\bullet \otimes_{\mathcal{O}_\mathcal{C}}
f^*\mathcal{M}^\bullet)
}
$$
est commutatif. Or, le côté gauche de ce diagramme est maintenant le côté
gauche du diagramme de l'énoncé, par notre choix de $\mathcal{P}^\bullet$
et $\mathcal{Q}^\bullet$ et par le Lemme
\ref{lemma-tensor-product-K-flat}.
\end{proof}









\section{Cohomologie des complexes non bornés}
\label{section-unbounded}

\noindent
Soit $(\mathcal{C}, \mathcal{O})$ un site annelé. La catégorie
$\textit{Mod}(\mathcal{O})$ est une catégorie abélienne de Grothendieck :
elle admet toutes les colimites, les colimites filtrantes y sont exactes,
et elle possède un générateur, à savoir
$$
\bigoplus\nolimits_{U \in \Ob(\mathcal{C})} j_{U!}\mathcal{O}_U,
$$
voir Modules sur les sites, Section \ref{sites-modules-section-kernels}, et
Lemmes \ref{sites-modules-lemma-j-shriek-flat} et
\ref{sites-modules-lemma-module-quotient-flat}. D'après Injectifs, Théorème
\ref{injectives-theorem-K-injective-embedding-grothendieck}, pour tout
complexe $\mathcal{F}^\bullet$ de $\mathcal{O}$-modules, il existe un
quasi-isomorphisme injectif
$\mathcal{F}^\bullet \to \mathcal{I}^\bullet$ vers un complexe K-injectif
de $\mathcal{O}$-modules ; de plus, cette inclusion peut être choisie
fonctoriellement en $\mathcal{F}^\bullet$. Il résulte de Catégories
dérivées, Lemme \ref{derived-lemma-enough-K-injectives-implies}, que
\begin{enumerate}
\item tout foncteur exact $F : K(\textit{Mod}(\mathcal{O})) \to \mathcal{D}$
à valeurs dans une catégorie triangulée $\mathcal{D}$ admet un foncteur
dérivé à droite
$RF : D(\mathcal{O}) \to \mathcal{D}$,
\item pour tout foncteur additif
$F : \textit{Mod}(\mathcal{O}) \to \mathcal{A}$
à valeurs dans une catégorie abélienne $\mathcal{A}$, nous considérons le
foncteur exact $F : K(\textit{Mod}(\mathcal{O})) \to K(\mathcal{A})$
induit par $F$ et obtenons un foncteur dérivé à droite
$RF : D(\mathcal{O}) \to D(\mathcal{A})$.
\end{enumerate}
Par construction, $RF(\mathcal{F}^\bullet) = F(\mathcal{I}^\bullet)$, où
$\mathcal{F}^\bullet \to \mathcal{I}^\bullet$ est comme ci-dessus.

\medskip\noindent
Voici quelques exemples de ce qui précède :
\begin{enumerate}
\item Le foncteur $\Gamma(\mathcal{C}, -) : \textit{Mod}(\mathcal{O}) \to
\text{Mod}_{\Gamma(\mathcal{C}, \mathcal{O})}$ donne naissance à
$$
R\Gamma(\mathcal{C}, -) :
D(\mathcal{O})
\longrightarrow
D(\Gamma(\mathcal{C}, \mathcal{O}))
$$
Nous noterons
$H^i(\mathcal{C}, K) = H^i(R\Gamma(\mathcal{C}, K))$ la cohomologie.
\item Pour un objet $U$ de $\mathcal{C}$, considérons le foncteur
$\Gamma(U, -) : \textit{Mod}(\mathcal{O}) \to
\text{Mod}_{\Gamma(U, \mathcal{O})}$. Il donne naissance à
$$
R\Gamma(U, -) : D(\mathcal{O}) \to D(\Gamma(U, \mathcal{O}))
$$
Nous noterons $H^i(U, K) = H^i(R\Gamma(U, K))$ la cohomologie.
\item Pour un morphisme de topos annelés
$f : (\Sh(\mathcal{C}), \mathcal{O}) \to (\Sh(\mathcal{D}), \mathcal{O}')$
considérons le foncteur
$f_* : \textit{Mod}(\mathcal{O}) \to \textit{Mod}(\mathcal{O}')$
qui donne naissance à l'image directe dérivée totale
$$
Rf_* : D(\mathcal{O}) \longrightarrow D(\mathcal{O}')
$$
sur les catégories dérivées non bornées.
\end{enumerate}

\begin{lemma}
\label{lemma-adjoint}
Soit $f : (\Sh(\mathcal{C}), \mathcal{O}) \to (\Sh(\mathcal{D}), \mathcal{O}')$
un morphisme de topos annelés. Le foncteur $Rf_*$ défini ci-dessus et le
foncteur $Lf^*$ défini dans le Lemme \ref{lemma-derived-base-change} sont
adjoints :
$$
\Hom_{D(\mathcal{O})}(Lf^*\mathcal{G}^\bullet, \mathcal{F}^\bullet)
=
\Hom_{D(\mathcal{O}')}(\mathcal{G}^\bullet, Rf_*\mathcal{F}^\bullet)
$$
bifonctoriellement en $\mathcal{F}^\bullet \in \Ob(D(\mathcal{O}))$ et
$\mathcal{G}^\bullet \in \Ob(D(\mathcal{O}'))$.
\end{lemma}

\begin{proof}
Cela résulte formellement de l'existence de $Rf_*$ et de $Lf^*$ ; voir
Catégories dérivées, Lemme
\ref{derived-lemma-derived-adjoint-functors}.
\end{proof}

\begin{lemma}
\label{lemma-derived-pushforward-composition}
Soient
$f : (\Sh(\mathcal{C}), \mathcal{O}_\mathcal{C}) \to
(\Sh(\mathcal{D}), \mathcal{O}_\mathcal{D})$
et
$g : (\Sh(\mathcal{D}), \mathcal{O}_\mathcal{D}) \to
(\Sh(\mathcal{E}), \mathcal{O}_\mathcal{E})$
des morphismes de topos annelés. Alors
$Rg_* \circ Rf_* = R(g \circ f)_*$ comme foncteurs
$D(\mathcal{O}_\mathcal{C}) \to D(\mathcal{O}_\mathcal{E})$.
\end{lemma}

\begin{proof}
D'après le Lemme \ref{lemma-adjoint}, $Rg_* \circ Rf_*$ est adjoint à
$Lf^* \circ Lg^*$. Or $Lf^* \circ Lg^* = L(g \circ f)^*$ d'après le
Lemme \ref{lemma-derived-pullback-composition} ; l'unicité des foncteurs
adjoints donne donc $Rg_* \circ Rf_* = R(g \circ f)_*$.
\end{proof}

\begin{remark}
\label{remark-base-change}
La construction des foncteurs dérivés non bornés $Lf^*$ et $Rf_*$ permet
de construire le morphisme de changement de base en toute généralité.
Supposons en effet que
$$
\xymatrix{
(\Sh(\mathcal{C}'), \mathcal{O}_{\mathcal{C}'})
\ar[r]_{g'} \ar[d]_{f'} &
(\Sh(\mathcal{C}), \mathcal{O}_\mathcal{C}) \ar[d]^f \\
(\Sh(\mathcal{D}'), \mathcal{O}_{\mathcal{D}'})
\ar[r]^g &
(\Sh(\mathcal{D}), \mathcal{O}_\mathcal{D})
}
$$
soit un diagramme commutatif de topos annelés. Soit $K$ un objet de
$D(\mathcal{O}_\mathcal{C})$. Il existe alors un morphisme canonique de
changement de base
$$
Lg^*Rf_*K \longrightarrow R(f')_*L(g')^*K
$$
dans $D(\mathcal{O}_{\mathcal{D}'})$. Ce morphisme est en effet adjoint à
un morphisme $L(f')^*Lg^*Rf_*K \to L(g')^*K$. Puisque
$L(f')^* \circ Lg^* = L(g')^* \circ Lf^*$, cela revient à donner un
morphisme $L(g')^*Lf^*Rf_*K \to L(g')^*K$, que l'on peut prendre égal à
l'image par $L(g')^*$ du morphisme d'adjonction $Lf^*Rf_*K \to K$.
\end{remark}

\begin{remark}
\label{remark-compose-base-change}
Considérons un diagramme commutatif
$$
\xymatrix{
(\Sh(\mathcal{B}'), \mathcal{O}_{\mathcal{B}'})
\ar[r]_k \ar[d]_{f'} &
(\Sh(\mathcal{B}), \mathcal{O}_\mathcal{B}) \ar[d]^f \\
(\Sh(\mathcal{C}'), \mathcal{O}_{\mathcal{C}'})
\ar[r]^l \ar[d]_{g'} &
(\Sh(\mathcal{C}), \mathcal{O}_\mathcal{C}) \ar[d]^g \\
(\Sh(\mathcal{D}'), \mathcal{O}_{\mathcal{D}'})
\ar[r]^m &
(\Sh(\mathcal{D}), \mathcal{O}_\mathcal{D}) \\
}
$$
de topos annelés. Les morphismes de changement de base de la Remarque
\ref{remark-base-change}, pour les deux carrés, se composent alors pour
donner le morphisme de changement de base du rectangle extérieur. Plus
précisément, la composée
\begin{align*}
Lm^* \circ R(g \circ f)_*
& =
Lm^* \circ Rg_* \circ Rf_* \\
& \to Rg'_* \circ Ll^* \circ Rf_* \\
& \to Rg'_* \circ Rf'_* \circ Lk^* \\
& = R(g' \circ f')_* \circ Lk^*
\end{align*}
est le morphisme de changement de base du rectangle. Nous omettons la
vérification.
\end{remark}

\begin{remark}
\label{remark-compose-base-change-horizontal}
Considérons un diagramme commutatif
$$
\xymatrix{
(\Sh(\mathcal{C}''), \mathcal{O}_{\mathcal{C}''})
\ar[r]_{g'} \ar[d]_{f''} &
(\Sh(\mathcal{C}'), \mathcal{O}_{\mathcal{C}'})
\ar[r]_g \ar[d]_{f'} &
(\Sh(\mathcal{C}), \mathcal{O}_\mathcal{C}) \ar[d]^f \\
(\Sh(\mathcal{D}''), \mathcal{O}_{\mathcal{D}''})
\ar[r]^{h'} &
(\Sh(\mathcal{D}'), \mathcal{O}_{\mathcal{D}'})
\ar[r]^h &
(\Sh(\mathcal{D}), \mathcal{O}_\mathcal{D})
}
$$
de topos annelés. Les morphismes de changement de base de la Remarque
\ref{remark-base-change}, pour les deux carrés, se composent alors pour
donner le morphisme de changement de base du rectangle extérieur. Plus
précisément, la composée
\begin{align*}
L(h \circ h')^* \circ Rf_*
& =
L(h')^* \circ Lh^* \circ Rf_* \\
& \to L(h')^* \circ Rf'_* \circ Lg^* \\
& \to Rf''_* \circ L(g')^* \circ Lg^* \\
& = Rf''_* \circ L(g \circ g')^*
\end{align*}
est le morphisme de changement de base du rectangle. Nous omettons la
vérification.
\end{remark}



\begin{lemma}
\label{lemma-adjoints-push-pull-compatibility}
Soit $f : (\Sh(\mathcal{C}), \mathcal{O}_\mathcal{C}) \to
(\Sh(\mathcal{D}), \mathcal{O}_\mathcal{D})$ un morphisme de topos annelés.
Soit $\mathcal{K}^\bullet$ un complexe de
$\mathcal{O}_\mathcal{C}$-modules. Le diagramme
$$
\xymatrix{
Lf^*f_*\mathcal{K}^\bullet \ar[r] \ar[d] &
f^*f_*\mathcal{K}^\bullet \ar[d] \\
Lf^*Rf_*\mathcal{K}^\bullet \ar[r] &
\mathcal{K}^\bullet
}
$$
provenant de $Lf^* \to f^*$ sur les complexes, de $f_* \to Rf_*$ sur les
complexes et de l'adjonction $Lf^* \circ Rf_* \to \text{id}$ est
commutatif dans $D(\mathcal{O}_\mathcal{C})$.
\end{lemma}

\begin{proof}
Nous utiliserons l'existence de résolutions K-plates et de résolutions
K-injectives ; voir les Lemmes \ref{lemma-K-flat-resolution},
\ref{lemma-derived-base-change} et \ref{lemma-pullback-K-flat}, ainsi que
la discussion ci-dessus. Choisissons un quasi-isomorphisme
$\mathcal{K}^\bullet \to \mathcal{I}^\bullet$, où $\mathcal{I}^\bullet$
est K-injectif comme complexe de $\mathcal{O}_\mathcal{C}$-modules.
Choisissons un quasi-isomorphisme
$\mathcal{Q}^\bullet \to f_*\mathcal{I}^\bullet$, où
$\mathcal{Q}^\bullet$ est un complexe K-plat à termes plats de
$\mathcal{O}_\mathcal{D}$-modules. Nous pouvons choisir un complexe K-plat de
$\mathcal{O}_\mathcal{D}$-modules, noté $\mathcal{P}^\bullet$, à termes plats,
ainsi qu'un diagramme de morphismes de complexes
$$
\xymatrix{
\mathcal{P}^\bullet \ar[r] \ar[d] &
f_*\mathcal{K}^\bullet \ar[d] \\
\mathcal{Q}^\bullet \ar[r] & f_*\mathcal{I}^\bullet
}
$$
commutatif à homotopie près, dans lequel la flèche horizontale supérieure
est un quasi-isomorphisme. En effet, nous pouvons d'abord choisir un tel
diagramme pour un certain complexe $\mathcal{P}^\bullet$, car les
quasi-isomorphismes forment un système multiplicatif dans la catégorie
homotopique des complexes, puis choisir une résolution de
$\mathcal{P}^\bullet$ par un complexe K-plat à termes plats. En prenant les
images inverses, nous obtenons un diagramme de morphismes de complexes
$$
\xymatrix{
f^*\mathcal{P}^\bullet \ar[r] \ar[d] &
f^*f_*\mathcal{K}^\bullet \ar[d] \ar[r] &
\mathcal{K}^\bullet \ar[d] \\
f^*\mathcal{Q}^\bullet \ar[r] &
f^*f_*\mathcal{I}^\bullet \ar[r] &
\mathcal{I}^\bullet
}
$$
commutatif à homotopie près. Le rectangle extérieur établit l'assertion
du lemme.
\end{proof}



\begin{remark}
\label{remark-cup-product}
Soit $f : (\Sh(\mathcal{C}), \mathcal{O}_\mathcal{C}) \to
(\Sh(\mathcal{D}), \mathcal{O}_\mathcal{D})$ un morphisme de topos
annelés. L'adjonction de $Lf^*$ et $Rf_*$ permet de construire un cup-produit
relatif
$$
Rf_*K \otimes_{\mathcal{O}_\mathcal{D}}^\mathbf{L} Rf_*L
\longrightarrow
Rf_*(K \otimes_{\mathcal{O}_\mathcal{C}}^\mathbf{L} L)
$$
dans $D(\mathcal{O}_\mathcal{D})$, pour tous $K, L$ dans
$D(\mathcal{O}_\mathcal{C})$. Ce morphisme est en effet adjoint à un
morphisme
$Lf^*(Rf_*K \otimes_{\mathcal{O}_\mathcal{D}}^\mathbf{L} Rf_*L) \to
K \otimes_{\mathcal{O}_\mathcal{C}}^\mathbf{L} L$, que l'on peut prendre
égal à la composée de l'isomorphisme
$Lf^*(Rf_*K \otimes_{\mathcal{O}_\mathcal{D}}^\mathbf{L} Rf_*L) =
Lf^*Rf_*K \otimes_{\mathcal{O}_\mathcal{C}}^\mathbf{L} Lf^*Rf_*L$
(Lemme \ref{lemma-pullback-tensor-product}) avec le morphisme
$Lf^*Rf_*K \otimes_{\mathcal{O}_\mathcal{C}}^\mathbf{L} Lf^*Rf_*L
\to K \otimes_{\mathcal{O}_\mathcal{C}}^\mathbf{L} L$
provenant de la co-unité $Lf^* \circ Rf_* \to \text{id}$.
\end{remark}

\begin{lemma}
\label{lemma-torsion}
Soit $\mathcal{C}$ un site. Notons
$\mathcal{A} \subset \textit{Ab}(\mathcal{C})$ la sous-catégorie de Serre
constituée des faisceaux abéliens de torsion. Alors le foncteur
$D(\mathcal{A}) \to D_\mathcal{A}(\mathcal{C})$ est une équivalence.
\end{lemma}

\begin{proof}
Une observation essentielle est qu'un faisceau abélien injectif
$\mathcal{I}$ est divisible. En effet, si $s \in \mathcal{I}(U)$ est une
section locale, interprétons $s$ comme un morphisme
$s : j_{U!}\mathbf{Z} \to \mathcal{I}$ et appliquons la propriété
définissant un objet injectif au monomorphisme de faisceaux
$n : j_{U!}\mathbf{Z} \to j_{U!}\mathbf{Z}$. Il existe alors
$s' \in \mathcal{I}(U)$ tel que $ns' = s$.

\medskip\noindent
Pour un faisceau $\mathcal{F}$, notons $\mathcal{F}_{tor}$ son sous-faisceau
de torsion. Nous affirmons que, si $\mathcal{I}^\bullet$ est un complexe
de faisceaux abéliens injectifs dont les faisceaux de cohomologie sont de
torsion, alors
$$
\mathcal{I}^\bullet_{tor} \to \mathcal{I}^\bullet
$$
est un quasi-isomorphisme. En effet, la platitude de $\mathbf{Q}$ sur
$\mathbf{Z}$ donne
$$
H^p(\mathcal{I}^\bullet) \otimes_\mathbf{Z} \mathbf{Q} =
H^p(\mathcal{I}^\bullet \otimes_\mathbf{Z} \mathbf{Q})
$$
qui est nul car les faisceaux de cohomologie sont de torsion. Par
divisibilité (établie ci-dessus), le morphisme
$\mathcal{I}^\bullet \to \mathcal{I}^\bullet \otimes_\mathbf{Z} \mathbf{Q}$
est surjectif de noyau $\mathcal{I}^\bullet_{tor}$. L'affirmation résulte
alors de la suite exacte longue des faisceaux de cohomologie associée à la
suite exacte courte ainsi obtenue.

\medskip\noindent
Pour démontrer le lemme, construisons un adjoint à droite
$T : D(\mathcal{C}) \to D(\mathcal{A})$. Étant donné $K$ dans
$D(\mathcal{C})$, nous pouvons représenter $K$ par un complexe K-injectif
$\mathcal{I}^\bullet$ dont les faisceaux de cohomologie sont injectifs ;
voir Injectifs, Théorème
\ref{injectives-theorem-K-injective-embedding-grothendieck}. Posons alors
$T(K) = \mathcal{I}^\bullet_{tor}$ ; autrement dit, $T$ est le foncteur
dérivé à droite du foncteur de torsion. Le foncteur $T$ est adjoint à
droite de $i : D(\mathcal{A}) \to D_\mathcal{A}(\mathcal{C})$. Cela résulte
immédiatement de l'observation suivante : si $\mathcal{F}^\bullet$ est un
complexe de faisceaux de torsion, alors
$$
\Hom_{K(\mathcal{A})}(\mathcal{F}^\bullet, I^\bullet_{tor}) =
\Hom_{K(\textit{Ab}(\mathcal{C}))}(\mathcal{F}^\bullet, I^\bullet)
$$
en particulier, $\mathcal{I}^\bullet_{tor}$ est un complexe K-injectif de
$\mathcal{A}$. Nous omettons quelques détails ; cela résulte aussi, en cas
de doute, du résultat plus général Catégories dérivées, Lemme
\ref{derived-lemma-derived-adjoint-functors}. Notre affirmation ci-dessus
donne $L = T(i(L))$ pour $L$ dans $D(\mathcal{A})$, et $i(T(K)) = K$ si
$K$ appartient à $D_\mathcal{A}(\mathcal{C})$. Le résultat découle alors
de Catégories, Lemme \ref{categories-lemma-adjoint-fully-faithful}.
\end{proof}




\section{Quelques propriétés des complexes K-injectifs}
\label{section-properties-K-injective}

\noindent
Soit $(\mathcal{C}, \mathcal{O})$ un site annelé et soit $U$ un objet de
$\mathcal{C}$. Notons
$j : (\Sh(\mathcal{C}/U), \mathcal{O}_U) \to (\Sh(\mathcal{C}), \mathcal{O})$
le morphisme de localisation correspondant. Le foncteur image inverse
$j^*$ est exact, puisqu'il n'est autre que le foncteur de restriction.
Ainsi, sur tout complexe, l'image inverse dérivée $Lj^*$ se calcule par
simple restriction du complexe. Nous notons souvent simplement le foncteur
correspondant
$$
D(\mathcal{O}) \to D(\mathcal{O}_U),
\quad
E \mapsto j^*E = E|_U
$$
De même, le prolongement par zéro
$j_! : \textit{Mod}(\mathcal{O}_U) \to \textit{Mod}(\mathcal{O})$ (voir
Modules sur les sites, Définition
\ref{sites-modules-definition-localize-ringed-site}) est un foncteur exact
(Modules sur les sites, Lemme
\ref{sites-modules-lemma-extension-by-zero-exact}). Il induit donc un
foncteur
$$
j_! : D(\mathcal{O}_U) \to D(\mathcal{O}), \quad
F \mapsto j_!F
$$
en appliquant simplement $j_!$ à tout complexe représentant l'objet $F$.

\begin{lemma}
\label{lemma-restrict-K-injective-to-open}
Soit $(\mathcal{C}, \mathcal{O})$ un site annelé et soit $U$ un objet de
$\mathcal{C}$. La restriction d'un complexe K-injectif de
$\mathcal{O}$-modules à $\mathcal{C}/U$ est un complexe K-injectif de
$\mathcal{O}_U$-modules.
\end{lemma}

\begin{proof}
Cela résulte immédiatement de Catégories dérivées, Lemme
\ref{derived-lemma-adjoint-preserve-K-injectives}, et du fait que le
foncteur de restriction admet l'adjoint à gauche exact $j_!$. Voir la
discussion ci-dessus.
\end{proof}

\begin{lemma}
\label{lemma-unbounded-cohomology-of-open}
Soit $(\mathcal{C}, \mathcal{O})$ un site annelé. Soit
$U \in \Ob(\mathcal{C})$. Pour $K$ dans $D(\mathcal{O})$, nous avons
$H^p(U, K) = H^p(\mathcal{C}/U, K|_{\mathcal{C}/U})$.
\end{lemma}

\begin{proof}
Soit $\mathcal{I}^\bullet$ un complexe K-injectif de $\mathcal{O}$-modules
représentant $K$. Alors
$$
H^q(U, K) = H^q(\Gamma(U, \mathcal{I}^\bullet)) =
H^q(\Gamma(\mathcal{C}/U, \mathcal{I}^\bullet|_{\mathcal{C}/U}))
$$
par construction de la cohomologie. D'après le Lemme
\ref{lemma-restrict-K-injective-to-open}, le complexe
$\mathcal{I}^\bullet|_{\mathcal{C}/U}$ est un complexe K-injectif
représentant $K|_{\mathcal{C}/U}$, d'où le lemme.
\end{proof}

\begin{lemma}
\label{lemma-sheafification-cohomology}
Soit $(\mathcal{C}, \mathcal{O})$ un site annelé. Soit $K$ un objet de
$D(\mathcal{O})$. Le faisceau associé au préfaisceau
$$
U \mapsto H^q(U, K) = H^q(\mathcal{C}/U, K|_{\mathcal{C}/U})
$$
est le $q$-ième faisceau de cohomologie $H^q(K)$ de $K$.
\end{lemma}

\begin{proof}
L'égalité $H^q(U, K) = H^q(\mathcal{C}/U, K|_{\mathcal{C}/U})$ résulte du
Lemme \ref{lemma-unbounded-cohomology-of-open}. Choisissons un complexe
K-injectif $\mathcal{I}^\bullet$ représentant $K$. Alors
$$
H^q(U, K) =
\frac{\Ker(\mathcal{I}^q(U) \to \mathcal{I}^{q + 1}(U))}
{\Im(\mathcal{I}^{q - 1}(U) \to \mathcal{I}^q(U))}.
$$
par notre construction de la cohomologie. Comme
$H^q(K) = \Ker(\mathcal{I}^q \to \mathcal{I}^{q + 1})/
\Im(\mathcal{I}^{q - 1} \to \mathcal{I}^q)$, le résultat est clair.
\end{proof}

\begin{lemma}
\label{lemma-restrict-direct-image-open}
Soit $f : (\mathcal{C}, \mathcal{O}_\mathcal{C}) \to
(\mathcal{D}, \mathcal{O}_\mathcal{D})$ un morphisme de sites annelés
correspondant au foncteur continu $u : \mathcal{D} \to \mathcal{C}$.
Pour $V \in \mathcal{D}$, posons $U = u(V)$ et notons
$g : (\mathcal{C}/U, \mathcal{O}_U) \to (\mathcal{D}/V, \mathcal{O}_V)$
le morphisme induit de sites annelés (Modules sur les sites, Lemme
\ref{sites-modules-lemma-localize-morphism-ringed-sites}).
Alors $(Rf_*E)|_{\mathcal{D}/V} = Rg_*(E|_{\mathcal{C}/U})$ pour $E$ dans
$D(\mathcal{O}_\mathcal{C})$.
\end{lemma}

\begin{proof}
Représentons $E$ par un complexe K-injectif $\mathcal{I}^\bullet$ de
$\mathcal{O}_\mathcal{C}$-modules. Alors
$Rf_*(E) = f_*\mathcal{I}^\bullet$
et $Rg_*(E|_{\mathcal{C}/U}) = g_*(\mathcal{I}^\bullet|_{\mathcal{C}/U})$
d'après le Lemme \ref{lemma-restrict-K-injective-to-open}. Comme il est
clair que
$(f_*\mathcal{F})|_{\mathcal{D}/V} = g_*(\mathcal{F}|_{\mathcal{C}/U})$
pour tout faisceau $\mathcal{F}$ sur $\mathcal{C}$ (voir Modules sur les
sites, Lemme \ref{sites-modules-lemma-localize-morphism-ringed-sites}, ou
le résultat plus élémentaire Sites, Lemme
\ref{sites-lemma-localize-morphism}), le résultat s'ensuit.
\end{proof}

\begin{lemma}
\label{lemma-Leray-unbounded}
Soit $f : (\mathcal{C}, \mathcal{O}_\mathcal{C}) \to
(\mathcal{D}, \mathcal{O}_\mathcal{D})$ un morphisme de sites annelés
correspondant au foncteur continu $u : \mathcal{D} \to \mathcal{C}$.
Alors $R\Gamma(\mathcal{D}, -) \circ Rf_* = R\Gamma(\mathcal{C}, -)$ comme
foncteurs $D(\mathcal{O}_\mathcal{C}) \to D(\Gamma(\mathcal{O}_\mathcal{D}))$.
Plus généralement, pour $V \in \mathcal{D}$ et $U = u(V)$, nous avons
$R\Gamma(U, -) = R\Gamma(V, -) \circ Rf_*$.
\end{lemma}

\begin{proof}
Considérons le topos ponctuel $pt$ muni de $\mathcal{O}_{pt}$ défini par
l'anneau $\Gamma(\mathcal{O}_\mathcal{D})$. Il existe un morphisme canonique
$(\mathcal{D}, \mathcal{O}_\mathcal{D}) \to (pt, \mathcal{O}_{pt})$
de topos annelés induisant l'identification sur les sections globales des
faisceaux structuraux. Alors
$D(\mathcal{O}_{pt}) = D(\Gamma(\mathcal{O}_\mathcal{D}))$.
L'assertion
$R\Gamma(\mathcal{D}, -) \circ Rf_* = R\Gamma(\mathcal{C}, -)$
résulte du Lemme \ref{lemma-derived-pushforward-composition} appliqué à
$$
(\mathcal{C}, \mathcal{O}_\mathcal{C}) \to
(\mathcal{D}, \mathcal{O}_\mathcal{D}) \to (pt, \mathcal{O}_{pt})
$$
La seconde assertion, plus générale, résulte de la première après
application du Lemme \ref{lemma-restrict-direct-image-open}.
\end{proof}

\begin{lemma}
\label{lemma-unbounded-describe-higher-direct-images}
Soit $f : (\mathcal{C}, \mathcal{O}_\mathcal{C}) \to
(\mathcal{D}, \mathcal{O}_\mathcal{D})$ un morphisme de sites annelés
correspondant au foncteur continu $u : \mathcal{D} \to \mathcal{C}$.
Soit $K$ dans $D(\mathcal{O}_\mathcal{C})$. Alors $H^i(Rf_*K)$ est le
faisceau associé au préfaisceau
$$
V \mapsto H^i(u(V), K) = H^i(V, Rf_*K)
$$
\end{lemma}

\begin{proof}
L'égalité $H^i(u(V), K) = H^i(V, Rf_*K)$ s'obtient en prenant la
cohomologie dans la seconde assertion du Lemme
\ref{lemma-Leray-unbounded}. L'assertion concernant le faisceau associé
résulte alors du Lemme \ref{lemma-sheafification-cohomology}.
\end{proof}

\begin{lemma}
\label{lemma-modules-abelian-unbounded}
Soit $(\mathcal{C}, \mathcal{O}_\mathcal{C})$ un site annelé. Soit $K$ un
objet de $D(\mathcal{O}_\mathcal{C})$ et notons $K_{ab}$ son image dans
$D(\underline{\mathbf{Z}}_\mathcal{C})$.
\begin{enumerate}
\item Il existe un morphisme canonique
$R\Gamma(\mathcal{C}, K) \to R\Gamma(\mathcal{C}, K_{ab})$
qui est un isomorphisme dans $D(\textit{Ab})$.
\item Pour tout $U \in \mathcal{C}$, il existe un morphisme canonique
$R\Gamma(U, K) \to R\Gamma(U, K_{ab})$
qui est un isomorphisme dans $D(\textit{Ab})$.
\item Soit $f : (\mathcal{C}, \mathcal{O}_\mathcal{C}) \to
(\mathcal{D}, \mathcal{O}_\mathcal{D})$ un morphisme de sites annelés.
Il existe un morphisme canonique $Rf_*K \to Rf_*(K_{ab})$ qui est un
isomorphisme dans $D(\underline{\mathbf{Z}}_\mathcal{D})$.
\end{enumerate}
\end{lemma}

\begin{proof}
Le morphisme se construit comme suit. Choisissons un complexe K-injectif
$\mathcal{I}^\bullet$ représentant $K$. Choisissons un quasi-isomorphisme
$\mathcal{I}^\bullet \to \mathcal{J}^\bullet$, où $\mathcal{J}^\bullet$
est un complexe K-injectif de groupes abéliens. Le morphisme de (1) est
alors donné par
$\Gamma(\mathcal{C}, \mathcal{I}^\bullet) \to
\Gamma(\mathcal{C}, \mathcal{J}^\bullet)$
(2) est donné par
$\Gamma(U, \mathcal{I}^\bullet) \to \Gamma(U, \mathcal{J}^\bullet)$
et le morphisme de (3) est donné par
$f_*\mathcal{I}^\bullet \to f_*\mathcal{J}^\bullet$.
Pour montrer que ces morphismes sont des isomorphismes, il suffit de
prouver qu'ils induisent des isomorphismes sur les groupes de cohomologie
et les faisceaux de cohomologie. D'après les Lemmes
\ref{lemma-unbounded-cohomology-of-open} et
\ref{lemma-unbounded-describe-higher-direct-images}, il suffit de montrer
que le morphisme
$$
H^0(\mathcal{C}, K) \longrightarrow H^0(\mathcal{C}, K_{ab})
$$
est un isomorphisme. Observons que
$$
H^0(\mathcal{C}, K) =
\Hom_{D(\mathcal{O}_\mathcal{C})}(\mathcal{O}_\mathcal{C}, K)
$$
et de même pour l'autre groupe. Choisissons un complexe quelconque
$\mathcal{K}^\bullet$ de $\mathcal{O}_\mathcal{C}$-modules représentant
$K$. Par construction de la catégorie dérivée comme localisation, nous
avons
$$
\Hom_{D(\mathcal{O}_\mathcal{C})}(\mathcal{O}_\mathcal{C}, K) =
\colim_{s : \mathcal{F}^\bullet \to \mathcal{O}_\mathcal{C}}
\Hom_{K(\mathcal{O}_\mathcal{C})}(\mathcal{F}^\bullet, \mathcal{K}^\bullet)
$$
où la colimite porte sur les quasi-isomorphismes $s$ de complexes de
$\mathcal{O}_\mathcal{C}$-modules. De même, nous avons
$$
\Hom_{D(\underline{\mathbf{Z}}_\mathcal{C})}
(\underline{\mathbf{Z}}_\mathcal{C}, K) =
\colim_{s : \mathcal{G}^\bullet \to \underline{\mathbf{Z}}_\mathcal{C}}
\Hom_{K(\underline{\mathbf{Z}}_\mathcal{C})}
(\mathcal{G}^\bullet, \mathcal{K}^\bullet)
$$
Observons ensuite que les quasi-isomorphismes
$s : \mathcal{G}^\bullet \to \underline{\mathbf{Z}}_\mathcal{C}$
pour lesquels $\mathcal{G}^\bullet$ est un complexe borné supérieurement
de $\underline{\mathbf{Z}}_\mathcal{C}$-modules plats forment un système
cofinal. (Cela résulte de Modules sur les sites, Lemme
\ref{sites-modules-lemma-module-quotient-flat}, et de Catégories dérivées,
Lemme \ref{derived-lemma-subcategory-left-resolution} ; voir la discussion
de la Section \ref{section-flat}.) Nous pouvons donc construire un inverse
du morphisme
$H^0(\mathcal{C}, K) \longrightarrow H^0(\mathcal{C}, K_{ab})$
en représentant un élément $\xi \in H^0(\mathcal{C}, K_{ab})$ par une paire
$$
(s : \mathcal{G}^\bullet \to \underline{\mathbf{Z}}_\mathcal{C},
a : \mathcal{G}^\bullet \to \mathcal{K}^\bullet)
$$
où $\mathcal{G}^\bullet$ est un complexe borné supérieurement de
$\underline{\mathbf{Z}}_\mathcal{C}$-modules plats, et en envoyant cette
paire sur
$$
(\mathcal{G}^\bullet \otimes_{\underline{\mathbf{Z}}_\mathcal{C}}
\mathcal{O}_\mathcal{C}
\to \mathcal{O}_\mathcal{C},
\mathcal{G}^\bullet  \otimes_{\underline{\mathbf{Z}}_\mathcal{C}}
\mathcal{O}_\mathcal{C}
\to \mathcal{K}^\bullet)
$$
Il suffit ici de remarquer que la première flèche est un
quasi-isomorphisme d'après les Lemmes
\ref{lemma-derived-tor-quasi-isomorphism-other-side} et
\ref{lemma-bounded-flat-K-flat}. Nous omettons la vérification détaillée du
fait que cette construction donne bien un inverse.
\end{proof}

\begin{lemma}
\label{lemma-adjoint-lower-shriek-restrict}
Soit $(\mathcal{C}, \mathcal{O})$ un site annelé et soit $U$ un objet de
$\mathcal{C}$. Notons
$j : (\Sh(\mathcal{C}/U), \mathcal{O}_U) \to (\Sh(\mathcal{C}), \mathcal{O})$
le morphisme de localisation correspondant. Le foncteur de restriction
$D(\mathcal{O}) \to D(\mathcal{O}_U)$ est adjoint à droite du prolongement
par zéro $j_! : D(\mathcal{O}_U) \to D(\mathcal{O})$.
\end{lemma}

\begin{proof}
Nous devons montrer que
$$
\Hom_{D(\mathcal{O})}(j_!E, F) = \Hom_{D(\mathcal{O}_U)}(E, F|_U)
$$
Choisissons un complexe $\mathcal{E}^\bullet$ de $\mathcal{O}_U$-modules
représentant $E$ et un complexe K-injectif $\mathcal{I}^\bullet$
représentant $F$. D'après le Lemme
\ref{lemma-restrict-K-injective-to-open}, le complexe
$\mathcal{I}^\bullet|_U$ est lui aussi K-injectif. La formule ci-dessus
devient donc
$$
\Hom_{D(\mathcal{O})}(j_!\mathcal{E}^\bullet, \mathcal{I}^\bullet) =
\Hom_{D(\mathcal{O}_U)}(\mathcal{E}^\bullet, \mathcal{I}^\bullet|_U)
$$
ce qui est vrai puisque $|_U$ et $j_!$ sont des foncteurs adjoints
(Modules sur les sites, Lemme
\ref{sites-modules-lemma-extension-by-zero}) et d'après Catégories dérivées,
Lemme \ref{derived-lemma-K-injective}.
\end{proof}

\begin{lemma}
\label{lemma-j-shriek-and-tensor}
Soit $(\mathcal{C}, \mathcal{O})$ un site annelé. Soit
$U \in \Ob(\mathcal{C})$. Pour $L$ dans $D(\mathcal{O}_U)$ et $K$ dans
$D(\mathcal{O})$, nous avons
$j_!L \otimes_\mathcal{O}^\mathbf{L} K =
j_!(L \otimes_{\mathcal{O}_U}^\mathbf{L} K|_U)$.
\end{lemma}

\begin{proof}
Représentons $L$ par un complexe de $\mathcal{O}_U$-modules et $K$ par un
complexe K-plat de $\mathcal{O}$-modules, puis appliquons Modules sur les
sites, Lemme \ref{sites-modules-lemma-j-shriek-and-tensor}. Nous omettons
les détails.
\end{proof}

\begin{lemma}
\label{lemma-K-injective-flat}
Soit $f : (\Sh(\mathcal{C}), \mathcal{O}_\mathcal{C}) \to
(\Sh(\mathcal{D}), \mathcal{O}_\mathcal{D})$ un morphisme plat de topos
annelés. Si $\mathcal{I}^\bullet$ est un complexe K-injectif de
$\mathcal{O}_\mathcal{C}$-modules, alors $f_*\mathcal{I}^\bullet$ est
K-injectif comme complexe de $\mathcal{O}_\mathcal{D}$-modules.
\end{lemma}

\begin{proof}
Cela est vrai car
$$
\Hom_{K(\mathcal{O}_\mathcal{D})}(\mathcal{F}^\bullet, f_*\mathcal{I}^\bullet)
=
\Hom_{K(\mathcal{O}_\mathcal{C})}(f^*\mathcal{F}^\bullet, \mathcal{I}^\bullet)
$$
d'après Modules sur les sites, Lemme
\ref{sites-modules-lemma-adjoint-pullback-pushforward-modules}, et car
$f^*$ est exact puisque $f$ est supposé plat.
\end{proof}

\begin{lemma}
\label{lemma-hom-K-injective}
Soit $\mathcal{C}$ un site. Soit $\mathcal{O} \to \mathcal{O}'$ un
morphisme de faisceaux d'anneaux. Si $\mathcal{I}^\bullet$ est un complexe
K-injectif de $\mathcal{O}$-modules, alors
$\SheafHom_\mathcal{O}(\mathcal{O}', \mathcal{I}^\bullet)$
est un complexe K-injectif de $\mathcal{O}'$-modules.
\end{lemma}

\begin{proof}
Cela est vrai car
$\Hom_{K(\mathcal{O}')}(\mathcal{G}^\bullet,
\Hom_\mathcal{O}(\mathcal{O}', \mathcal{I}^\bullet)) =
\Hom_{K(\mathcal{O})}(\mathcal{G}^\bullet, \mathcal{I}^\bullet)$
d'après Modules sur les sites, Lemme
\ref{sites-modules-lemma-adjoint-hom-restrict}.
\end{proof}






\section{Localisation et cohomologie}
\label{section-localization}

\noindent
Soit $\mathcal{C}$ un site et soit $f : X \to Y$ un morphisme de
$\mathcal{C}$. Nous obtenons alors un morphisme de topos
$$
j_{X/Y} : \Sh(\mathcal{C}/X) \longrightarrow \Sh(\mathcal{C}/Y)
$$
Voir Sites, Sections \ref{sites-section-localize} et
\ref{sites-section-more-localize}. Certaines questions de cohomologie sont
plus simples pour ce type de morphismes de topos. Voici un exemple où l'on
obtient une forme élémentaire du théorème de changement de base.

\begin{lemma}
\label{lemma-localize-cartesian-square}
Soit $\mathcal{C}$ un site. Soit
$$
\xymatrix{
X' \ar[d] \ar[r] & X \ar[d] \\
Y' \ar[r] & Y
}
$$
un diagramme cartésien de $\mathcal{C}$. Alors nous avons
$j_{Y'/Y}^{-1} \circ Rj_{X/Y, *} = Rj_{X'/Y', *} \circ j_{X'/X}^{-1}$
comme foncteurs $D(\mathcal{C}/X) \to D(\mathcal{C}/Y')$.
\end{lemma}

\begin{proof}
Soit $E \in D(\mathcal{C}/X)$. Choisissons un complexe K-injectif
$\mathcal{I}^\bullet$ de faisceaux abéliens sur $\mathcal{C}/X$
représentant $E$. D'après le Lemme
\ref{lemma-restrict-K-injective-to-open},
$j_{X'/X}^{-1}\mathcal{I}^\bullet$ est lui aussi K-injectif. Nous pouvons
donc calculer $Rj_{X'/Y'}(j_{X'/X}^{-1}E)$ par
$j_{X'/Y', *}j_{X'/X}^{-1}\mathcal{I}^\bullet$. L'égalité résulte alors de
Sites, Lemme \ref{sites-lemma-localize-cartesian-square}.
\end{proof}

\noindent
Si nous avons un site annelé $(\mathcal{C}, \mathcal{O})$ et un morphisme
$f : X \to Y$ de $\mathcal{C}$, alors $j_{X/Y}$ devient un morphisme de
topos annelés
$$
j_{X/Y} :
(\Sh(\mathcal{C}/X), \mathcal{O}_X)
\longrightarrow
(\Sh(\mathcal{C}/Y), \mathcal{O}_Y)
$$
Voir Modules sur les sites, Lemme
\ref{sites-modules-lemma-relocalize}.

\begin{lemma}
\label{lemma-localize-cartesian-square-modules}
Soit $(\mathcal{C}, \mathcal{O})$ un site annelé. Soit
$$
\xymatrix{
X' \ar[d] \ar[r] & X \ar[d] \\
Y' \ar[r] & Y
}
$$
un diagramme cartésien de $\mathcal{C}$. Alors nous avons
$j_{Y'/Y}^* \circ Rj_{X/Y, *} = Rj_{X'/Y', *} \circ j_{X'/X}^*$
comme foncteurs
$D(\mathcal{O}_X) \to D(\mathcal{O}_{Y'})$.
\end{lemma}

\begin{proof}
Puisque $j_{Y'/Y}^{-1}\mathcal{O}_Y = \mathcal{O}_{Y'}$, nous avons
$j_{Y'/Y}^* = Lj_{Y'/Y}^* = j_{Y'/Y}^{-1}$. De même,
$j_{X'/X}^* = Lj_{X'/X}^* = j_{X'/X}^{-1}$. D'après le Lemme
\ref{lemma-modules-abelian-unbounded}, il suffit donc de démontrer le
résultat dans les catégories dérivées de faisceaux abéliens, ce que nous
avons fait dans le Lemme \ref{lemma-localize-cartesian-square}.
\end{proof}












\section{Systèmes projectifs et cohomologie}
\label{section-inverse-systems}

\noindent
Nous démontrons quelques résultats sur les systèmes projectifs de faisceaux
de modules.

\begin{lemma}
\label{lemma-ML-general}
Soit $I$ un idéal d'un anneau $A$. Soit $\mathcal{C}$ un site. Soit
$$
\ldots \to \mathcal{F}_3 \to \mathcal{F}_2 \to \mathcal{F}_1
$$
un système projectif de faisceaux de $A$-modules sur $\mathcal{C}$ tel que
$\mathcal{F}_n = \mathcal{F}_{n + 1}/I^n\mathcal{F}_{n + 1}$. Soit
$p \geq 0$. Supposons que
$$
\bigoplus\nolimits_{n \geq 0} H^{p + 1}(\mathcal{C}, I^n\mathcal{F}_{n + 1})
$$
satisfasse la condition de chaîne ascendante comme
$\bigoplus_{n \geq 0} I^n/I^{n + 1}$-module gradué. Alors le système
projectif $M_n = H^p(\mathcal{C}, \mathcal{F}_n)$ satisfait la condition de
Mittag-Leffler\footnote{En fait, il existe $c \geq 0$ tel que
$\Im(M_n \to M_{n - c})$ soit l'image stable pour tout $n \geq c$.}.
\end{lemma}

\begin{proof}
Posons $N_n = H^{p + 1}(\mathcal{C}, I^n\mathcal{F}_{n + 1})$ et soit
$\delta_n : M_n \to N_n$ le morphisme de connexion en cohomologie provenant
de la suite exacte courte
$0 \to I^n\mathcal{F}_{n + 1} \to \mathcal{F}_{n + 1} \to \mathcal{F}_n \to 0$.
Alors $\bigoplus \Im(\delta_n) \subset \bigoplus N_n$ est un sous-module
gradué. En effet, si $s \in M_n$ et $f \in I^m$, nous avons un diagramme
commutatif
$$
\xymatrix{
0 \ar[r] &
I^n\mathcal{F}_{n + 1} \ar[d]_f \ar[r] &
\mathcal{F}_{n + 1} \ar[d]_f \ar[r] &
\mathcal{F}_n \ar[d]_f \ar[r] & 0 \\
0 \ar[r] &
I^{n + m}\mathcal{F}_{n + m + 1} \ar[r] &
\mathcal{F}_{n + m + 1} \ar[r] &
\mathcal{F}_{n + m} \ar[r] & 0
}
$$
Le morphisme vertical médian est défini en relevant une section locale de
$\mathcal{F}_{n + 1}$ en une section de $\mathcal{F}_{n + m + 1}$, puis en
la multipliant par $f$ ; les autres flèches verticales se définissent de
façon analogue. Nous en déduisons que
$\delta_{n + m}(fs) = f \delta_n(s)$. Par hypothèse, nous pouvons trouver
$s_j \in M_{n_j}$, $j = 1, \ldots, N$, tels que les
$\delta_{n_j}(s_j)$ engendrent $\bigoplus \Im(\delta_n)$ comme module
gradué. Soit $n > c = \max(n_j)$ et soit $s \in M_n$. Il existe alors
$f_j \in I^{n - n_j}$ tels que
$\delta_n(s) = \sum f_j \delta_{n_j}(s_j)$. Nous en déduisons que
$\delta(s - \sum f_j s_j) = 0$, c'est-à-dire qu'il existe
$s' \in M_{n + 1}$ dont l'image est $s - \sum f_js_j$ dans $M_n$. Il
s'ensuit que
$$
\Im(M_{n + 1} \to M_{n - c}) = \Im(M_n \to M_{n - c})
$$
En effet, les éléments $f_js_j$ ont une image nulle dans $M_{n - c}$.
Cela démontre le lemme.
\end{proof}

\begin{lemma}
\label{lemma-ML-general-better}
Soit $I$ un idéal d'un anneau $A$. Soit $\mathcal{C}$ un site. Soit
$$
\ldots \to \mathcal{F}_3 \to \mathcal{F}_2 \to \mathcal{F}_1
$$
un système projectif de $A$-modules sur $\mathcal{C}$ tel que
$\mathcal{F}_n = \mathcal{F}_{n + 1}/I^n\mathcal{F}_{n + 1}$. Soit
$p \geq 0$. Pour tout $n$, posons
$$
N_n =
\bigcap\nolimits_{m \geq n}
\Im\left(
H^{p + 1}(\mathcal{C}, I^n\mathcal{F}_{m + 1}) \to
H^{p + 1}(\mathcal{C}, I^n\mathcal{F}_{n + 1})
\right)
$$
Si $\bigoplus N_n$ satisfait la condition de chaîne ascendante comme
$\bigoplus_{n \geq 0} I^n/I^{n + 1}$-module gradué, alors le système
projectif $M_n = H^p(\mathcal{C}, \mathcal{F}_n)$ satisfait la condition de
Mittag-Leffler\footnote{En fait, il existe $c \geq 0$ tel que
$\Im(M_n \to M_{n - c})$ soit l'image stable pour tout $n \geq c$.}.
\end{lemma}

\begin{proof}
La démonstration est exactement la même que celle du Lemme
\ref{lemma-ML-general}. En effet, les arguments qui y sont donnés
établiront le résultat dès que nous aurons montré que $\bigoplus N_n$ est
un sous-$\bigoplus_{n \geq 0} I^n/I^{n + 1}$-module gradué de
$\bigoplus H^{p + 1}(\mathcal{C}, I^n\mathcal{F}_{n + 1})$ et que l'image
des morphismes de connexion
$\delta_n : M_n \to H^{p + 1}(\mathcal{C}, I^n\mathcal{F}_{n + 1})$ est
contenue dans $N_n$.

\medskip\noindent
Supposons que $\xi \in N_n$ et $f \in I^k$. Choisissons
$m \gg n + k$, puis un élément
$\xi' \in H^{p + 1}(\mathcal{C}, I^n\mathcal{F}_{m + 1})$ relevant
$\xi$. Considérons le diagramme
$$
\xymatrix{
0 \ar[r] &
I^n\mathcal{F}_{m + 1} \ar[d]_f \ar[r] &
\mathcal{F}_{m + 1} \ar[d]_f \ar[r] &
\mathcal{F}_n \ar[d]_f \ar[r] & 0 \\
0 \ar[r] &
I^{n + k}\mathcal{F}_{m + 1} \ar[r] &
\mathcal{F}_{m + 1} \ar[r] &
\mathcal{F}_{n + k} \ar[r] & 0
}
$$
construit comme dans la démonstration du Lemme \ref{lemma-ML-general}.
Nous obtenons un morphisme induit en cohomologie et voyons que
$f \xi' \in H^{p + 1}(\mathcal{C}, I^{n + k}\mathcal{F}_{m + 1})$ a pour
image $f \xi$. Comme cela vaut pour tout $m \gg n + k$, nous voyons que
$f\xi$ appartient à $N_{n + k}$, comme voulu.

\medskip\noindent
Pour voir que l'image des morphismes de connexion $\delta_n$ est contenue
dans $N_n$, considérons les diagrammes
$$
\xymatrix{
0 \ar[r] &
I^n\mathcal{F}_{m + 1} \ar[d] \ar[r] &
\mathcal{F}_{m + 1} \ar[d] \ar[r] &
\mathcal{F}_n \ar[d] \ar[r] & 0 \\
0 \ar[r] &
I^n\mathcal{F}_{n + 1} \ar[r] &
\mathcal{F}_{n + 1} \ar[r] &
\mathcal{F}_n \ar[r] & 0
}
$$
pour $m \geq n$. L'examen des morphismes induits en cohomologie permet de
conclure.
\end{proof}

\begin{lemma}
\label{lemma-topology-I-adic-general}
Soit $I$ un idéal d'un anneau $A$. Soit $\mathcal{C}$ un site. Soit
$$
\ldots \to \mathcal{F}_3 \to \mathcal{F}_2 \to \mathcal{F}_1
$$
un système projectif de faisceaux de $A$-modules sur $\mathcal{C}$ tel que
$\mathcal{F}_n = \mathcal{F}_{n + 1}/I^n\mathcal{F}_{n + 1}$. Soit
$p \geq 0$. Supposons que
$$
\bigoplus\nolimits_{n \geq 0} H^p(\mathcal{C}, I^n\mathcal{F}_{n + 1})
$$
satisfasse la condition de chaîne ascendante comme
$\bigoplus_{n \geq 0} I^n/I^{n + 1}$-module gradué. Alors la topologie
limite sur $M = \lim H^p(\mathcal{C}, \mathcal{F}_n)$ est la topologie
$I$-adique.
\end{lemma}

\begin{proof}
Posons $F^n = \Ker(M \to H^p(\mathcal{C}, \mathcal{F}_n))$ pour
$n \geq 1$ et $F^0 = M$. Observons que
$I F^n \subset F^{n + 1}$. En particulier, $I^n M \subset F^n$. La
topologie $I$-adique est donc plus fine que la topologie limite. Pour la
réciproque, nous montrerons que, pour tout $n$, il existe $m \geq n$ tel
que $F^m \subset I^nM$\footnote{En fait, il existe $c \geq 0$ tel que
$F^{n + c} \subset I^nM$ pour tout $n$.}. Nous avons des morphismes
injectifs
$$
F^n/F^{n + 1} \longrightarrow H^p(\mathcal{C}, \mathcal{F}_{n + 1})
$$
dont l'image est contenue dans celle de
$H^p(\mathcal{C}, I^n\mathcal{F}_{n + 1}) \to
H^p(\mathcal{C}, \mathcal{F}_{n + 1})$.
Notons
$$
E_n \subset H^p(\mathcal{C}, I^n\mathcal{F}_{n + 1})
$$
l'image inverse de $F^n/F^{n + 1}$. Alors $\bigoplus E_n$ est un
sous-$\bigoplus I^n/I^{n + 1}$-module gradué de
$\bigoplus H^p(\mathcal{C}, I^n\mathcal{F}_{n + 1})$, et
$\bigoplus E_n \to \bigoplus F^n/F^{n + 1}$ est un homomorphisme de
modules gradués ; nous omettons les détails. Par hypothèse,
$\bigoplus E_n$ est engendré par un nombre fini d'éléments homogènes sur
$\bigoplus I^n/I^{n + 1}$. Puisque $E_n \to F^n/F^{n + 1}$ est surjectif,
il en va de même de $\bigoplus F^n/F^{n + 1}$. Nous pouvons donc trouver
$r$, des entiers $c_1, \ldots, c_r \geq 0$ et des $a_i \in F^{c_i}$ dont
les images engendrent $\bigoplus F^n/F^{n + 1}$. Posons
$c = \max(c_i)$.

\medskip\noindent
Pour $n \geq c$, nous affirmons que $I F^n = F^{n + 1}$. Cette
affirmation donne $F^{n + c} = I^nF^c \subset I^nM$, comme voulu. Pour la
démontrer, supposons $a \in F^{n + 1}$. L'image de $a$ dans
$F^{n + 1}/F^{n + 2}$ est une combinaison linéaire des $a_i$. Ainsi,
$a  - \sum f_i a_i \in F^{n + 2}$ pour certains
$f_i \in I^{n + 1 - c_i}$. Comme
$I^{n + 1 - c_i} = I \cdot I^{n - c_i}$ puisque $n \geq c_i$, nous
pouvons écrire $f_i = \sum g_{i, j} h_{i, j}$ avec $g_{i, j} \in I$ et
$h_{i, j}a_i \in F^n$. Nous voyons donc que
$F^{n + 1} = F^{n + 2} + IF^n$. Une récurrence immédiate donne
$F^{n + 1} = F^{n + e} + IF^n$ pour tout $e > 0$. Il en résulte que
$IF^n$ est dense dans $F^{n + 1}$. Choisissons des générateurs
$k_1, \ldots, k_r$ de $I$ et considérons l'application continue
$$
u : (F^n)^{\oplus r} \longrightarrow F^{n + 1},\quad
(x_1, \ldots, x_r) \mapsto \sum k_i x_i
$$
(pour la topologie limite). D'après ce qui précède, l'image de
$(F^m)^{\oplus r}$ par $u$ est dense dans $F^{m + 1}$ pour tout
$m \geq n$. Le lemme de l'application ouverte (Compléments d'algèbre,
Lemme \ref{more-algebra-lemma-open-mapping}) montre que $u$ est ouverte.
Ainsi, $u$ est surjective. Par conséquent, $IF^n = F^{n + 1}$ pour
$n \geq c$, ce
qui achève la démonstration.
\end{proof}

\begin{lemma}
\label{lemma-topology-I-adic-general-better}
Soit $I$ un idéal d'un anneau $A$. Soit $\mathcal{C}$ un site. Soit
$$
\ldots \to \mathcal{F}_3 \to \mathcal{F}_2 \to \mathcal{F}_1
$$
un système projectif de faisceaux de $A$-modules sur $\mathcal{C}$ tel que
$\mathcal{F}_n = \mathcal{F}_{n + 1}/I^n\mathcal{F}_{n + 1}$. Soit
$p \geq 0$. Pour tout $n$, posons
$$
N_n =
\bigcap\nolimits_{m \geq n}
\Im\left(
H^p(\mathcal{C}, I^n\mathcal{F}_{m + 1}) \to
H^p(\mathcal{C}, I^n\mathcal{F}_{n + 1})
\right)
$$
Si $\bigoplus N_n$ satisfait la condition de chaîne ascendante comme
$\bigoplus_{n \geq 0} I^n/I^{n + 1}$-module gradué, alors la topologie
limite sur $M = \lim H^p(\mathcal{C}, \mathcal{F}_n)$ est la topologie
$I$-adique.
\end{lemma}

\begin{proof}
La démonstration est exactement la même que celle du Lemme
\ref{lemma-topology-I-adic-general}. En effet, les arguments qui y sont
donnés établiront le résultat dès que nous aurons montré que
$\bigoplus N_n$ est un sous-$\bigoplus_{n \geq 0} I^n/I^{n + 1}$-module
gradué de $\bigoplus H^{p + 1}(\mathcal{C}, I^n\mathcal{F}_{n + 1})$ et que
$F^n/F^{n + 1} \subset H^p(\mathcal{C}, \mathcal{F}_{n + 1})$ est contenu
dans l'image de
$N_n \to H^p(\mathcal{C}, \mathcal{F}_{n + 1})$. Dans la démonstration du
Lemme \ref{lemma-ML-general-better}, nous avons établi l'assertion sur la
structure de module.

\medskip\noindent
Soit $t \in F^n$. Choisissons un élément
$s \in H^p(\mathcal{C}, I^n\mathcal{F}_{n + 1})$
dont l'image est l'image de $t$ dans
$H^p(\mathcal{C}, \mathcal{F}_{n + 1})$. Nous devons montrer que $s$
appartient à $N_n$. Or $F^n$ est le noyau du morphisme
$M \to H^p(\mathcal{C}, \mathcal{F}_n)$ ; pour tout $m \geq n$, nous
pouvons donc envoyer $t$ sur un élément
$t_m \in H^p(\mathcal{C}, \mathcal{F}_{m + 1})$ dont l'image dans
$H^p(\mathcal{C}, \mathcal{F}_n)$ est nulle. Considérons la suite de
cohomologie
$$
H^{p - 1}(\mathcal{C}, \mathcal{F}_n) \to
H^p(\mathcal{C}, I^n\mathcal{F}_{m + 1}) \to
H^p(\mathcal{C}, \mathcal{F}_{m + 1}) \to
H^p(\mathcal{C}, \mathcal{F}_n)
$$
provenant de la suite exacte courte
$0 \to I^n\mathcal{F}_{m + 1} \to \mathcal{F}_{m + 1} \to \mathcal{F}_n \to 0$.
Nous pouvons choisir
$s_m \in H^p(\mathcal{C}, I^n\mathcal{F}_{m + 1})$ d'image $t_m$. En
comparant la suite ci-dessus à celle obtenue pour $m = n$, nous voyons que
l'image de $s_m$ est égale à $s$ à un élément de l'image de
$H^{p - 1}(\mathcal{C}, \mathcal{F}_n) \to
H^p(\mathcal{C}, I^n\mathcal{F}_{n + 1})$.
près. Cependant, ce morphisme se factorise par le morphisme
$H^p(\mathcal{C}, I^n\mathcal{F}_{m + 1}) \to
H^p(\mathcal{C}, I^n\mathcal{F}_{n + 1})$
et nous voyons que $s$ appartient à l'image, comme voulu.
\end{proof}












\section{Limites dérivées et homotopiques}
\label{section-derived-limits}

\noindent
Soit $\mathcal{C}$ un site. Considérons la catégorie
$\mathcal{C} \times \mathbf{N}$ pour laquelle
$\Mor((U, n), (V, m)) = \emptyset$ si $n > m$ et
$\Mor((U, n), (V, m)) = \Mor(U, V)$ sinon. Munissons-la d'une structure
de site en prenant pour recouvrements les familles
$\{(U_i, n) \to (U, n)\}$ telles que $\{U_i \to U\}$ soit un recouvrement
de $\mathcal{C}$. Le lecteur vérifie alors immédiatement que les faisceaux
sur $\mathcal{C} \times \mathbf{N}$ s'identifient aux systèmes projectifs
de faisceaux sur $\mathcal{C}$. En particulier,
$\textit{Ab}(\mathcal{C} \times \mathbf{N})$ est la catégorie des systèmes
projectifs de faisceaux abéliens sur $\mathcal{C}$. Considérons maintenant
le foncteur
$$
\lim : \textit{Ab}(\mathcal{C} \times \mathbf{N}) \to \textit{Ab}(\mathcal{C})
$$
qui associe à un système projectif sa limite. Ce n'est autre que $g_*$, où
$g : \Sh(\mathcal{C} \times \mathbf{N}) \to \Sh(\mathcal{C})$ est le
morphisme de topos associé au foncteur continu et cocontinu
$\mathcal{C} \times \mathbf{N} \to \mathcal{C}$. (Observons que $g^{-1}$
associe à un faisceau sur $\mathcal{C}$ le système projectif constant
correspondant.)

\medskip\noindent
Le formalisme général exposé ci-dessus fournit un foncteur dérivé
$$
R\lim = Rg_* : D(\mathcal{C} \times \mathbf{N}) \to D(\mathcal{C}).
$$
Comme indiqué, ce foncteur est souvent noté $R\lim$.

\medskip\noindent
Par ailleurs, les foncteurs continus et cocontinus
$\mathcal{C} \to \mathcal{C} \times \mathbf{N}$,
$U \mapsto (U, n)$ définissent des morphismes de topos
$i_n : \Sh(\mathcal{C}) \to \Sh(\mathcal{C} \times \mathbf{N})$.
Bien entendu, $i_n^{-1}$ est le foncteur qui extrait le $n$-ième terme du
système projectif. Il existe donc des transformations de foncteurs
$i_{n + 1}^{-1} \to i_n^{-1}$. Ainsi, pour
$K \in D(\mathcal{C} \times \mathbf{N})$, nous obtenons
$K_n = i_n^{-1}K \in D(\mathcal{C})$ et des morphismes
$K_{n + 1} \to K_n$. Dans Catégories dérivées, Définition
\ref{derived-definition-derived-limit}, nous avons défini la notion de
limite homotopique
$$
R\lim K_n \in D(\mathcal{C})
$$
Nous affirmons que les deux notions coïncident (dans la mesure où cette
assertion a un sens).

\begin{lemma}
\label{lemma-derived-limit-is-ok}
Soit $\mathcal{C}$ un site. Soit $K$ un objet de
$D(\mathcal{C} \times \mathbf{N})$. Posons $K_n = i_n^{-1}K$ comme
ci-dessus. Alors
$$
R\lim K \cong R\lim K_n
$$
dans $D(\mathcal{C})$.
\end{lemma}

\begin{proof}
Pour calculer $R\lim$ sur un objet $K$ de
$D(\mathcal{C} \times \mathbf{N})$, choisissons un représentant K-injectif
$\mathcal{I}^\bullet$ dont les termes sont des objets injectifs de
$\textit{Ab}(\mathcal{C} \times \mathbf{N})$ ; voir Injectifs, Théorème
\ref{injectives-theorem-K-injective-embedding-grothendieck}. Nous pouvons,
et allons, considérer $\mathcal{I}^\bullet$ comme un système projectif de
complexes $(\mathcal{I}_n^\bullet)$. Nous voyons alors que
$$
R\lim K = \lim \mathcal{I}_n^\bullet
$$
où le membre de droite est la limite projective prise terme à terme.

\medskip\noindent
Soit $\mathcal{J} = (\mathcal{J}_n)$ un objet injectif de
$\textit{Ab}(\mathcal{C} \times \mathbf{N})$. Les morphismes
$(U, n) \to (U, n + 1)$ sont des monomorphismes de
$\mathcal{C} \times \mathbf{N}$ ; par conséquent,
$\mathcal{J}(U, n + 1) \to \mathcal{J}(U, n)$ est surjectif (Lemme
\ref{lemma-restriction-along-monomorphism-surjective}). Il s'ensuit que
$\mathcal{J}_{n + 1} \to \mathcal{J}_n$ est surjectif comme morphisme de
préfaisceaux.

\medskip\noindent
Notons que le foncteur $i_n^{-1}$ admet un adjoint à gauche exact
$i_{n, !}$. En effet, $i_{n, !}\mathcal{F}$ est le système projectif
$\ldots 0 \to 0 \to \mathcal{F} \to \ldots \to \mathcal{F}$. Les
complexes $i_n^{-1}\mathcal{I}^\bullet = \mathcal{I}_n^\bullet$ sont donc
K-injectifs d'après Catégories dérivées, Lemme
\ref{derived-lemma-adjoint-preserve-K-injectives}.

\medskip\noindent
Comme nous avons choisi notre complexe K-injectif à termes injectifs, nous
en déduisons que
$$
0 \to  \lim \mathcal{I}_n^\bullet \to \prod \mathcal{I}_n^\bullet
\to \prod \mathcal{I}_n^\bullet \to 0
$$
est une suite exacte courte de complexes de faisceaux abéliens, puisqu'elle
est une suite exacte courte de complexes de préfaisceaux abéliens. De plus,
les produits du milieu et de droite représentent les produits dans
$D(\mathcal{C})$ ; voir Injectifs, Lemme
\ref{injectives-lemma-derived-products}, et sa démonstration (c'est ici que
nous utilisons la K-injectivité de $\mathcal{I}_n^\bullet$). Ainsi,
$R\lim K$ est une limite homotopique du système projectif $(K_n)$, par
définition des limites homotopiques dans les catégories triangulées.
\end{proof}

\begin{lemma}
\label{lemma-RGamma-commutes-with-Rlim}
Soit $(\mathcal{C}, \mathcal{O})$ un site annelé. Les foncteurs
$R\Gamma(\mathcal{C}, -)$ et $R\Gamma(U, -)$, pour
$U \in \Ob(\mathcal{C})$, commutent à $R\lim$. De plus, il existe des
suites exactes courtes
$$
0 \to
R^1\lim H^{m - 1}(U, K_n) \to H^m(U, R\lim K_n) \to
\lim H^m(U, K_n) \to 0
$$
pour tout système projectif $(K_n)$ dans $D(\mathcal{O})$ et tout
$m \in \mathbf{Z}$. Il en va de même pour
$H^m(\mathcal{C}, R\lim K_n)$.
\end{lemma}

\begin{proof}
La première assertion résulte d'Injectifs, Lemme
\ref{injectives-lemma-RF-commutes-with-Rlim}. Nous pouvons ensuite appliquer
Compléments d'algèbre, Remarque
\ref{more-algebra-remark-compare-derived-limit}, à
$R\lim R\Gamma(U, K_n) = R\Gamma(U, R\lim K_n)$ pour obtenir les suites
exactes courtes.
\end{proof}

\begin{lemma}
\label{lemma-Rf-commutes-with-Rlim}
Soit $f : (\Sh(\mathcal{C}), \mathcal{O}) \to (\Sh(\mathcal{C}'), \mathcal{O}')$
un morphisme de topos annelés. Alors $Rf_*$ commute à $R\lim$, c'est-à-dire
que $Rf_*$ commute aux limites dérivées.
\end{lemma}

\begin{proof}
Soit $(K_n)$ un système projectif d'objets de $D(\mathcal{O})$. Par
récurrence sur $n$, nous pouvons choisir des complexes
$\mathcal{K}_n^\bullet$ de $\mathcal{O}$-modules et des morphismes de
complexes $\mathcal{K}_{n + 1}^\bullet \to \mathcal{K}_n^\bullet$
représentant les morphismes $K_{n + 1} \to K_n$ dans $D(\mathcal{O})$.
Autrement dit, il existe un objet $K$ de
$D(\mathcal{C} \times \mathbf{N})$ dont le système projectif associé est
celui qui nous est donné. Considérons ensuite le diagramme commutatif
$$
\xymatrix{
\Sh(\mathcal{C} \times \mathbf{N}) \ar[r]_g \ar[d]_{f \times 1} &
\Sh(\mathcal{C}) \ar[d]_f \\
\Sh(\mathcal{C}' \times \mathbf{N}) \ar[r]^{g'} &
\Sh(\mathcal{C}')
}
$$
de morphismes de topos. Il en résulte que
$R\lim R(f \times 1)_*K = Rf_* R\lim K$. En explicitant les définitions et
en utilisant le Lemme \ref{lemma-derived-limit-is-ok}, nous obtenons
$R\lim (Rf_*K_n) = Rf_*(R\lim K_n)$.

\medskip\noindent
Autre démonstration lorsque $\mathcal{C}$ possède suffisamment de points.
Considérons le triangle distingué qui définit
$$
R\lim K_n \to \prod K_n \to \prod K_n
$$
dans $D(\mathcal{O})$. En appliquant le foncteur exact $Rf_*$, nous obtenons
le triangle distingué
$$
Rf_*(R\lim K_n) \to Rf_*\left(\prod K_n\right) \to Rf_*\left(\prod K_n\right)
$$
dans $D(\mathcal{O}')$. Il suffit donc de prouver que $Rf_*$ commute aux
produits dans la catégorie dérivée (lesquels ne sont pas simplement donnés
par les produits de complexes ; voir Injectifs, Lemme
\ref{injectives-lemma-derived-products}). Or, comme $Rf_*$ est un adjoint à
droite d'après le Lemme \ref{lemma-adjoint}, cela résulte formellement de
Catégories, Lemme \ref{categories-lemma-adjoint-exact}. Attention : nous ne
pouvons pas appliquer directement Catégories, Lemme
\ref{categories-lemma-adjoint-exact}, car $R\lim K_n$ n'est pas une limite
dans $D(\mathcal{O})$.
\end{proof}

\begin{remark}
\label{remark-discuss-derived-limit}
Soit $(\mathcal{C}, \mathcal{O})$ un site annelé. Soit $(K_n)$ un système
projectif dans $D(\mathcal{O})$. Posons $K = R\lim K_n$. Pour chaque $n$
et $m$, soit $\mathcal{H}^m_n = H^m(K_n)$ le $m$-ième faisceau de
cohomologie de $K_n$, et posons de même $\mathcal{H}^m = H^m(K)$. Notons
$\underline{\mathcal{H}}^m_n$ le préfaisceau
$$
U \longmapsto \underline{\mathcal{H}}^m_n(U) = H^m(U, K_n)
$$
De même, posons $\underline{\mathcal{H}}^m(U) = H^m(U, K)$. D'après le
Lemme \ref{lemma-sheafification-cohomology}, $\mathcal{H}^m_n$ est le
faisceau associé à $\underline{\mathcal{H}}^m_n$, et $\mathcal{H}^m$ est
le faisceau associé à $\underline{\mathcal{H}}^m$. Nous avons le diagramme
$$
\xymatrix{
K \ar@{=}[d] &
\underline{\mathcal{H}}^m \ar[d] \ar[r] & 
\mathcal{H}^m \ar[d] \\
R\lim K_n &
\lim \underline{\mathcal{H}}^m_n \ar[r] & 
\lim \mathcal{H}^m_n
}
$$
En général, $\lim \mathcal{H}^m_n$ n'est pas nécessairement le faisceau
associé à $\lim \underline{\mathcal{H}}^m_n$. Si
$U \in \mathcal{C}$, nous avons les suites exactes courtes
\begin{equation}
\label{equation-ses-Rlim-over-U}
0 \to
R^1\lim \underline{\mathcal{H}}^{m - 1}_n(U) \to
\underline{\mathcal{H}}^m(U) \to
\lim \underline{\mathcal{H}}^m_n(U) \to 0
\end{equation}
d'après le Lemme \ref{lemma-RGamma-commutes-with-Rlim}.
\end{remark}

\noindent
Le lemme suivant s'applique à un système projectif de modules
quasi-cohérents dont les morphismes de transition sont surjectifs sur un
espace algébrique ou un champ algébrique.

\begin{lemma}
\label{lemma-inverse-limit-is-derived-limit}
Soit $(\mathcal{C}, \mathcal{O})$ un site annelé. Soit $(\mathcal{F}_n)$ un
système projectif de $\mathcal{O}$-modules. Soit
$\mathcal{B} \subset \Ob(\mathcal{C})$ une partie. Supposons que
\begin{enumerate}
\item tout objet de $\mathcal{C}$ possède un recouvrement dont les membres
appartiennent à $\mathcal{B}$,
\item $H^p(U, \mathcal{F}_n) = 0$ pour $p > 0$ et $U \in \mathcal{B}$,
\item le système projectif $\mathcal{F}_n(U)$ a un $R^1\lim$ nul pour
$U \in \mathcal{B}$.
\end{enumerate}
Alors $R\lim \mathcal{F}_n = \lim \mathcal{F}_n$ et nous avons
$H^p(U, \lim \mathcal{F}_n) = 0$ pour $p > 0$ et $U \in \mathcal{B}$.
\end{lemma}

\begin{proof}
Posons $K_n = \mathcal{F}_n$ et $K = R\lim \mathcal{F}_n$. Avec les
notations de la Remarque \ref{remark-discuss-derived-limit}, l'hypothèse
(2) donne, pour $U \in \mathcal{B}$,
$\underline{\mathcal{H}}_n^m(U) = 0$ si $m \not = 0$ et
$\underline{\mathcal{H}}_n^0(U) = \mathcal{F}_n(U)$. L'Équation
(\ref{equation-ses-Rlim-over-U}) et l'hypothèse (3) donnent
$\underline{\mathcal{H}}^m(U) = 0$ si $m \not = 0$, et ce groupe est égal à
$\lim \mathcal{F}_n(U)$ si $m = 0$. En prenant les faisceaux associés et en
utilisant (1), nous obtenons $\mathcal{H}^m = 0$ si $m \not = 0$, et ce
faisceau est égal à $\lim \mathcal{F}_n$ si $m = 0$. Ainsi
$K = \lim \mathcal{F}_n$. Comme
$H^m(U, K) = \underline{\mathcal{H}}^m(U) = 0$ pour $m > 0$ (voir
ci-dessus), la seconde assertion est satisfaite.
\end{proof}

\begin{lemma}
\label{lemma-cohomology-derived-limit-injective}
Soit $(\mathcal{C}, \mathcal{O})$ un site annelé. Soit $(K_n)$ un système
projectif dans $D(\mathcal{O})$. Soient $V \in \Ob(\mathcal{C})$ et
$m \in \mathbf{Z}$. Supposons qu'il existe un entier $n(V)$ et un système
cofinal $\text{Cov}_V$ de recouvrements de $V$ tels que, pour
$\{V_i \to V\} \in \text{Cov}_V$,
\begin{enumerate}
\item $R^1\lim H^{m - 1}(V_i, K_n) = 0$, et
\item $H^m(V_i, K_n) \to H^m(V_i, K_{n(V)})$ soit injectif pour
$n \geq n(V)$.
\end{enumerate}
Alors le morphisme sur les sections
$H^m(R\lim K_n)(V) \to H^m(K_{n(V)})(V)$ est injectif.
\end{lemma}

\begin{proof}
Soit $\gamma \in H^m(R\lim K_n)(V)$ dont l'image dans
$H^m(K_{n(V)})(V)$ est nulle. Comme $H^m(R\lim K_n)$ est le faisceau
associé au préfaisceau $U \mapsto H^m(U, R\lim K_n)$ (d'après le Lemme
\ref{lemma-sheafification-cohomology}), nous pouvons choisir
$\{V_i \to V\} \in \text{Cov}_V$ et des éléments
$\tilde\gamma_i \in H^m(V_i, R\lim K_n)$ dont les images sont les
$\gamma|_{V_i}$. L'image de $\tilde\gamma_i$ est alors un élément
$\tilde\gamma_{i, n(V)} \in H^m(V_i, K_{n(V)})$. Puisque
$H^m(K_{n(V)})$ est le faisceau associé à
$U \mapsto H^m(U, K_{n(V)})$ (à nouveau d'après le Lemme
\ref{lemma-sheafification-cohomology}), nous pouvons, après avoir remplacé
$\{V_i \to V\}$ par un raffinement, supposer que
$\tilde\gamma_{i, n(V)} = 0$ pour tout $i$. Pour ce recouvrement,
considérons les suites exactes courtes
$$
0 \to
R^1\lim H^{m - 1}(V_i, K_n) \to H^m(V_i, R\lim K_n) \to
\lim H^m(V_i, K_n) \to 0
$$
du Lemme \ref{lemma-RGamma-commutes-with-Rlim}. D'après l'hypothèse (1),
le groupe de gauche est nul et, d'après l'hypothèse (2), le groupe de
droite s'injecte dans $H^m(V_i, K_{n(V)})$. Nous en déduisons que
$\tilde\gamma_i = 0$, puis que $\gamma = 0$, comme voulu.
\end{proof}

\begin{lemma}
\label{lemma-is-limit-per-object}
Soit $(\mathcal{C}, \mathcal{O})$ un site annelé. Soit
$E \in D(\mathcal{O})$. Soit $\mathcal{B} \subset \Ob(\mathcal{C})$ une
partie. Supposons que
\begin{enumerate}
\item tout objet de $\mathcal{C}$ possède un recouvrement dont les membres
appartiennent à $\mathcal{B}$, et
\item pour tout $V \in \mathcal{B}$, il existe une fonction
$p(V, -) : \mathbf{Z} \to \mathbf{Z}$ et un système cofinal
$\text{Cov}_V$ de recouvrements de $V$ tels que
$$
H^p(V_i, H^{m - p}(E)) = 0
$$
pour tout $\{V_i \to V\} \in \text{Cov}_V$ et tous entiers $p, m$
vérifiant $p > p(V, m)$.
\end{enumerate}
Alors le morphisme $E \to R\lim \tau_{\geq -n} E$ de Catégories dérivées,
Remarque
\ref{derived-remark-map-into-derived-limit-truncations}
est un isomorphisme dans $D(\mathcal{O})$.
\end{lemma}

\begin{proof}
Posons $K_n = \tau_{\geq -n}E$ et $K = R\lim K_n$. Le morphisme canonique
$E \to K$ provient des morphismes canoniques
$E \to K_n = \tau_{\geq -n}E$. Nous devons montrer que $E \to K$ induit
un isomorphisme $H^m(E) \to H^m(K)$ de faisceaux de cohomologie. Dans la
suite de la démonstration, fixons $m$. Si $n \geq -m$, le morphisme
$E \to \tau_{\geq -n}E = K_n$ induit un isomorphisme
$H^m(E) \to H^m(K_n)$. Pour achever la démonstration, il suffit de montrer
que, pour tout $V \in \mathcal{B}$, il existe un entier $n(V) \geq -m$ tel
que le morphisme $H^m(K)(V) \to H^m(K_{n(V)})(V)$ soit injectif. En effet,
la composée
$$
H^m(E)(V) \to H^m(K)(V) \to H^m(K_{n(V)})(V)
$$
est alors une bijection et la seconde flèche est injective ; la première
flèche est donc bijective. La propriété (1) implique alors que
$H^m(E) \to H^m(K)$ est un isomorphisme. Posons
$$
n(V) = 1 + \max\{-m, p(V, m - 1) - m, -1 + p(V, m) - m, -2 + p(V, m + 1) - m\}.
$$
de sorte que, dans tous les cas, $n(V) \geq -m$. Affirmation : les
morphismes
$$
H^{m - 1}(V_i, K_{n + 1}) \to H^{m - 1}(V_i, K_n)
\quad\text{et}\quad
H^m(V_i, K_{n + 1}) \to H^m(V_i, K_n)
$$
sont des isomorphismes pour $n \geq n(V)$ et
$\{V_i \to V\} \in \text{Cov}_V$. L'affirmation implique que les
conditions (1) et (2) du Lemme
\ref{lemma-cohomology-derived-limit-injective} sont satisfaites, et donc
l'injectivité voulue. Rappelons (Catégories dérivées, Remarque
\ref{derived-remark-truncation-distinguished-triangle})
que nous avons des triangles distingués
$$
H^{-n - 1}(E)[n + 1] \to
K_{n + 1} \to K_n \to H^{-n - 1}(E)[n + 2]
$$
En examinant la suite exacte longue de cohomologie associée, l'affirmation
résulte de l'annulation de
$$
H^{m + n}(V_i, H^{-n - 1}(E)),\quad
H^{m + n + 1}(V_i, H^{-n - 1}(E)),\quad
H^{m + n + 2}(V_i, H^{-n - 1}(E))
$$
pour $n \geq n(V)$ et $\{V_i \to V\} \in \text{Cov}_V$. Cela découle de
notre choix de $n(V)$ et de l'hypothèse du lemme.
\end{proof}

\begin{lemma}
\label{lemma-is-limit-spaltenstein}
Soit $(\mathcal{C}, \mathcal{O})$ un site annelé. Soit
$E \in D(\mathcal{O})$. Soit $\mathcal{B} \subset \Ob(\mathcal{C})$ une
partie. Supposons que
\begin{enumerate}
\item tout objet de $\mathcal{C}$ possède un recouvrement dont les membres
appartiennent à $\mathcal{B}$, et
\item pour tout $V \in \mathcal{B}$, il existe un entier $d_V \geq 0$ et un
système cofinal $\text{Cov}_V$ de recouvrements de $V$ tels que
$$
H^p(V_i, H^q(E)) = 0 \text{ pour }
\{V_i \to V\} \in \text{Cov}_V,\ p > d_V, \text{ et }q < 0
$$
\end{enumerate}
Alors le morphisme $E \to R\lim \tau_{\geq -n} E$ de Catégories dérivées,
Remarque
\ref{derived-remark-map-into-derived-limit-truncations}
est un isomorphisme dans $D(\mathcal{O})$.
\end{lemma}

\begin{proof}
Cela résulte du Lemme \ref{lemma-is-limit-per-object}, avec
$p(V, m) = d_V + \max(0, m)$.
\end{proof}

\begin{lemma}
\label{lemma-is-limit}
Soit $(\mathcal{C}, \mathcal{O})$ un site annelé. Soit
$E \in D(\mathcal{O})$. Supposons qu'il existe une fonction
$p(-) : \mathbf{Z} \to \mathbf{Z}$ et une partie
$\mathcal{B} \subset \Ob(\mathcal{C})$ telles que
\begin{enumerate}
\item tout objet de $\mathcal{C}$ possède un recouvrement dont les membres
appartiennent à $\mathcal{B}$,
\item $H^p(V, H^{m - p}(E)) = 0$ pour $p > p(m)$ et
$V \in \mathcal{B}$.
\end{enumerate}
Alors le morphisme $E \to R\lim \tau_{\geq -n} E$ de Catégories dérivées,
Remarque
\ref{derived-remark-map-into-derived-limit-truncations}
est un isomorphisme dans $D(\mathcal{O})$.
\end{lemma}

\begin{proof}
Appliquons le Lemme \ref{lemma-is-limit-per-object} avec
$p(V, m) = p(m)$ et $\text{Cov}_V$ égal à l'ensemble des recouvrements
$\{V_i \to V\}$ tels que $V_i \in \mathcal{B}$ pour tout $i$.
\end{proof}

\begin{lemma}
\label{lemma-is-limit-dimension}
Soit $(\mathcal{C}, \mathcal{O})$ un site annelé. Soit
$E \in D(\mathcal{O})$. Supposons qu'il existe un entier $d \geq 0$ et
une partie $\mathcal{B} \subset \Ob(\mathcal{C})$ tels que
\begin{enumerate}
\item tout objet de $\mathcal{C}$ possède un recouvrement dont les membres
appartiennent à $\mathcal{B}$,
\item $H^p(V, H^q(E)) = 0$ pour $p > d$, $q < 0$ et
$V \in \mathcal{B}$.
\end{enumerate}
Alors le morphisme $E \to R\lim \tau_{\geq -n} E$ de Catégories dérivées,
Remarque
\ref{derived-remark-map-into-derived-limit-truncations}
est un isomorphisme dans $D(\mathcal{O})$.
\end{lemma}

\begin{proof}
Appliquons le Lemme \ref{lemma-is-limit-spaltenstein} avec $d_V = d$ et
$\text{Cov}_V$ égal à l'ensemble des recouvrements $\{V_i \to V\}$ tels
que $V_i \in \mathcal{B}$ pour tout $i$.
\end{proof}

\noindent
Les lemmes précédents permettent de calculer la cohomologie dans certaines
situations.

\begin{lemma}
\label{lemma-cohomology-over-U-trivial}
Soit $(\mathcal{C}, \mathcal{O})$ un site annelé. Soit $K$ un objet de
$D(\mathcal{O})$. Soit $\mathcal{B} \subset \Ob(\mathcal{C})$ une partie.
Supposons que
\begin{enumerate}
\item tout objet de $\mathcal{C}$ possède un recouvrement dont les membres
appartiennent à $\mathcal{B}$,
\item $H^p(U, H^q(K)) = 0$ pour tous $p > 0$, $q \in \mathbf{Z}$ et
$U \in \mathcal{B}$.
\end{enumerate}
Alors $H^q(U, K) = H^0(U, H^q(K))$ pour $q \in \mathbf{Z}$ et
$U \in \mathcal{B}$.
\end{lemma}

\begin{proof}
Observons que $K = R\lim \tau_{\geq -n} K$ d'après le Lemme
\ref{lemma-is-limit-dimension}, avec $d = 0$. Soit
$U \in \mathcal{B}$. L'Équation (\ref{equation-ses-Rlim-over-U}) donne une
suite exacte courte
$$
0 \to R^1\lim H^{q - 1}(U, \tau_{\geq -n}K) \to
H^q(U, K) \to \lim H^q(U, \tau_{\geq -n}K) \to 0
$$
La condition (2) implique
$H^q(U, \tau_{\geq -n} K) = H^0(U, H^q(\tau_{\geq -n} K))$
pour tout $q$, en utilisant la suite spectrale de Catégories dérivées,
Lemme \ref{derived-lemma-two-ss-complex-functor}. Cette suite spectrale
converge car $\tau_{\geq -n}K$ est borné inférieurement. Si $n > -q$,
nous avons $H^q(\tau_{\geq -n}K) = H^q(K)$. Ainsi, les systèmes à gauche
et à droite dans la suite exacte courte affichée sont stationnaires, de
valeurs respectives $H^0(U, H^{q - 1}(K))$ et $H^0(U, H^q(K))$, d'où le
lemme.
\end{proof}

\noindent
Voici un autre cas où nous pouvons décrire la limite dérivée.

\begin{lemma}
\label{lemma-derived-limit-suitable-system}
Soit $(\mathcal{C}, \mathcal{O})$ un site annelé. Soit $(K_n)$ un système
projectif d'objets de $D(\mathcal{O})$. Soit
$\mathcal{B} \subset \Ob(\mathcal{C})$ une partie. Supposons que
\begin{enumerate}
\item tout objet de $\mathcal{C}$ possède un recouvrement dont les membres
appartiennent à $\mathcal{B}$,
\item pour tous $U \in \mathcal{B}$ et $q \in \mathbf{Z}$, nous avons
\begin{enumerate}
\item $H^p(U, H^q(K_n)) = 0$ pour $p > 0$,
\item le système projectif $H^0(U, H^q(K_n))$ a un $R^1\lim$ nul.
\end{enumerate}
\end{enumerate}
Alors $H^q(R\lim K_n) = \lim H^q(K_n)$ pour $q \in \mathbf{Z}$.
\end{lemma}

\begin{proof}
Posons $K = R\lim K_n$. Nous utiliserons les notations de la Remarque
\ref{remark-discuss-derived-limit}. Soit $U \in \mathcal{B}$. D'après le
Lemme \ref{lemma-cohomology-over-U-trivial} et (2)(a), nous avons
$H^q(U, K_n) = H^0(U, H^q(K_n))$. Comme le foncteur $R\Gamma(U, -)$
commute aux limites dérivées, nous avons
$$
H^q(U, K) = H^q(R\lim R\Gamma(U, K_n)) = \lim H^0(U, H^q(K_n))
$$
où la dernière égalité résulte de Compléments d'algèbre, Remarque
\ref{more-algebra-remark-compare-derived-limit}, et de l'hypothèse (2)(b).
Ainsi, $H^q(U, K)$ est la limite projective des sections des faisceaux
$H^q(K_n)$ sur $U$. Comme $\lim H^q(K_n)$ est un faisceau, l'hypothèse
(1) montre que $H^q(K)$, qui est le faisceau associé au préfaisceau
$U \mapsto H^q(U, K)$, est égal à $\lim H^q(K_n)$. Cela démontre le lemme.
\end{proof}







\section{Construction de résolutions K-injectives}
\label{section-K-injective}

\noindent
Soit $(\mathcal{C}, \mathcal{O})$ un site annelé. Soit
$\mathcal{F}^\bullet$ un complexe de $\mathcal{O}$-modules. La catégorie
$\textit{Mod}(\mathcal{O})$ possède suffisamment d'injectifs ; nous pouvons
donc utiliser Catégories dérivées, Lemme
\ref{derived-lemma-special-inverse-system}, pour construire un diagramme
$$
\xymatrix{
\ldots \ar[r] &
\tau_{\geq -2}\mathcal{F}^\bullet \ar[r] \ar[d] &
\tau_{\geq -1}\mathcal{F}^\bullet \ar[d] \\
\ldots \ar[r] & \mathcal{I}_2^\bullet \ar[r] & \mathcal{I}_1^\bullet
}
$$
dans la catégorie des complexes de $\mathcal{O}$-modules, tel que
\begin{enumerate}
\item les flèches verticales soient des quasi-isomorphismes,
\item $\mathcal{I}_n^\bullet$ soit un complexe borné inférieurement
d'injectifs,
\item les flèches $\mathcal{I}_{n + 1}^\bullet \to \mathcal{I}_n^\bullet$
soient des surjections scindées terme à terme.
\end{enumerate}
La catégorie des $\mathcal{O}$-modules admet des limites (elles se calculent
au niveau des préfaisceaux) ; nous pouvons donc former la limite terme à
terme $\mathcal{I}^\bullet = \lim_n \mathcal{I}_n^\bullet$. D'après
Catégories dérivées, Lemmes
\ref{derived-lemma-bounded-below-injectives-K-injective} et
\ref{derived-lemma-limit-K-injectives}
il s'agit d'un complexe K-injectif. En général, le morphisme canonique
\begin{equation}
\label{equation-into-candidate-K-injective}
\mathcal{F}^\bullet \to \mathcal{I}^\bullet
\end{equation}
n'est pas nécessairement un quasi-isomorphisme. Le lemme suivant donne des
conditions qui assurent qu'il en est un.

\begin{lemma}
\label{lemma-K-injective}
Plaçons-nous dans la situation décrite ci-dessus. Notons
$\mathcal{H}^m = H^m(\mathcal{F}^\bullet)$ le $m$-ième faisceau de
cohomologie. Soit $\mathcal{B} \subset \Ob(\mathcal{C})$ une partie et
soit $d \in \mathbf{N}$. Supposons que
\begin{enumerate}
\item tout objet de $\mathcal{C}$ possède un recouvrement dont les membres
appartiennent à $\mathcal{B}$,
\item pour tout $U \in \mathcal{B}$, nous ayons
$H^p(U, \mathcal{H}^q) = 0$ pour $p > d$ et $q < 0$\footnote{Il suffit
que $\forall m$, $\exists p(m)$,
$H^p(U, \mathcal{H}^{m - p}) = 0$ pour $p > p(m)$ ; voir le Lemme
\ref{lemma-is-limit}.}.
\end{enumerate}
Alors (\ref{equation-into-candidate-K-injective}) est un
quasi-isomorphisme.
\end{lemma}

\begin{proof}
Il suffit, d'après Catégories dérivées, Lemme
\ref{derived-lemma-difficulty-K-injectives}, de montrer que le morphisme
$\mathcal{F}^\bullet \to R\lim \tau_{\geq -n} \mathcal{F}^\bullet$
est un isomorphisme. Cela résulte du Lemme
\ref{lemma-is-limit-dimension}.
\end{proof}

\noindent
Voici un lemme technique sur les faisceaux de cohomologie des limites terme
à terme de systèmes projectifs de complexes de modules. Il convient
d'éviter autant que possible d'utiliser ce lemme et de privilégier les
arguments faisant intervenir des limites projectives dérivées.

\begin{lemma}
\label{lemma-inverse-limit-complexes}
Soit $(\mathcal{C}, \mathcal{O})$ un site annelé. Soit
$(\mathcal{F}_n^\bullet)$ un système projectif de complexes de
$\mathcal{O}$-modules. Soit $m \in \mathbf{Z}$. Supposons donnés une partie
$\mathcal{B} \subset \Ob(\mathcal{C})$ et un entier $n_0$ tels que
\begin{enumerate}
\item tout objet de $\mathcal{C}$ possède un recouvrement dont les membres
appartiennent à $\mathcal{B}$,
\item pour tout $U \in \mathcal{B}$,
\begin{enumerate}
\item les systèmes projectifs de groupes abéliens
$\mathcal{F}_n^{m - 2}(U)$ et $\mathcal{F}_n^{m - 1}(U)$
aient un $R^1\lim$ nul (par exemple, qu'ils satisfassent la condition de
Mittag-Leffler),
\item le système projectif de groupes abéliens
$H^{m - 1}(\mathcal{F}_n^\bullet(U))$ ait un $R^1\lim$ nul (par exemple,
qu'il satisfasse la condition de Mittag-Leffler), et
\item nous ayons
$H^m(\mathcal{F}_n^\bullet(U)) = H^m(\mathcal{F}_{n_0}^\bullet(U))$
pour tout $n \geq n_0$.
\end{enumerate}
\end{enumerate}
Alors les morphismes
$H^m(\mathcal{F}^\bullet) \to \lim H^m(\mathcal{F}_n^\bullet)
\to H^m(\mathcal{F}_{n_0}^\bullet)$ sont des isomorphismes de faisceaux,
où $\mathcal{F}^\bullet = \lim \mathcal{F}_n^\bullet$ est la limite
projective terme à terme.
\end{lemma}

\begin{proof}
Soit $U \in \mathcal{B}$. Notons que
$H^m(\mathcal{F}^\bullet(U))$ est la cohomologie de
$$
\lim_n \mathcal{F}_n^{m - 2}(U) \to
\lim_n \mathcal{F}_n^{m - 1}(U) \to
\lim_n \mathcal{F}_n^m(U) \to
\lim_n \mathcal{F}_n^{m + 1}(U)
$$
au troisième terme en partant de la gauche. D'après les hypothèses (2)(a)
et (2)(b), nous pouvons appliquer Compléments d'algèbre, Lemme
\ref{more-algebra-lemma-apply-Mittag-Leffler-again}, pour conclure que
$$
H^m(\mathcal{F}^\bullet(U)) = \lim H^m(\mathcal{F}_n^\bullet(U))
$$
D'après l'hypothèse (2)(c), nous en déduisons
$$
H^m(\mathcal{F}^\bullet(U)) = H^m(\mathcal{F}_n^\bullet(U))
$$
pour tout $n \geq n_0$. L'hypothèse (1) montre que le faisceau associé à
$U \mapsto H^m(\mathcal{F}^\bullet(U))$ est égal au faisceau associé à
$U \mapsto H^m(\mathcal{F}_n^\bullet(U))$ pour tout $n \geq n_0$. Le
système projectif de faisceaux $H^m(\mathcal{F}_n^\bullet)$ est donc
constant pour $n \geq n_0$, de valeur $H^m(\mathcal{F}^\bullet)$, ce qui
démontre le lemme.
\end{proof}







\section{Dimension cohomologique bornée}
\label{section-bounded}

\noindent
Dans cette section, nous cherchons à déterminer quand un foncteur $Rf_*$
est de dimension cohomologique bornée. Cette question est assez délicate
pour les complexes non bornés.

\begin{situation}
\label{situation-olsson-laszlo}
Soit $\mathcal{C}$ un site. Soit $\mathcal{O}$ un faisceau d'anneaux sur
$\mathcal{C}$. Soit $\mathcal{A} \subset \textit{Mod}(\mathcal{O})$ une
sous-catégorie de Serre faible. Nous faisons l'hypothèse suivante : il
existe une partie $\mathcal{B} \subset \Ob(\mathcal{C})$ telle que
\begin{enumerate}
\item tout objet de $\mathcal{C}$ possède un recouvrement dont les membres
appartiennent à $\mathcal{B}$, et
\item pour tout $V \in \mathcal{B}$, il existe un entier $d_V$ et un
système cofinal $\text{Cov}_V$ de recouvrements de $V$ tels que
$$
H^p(V_i, \mathcal{F}) = 0 \text{ pour }
\{V_i \to V\} \in \text{Cov}_V,\ p > d_V, \text{ et }
\mathcal{F} \in \Ob(\mathcal{A})
$$
\end{enumerate}
\end{situation}

\begin{lemma}
\label{lemma-olsson-laszlo}
\begin{reference}
C'est \cite[Proposition 2.1.4]{six-I}, avec des hypothèses légèrement
modifiées ; il s'agit de l'analogue pour les sites de
\cite[Proposition 3.13]{Spaltenstein}.
\end{reference}
Dans la Situation \ref{situation-olsson-laszlo}, pour tout
$E \in D_\mathcal{A}(\mathcal{O})$, le morphisme
$E \to R\lim \tau_{\geq -n} E$
de Catégories dérivées, Remarque
\ref{derived-remark-map-into-derived-limit-truncations}
est un isomorphisme dans $D(\mathcal{O})$.
\end{lemma}

\begin{proof}
Cela résulte immédiatement du Lemme \ref{lemma-is-limit-spaltenstein}.
\end{proof}

\begin{lemma}
\label{lemma-olsson-laszlo-modified}
Dans la Situation \ref{situation-olsson-laszlo}, soit $(K_n)$ un système
projectif dans $D_\mathcal{A}^+(\mathcal{O})$. Supposons que, pour tout
$j$, le système projectif $(H^j(K_n))$ dans $\mathcal{A}$ soit
stationnaire, de valeur $\mathcal{H}^j$. Alors
$H^j(R\lim K_n) = \mathcal{H}^j$ pour tout $j$.
\end{lemma}

\begin{proof}
Soit $V \in \mathcal{B}$. Soit $\{V_i \to V\}$ un élément de l'ensemble
$\text{Cov}_V$ de la Situation \ref{situation-olsson-laszlo}. Comme $K_n$
est borné inférieurement, il existe une suite spectrale
$$
E_2^{p, q} = H^p(V_i, H^q(K_n))
$$
convergeant vers $H^{p + q}(V_i, K_n)$. Voir Catégories dérivées, Lemme
\ref{derived-lemma-two-ss-complex-functor}. Observons que, par hypothèse,
$E_2^{p, q} = 0$ pour $p > d_V$. Choisissons $n_0$ tel que
$$
\begin{matrix}
\mathcal{H}^{j + 1} & = & H^{j + 1}(K_n), \\
\mathcal{H}^j & = & H^j(K_n), \\
\ldots, \\
\mathcal{H}^{j - d_V - 2} & = & H^{j - d_V - 2}(K_n)
\end{matrix}
$$
pour tout $n \geq n_0$. En comparant les suites spectrales ci-dessus pour
$K_n$ et $K_{n_0}$, nous voyons que, pour $n \geq n_0$, les groupes de
cohomologie $H^{j - 1}(V_i, K_n)$ et $H^j(V_i, K_n)$ sont indépendants de
$n$. Il s'ensuit que le morphisme sur les sections
$H^j(R\lim K_n)(V) \to H^j(K_n)(V)$ est injectif pour $n$ assez grand
(en fonction de $V$) ; voir le Lemme
\ref{lemma-cohomology-derived-limit-injective}. Comme tout objet de
$\mathcal{C}$ peut être recouvert par des éléments de $\mathcal{B}$, nous
en déduisons que le morphisme
$H^j(R\lim K_n) \to \mathcal{H}^j$ est injectif.

\medskip\noindent
La surjectivité se démontre de manière analogue. Choisissons en effet
$U \in \Ob(\mathcal{C})$ et $\gamma \in \mathcal{H}^j(U)$. Nous voulons
relever $\gamma$ en une section de $H^j(R\lim K_n)$ après avoir remplacé
$U$ par les membres d'un recouvrement. D'après la propriété (1) de la
Situation \ref{situation-olsson-laszlo}, nous pouvons donc supposer
$U = V \in \mathcal{B}$. Choisissons $n_0$ tel que
$$
\begin{matrix}
\mathcal{H}^{j + 1} & = & H^{j + 1}(K_n), \\
\mathcal{H}^j & = & H^j(K_n), \\
\ldots, \\
\mathcal{H}^{j - d_V - 2} & = & H^{j - d_V - 2}(K_n)
\end{matrix}
$$
pour tout $n \geq n_0$. Choisissons un élément $\{V_i \to V\}$ de
$\text{Cov}_V$ tel que
$\gamma|_{V_i} \in \mathcal{H}^j(V_i) = H^j(K_{n_0})(V_i)$
se relève en un élément $\gamma_{n_0, i} \in H^j(V_i, K_{n_0})$. Cela est
possible car $H^j(K_{n_0})$ est le faisceau associé au préfaisceau
$U \mapsto H^j(U, K_{n_0})$ d'après le Lemme
\ref{lemma-sheafification-cohomology}. La discussion du premier paragraphe
de la démonstration montre que $H^{j - 1}(V_i, K_n)$ et
$H^j(V_i, K_n)$ sont indépendants de $n \geq n_0$. Ainsi,
$\gamma_{n_0, i}$ se relève en un élément
$\gamma_i \in H^j(V_i, R\lim K_n)$ d'après le Lemme
\ref{lemma-RGamma-commutes-with-Rlim}. Cela achève la démonstration.
\end{proof}

\begin{lemma}
\label{lemma-olsson-laszlo-map-version-one}
\begin{reference}
Il s'agit d'une version de \cite[Lemma 2.1.10]{six-I} avec des hypothèses
légèrement modifiées.
\end{reference}
Soit $f : (\Sh(\mathcal{C}), \mathcal{O}) \to (\Sh(\mathcal{C}'), \mathcal{O}')$
un morphisme de topos annelés. Soient
$\mathcal{A} \subset \textit{Mod}(\mathcal{O})$ et
$\mathcal{A}' \subset \textit{Mod}(\mathcal{O}')$ des sous-catégories de
Serre faibles. Supposons qu'il existe un entier $N$ tel que
\begin{enumerate}
\item $\mathcal{C}, \mathcal{O}, \mathcal{A}$ satisfassent l'hypothèse de
la Situation \ref{situation-olsson-laszlo},
\item $\mathcal{C}', \mathcal{O}', \mathcal{A}'$ satisfassent l'hypothèse
de la Situation \ref{situation-olsson-laszlo},
\item $R^pf_*\mathcal{F} \in \Ob(\mathcal{A}')$ pour
$p \geq 0$ et $\mathcal{F} \in \Ob(\mathcal{A})$,
\item $R^pf_*\mathcal{F} = 0$ pour
$p > N$ et $\mathcal{F} \in \Ob(\mathcal{A})$,
\end{enumerate}
Alors, pour $K$ dans $D_\mathcal{A}(\mathcal{O})$, nous avons
\begin{enumerate}
\item[(a)] $Rf_*K$ appartient à $D_{\mathcal{A}'}(\mathcal{O}')$,
\item[(b)] le morphisme
$H^j(Rf_*K) \to H^j(Rf_*(\tau_{\geq -n}K))$ est un isomorphisme pour
$j \geq N - n$.
\end{enumerate}
\end{lemma}

\begin{proof}
D'après le Lemme \ref{lemma-olsson-laszlo}, nous avons
$K = R\lim \tau_{\geq -n}K$. D'après le Lemme
\ref{lemma-Rf-commutes-with-Rlim}, nous avons
$Rf_*K = R\lim Rf_*\tau_{\geq -n}K$. Les complexes
$Rf_*\tau_{\geq -n}K$ sont bornés inférieurement. La suite spectrale
$$
E_2^{p, q} = R^pf_*H^q(\tau_{\geq -n}K)
$$
convergeant vers $H^{p + q}(Rf_*\tau_{\geq -n}K)$ (Catégories dérivées,
Lemme \ref{derived-lemma-two-ss-complex-functor}) et l'hypothèse (3)
montrent que $Rf_*\tau_{\geq -n}K$ appartient à
$D^+_{\mathcal{A}'}(\mathcal{O}')$ ; voir Homologie, Lemme
\ref{homology-lemma-biregular-ss-converges}. Observons que, pour
$m \geq n$, le morphisme
$$
Rf_*(\tau_{\geq -m}K) \longrightarrow Rf_*(\tau_{\geq -n}K)
$$
induit un isomorphisme sur les faisceaux de cohomologie en degrés
$j \geq -n + N$, d'après les suites spectrales ci-dessus. Nous pouvons
donc conclure en appliquant le Lemme
\ref{lemma-olsson-laszlo-modified}.
\end{proof}

\noindent
Il s'avère que nous avons parfois besoin d'une variante du lemme précédent
où les hypothèses sont légèrement différentes.

\begin{situation}
\label{situation-olsson-laszlo-prime}
Soit
$f : (\mathcal{C}, \mathcal{O}) \to (\mathcal{C}', \mathcal{O}')$
un morphisme de sites annelés. Soit $u : \mathcal{C}' \to \mathcal{C}$ le
foncteur continu de sites correspondant. Soit
$\mathcal{A} \subset \textit{Mod}(\mathcal{O})$ une sous-catégorie de Serre
faible. Nous faisons l'hypothèse suivante : il existe une partie
$\mathcal{B}' \subset \Ob(\mathcal{C}')$ telle que
\begin{enumerate}
\item tout objet de $\mathcal{C}'$ possède un recouvrement dont les membres
appartiennent à $\mathcal{B}'$, et
\item pour tout $V' \in \mathcal{B}'$, il existe un entier $d_{V'}$ et un
système cofinal $\text{Cov}_{V'}$ de recouvrements de $V'$ tels que
$$
H^p(u(V'_i), \mathcal{F}) = 0 \text{ pour }
\{V'_i \to V'\} \in \text{Cov}_{V'},\ p > d_{V'}, \text{ et }
\mathcal{F} \in \Ob(\mathcal{A})
$$
\end{enumerate}
\end{situation}

\begin{lemma}
\label{lemma-olsson-laszlo-map-version-two}
\begin{reference}
Il s'agit d'une version de \cite[Lemma 2.1.10]{six-I} avec des hypothèses
légèrement modifiées.
\end{reference}
Soit $f : (\mathcal{C}, \mathcal{O}) \to (\mathcal{C}', \mathcal{O}')$ un
morphisme de sites annelés. Supposons de plus qu'il existe un entier $N$
tel que
\begin{enumerate}
\item $\mathcal{C}, \mathcal{O}, \mathcal{A}$ satisfassent l'hypothèse de
la Situation \ref{situation-olsson-laszlo},
\item $f : (\mathcal{C}, \mathcal{O}) \to (\mathcal{C}', \mathcal{O}')$
et $\mathcal{A}$ satisfassent l'hypothèse de la Situation
\ref{situation-olsson-laszlo-prime},
\item $R^pf_*\mathcal{F} = 0$ pour
$p > N$ et $\mathcal{F} \in \Ob(\mathcal{A})$,
\end{enumerate}
Alors, pour $K$ dans $D_\mathcal{A}(\mathcal{O})$, le morphisme
$H^j(Rf_*K) \to H^j(Rf_*(\tau_{\geq -n}K))$ est un isomorphisme pour
$j \geq N - n$.
\end{lemma}

\begin{proof}
Soit $K$ dans $D_\mathcal{A}(\mathcal{O})$. D'après le Lemme
\ref{lemma-olsson-laszlo}, nous avons
$K = R\lim \tau_{\geq -n}K$. D'après le Lemme
\ref{lemma-Rf-commutes-with-Rlim}, nous avons
$Rf_*K = R\lim Rf_*(\tau_{\geq -n}K)$. Soit $V' \in \mathcal{B}'$ et soit
$\{V'_i \to V'\}$ un élément de $\text{Cov}_{V'}$. Considérons
$$
H^j(V'_i, Rf_*K) = H^j(u(V'_i), K)
\quad\text{et}\quad
H^j(V'_i, Rf_*(\tau_{\geq -n}K)) = H^j(u(V'_i), \tau_{\geq -n}K)
$$
L'hypothèse de la Situation \ref{situation-olsson-laszlo-prime} implique
que le dernier groupe est indépendant de $n$ pour $n$ assez grand, en
fonction de $j$ et de $d_{V'}$. Nous omettons certains détails. En
appliquant ceci à $j$ et $j - 1$, puis en utilisant le Lemme
\ref{lemma-RGamma-commutes-with-Rlim}, nous obtenons
$$
H^j(V'_i, Rf_*K) = \lim H^j(V'_i, Rf_*(\tau_{\geq -n} K))
$$
et le système de droite est constant pour $n$ supérieur à une constante ne
dépendant que de $d_{V'}$ et de $j$. Le Lemme
\ref{lemma-cohomology-derived-limit-injective} implique donc que
$$
H^j(Rf_*K)(V') \longrightarrow
\left(\lim H^j(Rf_*(\tau_{\geq -n}K))\right)(V')
$$
est injectif. Comme les éléments $V' \in \mathcal{B}'$ recouvrent tout
objet de $\mathcal{C}'$, nous en déduisons que le morphisme
$H^j(Rf_*K) \to \lim H^j(Rf_*(\tau_{\geq -n}K))$ est injectif. La suite
spectrale
$$
E_2^{p, q} = R^pf_*H^q(\tau_{\geq -n}K)
$$
convergeant vers $H^{p + q}(Rf_*(\tau_{\geq -n}K))$ (Catégories dérivées,
Lemme \ref{derived-lemma-two-ss-complex-functor}) et l'hypothèse (3)
montrent que $H^j(Rf_*(\tau_{\geq -n}K))$ est constant pour
$n \geq N - j$. Ainsi,
$H^j(Rf_*K) \to H^j(Rf_*(\tau_{\geq -n}K))$ est injectif pour
$j \geq N - n$.

\medskip\noindent
Nous avons ainsi démontré le lemme en remplaçant « isomorphisme » par
« injectif » dans sa dernière ligne. Choisissons maintenant $j$ et $n$
tels que $j \geq N - n$, puis considérons le triangle distingué
$$
\tau_{\leq -n - 1}K \to K \to \tau_{\geq -n}K \to (\tau_{\leq -n - 1}K)[1]
$$
Voir Catégories dérivées, Remarque
\ref{derived-remark-truncation-distinguished-triangle}.
Comme $\tau_{\geq -n}\tau_{\leq -n -1}K = 0$, l'injectivité déjà démontrée
pour $\tau_{-n - 1}K$ implique
$$
0 = H^j(Rf_*(\tau_{\leq -n - 1}K)) =
H^{j + 1}(Rf_*(\tau_{\leq -n - 1}K)) =
H^{j + 2}(Rf_*(\tau_{\leq -n - 1}K)) = \ldots
$$
La suite exacte longue de cohomologie associée au triangle distingué
$$
Rf_*(\tau_{\leq -n - 1}K) \to Rf_*K \to Rf_*(\tau_{\geq -n}K) \to
Rf_*(\tau_{\leq -n - 1}K)[1]
$$
implique alors que $H^j(Rf_*K) \to H^j(Rf_*(\tau_{\geq -n}K))$ est un
isomorphisme.
\end{proof}








\section{Mayer--Vietoris}
\label{section-mayer-vietoris}

\noindent
Pour l'énoncé et la démonstration usuels de Mayer--Vietoris, voir
Cohomologie, Section \ref{cohomology-section-mayer-vietoris}.

\medskip\noindent
Soit $(\mathcal{C}, \mathcal{O})$ un site annelé. Considérons un diagramme
commutatif
$$
\xymatrix{
E \ar[d] \ar[r] & Y \ar[d] \\
Z \ar[r] & X
}
$$
dans la catégorie $\mathcal{C}$. Dans cette situation, étant donné un objet
$K$ de $D(\mathcal{O})$, nous obtenons ce qui ressemble au début d'un
triangle distingué
$$
R\Gamma(X, K) \to
R\Gamma(Z, K) \oplus
R\Gamma(Y, K) \to
R\Gamma(E, K)
$$
Le lemme suivant précise cette observation.

\begin{lemma}
\label{lemma-c-square}
Dans la situation ci-dessus, choisissons un complexe K-injectif
$\mathcal{I}^\bullet$ de $\mathcal{O}$-modules représentant $K$. En
multipliant par $-1$ le morphisme canonique pour l'une des quatre flèches,
nous obtenons des morphismes de complexes
$$
\mathcal{I}^\bullet(X) \xrightarrow{\alpha}
\mathcal{I}^\bullet(Z) \oplus
\mathcal{I}^\bullet(Y) \xrightarrow{\beta}
\mathcal{I}^\bullet(E)
$$
avec $\beta \circ \alpha = 0$. Nous obtenons donc un morphisme canonique
$$
c^K_{X, Z, Y, E} :
\mathcal{I}^\bullet(X)
\longrightarrow
C(\beta)^\bullet[-1]
$$
Ce morphisme est canonique au sens où un autre choix de complexe
K-injectif représentant $K$ détermine une flèche isomorphe dans la
catégorie dérivée des groupes abéliens. Si $c^K_{X, Z, Y, E}$ est un
isomorphisme, son inverse fournit un triangle distingué canonique
$$
R\Gamma(X, K) \to
R\Gamma(Z, K) \oplus
R\Gamma(Y, K) \to
R\Gamma(E, K) \to
R\Gamma(X, K)[1]
$$
Toutes ces constructions sont fonctorielles en $K$.
\end{lemma}

\begin{proof}
Le lemme se démontre de lui-même. Par exemple, si
$\mathcal{J}^\bullet$ est un second complexe K-injectif représentant $K$,
nous pouvons choisir un quasi-isomorphisme
$\mathcal{I}^\bullet \to \mathcal{J}^\bullet$ qui détermine des
quasi-isomorphismes entre tous les complexes en jeu. Nous omettons les
détails. Pour la construction des cônes et leur relation avec les triangles
distingués, voir Catégories dérivées, Sections
\ref{derived-section-cones} et
\ref{derived-section-homotopy-triangulated}.
\end{proof}

\begin{lemma}
\label{lemma-two-out-of-three-blow-up-square}
Dans la situation ci-dessus, soit $K_1 \to K_2 \to K_3 \to K_1[1]$ un
triangle distingué dans $D(\mathcal{O})$. Si $c^{K_i}_{X, Z, Y, E}$ est un
quasi-isomorphisme pour deux valeurs de $i$ parmi $\{1, 2, 3\}$, il en est
de même pour la troisième valeur de $i$.
\end{lemma}

\begin{proof}
En faisant tourner le triangle, nous pouvons supposer que
$c^{K_1}_{X, Z, Y, E}$ et $c^{K_2}_{X, Z, Y, E}$ sont des
quasi-isomorphismes. Choisissons un morphisme
$f : \mathcal{I}^\bullet_1 \to \mathcal{I}^\bullet_2$ de complexes
K-injectifs de $\mathcal{O}$-modules représentant $K_1 \to K_2$. Alors
$K_3$ est représenté par le complexe K-injectif $C(f)^\bullet$ ; voir
Catégories dérivées, Lemme \ref{derived-lemma-triangle-K-injective}. Le
morphisme $c^{K_3}_{X, Z, Y, E}$ est donc un isomorphisme, car il est le
troisième côté d'un morphisme de triangles distingués dans
$K(\textit{Ab})$ dont les deux autres côtés sont des quasi-isomorphismes.
Nous omettons certains détails ; utiliser Catégories dérivées, Lemme
\ref{derived-lemma-third-isomorphism-triangle}.
\end{proof}

\noindent
Donnons un critère assurant que cette construction produit effectivement
un triangle distingué.

\begin{lemma}
\label{lemma-square-triangle}
Dans la situation ci-dessus, supposons que
\begin{enumerate}
\item $h_X^\# = h_Y^\# \amalg_{h_E^\#} h_Z^\#$, et
\item $h_E^\# \to h_Y^\#$ soit injectif.
\end{enumerate}
Alors la construction du Lemme \ref{lemma-c-square} produit un triangle
distingué
$$
R\Gamma(X, K) \to
R\Gamma(Z, K) \oplus
R\Gamma(Y, K) \to
R\Gamma(E, K) \to R\Gamma(X, K)[1]
$$
fonctoriel en $K$ dans $D(\mathcal{C})$.
\end{lemma}

\begin{proof}
Nous pouvons représenter $K$ par un complexe K-injectif dont les termes
sont des faisceaux abéliens injectifs ; voir la Section
\ref{section-unbounded}. Il suffit donc de montrer que, si $\mathcal{I}$
est un faisceau abélien injectif, alors
$$
0 \to \mathcal{I}(X) \to
\mathcal{I}(Z) \oplus \mathcal{I}(Y) \to
\mathcal{I}(E) \to 0
$$
est une suite exacte courte. La première flèche est injective car, d'après
la condition (1), le morphisme $h_Y \amalg h_Z \to h_X$ devient surjectif
après faisceautisation, ce qui signifie que
$\{Y \to X, Z \to X\}$ peut être raffiné par un recouvrement de $X$. La
dernière flèche est surjective car
$\mathcal{I}(Y) \to \mathcal{I}(E)$ est surjectif. En effet, nous avons
$\mathcal{I}(E) = \Hom(\mathbf{Z}_E^\#, \mathcal{I})$,
$\mathcal{I}(Y) = \Hom(\mathbf{Z}_Y^\#, \mathcal{I})$,
le morphisme $\mathbf{Z}_E^\# \to \mathbf{Z}_Y^\#$ est injectif d'après
(2), et $\mathcal{I}$ est un faisceau abélien injectif. Comparer à Modules
sur les sites, Section
\ref{sites-modules-section-free-abelian-sheaf}.
Enfin, supposons donnés $s \in \mathcal{I}(Y)$ et
$t \in \mathcal{F}(Z)$ ayant la même image dans $\mathcal{I}(E)$. Alors
$s$ et $t$ définissent un morphisme
$$
s \amalg t : h_Y^\# \amalg h_Z^\# \longrightarrow \mathcal{I}
$$
qui, par hypothèse, se factorise par
$h_Y^\# \amalg_{h_E^\#} h_Z^\#$. L'hypothèse (1) fournit donc un unique
morphisme $h_X^\# \to \mathcal{I}$, correspondant à un élément de
$\mathcal{I}(X)$ qui se restreint à $s$ sur $Y$ et à $t$ sur $Z$.
\end{proof}

\begin{lemma}
\label{lemma-square-triangle-general}
Soit $\mathcal{C}$ un site. Considérons un diagramme commutatif
$$
\xymatrix{
\mathcal{D} \ar[r] \ar[d] & \mathcal{F} \ar[d] \\
\mathcal{E} \ar[r] & \mathcal{G}
}
$$
de préfaisceaux d'ensembles sur $\mathcal{C}$ et supposons que
\begin{enumerate}
\item $\mathcal{G}^\# =
\mathcal{E}^\# \amalg_{\mathcal{D}^\#} \mathcal{F}^\#$, et
\item $\mathcal{D}^\# \to \mathcal{F}^\#$ soit injectif.
\end{enumerate}
Alors il existe un triangle distingué canonique
$$
R\Gamma(\mathcal{G}, K) \to
R\Gamma(\mathcal{E}, K) \oplus
R\Gamma(\mathcal{F}, K) \to
R\Gamma(\mathcal{D}, K) \to
R\Gamma(\mathcal{G}, K)[1]
$$
fonctoriel en $K \in D(\mathcal{C})$, où $R\Gamma(\mathcal{G}, -)$ désigne
la cohomologie étudiée dans la Section \ref{section-limp}.
\end{lemma}

\begin{proof}
Comme la faisceautisation est exacte et que
$R\Gamma(\mathcal{G}, -) = R\Gamma(\mathcal{G}^\#, -)$, nous pouvons
supposer que $\mathcal{D}, \mathcal{E}, \mathcal{F}, \mathcal{G}$ sont des
faisceaux d'ensembles. De plus, la cohomologie
$R\Gamma(\mathcal{G}, -)$ ne dépend que du topos, et non du site qui le
présente. D'après Sites, Lemme \ref{sites-lemma-topos-good-site}, nous
pouvons donc remplacer $\mathcal{C}$ par un site « plus grand », muni d'une
topologie sous-canonique, tel que $\mathcal{G} = h_X$,
$\mathcal{F} = h_Y$, $\mathcal{E} = h_Z$ et $\mathcal{D} = h_E$ pour
certains objets $X, Y, Z, E$ de $\mathcal{C}$. Dans ce cas, le résultat
découle du Lemme \ref{lemma-square-triangle}.
\end{proof}







\section{Comparaison de deux topologies}
\label{section-compare}

\noindent
Soit $\mathcal{C}$ une catégorie. Soient
$\text{Cov}(\mathcal{C}) \supset \text{Cov}'(\mathcal{C})$
deux manières de munir $\mathcal{C}$ d'une structure de site. Notons
$\tau$ la topologie correspondant à $\text{Cov}(\mathcal{C})$ et $\tau'$
celle correspondant à $\text{Cov}'(\mathcal{C})$. Le foncteur identité de
$\mathcal{C}$ définit alors un morphisme de sites
$$
\epsilon : \mathcal{C}_\tau \longrightarrow \mathcal{C}_{\tau'}
$$
où $\epsilon_*$ est le foncteur identité sur les préfaisceaux sous-jacents,
et où $\epsilon^{-1}$ est la faisceautisation pour $\tau$ d'un
$\tau'$-faisceau. Voir Sites, Exemples
\ref{sites-example-finer-topology} et
\ref{sites-example-finer-topology-bis}. Dans la situation ci-dessus, nous
avons les propriétés suivantes :
\begin{enumerate}
\item $\epsilon_* : \Sh(\mathcal{C}_\tau) \to \Sh(\mathcal{C}_{\tau'})$
est pleinement fidèle et $\epsilon^{-1} \circ \epsilon_* = \text{id}$,
\item $\epsilon_* : \textit{Ab}(\mathcal{C}_\tau) \to
\textit{Ab}(\mathcal{C}_{\tau'})$
est pleinement fidèle et $\epsilon^{-1} \circ \epsilon_* = \text{id}$,
\item $R\epsilon_* : D(\mathcal{C}_\tau) \to D(\mathcal{C}_{\tau'})$
est pleinement fidèle et $\epsilon^{-1} \circ R\epsilon_* = \text{id}$,
\item si $\mathcal{O}$ est un faisceau d'anneaux pour la topologie $\tau$,
alors $\mathcal{O}$ est également un faisceau pour la topologie $\tau'$ et
$\epsilon$ devient un morphisme plat de sites annelés
$$
\epsilon :
(\mathcal{C}_\tau, \mathcal{O}_\tau)
\longrightarrow
(\mathcal{C}_{\tau'}, \mathcal{O}_{\tau'})
$$
\item $\epsilon_* : \textit{Mod}(\mathcal{O}_\tau) \to
\textit{Mod}(\mathcal{O}_{\tau'})$ est pleinement fidèle et
$\epsilon^* \circ \epsilon_* = \text{id}$,
\item $R\epsilon_* : D(\mathcal{O}_\tau) \to D(\mathcal{O}_{\tau'})$
est pleinement fidèle et $\epsilon^* \circ R\epsilon_* = \text{id}$.
\end{enumerate}
Voici quelques explications.

\medskip\noindent
Pour (1). Soit $\mathcal{F}$ un faisceau d'ensembles pour la topologie
$\tau$. Alors $\epsilon_*\mathcal{F}$ n'est autre que $\mathcal{F}$
considéré comme faisceau pour la topologie $\tau'$. Appliquer
$\epsilon^{-1}$ revient à prendre le $\tau$-faisceau associé à
$\mathcal{F}$, ce qui ne change rien puisque $\mathcal{F}$ est déjà un
$\tau$-faisceau. Ainsi,
$\epsilon^{-1}(\epsilon_*\mathcal{F})) = \mathcal{F}$. La pleine fidélité
résulte de Catégories, Lemme
\ref{categories-lemma-adjoint-fully-faithful}.

\medskip\noindent
Pour (2). C'est une conséquence de (1), car prendre l'image inverse ou
l'image directe de faisceaux abéliens revient à effectuer ces opérations
sur les faisceaux d'ensembles sous-jacents.

\medskip\noindent
Pour (3). Soit $K$ un objet de $D(\mathcal{C}_\tau)$. Pour calculer
$R\epsilon_*K$, choisissons un complexe K-injectif
$\mathcal{I}^\bullet$ représentant $K$ et posons
$R\epsilon_*K = \epsilon_*\mathcal{I}^\bullet$. Comme
$\epsilon^{-1} : D(\mathcal{C}_{\tau'}) \to D(\mathcal{C}_\tau)$ se
calcule sur un objet $L$ en appliquant le foncteur exact $\epsilon^{-1}$ à
n'importe quel complexe de faisceaux abéliens représentant $L$, nous
voyons que $\epsilon^{-1}R\epsilon_*K$ est représenté par
$\epsilon^{-1}\epsilon_*\mathcal{I}^\bullet$. D'après (1), nous avons
$\mathcal{I}^\bullet = \epsilon^{-1}\epsilon_*\mathcal{I}^\bullet$.
Autrement dit, $\epsilon^{-1} \circ R\epsilon_* = \text{id}$, et nous
concluons comme précédemment.

\medskip\noindent
Pour (4). Observons que
$\epsilon^{-1}\mathcal{O}_{\tau'} = \mathcal{O}_\tau$ ; voir la discussion
de (1). Ainsi, $\epsilon$ est un morphisme plat de sites annelés ; voir
Modules sur les sites, Définition
\ref{sites-modules-definition-flat-morphism}. De plus, il est clair que
$\epsilon^* = \epsilon^{-1}$ sur les $\mathcal{O}_{\tau'}$-modules
(l'image inverse comme module a le même faisceau abélien sous-jacent que
l'image inverse du faisceau abélien sous-jacent).

\medskip\noindent
Pour (5). Cela résulte clairement de (2) et de ce que nous avons dit en (4).

\medskip\noindent
Pour (6). L'argument est analogue à celui de (3). Nous omettons les détails.












\section{Formalismes de la descente cohomologique}
\label{section-formal-cohomological-descent}

\noindent
Dans cette section, nous examinons seulement dans quelle mesure un
morphisme de topos annelés détermine un plongement de la catégorie dérivée
de la base dans celle de la source. Voici un résultat typique.

\begin{lemma}
\label{lemma-downstairs}
Soit $f : (\Sh(\mathcal{C}), \mathcal{O}_\mathcal{C}) \to
(\Sh(\mathcal{D}), \mathcal{O}_\mathcal{D})$ un morphisme de topos annelés.
Considérons la sous-catégorie pleine
$D' \subset D(\mathcal{O}_\mathcal{D})$ constituée des objets $K$ tels que
$$
K \longrightarrow Rf_*Lf^*K
$$
soit un isomorphisme. Alors $D'$ est une sous-catégorie strictement pleine,
triangulée et saturée de $D(\mathcal{O}_\mathcal{D})$, et le foncteur
$Lf^* : D' \to D(\mathcal{O}_\mathcal{C})$ est pleinement fidèle.
\end{lemma}

\begin{proof}
Voir Catégories dérivées, Définition \ref{derived-definition-saturated},
pour la définition de « saturée » dans ce contexte, et Catégories dérivées,
Lemme \ref{derived-lemma-triangulated-subcategory}, pour une discussion des
sous-catégories triangulées. Le morphisme canonique du lemme est l'unité
du couple de foncteurs adjoints $(Lf^*, Rf_*)$ ; voir le Lemme
\ref{lemma-adjoint}. Cela étant dit, nous omettons la démonstration du fait
que $D'$ est une sous-catégorie triangulée saturée : elle résulte
formellement de l'exactitude des foncteurs $Lf^*$ et $Rf_*$ entre
catégories triangulées. La dernière assertion résulte formellement de
l'adjonction de $Lf^*$ et $Rf_*$ ; comparer à Catégories, Lemme
\ref{categories-lemma-adjoint-fully-faithful}.
\end{proof}

\begin{lemma}
\label{lemma-upstairs}
Soit $f : (\Sh(\mathcal{C}), \mathcal{O}_\mathcal{C}) \to
(\Sh(\mathcal{D}), \mathcal{O}_\mathcal{D})$ un morphisme de topos annelés.
Considérons la sous-catégorie pleine
$D' \subset D(\mathcal{O}_\mathcal{C})$ constituée des objets $K$ tels que
$$
Lf^*Rf_*K \longrightarrow K
$$
soit un isomorphisme. Alors $D'$ est une sous-catégorie strictement pleine,
triangulée et saturée de $D(\mathcal{O}_\mathcal{C})$, et le foncteur
$Rf_* : D' \to D(\mathcal{O}_\mathcal{D})$ est pleinement fidèle.
\end{lemma}

\begin{proof}
Voir Catégories dérivées, Définition \ref{derived-definition-saturated},
pour la définition de « saturée » dans ce contexte, et Catégories dérivées,
Lemme \ref{derived-lemma-triangulated-subcategory}, pour une discussion des
sous-catégories triangulées. Le morphisme canonique du lemme est la
co-unité du couple de foncteurs adjoints $(Lf^*, Rf_*)$ ; voir le Lemme
\ref{lemma-adjoint}. Cela étant dit, nous omettons la démonstration du fait
que $D'$ est une sous-catégorie triangulée saturée : elle résulte
formellement de l'exactitude des foncteurs $Lf^*$ et $Rf_*$ entre
catégories triangulées. La dernière assertion résulte formellement de
l'adjonction de $Lf^*$ et $Rf_*$ ; comparer à Catégories, Lemme
\ref{categories-lemma-adjoint-fully-faithful}.
\end{proof}

\begin{lemma}
\label{lemma-bounded-in-image-upstairs}
Soit $f : (\Sh(\mathcal{C}), \mathcal{O}_\mathcal{C}) \to
(\Sh(\mathcal{D}), \mathcal{O}_\mathcal{D})$ un morphisme de topos annelés.
Soit $K$ un objet de $D(\mathcal{O}_\mathcal{C})$. Supposons que
\begin{enumerate}
\item $f$ soit plat,
\item $K$ soit borné inférieurement,
\item $f^*Rf_*H^q(K) \to H^q(K)$ soit un isomorphisme.
\end{enumerate}
Alors $f^*Rf_*K \to K$ est un isomorphisme.
\end{lemma}

\begin{proof}
Observons que $f^*Rf_*K \to K$ est un isomorphisme si et seulement s'il
induit un isomorphisme sur les faisceaux de cohomologie $H^j$. Observons
que
$H^j(f^*Rf_*K) = f^*H^j(Rf_*K) = f^*H^j(Rf_*\tau_{\leq j}K) =
H^j(f^*Rf_*\tau_{\leq j}K)$.
Nous pouvons donc supposer que $K$ est borné. La propriété (3) montre alors
que les faisceaux de cohomologie appartiennent à la sous-catégorie
triangulée $D' \subset D(\mathcal{O}_\mathcal{C})$ du Lemme
\ref{lemma-upstairs}. Par conséquent, $K$ y appartient également.
\end{proof}

\begin{lemma}
\label{lemma-bounded-in-image-downstairs}
Soit $f : (\Sh(\mathcal{C}), \mathcal{O}_\mathcal{C}) \to
(\Sh(\mathcal{D}), \mathcal{O}_\mathcal{D})$ un morphisme de topos annelés.
Soit $K$ un objet de $D(\mathcal{O}_\mathcal{D})$. Supposons que
\begin{enumerate}
\item $f$ soit plat,
\item $K$ soit borné inférieurement,
\item $H^q(K) \to Rf_*f^*H^q(K)$ soit un isomorphisme.
\end{enumerate}
Alors $K \to Rf_*f^*K$ est un isomorphisme.
\end{lemma}

\begin{proof}
Observons que $K \to Rf_*f^*K$ est un isomorphisme si et seulement s'il
induit un isomorphisme sur les faisceaux de cohomologie $H^j$. Observons
que
$H^j(Rf_*f^*K) = H^j(Rf_*\tau_{\leq j}f^*K) = H^j(Rf_*f^*\tau_{\leq j}K)$.
Nous pouvons donc supposer que $K$ est borné. La propriété (3) montre alors
que les faisceaux de cohomologie appartiennent à la sous-catégorie
triangulée $D' \subset D(\mathcal{O}_\mathcal{D})$ du Lemme
\ref{lemma-downstairs}. Par conséquent, $K$ y appartient également.
\end{proof}

\begin{lemma}
\label{lemma-equivalence-bounded}
Soit $f : (\Sh(\mathcal{C}), \mathcal{O}) \to (\Sh(\mathcal{C}'), \mathcal{O}')$
un morphisme de topos annelés. Soient
$\mathcal{A} \subset \textit{Mod}(\mathcal{O})$ et
$\mathcal{A}' \subset \textit{Mod}(\mathcal{O}')$ des sous-catégories de
Serre faibles. Supposons que
\begin{enumerate}
\item $f$ soit plat,
\item $f^*$ induise une équivalence de catégories
$\mathcal{A}' \to \mathcal{A}$,
\item $\mathcal{F}' \to Rf_*f^*\mathcal{F}'$ soit un isomorphisme pour
$\mathcal{F}' \in \Ob(\mathcal{A}')$.
\end{enumerate}
Alors
$f^* : D_{\mathcal{A}'}^+(\mathcal{O}') \to D_\mathcal{A}^+(\mathcal{O})$
est une équivalence de catégories dont un quasi-inverse est
$Rf_* : D_\mathcal{A}^+(\mathcal{O}) \to D_{\mathcal{A}'}^+(\mathcal{O}')$.
\end{lemma}

\begin{proof}
D'après les hypothèses (2) et (3) et les Lemmes
\ref{lemma-bounded-in-image-downstairs} et \ref{lemma-downstairs}, nous
voyons que
$f^* : D_{\mathcal{A}'}^+(\mathcal{O}') \to D_\mathcal{A}^+(\mathcal{O})$
est pleinement fidèle. Soit $\mathcal{F} \in \Ob(\mathcal{A})$. Nous
pouvons écrire $\mathcal{F} = f^*\mathcal{F}'$. Alors
$Rf_*\mathcal{F} = Rf_* f^*\mathcal{F}' = \mathcal{F}'$.
En particulier, $R^pf_*\mathcal{F} = 0$ pour $p > 0$ et
$f_*\mathcal{F} \in \Ob(\mathcal{A}')$. Ainsi, pour tout
$K \in D^+_\mathcal{A}(\mathcal{O})$, la suite spectrale
$E_2^{p, q} = R^pf_*H^q(K)$ convergeant vers $R^{p + q}f_*K$ montre que
$Rf_*K$ appartient à $D^+_{\mathcal{A}'}(\mathcal{O}')$. Les Lemmes
\ref{lemma-bounded-in-image-upstairs} et \ref{lemma-upstairs} montrent
également que
$Rf_* : D_\mathcal{A}^+(\mathcal{O}) \to D_{\mathcal{A}'}^+(\mathcal{O}')$
est pleinement fidèle. Comme $f^*$ et $Rf_*$ sont adjoints, nous obtenons
alors le résultat du lemme, par exemple à l'aide de Catégories, Lemme
\ref{categories-lemma-adjoint-fully-faithful}.
\end{proof}

\begin{lemma}
\label{lemma-equivalence-unbounded-one}
\begin{reference}
C'est l'analogue de \cite[Theorem 2.2.3]{six-I}.
\end{reference}
Soit $f : (\Sh(\mathcal{C}), \mathcal{O}) \to (\Sh(\mathcal{C}'), \mathcal{O}')$
un morphisme de topos annelés. Soient
$\mathcal{A} \subset \textit{Mod}(\mathcal{O})$ et
$\mathcal{A}' \subset \textit{Mod}(\mathcal{O}')$ des sous-catégories de
Serre faibles. Supposons que
\begin{enumerate}
\item $f$ soit plat,
\item $f^*$ induise une équivalence de catégories
$\mathcal{A}' \to \mathcal{A}$,
\item $\mathcal{F}' \to Rf_*f^*\mathcal{F}'$ soit un isomorphisme pour
$\mathcal{F}' \in \Ob(\mathcal{A}')$,
\item $\mathcal{C}, \mathcal{O}, \mathcal{A}$ satisfassent l'hypothèse de
la Situation \ref{situation-olsson-laszlo},
\item $\mathcal{C}', \mathcal{O}', \mathcal{A}'$ satisfassent l'hypothèse
de la Situation \ref{situation-olsson-laszlo}.
\end{enumerate}
Alors $f^* : D_{\mathcal{A}'}(\mathcal{O}') \to D_\mathcal{A}(\mathcal{O})$
est une équivalence de catégories dont un quasi-inverse est
$Rf_* : D_\mathcal{A}(\mathcal{O}) \to D_{\mathcal{A}'}(\mathcal{O}')$.
\end{lemma}

\begin{proof}
Comme $f^*$ est exact, il est clair que $f^*$ définit un foncteur
$f^* : D_{\mathcal{A}'}(\mathcal{O}') \to D_\mathcal{A}(\mathcal{O})$
comme dans l'énoncé du lemme et que ce foncteur commute aux foncteurs de
troncation $\tau_{\geq -n}$. Nous savons déjà que $f^*$ et $Rf_*$ sont des
équivalences quasi-inverses sur les catégories bornées inférieurement
correspondantes ; voir le Lemme \ref{lemma-equivalence-bounded}. Le Lemme
\ref{lemma-olsson-laszlo-map-version-one}, avec $N = 0$, montre que $Rf_*$
définit bien un foncteur
$Rf_* : D_\mathcal{A}(\mathcal{O}) \to D_{\mathcal{A}'}(\mathcal{O}')$
et que ce foncteur commute en outre aux foncteurs de troncation
$\tau_{\geq -n}$. Ainsi, pour $K$ dans $D_\mathcal{A}(\mathcal{O})$, le
morphisme $f^*Rf_*K \to K$ est un isomorphisme, puisque c'est le cas sur
les troncations. De même, pour $K'$ dans
$D_{\mathcal{A}'}(\mathcal{O}')$, le morphisme $K' \to Rf_*f^*K'$ est un
isomorphisme, puisque c'est le cas sur les troncations. Cela achève la
démonstration.
\end{proof}

\begin{lemma}
\label{lemma-equivalence-unbounded-two}
\begin{reference}
C'est l'analogue de \cite[Theorem 2.2.3]{six-I}.
\end{reference}
Soit $f : (\mathcal{C}, \mathcal{O}) \to (\mathcal{C}', \mathcal{O}')$ un
morphisme de sites annelés. Soient
$\mathcal{A} \subset \textit{Mod}(\mathcal{O})$ et
$\mathcal{A}' \subset \textit{Mod}(\mathcal{O}')$ des sous-catégories de
Serre faibles. Supposons que
\begin{enumerate}
\item $f$ soit plat,
\item $f^*$ induise une équivalence de catégories
$\mathcal{A}' \to \mathcal{A}$,
\item $\mathcal{F}' \to Rf_*f^*\mathcal{F}'$ soit un isomorphisme pour
$\mathcal{F}' \in \Ob(\mathcal{A}')$,
\item $\mathcal{C}, \mathcal{O}, \mathcal{A}$ satisfassent l'hypothèse de
la Situation \ref{situation-olsson-laszlo},
\item $f : (\mathcal{C}, \mathcal{O}) \to (\mathcal{C}', \mathcal{O}')$
et $\mathcal{A}$ satisfassent l'hypothèse de la Situation
\ref{situation-olsson-laszlo-prime}.
\end{enumerate}
Alors $f^* : D_{\mathcal{A}'}(\mathcal{O}') \to D_\mathcal{A}(\mathcal{O})$
est une équivalence de catégories dont un quasi-inverse est
$Rf_* : D_\mathcal{A}(\mathcal{O}) \to D_{\mathcal{A}'}(\mathcal{O}')$.
\end{lemma}

\begin{proof}
La démonstration de ce lemme est exactement la même que celle du Lemme
\ref{lemma-equivalence-unbounded-one}, à ceci près que la référence au
Lemme \ref{lemma-olsson-laszlo-map-version-one} est remplacée par une
référence au Lemme \ref{lemma-olsson-laszlo-map-version-two}.
\end{proof}








\section{Comparaison de deux topologies, II}
\label{section-compare-II}

\noindent
Soit $\mathcal{C}$ une catégorie. Soient
$\text{Cov}(\mathcal{C}) \supset \text{Cov}'(\mathcal{C})$
deux manières de munir $\mathcal{C}$ d'une structure de site. Notons
$\tau$ la topologie correspondant à $\text{Cov}(\mathcal{C})$ et $\tau'$
celle correspondant à $\text{Cov}'(\mathcal{C})$. Le foncteur identité de
$\mathcal{C}$ définit alors un morphisme de sites
$$
\epsilon : \mathcal{C}_\tau \longrightarrow \mathcal{C}_{\tau'}
$$
où $\epsilon_*$ est le foncteur identité sur les préfaisceaux sous-jacents,
et où $\epsilon^{-1}$ est la faisceautisation pour $\tau$ d'un
$\tau'$-faisceau (elle est donc manifestement exacte). Soit
$\mathcal{O}$ un faisceau d'anneaux pour la topologie $\tau$. Alors
$\mathcal{O}$ est également un faisceau pour la topologie $\tau'$ et
$\epsilon$ devient un morphisme de sites annelés
$$
\epsilon :
(\mathcal{C}_\tau, \mathcal{O}_\tau)
\longrightarrow
(\mathcal{C}_{\tau'}, \mathcal{O}_{\tau'})
$$
Pour plus de détails, voir la Section \ref{section-compare}.

\begin{lemma}
\label{lemma-compare-topologies-derived-adequate-modules}
Soit $\epsilon : (\mathcal{C}_\tau, \mathcal{O}_\tau) \to
(\mathcal{C}_{\tau'}, \mathcal{O}_{\tau'})$ comme ci-dessus. Soit
$\mathcal{B} \subset \Ob(\mathcal{C})$ une partie. Soit
$\mathcal{A} \subset \textit{PMod}(\mathcal{O})$ une sous-catégorie
pleine. Supposons que
\begin{enumerate}
\item tout objet de $\mathcal{A}$ soit un faisceau pour la topologie $\tau$,
\item $\mathcal{A}$ soit une sous-catégorie de Serre faible de
$\textit{Mod}(\mathcal{O}_\tau)$,
\item tout objet de $\mathcal{C}$ possède un $\tau'$-recouvrement dont les
membres appartiennent à $\mathcal{B}$, et
\item pour tout $U \in \mathcal{B}$, nous ayons
$H^p_\tau(U, \mathcal{F}) = 0$ pour $p > 0$ et tout
$\mathcal{F} \in \mathcal{A}$.
\end{enumerate}
Alors $\mathcal{A}$ est une sous-catégorie de Serre faible de
$\textit{Mod}(\mathcal{O}_{\tau'})$ et il existe une équivalence de
catégories triangulées
$D_\mathcal{A}(\mathcal{O}_\tau) = D_\mathcal{A}(\mathcal{O}_{\tau'})$
donnée par $\epsilon^*$ et $R\epsilon_*$.
\end{lemma}

\begin{proof}
Comme $\epsilon^{-1}\mathcal{O}_{\tau'} = \mathcal{O}_\tau$, nous voyons
que $\epsilon$ est un morphisme plat de sites annelés et que, de fait,
$\epsilon^{-1} = \epsilon^*$ sur les faisceaux de modules. La propriété
(1) permet de considérer tout objet de $\mathcal{A}$ à la fois comme un
faisceau de $\mathcal{O}_\tau$-modules et comme un faisceau de
$\mathcal{O}_{\tau'}$-modules. Autrement dit, nous avons des foncteurs
d'inclusion pleinement fidèles
$$
\mathcal{A} \to \textit{Mod}(\mathcal{O}_\tau) \to
\textit{Mod}(\mathcal{O}_{\tau'})
$$
Pour éviter toute confusion, nous noterons
$\mathcal{A}' \subset \textit{Mod}(\mathcal{O}_{\tau'})$
les images des objets de $\mathcal{A}$. Il est alors clair que
$\epsilon_* : \mathcal{A} \to \mathcal{A}'$ et
$\epsilon^* : \mathcal{A}' \to \mathcal{A}$ sont
des équivalences quasi-inverses (voir la discussion précédant le lemme et
utiliser le fait que les objets de $\mathcal{A}'$ sont des faisceaux pour
la topologie $\tau$).

\medskip\noindent
Les conditions (3) et (4) impliquent que
$R^p\epsilon_*\mathcal{F} = 0$ pour $p > 0$ et
$\mathcal{F} \in \Ob(\mathcal{A})$. En effet, $R^p\epsilon_*$ est le
faisceau associé au préfaisceau $U \mapsto H^p_\tau(U, \mathcal{F})$ ;
voir le Lemme \ref{lemma-higher-direct-images}. Ainsi, tout complexe exact
dans $\mathcal{A}$ (ce qui revient à un complexe exact dans
$\textit{Mod}(\mathcal{O}_\tau)$ dont les termes appartiennent à
$\mathcal{A}$ ; voir Homologie, Lemme
\ref{homology-lemma-characterize-weak-serre-subcategory}) reste exact après
application du foncteur $\epsilon_*$.

\medskip\noindent
Considérons une suite exacte
$$
\mathcal{F}'_0 \to \mathcal{F}'_1 \to
\mathcal{F}'_2 \to \mathcal{F}'_3 \to
\mathcal{F}'_4
$$
dans $\textit{Mod}(\mathcal{O}_{\tau'})$, où
$\mathcal{F}'_0, \mathcal{F}'_1, \mathcal{F}'_3, \mathcal{F}'_4$
appartiennent à $\mathcal{A}'$. L'application du foncteur exact
$\epsilon^*$ donne une suite exacte
$$
\epsilon^*\mathcal{F}'_0 \to \epsilon^*\mathcal{F}'_1 \to
\epsilon^*\mathcal{F}'_2 \to \epsilon^*\mathcal{F}'_3 \to
\epsilon^*\mathcal{F}'_4
$$
dans $\textit{Mod}(\mathcal{O}_\tau)$. Comme $\mathcal{A}$ est une
sous-catégorie de Serre faible et que
$\epsilon^*\mathcal{F}'_0, \epsilon^*\mathcal{F}'_1,
\epsilon^*\mathcal{F}'_3, \epsilon^*\mathcal{F}'_4$
appartiennent à $\mathcal{A}$, nous déduisons d'Homologie, Définition
\ref{homology-definition-serre-subcategory}, que
$\epsilon^*\mathcal{F}'_2$ appartient à $\mathcal{A}$. Considérons le
morphisme de suites
$$
\xymatrix{
\mathcal{F}'_0 \ar[r] \ar[d] &
\mathcal{F}'_1 \ar[r] \ar[d] &
\mathcal{F}'_2 \ar[r] \ar[d] &
\mathcal{F}'_3 \ar[r] \ar[d] &
\mathcal{F}'_4 \ar[d] \\
\epsilon_*\epsilon^*\mathcal{F}'_0 \ar[r] &
\epsilon_*\epsilon^*\mathcal{F}'_1 \ar[r] &
\epsilon_*\epsilon^*\mathcal{F}'_2 \ar[r] &
\epsilon_*\epsilon^*\mathcal{F}'_3 \ar[r] &
\epsilon_*\epsilon^*\mathcal{F}'_4
}
$$
La ligne inférieure est exacte d'après la discussion du paragraphe
précédent. Les flèches verticales d'indices $0$, $1$, $3$, $4$ sont des
isomorphismes d'après la discussion du premier paragraphe. Le lemme des $5$
(Homologie, Lemme \ref{homology-lemma-five-lemma}) donne
$\mathcal{F}'_2 \cong \epsilon_*\epsilon^*\mathcal{F}'_2$, et donc
$\mathcal{F}'_2$ appartient à $\mathcal{A}'$. Nous voyons ainsi que
$\mathcal{A}'$ est une sous-catégorie de Serre faible de
$\textit{Mod}(\mathcal{O}_{\tau'})$ ; voir Homologie, Définition
\ref{homology-definition-serre-subcategory}.

\medskip\noindent
À ce stade, il est légitime de parler des catégories dérivées
$D_\mathcal{A}(\mathcal{O}_\tau)$ et
$D_{\mathcal{A}'}(\mathcal{O}_{\tau'})$ ; voir Catégories dérivées,
Section \ref{derived-section-triangulated-sub}. Pour achever la
démonstration, montrons que les conditions (1) -- (5) du Lemme
\ref{lemma-equivalence-unbounded-two} s'appliquent. Nous avons déjà établi
(1), (2) et (3). Comme tout objet possède un $\tau'$-recouvrement par des
objets de $\mathcal{B}$, il possède a fortiori un $\tau$-recouvrement par
des objets de $\mathcal{B}$. La condition (4) du Lemme
\ref{lemma-equivalence-unbounded-two} est donc satisfaite. La condition (5)
l'est de même.
\end{proof}

\begin{lemma}
\label{lemma-descent-squares}
Soit $\epsilon : (\mathcal{C}_\tau, \mathcal{O}_\tau) \to
(\mathcal{C}_{\tau'}, \mathcal{O}_{\tau'})$ comme ci-dessus. Soit $A$ un
ensemble et, pour $\alpha \in A$, soit
$$
\xymatrix{
E_\alpha \ar[d] \ar[r] & Y_\alpha \ar[d] \\
Z_\alpha \ar[r] & X_\alpha
}
$$
un diagramme commutatif dans la catégorie $\mathcal{C}$. Supposons que
\begin{enumerate}
\item un $\tau'$-faisceau $\mathcal{F}'$ soit un $\tau$-faisceau dès que
$\mathcal{F}'(X_\alpha) =
\mathcal{F}'(Z_\alpha) \times_{\mathcal{F}'(E_\alpha)} 
\mathcal{F}'(Y_\alpha)$ pour tout $\alpha$,
\item pour $K'$ dans $D(\mathcal{O}_{\tau'})$ appartenant à l'image
essentielle de $R\epsilon_*$, les morphismes
$c^{K'}_{X_\alpha, Z_\alpha, Y_\alpha, E_\alpha}$ du Lemme
\ref{lemma-c-square} soient des isomorphismes pour tout $\alpha$.
\end{enumerate}
Alors $K' \in D^+(\mathcal{O}_{\tau'})$ appartient à l'image essentielle de
$R\epsilon_*$ si et seulement si les morphismes
$c^{K'}_{X_\alpha, Z_\alpha, Y_\alpha, E_\alpha}$ sont des isomorphismes
pour tout $\alpha$.
\end{lemma}

\begin{proof}
L'implication directe résulte de l'hypothèse (2). D'autre part, si $K'$
n'a qu'un seul faisceau de cohomologie non nul, l'implication réciproque
résulte de l'hypothèse (1). Dans le cas général, nous allons démontrer
l'implication réciproque par récurrence. Disons qu'un objet $K'$ de
$D^+(\mathcal{O}_{\tau'})$ satisfait (P) si les morphismes
$c^{K'}_{X_\alpha, Z_\alpha, Y_\alpha, E_\alpha}$ sont des isomorphismes
pour tout $\alpha \in A$.

\medskip\noindent
Soit donc $K'$ un objet de $D^+(\mathcal{O}_{\tau'})$ satisfaisant (P).
Choisissons un triangle distingué
$$
K' \to R\epsilon_*\epsilon^{-1}K' \to M' \to K'[1]
$$
dans $D^+(\mathcal{O}_{\tau'})$, où la première flèche est le morphisme
d'adjonction. D'après (2) et le Lemme
\ref{lemma-two-out-of-three-blow-up-square}, $M'$ satisfait (P). D'autre
part, en appliquant $\epsilon^{-1}$ et en utilisant l'égalité
$\epsilon^{-1}R\epsilon_* = \text{id}$ de la Section
\ref{section-compare}, nous obtenons $\epsilon^{-1}M' = 0$. Nous allons
montrer dans le paragraphe suivant que $M' = 0$, ce qui achèvera la
démonstration.

\medskip\noindent
Soit $K'$ un objet de $D^+(\mathcal{O}_{\tau'})$ satisfaisant (P) et tel
que $\epsilon^{-1}K' = 0$. Montrons que $K' = 0$. Étant donné
$n \in \mathbf{Z}$ tel que $H^i(K') = 0$ pour $i < n$, nous allons montrer
que $H^n(K') = 0$. Pour $\alpha \in A$, nous avons un triangle distingué
$$
R\Gamma_{\tau'}(X_\alpha, K') \to
R\Gamma_{\tau'}(Z_\alpha, K') \oplus
R\Gamma_{\tau'}(Y_\alpha, K') \to
R\Gamma_{\tau'}(E_\alpha, K') \to
R\Gamma_{\tau'}(X_\alpha, K')[1]
$$
d'après le Lemme \ref{lemma-c-square}. En prenant la cohomologie en degré
$n$ et en utilisant l'annulation supposée des faisceaux de cohomologie de
$K'$, nous obtenons une suite exacte
$$
0 \to
H^n_{\tau'}(X_\alpha, K') \to
H^n_{\tau'}(Z_\alpha, K') \oplus
H^n_{\tau'}(Y_\alpha, K') \to
H^n_{\tau'}(E_\alpha, K')
$$
qui n'est autre que la suite exacte
$$
0 \to
\Gamma(X_\alpha, H^n(K')) \to
\Gamma(Z_\alpha, H^n(K')) \oplus
\Gamma(Y_\alpha, H^n(K')) \to
\Gamma(E_\alpha, H^n(K'))
$$
L'hypothèse (1) montre que $H^n(K')$ est un $\tau$-faisceau. Or son
$\tau$-faisceau associé
$\epsilon^{-1}H^n(K') = H^n(\epsilon^{-1}K')$ est égal à $0$ puisque
$\epsilon^{-1}K' = 0$. Nous en déduisons que $H^n(K') = 0$, comme voulu.
\end{proof}

\begin{lemma}
\label{lemma-descent-squares-helper}
Soit $\epsilon : (\mathcal{C}_\tau, \mathcal{O}_\tau) \to
(\mathcal{C}_{\tau'}, \mathcal{O}_{\tau'})$ comme ci-dessus. Soit
$$
\xymatrix{
E \ar[d] \ar[r] & Y \ar[d] \\
Z \ar[r] & X
}
$$
un diagramme commutatif dans la catégorie $\mathcal{C}$ tel que
\begin{enumerate}
\item $h_X^\# = h_Y^\# \amalg_{h_E^\#} h_Z^\#$, et
\item $h_E^\# \to h_Y^\#$ soit injectif
\end{enumerate}
où ${}^\#$ désigne la $\tau$-faisceautisation. Alors, pour
$K' \in D(\mathcal{O}_{\tau'})$ appartenant à l'image essentielle de
$R\epsilon_*$, le morphisme $c^{K'}_{X, Z, Y, E}$ du Lemme
\ref{lemma-c-square} (construit avec la topologie $\tau'$) est un
isomorphisme.
\end{lemma}

\begin{proof}
Ce lemme auxiliaire est une conséquence presque immédiate du Lemme
\ref{lemma-square-triangle}, et nous recommandons vivement au lecteur d'en
omettre la démonstration. Écrivons $K' = R\epsilon_*K$. Choisissons un
complexe K-injectif de $\mathcal{O}_\tau$-modules, noté
$\mathcal{J}^\bullet$, représentant $K$. Alors
$\epsilon_*\mathcal{J}^\bullet$ est un complexe K-injectif de
$\mathcal{O}_{\tau'}$-modules représentant $K'$ ; voir le Lemme
\ref{lemma-K-injective-flat}. Ensuite,
$$
0 \to
\mathcal{J}^\bullet(X) \xrightarrow{\alpha}
\mathcal{J}^\bullet(Z) \oplus
\mathcal{J}^\bullet(Y) \xrightarrow{\beta}
\mathcal{J}^\bullet(E)
\to 0
$$
est une suite exacte courte de complexes de groupes abéliens ; voir le
Lemme \ref{lemma-square-triangle} et sa démonstration. Comme il s'agit de
la même suite de complexes de groupes abéliens que celle utilisée pour
définir $c^{K'}_{X, Z, Y, E}$, nous concluons.
\end{proof}
\section{Comparaison des cohomologies}
\label{section-compare-general}

\noindent
Nous développons une théorie générale qui nous aidera à comparer
les cohomologies pour différentes topologies. Étant donnés
$\mathcal{C}$, $\tau$ et $\tau'$ comme à la section \ref{section-compare},
ainsi qu'un morphisme $f : X \to Y$ de $\mathcal{C}$, nous obtenons un
diagramme commutatif de morphismes de topos
\begin{equation}
\label{equation-commutative-epsilon}
\vcenter{
\xymatrix{
\Sh(\mathcal{C}_\tau/X) \ar[r]_{f_\tau} \ar[d]_{\epsilon_X} &
\Sh(\mathcal{C}_\tau/Y) \ar[d]^{\epsilon_Y} \\
\Sh(\mathcal{C}_{\tau'}/X) \ar[r]^{f_{\tau'}} &
\Sh(\mathcal{C}_{\tau'}/Y)
}
}
\end{equation}
Ici, le morphisme $\epsilon_X$, resp.\ $\epsilon_Y$, est le morphisme de
comparaison de la section \ref{section-compare} pour la catégorie
$\mathcal{C}/X$ munie des deux topologies $\tau$ et $\tau'$.
Les morphismes $f_\tau$ et $f_{\tau'}$ sont des morphismes de
«~relocalisation~» (Sites, lemme \ref{sites-lemma-relocalize}).
La commutativité du diagramme est un cas particulier de
Sites, lemme \ref{sites-lemma-localize-morphism}
(appliqué avec $\mathcal{C} = \mathcal{C}_\tau/Y$,
$\mathcal{D} = \mathcal{C}_{\tau'}/Y$,
$u = \text{id}$, $U = X$ et $V = X$). Nous obtenons aussi
$\epsilon_{X, *} \circ f_\tau^{-1} = f_{\tau'}^{-1} \circ \epsilon_{Y, *}$
soit par ce lemme, soit directement.

\begin{situation}
\label{situation-compare}
On se place avec $\mathcal{C}$, $\tau$ et $\tau'$ comme à la section
\ref{section-compare}. Supposons donnés un sous-ensemble
$\mathcal{P} \subset \text{Flèches}(\mathcal{C})$ et, pour tout objet $X$ de
$\mathcal{C}$, une sous-catégorie de Serre faible
$\mathcal{A}'_X \subset \textit{Ab}(\mathcal{C}_{\tau'}/X)$.
Nous faisons les hypothèses suivantes :
\begin{enumerate}
\item
\label{item-base-change-P}
si $f : X \to Y$ appartient à $\mathcal{P}$ et si $Y' \to Y$ est quelconque,
alors $X \times_Y Y'$ existe et $X \times_Y Y' \to Y'$ appartient à $\mathcal{P}$,
\item
\label{item-restriction-A}
$f_{\tau'}^{-1}$ envoie $\mathcal{A}'_Y$ dans $\mathcal{A}'_X$
pour tout morphisme $f : X \to Y$ de $\mathcal{C}$,
\item
\label{item-A-sheaf}
si $X$ est dans $\mathcal{C}$ et $\mathcal{F}'$ dans $\mathcal{A}'_X$, alors
$\mathcal{F}'$ satisfait la condition de faisceau pour les recouvrements
$\tau$, c'est-à-dire
$\mathcal{F}' = \epsilon_{X, *}\epsilon_X^{-1}\mathcal{F}'$,
\item
\label{item-A-and-P}
si $f : X \to Y$ appartient à $\mathcal{P}$ et si
$\mathcal{F}' \in \Ob(\mathcal{A}'_X)$, alors
$R^if_{\tau', *}\mathcal{F}' \in \Ob(\mathcal{A}'_Y)$
pour $i \geq 0$.
\item
\label{item-refine-tau-by-P}
si $\{U_i \to U\}_{i \in I}$ est un recouvrement $\tau$, alors il existe
\begin{enumerate}
\item un recouvrement $\tau'$ $\{V_j \to U\}_{j \in J}$,
\item un recouvrement $\tau$ $\{f_j : W_j \to V_j\}$ réduit à un seul
morphisme $f_j \in \mathcal{P}$, et
\item un recouvrement $\tau'$ $\{W_{jk} \to W_j\}_{k \in K_j}$
\end{enumerate}
tels que $\{W_{jk} \to U\}_{j \in J, k \in K_j}$ soit un raffinement de
$\{U_i \to U\}_{i \in I}$.
\end{enumerate}
\end{situation}

\begin{lemma}
\label{lemma-A}
Dans la situation \ref{situation-compare}, pour $X$ dans $\mathcal{C}$,
notons $\mathcal{A}_X$ la classe des objets de
$\textit{Ab}(\mathcal{C}_\tau/X)$ de la forme
$\epsilon_X^{-1}\mathcal{F}'$, avec $\mathcal{F}'$ dans $\mathcal{A}'_X$.
Alors
\begin{enumerate}
\item pour $\mathcal{F}$ dans $\textit{Ab}(\mathcal{C}_\tau/X)$,
nous avons $\mathcal{F} \in \mathcal{A}_X \Leftrightarrow
\epsilon_{X, *}\mathcal{F} \in \mathcal{A}'_X$, et
\item $f_\tau^{-1}$ envoie $\mathcal{A}_Y$ dans $\mathcal{A}_X$
pour tout morphisme $f : X \to Y$ de $\mathcal{C}$.
\end{enumerate}
\end{lemma}

\begin{proof}
L'assertion (1) résulte de (\ref{item-A-sheaf}) et l'assertion (2)
résulte de (\ref{item-restriction-A}) et de la commutativité de
(\ref{equation-commutative-epsilon}), qui donne
$\epsilon_X^{-1} \circ f_{\tau'}^{-1} = f_\tau^{-1} \circ \epsilon_Y^{-1}$.
\end{proof}

\noindent
Notre prochain objectif est de démontrer les lemmes
\ref{lemma-compare-cohomology-general} et
\ref{lemma-cohomological-descent-general}. Nous procéderons par récurrence
en utilisant l'hypothèse de récurrence suivante.

\medskip\noindent
$(V_n)$ Pour $X$ dans $\mathcal{C}$ et $\mathcal{F}$ dans $\mathcal{A}_X$,
nous avons $R^i\epsilon_{X, *}\mathcal{F} = 0$ pour $1 \leq i \leq n$.

\begin{lemma}
\label{lemma-V-implies-C-general}
Dans la situation \ref{situation-compare}, supposons $(V_n)$ vérifiée.
Pour $f : X \to Y$ dans $\mathcal{P}$ et $\mathcal{F}$ dans $\mathcal{A}_X$,
nous avons $R^if_{\tau', *}\epsilon_{X, *}\mathcal{F} =
\epsilon_{Y, *}R^if_{\tau, *}\mathcal{F}$ pour $i \leq n$.
\end{lemma}

\begin{proof}
Nous utiliserons le diagramme commutatif
(\ref{equation-commutative-epsilon}) sans plus le mentionner. En particulier,
nous avons
$$
Rf_{\tau', *}R\epsilon_{X, *}\mathcal{F} =
R\epsilon_{Y, *}Rf_{\tau, *}\mathcal{F}
$$
L'hypothèse $(V_n)$ nous dit que
$\epsilon_{X, *}\mathcal{F} \to R\epsilon_{X, *}\mathcal{F}$
est un isomorphisme en degrés $\leq n$. Par conséquent,
$Rf_{\tau', *}\epsilon_{X, *}\mathcal{F} \to
Rf_{\tau', *}R\epsilon_{X, *}\mathcal{F}$ est un isomorphisme
en degrés $\leq n$. Nous en concluons que
$$
R^if_{\tau', *}\epsilon_{X, *}\mathcal{F} \to
H^i(R\epsilon_{Y, *}Rf_{\tau, *}\mathcal{F})
$$
est un isomorphisme pour $i \leq n$. Nous démontrerons le lemme en examinant
la deuxième page de la suite spectrale du
lemme \ref{lemma-relative-Leray} pour
$R\epsilon_{Y, *}Rf_{\tau, *}\mathcal{F}$. En voici une représentation :
$$
\begin{matrix}
\ldots &
\ldots &
\ldots &
\ldots \\
\epsilon_{Y, *}R^2f_{\tau, *}\mathcal{F} &
R^1\epsilon_{Y, *}R^2f_{\tau, *}\mathcal{F} &
R^2\epsilon_{Y, *}R^2f_{\tau, *}\mathcal{F} &
\ldots \\
\epsilon_{Y, *}R^1f_{\tau, *}\mathcal{F} &
R^1\epsilon_{Y, *}R^1f_{\tau, *}\mathcal{F} &
R^2\epsilon_{Y, *}R^1f_{\tau, *}\mathcal{F} &
\ldots \\
\epsilon_{Y, *}f_{\tau, *}\mathcal{F} &
R^1\epsilon_{Y, *}f_{\tau, *}\mathcal{F} &
R^2\epsilon_{Y, *}f_{\tau, *}\mathcal{F} &
\ldots
\end{matrix}
$$
Notons $(C_m)$ l'hypothèse : $R^if_{\tau', *}\epsilon_{X, *}\mathcal{F} =
\epsilon_{Y, *}R^if_{\tau, *}\mathcal{F}$ pour $i \leq m$. Observons que
$(C_0)$ est vérifiée. Nous allons montrer que
$(C_{m - 1}) \Rightarrow (C_m)$ pour $m < n$. En effet, si
$(C_{m - 1})$ est vérifiée, alors, pour
$n \geq p > 0$ et $q \leq m - 1$, nous avons
\begin{align*}
R^p\epsilon_{Y, *}R^qf_{\tau, *}\mathcal{F}
& =
R^p\epsilon_{Y, *}
\epsilon_Y^{-1} \epsilon_{Y, *} R^qf_{\tau, *}\mathcal{F} \\
& =
R^p\epsilon_{Y, *}
\epsilon_Y^{-1}R^qf_{\tau', *}\epsilon_{X, *}\mathcal{F} = 0
\end{align*}
La première égalité vient de
$\epsilon_Y^{-1}\epsilon_{Y, *} = \text{id}$, la deuxième de
$(C_{m - 1})$, et la dernière de $(V_n)$, car
$\epsilon_Y^{-1}R^qf_{\tau', *}\epsilon_{X, *}\mathcal{F}$
appartient à $\mathcal{A}_Y$ par (\ref{item-A-and-P}).
En examinant la suite spectrale, nous voyons que
$E_2^{0, m} = \epsilon_{Y, *}R^mf_{\tau, *}\mathcal{F}$
est le seul terme $E_2^{p, q}$ non nul avec $p + q = m$.
Rappelons que $\text{d}_r^{p, q} : E_r^{p, q} \to E_r^{p + r, q - r + 1}$.
Il n'y a donc aucune différentielle non nulle
$\text{d}_r^{p, q}$, $r \geq 2$, partant de cette position ou y arrivant.
Nous en concluons que
$H^m(R\epsilon_{Y, *}Rf_{\tau, *}\mathcal{F}) =
\epsilon_{Y, *}R^mf_{\tau, *}\mathcal{F}$, ce qui implique
$(C_m)$ d'après la discussion précédente.

\medskip\noindent
Enfin, supposons $(C_{n - 1})$. La même analyse montre que
$E_2^{0, n} = \epsilon_{Y, *}R^nf_{\tau, *}\mathcal{F}$
est le seul terme $E_2^{p, q}$ non nul avec $p + q = n$.
Il n'y a toujours aucune différentielle non nulle arrivant à cette position,
mais une différentielle non nulle peut en partir, à savoir le morphisme
$d_{n + 1}^{0, n} : \epsilon_{Y, *}R^nf_{\tau, *}\mathcal{F} \to
R^{n + 1}\epsilon_{Y, *}f_{\tau, *}\mathcal{F}$.
Nous en concluons qu'il existe une suite exacte
$$
0 \to R^nf_{\tau', *}\epsilon_{X, *}\mathcal{F} \to 
\epsilon_{Y, *}R^nf_{\tau, *}\mathcal{F} \to
R^{n + 1}\epsilon_{Y, *}f_{\tau, *}\mathcal{F}
$$
Par (\ref{item-A-and-P}) et (\ref{item-A-sheaf}), le faisceau
$R^nf_{\tau', *}\epsilon_{X, *}\mathcal{F}$
satisfait la condition de faisceau pour les recouvrements $\tau$, tout comme
$\epsilon_{Y, *}R^nf_{\tau, *}\mathcal{F}$
(utiliser la description de $\epsilon_*$ à la section \ref{section-compare}).
Cependant, le faisceautisé pour $\tau$ du faisceau pour $\tau'$
$R^{n + 1}\epsilon_{Y, *}f_{\tau, *}\mathcal{F}$
est nul (par localité de la cohomologie ; utiliser les
lemmes \ref{lemma-kill-cohomology-class-on-covering} et
\ref{lemma-higher-direct-images}).
Ainsi, $R^nf_{\tau', *}\epsilon_{X, *}\mathcal{F} \to 
\epsilon_{Y, *}R^nf_{\tau, *}\mathcal{F}$
doit donc être un isomorphisme, ce qui achève la démonstration.
\end{proof}

\noindent
Si $E'$, resp.\ $E$, est un objet de
$D(\mathcal{C}_{\tau'}/X)$, resp.\ $D(\mathcal{C}_\tau/X)$,
nous noterons
$H^n_{\tau'}(U, E')$, resp.\ $H^n_\tau(U, E)$,
la cohomologie de $E'$, resp.\ $E$,
au-dessus d'un objet $U$ de $\mathcal{C}/X$.

\begin{lemma}
\label{lemma-V-implies-cohomology-general}
Dans la situation \ref{situation-compare}, si $(V_n)$ est vérifiée, alors,
pour $X$ dans $\mathcal{C}$ et $L \in D(\mathcal{C}_{\tau'}/X)$
tels que $H^i(L) = 0$ pour $i < 0$ et que $H^i(L)$ appartienne à
$\mathcal{A}'_X$ pour $0 \leq i \leq n$, nous avons
$H^n_{\tau'}(X, L) = H^n_\tau(X, \epsilon_X^{-1}L)$.
\end{lemma}

\begin{proof}
Par le lemme \ref{lemma-Leray-unbounded}, nous avons
$H^n_\tau(X, \epsilon_X^{-1}L) =
H^n_{\tau'}(X, R\epsilon_{X, *}\epsilon_X^{-1}L)$.
Il existe une suite spectrale
$$
E_2^{p, q} = R^p\epsilon_{X, *}\epsilon_X^{-1}H^q(L)
$$
qui converge vers $H^{p + q}(R\epsilon_{X, *}\epsilon_X^{-1}L)$.
Par $(V_n)$, le terme $E_2^{p, q}$ est nul pour
$0 < p \leq n$ et $0 \leq q \leq n$. Ainsi, les termes
$E_2^{0, q} = \epsilon_{X, *}\epsilon_X^{-1}H^q(L) = H^q(L)$
sont les seuls termes $E_2^{p, q}$ non nuls avec $p + q \leq n$.
Il s'ensuit que le morphisme
$$
L \longrightarrow R\epsilon_{X, *}\epsilon_X^{-1}L
$$
est un isomorphisme en degrés $\leq n$ (on omet un détail mineur).
Nous obtenons donc
$H^i_{\tau'}(X, L) = H^i_{\tau'}(X, R\epsilon_{X, *}\epsilon_X^{-1}L)$
pour $i \leq n$. Le lemme est ainsi démontré.
\end{proof}

\begin{lemma}
\label{lemma-V-implies-cohomology-extra-general}
Dans la situation \ref{situation-compare}, si $(V_n)$ est vérifiée, alors,
pour $X$ dans $\mathcal{C}$ et $\mathcal{F}$ dans $\mathcal{A}_X$, le morphisme
$H^{n + 1}_{\tau'}(X, \epsilon_{X, *}\mathcal{F}) \to
H^{n + 1}_\tau(X, \mathcal{F})$
est injectif et son image est formée des classes qui deviennent nulles sur
un recouvrement $\tau'$ de $X$.
\end{lemma}

\begin{proof}
Rappelons que $\epsilon_X^{-1}\epsilon_{X, *}\mathcal{F} = \mathcal{F}$ ;
le morphisme est donc donné par l'image inverse des classes de cohomologie
par $\epsilon_X$. La suite spectrale de Leray
(lemme \ref{lemma-Leray})
$$
E_2^{p, q} = H^p_{\tau'}(X, R^q\epsilon_{X, *}\mathcal{F})
\Rightarrow
H^{p + q}_\tau(X, \mathcal{F})
$$
combinée à l'annulation supposée donne une suite exacte
$$
0 \to
H^{n + 1}_{\tau'}(X, \epsilon_{X, *}\mathcal{F}) \to
H^{n + 1}_\tau(X, \mathcal{F}) \to
H^0_{\tau'}(X, R^{n + 1}\epsilon_{X, *}\mathcal{F})
$$
C'est une reformulation du lemme.
\end{proof}

\begin{lemma}
\label{lemma-make-class-zero-general}
Dans la situation \ref{situation-compare}, soit $f : X \to Y$
un élément de $\mathcal{P}$ tel que $\{X \to Y\}$ soit un recouvrement
$\tau$. Soit $\mathcal{F}'$ dans $\mathcal{A}'_Y$. Si $n \geq 0$ et
$$
\theta \in
\text{Égaliseur}\left(
\xymatrix{
H^{n + 1}_{\tau'}(X, \mathcal{F}')
\ar@<1ex>[r] \ar@<-1ex>[r] &
H^{n + 1}_{\tau'}(X \times_Y X, \mathcal{F}')
}
\right)
$$
alors il existe un recouvrement $\tau'$ $\{Y_i \to Y\}$
tel que la restriction de $\theta$ soit nulle dans
$H^{n + 1}_{\tau'}(Y_i \times_Y X, \mathcal{F}')$.
\end{lemma}

\begin{proof}
Observons que $X \times_Y X$ existe d'après (\ref{item-base-change-P}).
Pour $Z$ dans $\mathcal{C}/Y$, notons $\mathcal{F}'|_Z$
la restriction de $\mathcal{F}'$ à $\mathcal{C}_{\tau'}/Z$.
Rappelons que $H^{n + 1}_{\tau'}(X, \mathcal{F}') =
H^{n + 1}(\mathcal{C}_{\tau'}/X, \mathcal{F}'|_X)$, voir le
lemme \ref{lemma-cohomology-of-open}.
Le lemme affirme que l'image
$\overline{\theta} \in H^0(Y, R^{n + 1}f_{\tau', *}\mathcal{F}'|_X)$
de $\theta$ est nulle. Considérons le diagramme cartésien
$$
\xymatrix{
X \times_Y X \ar[d]_{\text{pr}_1} \ar[r]_{\text{pr}_2} &
X \ar[d]^f \\
X \ar[r]^f & Y
}
$$
Par changement de base trivial (lemme \ref{lemma-localize-cartesian-square}),
nous avons
$$
f_{\tau'}^{-1}R^{n + 1}f_{\tau', *}(\mathcal{F}'|_X) =
R^{n + 1}\text{pr}_{1, \tau', *}(\mathcal{F}'|_{X \times_Y X})
$$
Si $\text{pr}_1^{-1}\theta = \text{pr}_2^{-1}\theta$,
alors la section $f_{\tau'}^{-1}\overline{\theta}$ de
$f_{\tau'}^{-1}R^{n + 1}f_{\tau', *}(\mathcal{F}'|_X)$ est nulle,
car il est clair que $\text{pr}_1^{-1}\theta$ s'envoie sur l'élément nul de
$H^0(X, R^{n + 1}\text{pr}_{1, \tau', *}(\mathcal{F}'|_{X \times_Y X}))$.
D'après (\ref{item-restriction-A}), $\mathcal{F}'|_X$ appartient à
$\mathcal{A}'_X$. Ainsi,
$\mathcal{G}' = R^{n + 1}f_{\tau', *}(\mathcal{F}'|_X)$
est un objet de $\mathcal{A}'_Y$ par (\ref{item-A-and-P}).
Par conséquent, $\mathcal{G}'$ satisfait la condition de faisceau pour les
recouvrements $\tau$ d'après (\ref{item-A-sheaf}).
Comme $\{X \to Y\}$ est un recouvrement $\tau$, l'application de restriction
$\mathcal{G}'(Y) \to \mathcal{G}'(X)$ est injective.
Il s'ensuit que $\overline{\theta}$ est nulle.
\end{proof}

\begin{lemma}
\label{lemma-induction-step-V-C-general}
Dans la situation \ref{situation-compare}, nous avons
$(V_n) \Rightarrow (V_{n + 1})$.
\end{lemma}

\begin{proof}
Soient $X$ dans $\mathcal{C}$ et $\mathcal{F}$ dans $\mathcal{A}_X$.
Soit $\xi \in H^{n + 1}_\tau(U, \mathcal{F})$ pour un certain $U/X$.
Nous devons montrer que la restriction de $\xi$ est nulle sur les membres
d'un recouvrement $\tau'$ de $U$, voir le
lemme \ref{lemma-higher-direct-images}.
Il en résulte que nous pouvons remplacer $U$ par les membres d'un
recouvrement $\tau'$ de $U$.

\medskip\noindent
Par localité de la cohomologie
(lemme \ref{lemma-kill-cohomology-class-on-covering}),
nous pouvons choisir un recouvrement $\tau$ $\{U_i \to U\}$
tel que la restriction de $\xi$ à $U_i$ soit nulle.
Choisissons $\{V_j \to V\}$, $\{f_j : W_j \to V_j\}$ et
$\{W_{jk} \to W_j\}$ comme dans (\ref{item-refine-tau-by-P}).
En remplaçant $U$ par $V_j$ et $\mathcal{F}$ par sa restriction à
$\mathcal{C}_\tau/V_j$, ce qui est permis par
(\ref{item-base-change-P}), nous nous ramenons au cas étudié au paragraphe
suivant.

\medskip\noindent
Ici, $f : X \to Y$ est un élément de $\mathcal{P}$ tel que
$\{X \to Y\}$ soit un recouvrement $\tau$,
$\mathcal{F}$ est un objet de $\mathcal{A}_Y$, et
$\xi \in H^{n + 1}_\tau(Y, \mathcal{F})$ est tel qu'il existe un
recouvrement $\tau'$ $\{X_i \to X\}_{i \in I}$ tel que $\xi$ se restreigne
à zéro sur $X_i$ pour tout $i \in I$.
Il s'agit de montrer que la restriction de $\xi$ est nulle sur un
recouvrement $\tau'$ de $Y$.

\medskip\noindent
Par le lemme \ref{lemma-V-implies-cohomology-extra-general}, il existe une
unique classe de cohomologie $\tau'$
$\theta \in H^{n + 1}_{\tau'}(X, \epsilon_{X, *}\mathcal{F})$
dont l'image est $\xi|_X$.
Comme les images inverses de $\xi|_X$ par les deux projections donnent la
même classe sur $X \times_Y X$, il en est de même pour $\theta$
(par unicité).
D'après le lemme \ref{lemma-make-class-zero-general}, après avoir remplacé
$Y$ par les membres d'un recouvrement $\tau'$, nous pouvons supposer que
$\theta = 0$. Nous pouvons donc supposer que $\xi|_X$ est nulle.

\medskip\noindent
Soit $f : X \to Y$ un élément de $\mathcal{P}$ tel que
$\{X \to Y\}$ soit un recouvrement $\tau$,
soit $\mathcal{F}$ un objet de $\mathcal{A}_Y$, et soit
$\xi \in H^{n + 1}_\tau(Y, \mathcal{F})$ dont l'image dans
$H^{n + 1}_\tau(X, \mathcal{F})$ est nulle.
Il faut montrer que la restriction de $\xi$ est nulle sur un recouvrement
$\tau'$ de $Y$.

\medskip\noindent
Les hypothèses nous disent que l'image de $\xi$ par le morphisme
$$
\mathcal{F} \longrightarrow Rf_{\tau, *}f_\tau^{-1}\mathcal{F}
$$
est nulle ; on utilise ici le lemme \ref{lemma-Leray-unbounded}.
Un argument simple fondé sur le triangle distingué de troncation
(Catégories dérivées, remarque
\ref{derived-remark-truncation-distinguished-triangle}) montre que
l'image de $\xi$ par le morphisme
$$
\mathcal{F} \longrightarrow \tau_{\leq n}Rf_{\tau, *}f_\tau^{-1}\mathcal{F}
$$
est nulle. Nous allons comparer celui-ci au morphisme
$\epsilon_{Y, *}\mathcal{F} \to K$, où
$$
K = \tau_{\leq n}Rf_{\tau', *}f_{\tau'}^{-1}\epsilon_{Y, *}\mathcal{F} =
\tau_{\leq n}Rf_{\tau', *}\epsilon_{X, *}f_{\tau}^{-1}\mathcal{F}
$$
L'égalité
$\epsilon_{X, *} f_\tau^{-1} = f_{\tau'}^{-1} \epsilon_{Y, *}$
est une propriété de (\ref{equation-commutative-epsilon}). Considérons le
morphisme
$$
Rf_{\tau', *}\epsilon_{X, *}f_{\tau}^{-1}\mathcal{F} \longrightarrow
Rf_{\tau', *}R\epsilon_{X, *}f_{\tau}^{-1}\mathcal{F} =
R\epsilon_{Y, *}Rf_{\tau, *}f_\tau^{-1}\mathcal{F}
$$
utilisé dans la démonstration du lemme \ref{lemma-V-implies-C-general} ;
il induit par adjonction un morphisme
$$
\epsilon_Y^{-1} Rf_{\tau', *}\epsilon_{X, *}f_{\tau}^{-1}\mathcal{F} \to
Rf_{\tau, *}f_\tau^{-1}\mathcal{F}
$$
En prenant les troncations, nous obtenons un morphisme
$$
\epsilon_Y^{-1}K
\longrightarrow
\tau_{\leq n}Rf_{\tau, *}f_\tau^{-1}\mathcal{F}
$$
qui est un isomorphisme par le lemme \ref{lemma-V-implies-C-general} ;
ce lemme s'applique, car $f_\tau^{-1}\mathcal{F}$ appartient à
$\mathcal{A}_X$ d'après le lemme \ref{lemma-A}. Choisissons un triangle
distingué
$$
\epsilon_{Y, *}\mathcal{F} \to K \to L \to \epsilon_{Y, *}\mathcal{F}[1]
$$
Le morphisme $\mathcal{F} \to f_{\tau, *}f_\tau^{-1}\mathcal{F}$
est injectif puisque $\{X \to Y\}$ est un recouvrement $\tau$. Ainsi,
$\epsilon_{Y, *}\mathcal{F} \to
\epsilon_{Y, *}f_{\tau, *}f_\tau^{-1}\mathcal{F} =
f_{\tau', *}f_{\tau'}^{-1}\epsilon_{Y, *}\mathcal{F}$
est lui aussi injectif. Par conséquent, les seuls faisceaux de cohomologie
non nuls de $L$ sont en degrés $0, \ldots, n$.
Comme $f_{\tau', *}f_{\tau'}^{-1}\epsilon_{Y, *}\mathcal{F}$
appartient à $\mathcal{A}'_Y$ d'après (\ref{item-restriction-A}) et
(\ref{item-A-and-P}), nous en concluons que
$$
H^0(L) = \Coker(\epsilon_{Y, *}\mathcal{F} \to
f_{\tau', *}f_{\tau'}^{-1}\epsilon_{Y, *}\mathcal{F})
$$
appartient à la sous-catégorie de Serre faible $\mathcal{A}'_Y$. Pour
$1 \leq i \leq n$, nous voyons que
$H^i(L) = R^if_{\tau', *}f_{\tau'}^{-1}\epsilon_{Y, *}\mathcal{F}$
appartient à $\mathcal{A}'_Y$ d'après (\ref{item-restriction-A}) et
(\ref{item-A-and-P}). En tirant le triangle distingué ci-dessus en arrière
par $\epsilon_Y$, nous obtenons le triangle distingué
$$
\mathcal{F} \to \tau_{\leq n}Rf_{\tau, *}f_\tau^{-1}\mathcal{F}
\to \epsilon_Y^{-1}L \to \mathcal{F}[1]
$$
Puisque l'image de $\xi$ dans le terme du milieu est nulle, nous voyons que
$\xi$ est l'image d'un élément
$\xi' \in H^n_\tau(Y, \epsilon_Y^{-1}L)$.
D'après le lemme \ref{lemma-V-implies-cohomology-general}, nous avons
$$
H^n_{\tau'}(Y, L) = H^n_\tau(Y, \epsilon_Y^{-1}L),
$$
Nous pouvons donc relever $\xi'$ en un élément de $H^n_{\tau'}(Y, L)$,
puis prendre son bord dans
$H^{n + 1}_{\tau'}(Y, \epsilon_{Y, *}\mathcal{F})$
pour voir que $\xi$ appartient à l'image du morphisme canonique
$H^{n + 1}_{\tau'}(Y, \epsilon_{Y, *}\mathcal{F}) \to
H^{n + 1}_\tau(Y, \mathcal{F})$.
La localité de la cohomologie pour
$H^{n + 1}_{\tau'}(Y,\epsilon_{Y, *}\mathcal{F})$, voir le
lemme \ref{lemma-kill-cohomology-class-on-covering}, permet de conclure.
\end{proof}

\begin{lemma}
\label{lemma-V-C-all-n-general}
Dans la situation \ref{situation-compare}, la propriété $(V_n)$ est vraie
pour tout $n$. De plus :
\begin{enumerate}
\item Pour $X$ dans $\mathcal{C}$ et
$K' \in D^+_{\mathcal{A}'_X}(\mathcal{C}_{\tau'}/X)$, le morphisme
$K' \to R\epsilon_{X, *}(\epsilon_X^{-1}K')$ est un isomorphisme.
\item Pour $f : X \to Y$ dans $\mathcal{P}$ et
$K' \in D^+_{\mathcal{A}'_X}(\mathcal{C}_{\tau'}/X)$, nous avons
$Rf_{\tau', *}K' \in D^+_{\mathcal{A}'_Y}(\mathcal{C}_{\tau'}/Y)$ et
$\epsilon_Y^{-1}(Rf_{\tau', *}K') = Rf_{\tau, *}(\epsilon_X^{-1}K')$.
\end{enumerate}
\end{lemma}

\begin{proof}
Observons que $(V_0)$ est vérifiée, car c'est la condition vide.
Nous obtenons alors $(V_n)$ pour tout $n$ par le
lemme \ref{lemma-induction-step-V-C-general}.

\medskip\noindent
Démonstration de (1). L'objet $K = \epsilon_X^{-1}K'$ a pour faisceaux de
cohomologie $H^i(K) = \epsilon_X^{-1}H^i(K')$, qui appartiennent à
$\mathcal{A}_X$. La suite spectrale
$$
E_2^{p, q} = R^p\epsilon_{X, *} H^q(K) \Rightarrow
H^{p + q}(R\epsilon_{X, *}K)
$$
dégénère donc, par $(V_n)$ pour tout $n$, et nous obtenons
$$
H^n(R\epsilon_{X, *}K) = \epsilon_{X, *}H^n(K) =
\epsilon_{X, *}\epsilon_X^{-1}H^n(K') = H^n(K').
$$
La dernière égalité vient encore de ce que $H^n(K')$ appartient à
$\mathcal{A}'_X$. Le morphisme canonique
$K' \to R\epsilon_{X, *}(\epsilon_X^{-1}K')$ est donc un isomorphisme.

\medskip\noindent
Démonstration de (2). En utilisant la suite spectrale
$$
E_2^{p, q} = R^pf_{\tau', *}H^q(K') \Rightarrow R^{p + q}f_{\tau', *}K'
$$
le fait que $R^pf_{\tau', *}H^q(K')$ appartient à
$\mathcal{A}'_Y$ par (\ref{item-A-and-P}), le fait que $\mathcal{A}'_Y$ est
une sous-catégorie de Serre faible de
$\textit{Ab}(\mathcal{C}_{\tau'}/Y)$, et Homologie, lemme
\ref{homology-lemma-biregular-ss-converges}, nous concluons que
$Rf_{\tau', *}K' \in D^+_{\mathcal{A}'_Y}(\mathcal{C}_{\tau'}/Y)$.
Pour achever la démonstration, il faut montrer que le morphisme de changement
de base
$$
\epsilon_Y^{-1}(Rf_{\tau', *}K')
\longrightarrow
Rf_{\tau, *}(\epsilon_X^{-1}K')
$$
est un isomorphisme. En comparant la suite spectrale ci-dessus à la suite
spectrale
$$
E_2^{p, q} = R^pf_{\tau, *}H^q(\epsilon_X^{-1}K')
\Rightarrow R^{p + q}f_{\tau, *}\epsilon_X^{-1}K'
$$
nous nous ramenons au cas où $K'$ possède un unique faisceau de cohomologie
non nul $\mathcal{F}'$, appartenant à $\mathcal{A}'_X$ ; nous omettons les
détails. Le lemme \ref{lemma-V-implies-C-general} donne alors
$\epsilon_Y^{-1}R^if_{\tau', *}\mathcal{F}' =
R^if_{\tau, *}\epsilon_X^{-1}\mathcal{F}'$ pour tout $i$,
ce qui achève la démonstration.
\end{proof}

\begin{lemma}
\label{lemma-cohomological-descent-general}
Dans la situation \ref{situation-compare}, pour tout $X$ dans
$\mathcal{C}$, la catégorie
$\mathcal{A}_X \subset \textit{Ab}(\mathcal{C}_\tau/X)$
est une sous-catégorie de Serre faible, et le foncteur
$$
R\epsilon_{X, *} :
D^+_{\mathcal{A}_X}(\mathcal{C}_\tau/X)
\longrightarrow
D^+_{\mathcal{A}'_X}(\mathcal{C}_{\tau'}/X)
$$
est une équivalence dont un quasi-inverse est donné par $\epsilon_X^{-1}$.
\end{lemma}

\begin{proof}
Nous devons vérifier pour $\mathcal{A}_X$ les conditions énumérées dans
Homologie, lemme \ref{homology-lemma-characterize-weak-serre-subcategory}.
Si $\varphi : \mathcal{F} \to \mathcal{G}$ est un morphisme dans
$\mathcal{A}_X$, alors $\epsilon_{X, *}\varphi :
\epsilon_{X, *}\mathcal{F} \to \epsilon_{X, *}\mathcal{G}$
est un morphisme dans $\mathcal{A}'_X$. Par conséquent,
$\Ker(\epsilon_{X, *}\varphi)$ et $\Coker(\epsilon_{X, *}\varphi)$ sont des
objets de $\mathcal{A}'_X$, puisque celle-ci est une sous-catégorie de Serre
faible de $\textit{Ab}(\mathcal{C}_{\tau'}/X)$. En appliquant
$\epsilon_X^{-1}$, nous obtenons une suite exacte
$$
0 \to
\epsilon_X^{-1}\Ker(\epsilon_{X, *}\varphi) \to
\mathcal{F} \to \mathcal{G} \to
\epsilon_X^{-1}\Coker(\epsilon_{X, *}\varphi) \to 0
$$
et nous voyons que $\Ker(\varphi)$ et $\Coker(\varphi)$ appartiennent à
$\mathcal{A}_X$. Enfin, supposons que
$$
0 \to \mathcal{F}_1 \to \mathcal{F}_2 \to \mathcal{F}_3 \to 0
$$
soit une suite exacte courte dans $\textit{Ab}(\mathcal{C}_\tau/X)$,
avec $\mathcal{F}_1$ et $\mathcal{F}_3$ dans $\mathcal{A}_X$.
En appliquant $\epsilon_{X, *}$, nous obtenons une suite exacte
$$
0 \to 
\epsilon_{X, *}\mathcal{F}_1 \to
\epsilon_{X, *}\mathcal{F}_2 \to
\epsilon_{X, *}\mathcal{F}_3 \to
R^1\epsilon_{X, *}\mathcal{F}_1 = 0
$$
L'annulation vient du lemme \ref{lemma-V-C-all-n-general}.
Ainsi, $\epsilon_{X, *}\mathcal{F}_2$ appartient à $\mathcal{A}'_X$, car
celle-ci est une sous-catégorie de Serre faible de
$\textit{Ab}(\mathcal{C}_{\tau'}/X)$. En tirant en arrière par $\epsilon_X$,
nous concluons que $\mathcal{F}_2$ appartient à $\mathcal{A}_X$.

\medskip\noindent
Ainsi, $\mathcal{A}_X$ est une sous-catégorie de Serre faible de
$\textit{Ab}(\mathcal{C}_\tau/X)$, et il est légitime de considérer la
catégorie $D^+_{\mathcal{A}_X}(\mathcal{C}_\tau/X)$.
Observons que $\epsilon_X^{-1} : \mathcal{A}'_X \to \mathcal{A}_X$
est une équivalence et que
$\mathcal{F}' \to R\epsilon_{X, *}\epsilon_X^{-1}\mathcal{F}'$
est un isomorphisme pour $\mathcal{F}'$ dans $\mathcal{A}'_X$, puisque
$(V_n)$ est vérifiée pour tout $n$ d'après le
lemme \ref{lemma-V-C-all-n-general}. Nous concluons donc par le
lemme \ref{lemma-equivalence-bounded}.
\end{proof}

\begin{lemma}
\label{lemma-compare-cohomology-general}
Dans la situation \ref{situation-compare}, soit $X$ dans $\mathcal{C}$.
\begin{enumerate}
\item pour $\mathcal{F}'$ dans $\mathcal{A}'_X$, nous avons
$H^n_{\tau'}(X, \mathcal{F}') = H^n_\tau(X, \epsilon_X^{-1}\mathcal{F}')$,
\item pour $K' \in D^+_{\mathcal{A}'_X}(\mathcal{C}_{\tau'}/X)$,
nous avons $H^n_{\tau'}(X, K') = H^n_\tau(X, \epsilon_X^{-1}K')$.
\end{enumerate}
\end{lemma}

\begin{proof}
Cela résulte du lemme \ref{lemma-V-C-all-n-general} et de la
remarque \ref{remark-before-Leray}.
\end{proof}
























\section{Cohomologie sur les espaces séparés et localement quasi-compacts}
\label{section-cohomology-LC}

\noindent
Nous continuons à suivre notre convention consistant à dire
«~séparé et localement quasi-compact~» plutôt que «~localement compact~»,
comme on le fait souvent dans la littérature. Notons $\textit{LC}$ la
catégorie dont les objets sont les espaces topologiques séparés et localement
quasi-compacts et dont les morphismes sont les applications continues.

\begin{lemma}
\label{lemma-LC-basic}
La catégorie $\textit{LC}$ admet des produits fibrés et un objet final, et
admet donc toutes les limites finies. Étant donnés des morphismes $X \to Z$
et $Y \to Z$ dans $\textit{LC}$ tels que $X$ et $Y$ soient quasi-compacts,
$X \times_Z Y$ est quasi-compact.
\end{lemma}

\begin{proof}
L'objet final est l'espace réduit à un point. Étant donnés des morphismes
$X \to Z$ et $Y \to Z$ de $\textit{LC}$, le produit fibré $X \times_Z Y$
est un sous-espace de $X \times Y$. Ainsi, $X \times_Z Y$ est séparé puisque
$X \times Y$ est séparé d'après
Topologie, section \ref{topology-section-Hausdorff}.

\medskip\noindent
Si $X$ et $Y$ sont quasi-compacts, alors $X \times Y$ est quasi-compact
d'après Topologie, théorème \ref{topology-theorem-tychonov}.
Comme $X \times_Z Y$ est une partie fermée de $X \times Y$
(Topologie, lemme \ref{topology-lemma-fibre-product-closed}),
on voit que $X \times_Z Y$ est quasi-compact d'après
Topologie, lemme \ref{topology-lemma-closed-in-quasi-compact}.

\medskip\noindent
Enfin, revenons au cas général. Si $x \in X$ et $y \in Y$, on peut choisir
des voisinages quasi-compacts $x \in E \subset X$ et $y \in F \subset Y$ ;
le résultat précédent montre alors que $E \times_Z F$ est un voisinage
quasi-compact de $(x, y)$. Ainsi, $X \times_Z Y$ est un objet de
$\textit{LC}$ d'après
Topologie, lemme \ref{topology-lemma-locally-quasi-compact-Hausdorff}.
\end{proof}

\noindent
On peut munir $\textit{LC}$ d'une topologie plus fine que la topologie usuelle.

\begin{definition}
\label{definition-covering-LC}
Soit $\{f_i : X_i \to X\}$ une famille de morphismes de but fixé dans la
catégorie $\textit{LC}$. On dit que cette famille est un
{\it recouvrement qc}\footnote{Cette notation n'est pas standard.
Nous l'avons choisie pour rappeler au lecteur les recouvrements fpqc
de schémas.}
si, pour tout $x \in X$, il existe $i_1, \ldots, i_n \in I$ et des parties
quasi-compactes $E_j \subset X_{i_j}$ telles que
$\bigcup f_{i_j}(E_j)$ soit un voisinage de $x$.
\end{definition}

\noindent
Remarquons qu'un recouvrement ouvert $X = \bigcup U_i$ d'un objet de
$\textit{LC}$ fournit un recouvrement qc $\{U_i \to X\}$, puisque $X$ est
localement quasi-compact. Commençons par le lemme obligé.

\begin{lemma}
\label{lemma-qc}
Soit $X$ un espace séparé et localement quasi-compact, autrement dit un
objet de $\textit{LC}$.
\begin{enumerate}
\item Si $X' \to X$ est un isomorphisme dans $\textit{LC}$, alors
$\{X' \to X\}$ est un recouvrement qc.
\item Si $\{f_i : X_i \to X\}_{i\in I}$ est un recouvrement qc et si, pour
chaque $i$, on a un recouvrement qc
$\{g_{ij} : X_{ij} \to X_i\}_{j\in J_i}$, alors
$\{X_{ij} \to X\}_{i \in I, j\in J_i}$ est un recouvrement qc.
\item Si $\{X_i \to X\}_{i\in I}$ est un recouvrement qc et si
$X' \to X$ est un morphisme de $\textit{LC}$, alors
$\{X' \times_X X_i \to X'\}_{i\in I}$ est un recouvrement qc.
\end{enumerate}
\end{lemma}

\begin{proof}
Le point (1) résulte de la remarque précédente selon laquelle les
recouvrements ouverts sont des recouvrements qc.

\medskip\noindent
Démonstration de (2). Soit $x \in X$. Choisissons
$i_1, \ldots, i_n \in I$ et des parties quasi-compactes
$E_a \subset X_{i_a}$ telles que $\bigcup f_{i_a}(E_a)$ soit un voisinage
de $x$. Pour tout $e \in E_a$, on peut trouver une partie finie
$J_e \subset J_{i_a}$ et des parties quasi-compactes
$F_{e, j} \subset X_{i_a j}$, $j \in J_e$, telles que
$\bigcup g_{i_a j}(F_{e, j})$ soit un voisinage de $e$. Comme $E_a$ est
quasi-compact, on trouve une famille finie
$e_1, \ldots, e_{m_a}$ telle que
$$
E_a \subset
\bigcup\nolimits_{k = 1, \ldots, m_a}
\bigcup\nolimits_{j \in J_{e_k}} g_{i_a j}(F_{e_k, j})
$$
On voit alors que
$$
\bigcup\nolimits_{a = 1, \ldots, n}
\bigcup\nolimits_{k = 1, \ldots, m_a}
\bigcup\nolimits_{j \in J_{e_k}} f_{i_a}(g_{i_a j}(F_{e_k, j}))
$$
est un voisinage de $x$.

\medskip\noindent
Démonstration de (3). Soit $x' \in X'$ un point, et soit $x \in X$ son
image. Choisissons $i_1, \ldots, i_n \in I$ et des parties quasi-compactes
$E_j \subset X_{i_j}$ telles que $\bigcup f_{i_j}(E_j)$ soit un voisinage
de $x$. Choisissons un voisinage quasi-compact $F \subset X'$ de $x'$ dont
l'image soit contenue dans le voisinage quasi-compact
$\bigcup f_{i_j}(E_j)$ de $x$. Alors
$F \times_X E_j \subset X' \times_X X_{i_j}$ est une partie quasi-compacte,
et $F$ est l'image de l'application
$\coprod F \times_X E_j \to F$. Le changement de base est donc un
recouvrement qc, ce qui achève la démonstration.
\end{proof}

\noindent
Puisque tous les objets de $\textit{LC}$ sont séparés, tout morphisme
$f : X \to Y$ de $\textit{LC}$ est une application continue séparée
d'espaces topologiques. Ainsi, $f$ est une application propre d'espaces
topologiques si et seulement si $f$ est universellement fermée. Voir la
discussion dans Topologie, section \ref{topology-section-proper}.

\begin{lemma}
\label{lemma-proper-surjective-is-qc-covering}
Soit $f : X \to Y$ un morphisme de $\textit{LC}$. Si $f$ est propre et
surjectif, alors $\{f : X \to Y\}$ est un recouvrement qc.
\end{lemma}

\begin{proof}
Soit $y \in Y$ un point. Pour tout $x \in X_y$, choisissons un voisinage
quasi-compact $E_x \subset X$, puis un ouvert $x \in U_x \subset E_x$.
Comme $f$ est propre, la fibre $X_y$ est quasi-compacte ; on trouve donc
$x_1, \ldots, x_n \in X_y$ tels que
$X_y \subset U_{x_1} \cup \ldots \cup U_{x_n}$.
Nous affirmons que $f(E_{x_1}) \cup \ldots \cup f(E_{x_n})$ est un voisinage
de $y$. En effet, comme $f$ est fermée
(Topologie, théorème \ref{topology-theorem-characterize-proper}),
on voit que
$Z = f(X \setminus (U_{x_1} \cup \ldots \cup U_{x_n}))$
est une partie fermée de $Y$ qui ne contient pas $y$. Comme $f$ est
surjective, on voit que $Y \setminus Z$ est contenu dans
$f(E_{x_1}) \cup \ldots \cup f(E_{x_n})$, comme voulu.
\end{proof}

\noindent
Mis à part quelques questions ensemblistes, le lemme \ref{lemma-qc} montre
que $\textit{LC}$, munie de la collection des recouvrements qc, forme un
site. Nous noterons $\textit{LC}_{qc}$ ce site (convenablement modifié pour
résoudre les questions ensemblistes).

\begin{remark}[Questions ensemblistes]
\label{remark-set-theoretic-LC}
La catégorie $\textit{LC}$ est une «~grande~» catégorie, car ses objets
forment une classe propre. De même, les recouvrements forment une classe
propre. Définissons la {\it taille} d'un espace topologique $X$ comme le
cardinal de l'ensemble des points de $X$. Choisissons une fonction $Bound$ sur
les cardinaux, par exemple comme dans
Ensembles, équation (\ref{sets-equation-bound}).
Enfin, soit $S_0$ un ensemble initial d'objets de $\textit{LC}$, par exemple
$S_0 = \{(\mathbf{R}, \text{topologie euclidienne})\}$.
Exactement comme dans Ensembles, lemme \ref{sets-lemma-construct-category},
on peut choisir un ordinal limite $\alpha$ tel que
$\textit{LC}_\alpha = \textit{LC} \cap V_\alpha$ contienne $S_0$ et soit
stable par toutes les limites et colimites dénombrables qui existent dans
$\textit{LC}$. De plus, si $X \in \textit{LC}_\alpha$, si
$Y \in \textit{LC}$ et si
$\text{size}(Y) \leq Bound(\text{size}(X))$, alors $Y$ est isomorphe à un
objet de $\textit{LC}_\alpha$.
Appliquons ensuite Ensembles, lemme \ref{sets-lemma-coverings-site}, pour
choisir un ensemble $\text{Cov}$ de recouvrements qc sur
$\textit{LC}_\alpha$ tel que tout recouvrement qc dans
$\textit{LC}_\alpha$ soit combinatoirement équivalent à un recouvrement
appartenant à cet ensemble. On obtient ainsi un site
$(\textit{LC}_\alpha, \text{Cov})$, que l'on notera $\textit{LC}_{qc}$.
\end{remark}

\noindent
Il existe une seconde topologie sur le site $\textit{LC}_{qc}$ de la
remarque \ref{remark-set-theoretic-LC}. Pour un objet $X$, on considère tous
les recouvrements $\{X_i \to X\}$ de $\textit{LC}_{qc}$ tels que
$X_i \to X$ soit une immersion ouverte. Nous notons ce site
$\textit{LC}_{Zar}$. Le foncteur identité
$\textit{LC}_{Zar} \to \textit{LC}_{qc}$ est continu et définit un morphisme
de sites
$$
\epsilon : \textit{LC}_{qc} \longrightarrow \textit{LC}_{Zar}
$$
Voir la section \ref{section-compare}.
Pour un espace topologique séparé et localement quasi-compact $X$, plus
précisément pour $X \in \Ob(\textit{LC}_{qc})$, nous notons le morphisme
induit
$$
\epsilon_X : \textit{LC}_{qc}/X \longrightarrow \textit{LC}_{Zar}/X
$$
(voir Sites, lemme \ref{sites-lemma-localize-morphism}).
Soit $X_{Zar}$ le site dont les objets sont les ouverts de $X$, voir
Sites, exemple \ref{sites-example-site-topological}.
Il existe un morphisme de sites
$$
\pi_X : \textit{LC}_{Zar}/X \longrightarrow X_{Zar}
$$
défini par le foncteur continu
$X_{Zar} \to \textit{LC}_{Zar}/X$, $U \mapsto U$.
En effet, $X_{Zar}$ possède des produits fibrés et un objet final, le
foncteur ci-dessus y commute, et Sites, proposition
\ref{sites-proposition-get-morphism}, s'applique.
Nous considérons souvent $\pi$ comme un morphisme de topos
$$
\pi_X : \Sh(\textit{LC}_{Zar}/X) \longrightarrow \Sh(X)
$$
en utilisant l'égalité $\Sh(X_{Zar}) = \Sh(X)$.

\begin{lemma}
\label{lemma-describe-pullback-pi}
Soit $X$ un objet de $\textit{LC}_{qc}$ et soit $\mathcal{F}$ un faisceau
sur $X$. La règle
$$
\textit{LC}_{qc}/X \longrightarrow \textit{Sets},\quad
(f : Y \to X) \longmapsto \Gamma(Y, f^{-1}\mathcal{F})
$$
est un faisceau et, a fortiori, aussi un faisceau sur
$\textit{LC}_{Zar}/X$. Ce faisceau est égal à
$\pi_X^{-1}\mathcal{F}$ sur $\textit{LC}_{Zar}/X$ et à
$\epsilon_X^{-1}\pi_X^{-1}\mathcal{F}$ sur $\textit{LC}_{qc}/X$.
\end{lemma}

\begin{proof}
Notons $\mathcal{G}$ le préfaisceau donné par la formule du lemme.
Bien entendu, l'image inverse $f^{-1}$ figurant dans la formule désigne
l'image inverse usuelle des faisceaux sur les espaces topologiques.
Il découle immédiatement des définitions que $\mathcal{G}$ est un faisceau
pour la topologie Zar.

\medskip\noindent
Soit $Y \to X$ un morphisme de $\textit{LC}_{qc}$ et soit
$\mathcal{V} = \{g_i : Y_i \to Y\}_{i \in I}$ un recouvrement qc.
Pour montrer que $\mathcal{G}$ est un faisceau pour la topologie qc, il
suffit de montrer que
$\mathcal{G}(Y) \to H^0(\mathcal{V}, \mathcal{G})$
est un isomorphisme, voir Sites, section
\ref{sites-section-sheafification}. Remarquons d'abord que le morphisme est
injectif, puisqu'un recouvrement qc est surjectif et que l'égalité des
sections se détecte sur les germes (utiliser Faisceaux, lemmes
\ref{sheaves-lemma-sheaf-subset-stalks} et
\ref{sheaves-lemma-stalk-pullback-presheaf}). Ainsi, $\mathcal{G}$ est un
préfaisceau séparé sur $\textit{LC}_{qc}$ ; il suffit donc de montrer que
tout élément
$(s_i) \in H^0(\mathcal{V}, \mathcal{G})$
s'envoie sur un élément de l'image de $\mathcal{G}(Y)$ après remplacement de
$\mathcal{V}$ par un raffinement
(Sites, théorème \ref{sites-theorem-plus}).

\medskip\noindent
En identifiant les faisceaux sur $Y_{i, Zar}$ aux faisceaux sur $Y_i$, nous
voyons que $\mathcal{G}|_{Y_{i, Zar}}$ est l'image inverse de
$f^{-1}\mathcal{F}$ par l'application continue $g_i : Y_i \to Y$. Nous
pouvons donc choisir un recouvrement ouvert $Y_i = \bigcup V_{ij}$ tel que,
pour chaque $j$, il existe un ouvert $W_{ij} \subset Y$ et une section
$t_{ij} \in \mathcal{G}(W_{ij})$, avec $V_{ij}$ envoyé dans $W_{ij}$ et
$s|_{V_{ij}}$ image inverse de $t_{ij}$. Autrement dit, après avoir raffiné
le recouvrement $\{Y_i \to Y\}$, nous pouvons supposer qu'il existe des
ouverts $W_i \subset Y$ tels que $Y_i \to Y$ se factorise par $W_i$, et des
sections $t_i$ de $\mathcal{G}$ sur $W_i$ dont les restrictions sont les
sections données $s_i$. De plus, si $y \in Y$ appartient aux images de
$Y_i \to Y$ et de $Y_j \to Y$, alors les images $t_{i, y}$ et $t_{j, y}$
dans le germe $f^{-1}\mathcal{F}_y$ sont égales (car $s_i$ et $s_j$ sont
égales sur $Y_i \times_Y Y_j$). Ainsi, pour $y \in Y$, il existe un élément
bien défini $t_y$ de $f^{-1}\mathcal{F}_y$, égal à $t_{i, y}$ dès que $y$
appartient à l'image de $Y_i \to Y$.
Nous allons montrer que l'élément $(t_y)$ provient d'une section globale de
$f^{-1}\mathcal{F}$ sur $Y$, ce qui achèvera la démonstration du lemme.

\medskip\noindent
Il suffit de le montrer localement sur $Y$, voir Faisceaux, section
\ref{sheaves-section-sheafification}. Soit $y_0 \in Y$. Choisissons
$i_1, \ldots, i_n \in I$ et des parties quasi-compactes
$E_j \subset Y_{i_j}$ telles que $\bigcup g_{i_j}(E_j)$ soit un voisinage de
$y_0$. Soit $V \subset Y$ un voisinage ouvert de $y_0$, contenu dans
$\bigcup g_{i_j}(E_j)$ et dans $W_{i_1} \cap \ldots \cap W_{i_n}$.
Comme $t_{i_1, y_0} = \ldots = t_{i_n, y_0}$, après avoir rétréci $V$, nous
pouvons supposer que les sections $t_{i_j}|_V$, $j = 1, \ldots, n$, de
$f^{-1}\mathcal{F}$ sont égales. Puisque
$V \subset \bigcup g_{i_j}(E_j)$, nous voyons que $(t_y)_{y \in V}$ provient
de cette section.

\medskip\noindent
Il reste à montrer que $\mathcal{G}$ est égal à
$\epsilon_X^{-1}\pi_X^{-1}\mathcal{F}$ sur $\textit{LC}_{qc}$,
resp.\ à $\pi_X^{-1}\mathcal{F}$ sur $\textit{LC}_{Zar}$.
Dans les deux cas, l'image inverse est définie en prenant le préfaisceau
$$
(f : Y \to X)
\longmapsto
\colim_{f(Y) \subset U \subset X} \mathcal{F}(U)
$$
puis en le faisceautisant. La faisceautisation pour la topologie Zar produit
exactement notre faisceau $\mathcal{G}$ ; le fait que $\mathcal{G}$ soit un
faisceau qc montre que cela fonctionne aussi pour la topologie qc.
\end{proof}

\noindent
Soit $X \in \Ob(\textit{LC}_{Zar})$ et soit $\mathcal{H}$ un faisceau
abélien sur $\textit{LC}_{Zar}/X$. Nous noterons
$H^n_{Zar}(U, \mathcal{H})$ la cohomologie de $\mathcal{H}$ au-dessus d'un
objet $U$ de $\textit{LC}_{Zar}/X$.

\begin{lemma}
\label{lemma-collect-true-things-Zar}
Soit $X$ un objet de $\textit{LC}_{Zar}$. Alors :
\begin{enumerate}
\item pour $\mathcal{F} \in \textit{Ab}(X)$, nous avons
$H^n_{Zar}(X, \pi_X^{-1}\mathcal{F}) = H^n(X, \mathcal{F})$,
\item $\pi_{X, *} : \textit{Ab}(\textit{LC}_{Zar}/X) \to \textit{Ab}(X)$
est exact,
\item l'unité $\text{id} \to \pi_{X, *} \circ \pi_X^{-1}$ de l'adjonction
est un isomorphisme, et
\item pour $K \in D(X)$, le morphisme canonique
$K \to R\pi_{X, *} \pi_X^{-1}K$ est un isomorphisme.
\end{enumerate}
Soit $f : X \to Y$ un morphisme de $\textit{LC}_{Zar}$. Alors :
\begin{enumerate}
\item[(5)] il existe un diagramme commutatif
$$
\xymatrix{
\Sh(\textit{LC}_{Zar}/X) \ar[r]_{f_{Zar}} \ar[d]_{\pi_X} &
\Sh(\textit{LC}_{Zar}/Y) \ar[d]^{\pi_Y} \\
\Sh(X_{Zar}) \ar[r]^f &
\Sh(Y_{Zar})
}
$$
de topos,
\item[(6)] pour $L \in D^+(Y)$, nous avons
$H^n_{Zar}(X, \pi_Y^{-1}L) = H^n(X, f^{-1}L)$,
\item[(7)] si $f$ est propre, alors :
\begin{enumerate}
\item $\pi_Y^{-1} \circ f_* = f_{Zar, *} \circ \pi_X^{-1}$ comme foncteurs
$\Sh(X) \to \Sh(\textit{LC}_{Zar}/Y)$,
\item $\pi_Y^{-1} \circ Rf_* = Rf_{Zar, *} \circ \pi_X^{-1}$ comme
foncteurs $D^+(X) \to D^+(\textit{LC}_{Zar}/Y)$.
\end{enumerate}
\end{enumerate}
\end{lemma}

\begin{proof}
Démonstration de (1).
L'égalité $H^n_{Zar}(X, \pi_X^{-1}\mathcal{F}) = H^n(X, \mathcal{F})$
est un fait général provenant de l'observation immédiate que les recouvrements
de $X$ dans $\textit{LC}_{Zar}$ sont exactement les recouvrements ouverts de
$X$. Pour une démonstration détaillée, on peut appliquer le
lemme \ref{lemma-cohomology-bigger-site} au foncteur
$X_{Zar} \to \textit{LC}_{Zar}$.

\medskip\noindent
Démonstration de (2). Cela est vrai parce que
$\pi_{X, *} = \tau_X^{-1}$ pour un certain morphisme de topos
$\tau_X : \Sh(X_{Zar}) \to \Sh(\textit{LC}_{Zar})$, comme il résulte de
Sites, lemme \ref{sites-lemma-bigger-site}, appliqué au foncteur
$X_{Zar} \to \textit{LC}_{Zar}/X$ utilisé pour définir $\pi_X$.

\medskip\noindent
Démonstration de (3). Cela est vrai parce que
$\tau_X^{-1} \circ \pi_X^{-1}$ est le foncteur identité d'après Sites,
lemme \ref{sites-lemma-bigger-site}. On peut aussi le déduire de la
description explicite de $\pi_X^{-1}$ dans le
lemme \ref{lemma-describe-pullback-pi}.

\medskip\noindent
Démonstration de (4). Appliquer (3) à un complexe de faisceaux abéliens
représentant $K$.

\medskip\noindent
Démonstration de (5). Le morphisme de topos $f_{Zar}$ provient d'une
application de Sites, lemme \ref{sites-lemma-relocalize}, et, dans notre cas,
du foncteur continu
$Z/Y \mapsto Z \times_Y X/X$ d'après Sites, lemme
\ref{sites-lemma-relocalize-given-fibre-products}.
Le diagramme commute simplement parce que les foncteurs continus
correspondants se composent correctement
(voir Sites, lemme \ref{sites-lemma-composition-morphisms-sites}).

\medskip\noindent
Démonstration de (6). Nous avons
$H^n_{Zar}(X, \pi_Y^{-1}\mathcal{G}) =
H^n_{Zar}(X, f_{Zar}^{-1}\pi_Y^{-1}\mathcal{G})$
pour $\mathcal{G}$ dans $\textit{Ab}(Y)$, voir le
lemme \ref{lemma-cohomology-of-open}. Ceci est égal à
$H^n_{Zar}(X, \pi_X^{-1}f^{-1}\mathcal{G})$ par commutativité du diagramme de
(5). Nous concluons donc par (1) lorsque $L$ est réduit à un seul faisceau en
degré $0$. Le cas général s'obtient en représentant $L$ par un complexe de
faisceaux abéliens borné inférieurement.

\medskip\noindent
Démonstration de (7a). Soit $\mathcal{F}$ un faisceau sur $X$.
Soit $g : Z \to Y$ un objet de $\textit{LC}_{Zar}/Y$. Considérons le produit
fibré
$$
\xymatrix{
Z' \ar[r]_{f'} \ar[d]_{g'} & Z \ar[d]^g \\
X \ar[r]^f & Y
}
$$
Nous avons alors
$$
(f_{Zar, *}\pi_X^{-1}\mathcal{F})(Z/Y) =
(\pi_X^{-1}\mathcal{F})(Z'/X)  =
\Gamma(Z', (g')^{-1}\mathcal{F})  =
\Gamma(Z, f'_*(g')^{-1}\mathcal{F})
$$
la deuxième égalité venant du lemme \ref{lemma-describe-pullback-pi}.
D'autre part,
$$
(\pi_Y^{-1}f_*\mathcal{F})(Z/Y) = \Gamma(Z, g^{-1}f_*\mathcal{F})
$$
encore d'après le lemme \ref{lemma-describe-pullback-pi}.
Le changement de base propre pour les faisceaux d'ensembles
(Cohomologie, lemme
\ref{cohomology-lemma-proper-base-change-sheaves-of-sets}) montre donc que
les deux ensembles sont canoniquement isomorphes. Cet isomorphisme est
compatible aux applications de restriction et définit un isomorphisme
$\pi_Y^{-1}f_*\mathcal{F} = f_{Zar, *}\pi_X^{-1}\mathcal{F}$.
et donc un isomorphisme de foncteurs
$\pi_Y^{-1} \circ f_* = f_{Zar, *} \circ \pi_X^{-1}$.

\medskip\noindent
Démonstration de (7b). Soit $K \in D^+(X)$. D'après le
lemme \ref{lemma-unbounded-describe-higher-direct-images}, le $n$-ième
faisceau de cohomologie de $Rf_{Zar, *}\pi_X^{-1}K$ est le faisceau associé
au préfaisceau
$$
(g : Z \to Y) \longmapsto H^n_{Zar}(Z', \pi_X^{-1}K)
$$
avec les notations ci-dessus. Observons que
\begin{align*}
H^n_{Zar}(Z', \pi_X^{-1}K)
& =
H^n(Z', (g')^{-1}K) \\
& =
H^n(Z, Rf'_*(g')^{-1}K) \\
& =
H^n(Z, g^{-1}Rf_*K) \\
& =
H^n_{Zar}(Z, \pi_Y^{-1}Rf_*K)
\end{align*}
La première égalité est (6) appliqué à $K$ et à $g' : Z' \to X$.
La deuxième est Leray pour $f' : Z' \to Z$
(Cohomologie, lemme \ref{cohomology-lemma-before-Leray}).
La troisième est le théorème de changement de base propre
(Cohomologie, théorème \ref{cohomology-theorem-proper-base-change}).
La quatrième est (6) appliqué à $g : Z \to Y$ et à $Rf_*K$.
Ainsi, $Rf_{Zar, *}\pi_X^{-1}K$ et $\pi_Y^{-1}Rf_*K$ ont les mêmes faisceaux
de cohomologie. Nous omettons la vérification que le morphisme canonique de
changement de base $\pi_Y^{-1}Rf_*K \to Rf_{Zar, *}\pi_X^{-1}K$ induit cet
isomorphisme.
\end{proof}

\noindent
Dans la situation du lemme \ref{lemma-describe-pullback-pi}, la composition
de $\epsilon$ et $\pi$ et l'égalité $\Sh(X) = \Sh(X_{Zar})$ déterminent un
morphisme de topos
$$
a_X : \Sh(\textit{LC}_{qc}/X) \longrightarrow \Sh(X)
$$

\begin{lemma}
\label{lemma-push-pull-LC}
Soit $f : X \to Y$ un morphisme de $\textit{LC}_{qc}$. Il existe alors des
diagrammes commutatifs de topos
$$
\vcenter{
\xymatrix{
\Sh(\textit{LC}_{qc}/X) \ar[r]_{f_{qc}} \ar[d]_{\epsilon_X} &
\Sh(\textit{LC}_{qc}/Y) \ar[d]^{\epsilon_Y} \\
\Sh(\textit{LC}_{Zar}/X) \ar[r]^{f_{Zar}} &
\Sh(\textit{LC}_{Zar}/Y)
}
}
\quad\text{et}\quad
\vcenter{
\xymatrix{
\Sh(\textit{LC}_{qc}/X) \ar[r]_{f_{qc}} \ar[d]_{a_X} &
\Sh(\textit{LC}_{qc}/Y) \ar[d]^{a_Y} \\
\Sh(X) \ar[r]^f &
\Sh(Y)
}
}
$$
avec $a_X = \pi_X \circ \epsilon_X$, $a_Y = \pi_Y \circ \epsilon_Y$.
Si $f$ est propre, alors $a_Y^{-1} \circ f_* = f_{qc, *} \circ a_X^{-1}$.
\end{lemma}

\begin{proof}
Le morphisme de topos $f_{qc}$ est celui de Sites, lemme
\ref{sites-lemma-relocalize}, qui, dans notre cas, provient du foncteur
continu $Z/Y \mapsto Z \times_Y X/X$, voir Sites, lemme
\ref{sites-lemma-relocalize-given-fibre-products}.
Le diagramme de gauche commute parce que les foncteurs continus
correspondants se composent correctement
(voir Sites, lemme \ref{sites-lemma-composition-morphisms-sites}).
Le diagramme de droite commute parce que celui de gauche commute et à cause
du point (5) du lemme \ref{lemma-collect-true-things-Zar}.

\medskip\noindent
Démonstration de la dernière assertion. On pourrait répéter la démonstration
du point (7a) du lemme \ref{lemma-collect-true-things-Zar} ; nous allons
plutôt l'en déduire. Comme $\epsilon_{Y, *}$ est le foncteur identité sur les
préfaisceaux sous-jacents, il reflète les isomorphismes. La description du
lemme \ref{lemma-describe-pullback-pi} montre que
$\epsilon_{Y, *} \circ a_Y^{-1} = \pi_Y^{-1}$, et de même pour $X$.
Pour montrer que le morphisme canonique
$a_Y^{-1}f_*\mathcal{F} \to f_{qc, *}a_X^{-1}\mathcal{F}$
est un isomorphisme, il suffit de montrer que
$$
\pi_Y^{-1}f_*\mathcal{F} =
\epsilon_{Y, *}a_Y^{-1}f_*\mathcal{F} \to
\epsilon_{Y, *}f_{qc, *}a_X^{-1}\mathcal{F} =
f_{Zar, *}\epsilon_{X, *}a_X^{-1}\mathcal{F} =
f_{Zar, *}\pi_X^{-1}\mathcal{F}
$$
est un isomorphisme. C'est le point (7a) du
lemme \ref{lemma-collect-true-things-Zar}.
\end{proof}

\begin{lemma}
\label{lemma-compare-qc-zar}
Considérons le morphisme de comparaison
$\epsilon : \textit{LC}_{qc} \to \textit{LC}_{Zar}$.
Notons $\mathcal{P}$ la classe des applications propres d'espaces
topologiques. Pour $X$ dans $\textit{LC}_{Zar}$, notons
$\mathcal{A}'_X \subset \textit{Ab}(\textit{LC}_{Zar}/X)$
la sous-catégorie pleine formée des faisceaux de la forme
$\pi_X^{-1}\mathcal{F}$ avec $\mathcal{F}$ dans $\textit{Ab}(X)$.
Alors les propriétés
(\ref{item-base-change-P}),
(\ref{item-restriction-A}),
(\ref{item-A-sheaf}),
(\ref{item-A-and-P}) et
(\ref{item-refine-tau-by-P})
de la situation \ref{situation-compare} sont vérifiées.
\end{lemma}

\begin{proof}
Nous montrons d'abord que
$\mathcal{A}'_X \subset \textit{Ab}(\textit{LC}_{Zar}/X)$ est une
sous-catégorie de Serre faible en vérifiant les conditions (1), (2), (3) et
(4) de Homologie, lemme
\ref{homology-lemma-characterize-weak-serre-subcategory}.
Les points (1), (2) et (3) sont immédiats, car $\pi_X^{-1}$ est exact et
pleinement fidèle d'après le point (3) du
lemme \ref{lemma-collect-true-things-Zar}. Si
$0 \to \pi_X^{-1}\mathcal{F} \to \mathcal{G} \to \pi_X^{-1}\mathcal{F}' \to 0$
est une suite exacte courte dans $\textit{Ab}(\textit{LC}_{Zar}/X)$, alors
$0 \to \mathcal{F} \to \pi_{X, *}\mathcal{G} \to \mathcal{F}' \to 0$ est
exacte d'après le point (2) du lemme
\ref{lemma-collect-true-things-Zar}. Par conséquent,
$\mathcal{G} = \pi_X^{-1}\pi_{X, *}\mathcal{G}$ appartient à
$\mathcal{A}'_X$, ce qui vérifie la dernière condition.

\medskip\noindent
La propriété (\ref{item-base-change-P}) résulte du lemme
\ref{lemma-LC-basic} et du fait que le changement de base d'une application
propre est propre (voir Topologie, théorème
\ref{topology-theorem-characterize-proper} et lemme
\ref{topology-lemma-base-change-separated}).

\medskip\noindent
La propriété (\ref{item-restriction-A}) résulte du diagramme commutatif (5)
du lemme \ref{lemma-collect-true-things-Zar}.

\medskip\noindent
La propriété (\ref{item-A-sheaf}) est le
lemme \ref{lemma-describe-pullback-pi}.

\medskip\noindent
La propriété (\ref{item-A-and-P}) est le point (7)(b) du
lemme \ref{lemma-collect-true-things-Zar}.

\medskip\noindent
Démonstration de (\ref{item-refine-tau-by-P}). Soit donné un recouvrement qc
$\{U_i \to U\}$. Pour $u \in U$, choisissons $i_1, \ldots, i_m \in I$ et
des parties quasi-compactes $E_j \subset U_{i_j}$ telles que
$\bigcup f_{i_j}(E_j)$ soit un voisinage de $u$.
Observons que $Y = \coprod_{j = 1, \ldots, m} E_j \to U$ est propre comme
application continue entre espaces séparés quasi-compacts
(Topologie, lemme \ref{topology-lemma-closed-map}).
Choisissons un voisinage ouvert $u \in V$ contenu dans
$\bigcup f_{i_j}(E_j)$. Alors $Y \times_U V \to V$ est un morphisme propre
surjectif, donc un recouvrement $qc$ d'après le
lemme \ref{lemma-proper-surjective-is-qc-covering}.
Comme nous pouvons le faire pour tout $u \in U$, la propriété
(\ref{item-refine-tau-by-P}) est vérifiée.
\end{proof}

\begin{lemma}
\label{lemma-V-C-all-n}
Avec les notations ci-dessus :
\begin{enumerate}
\item Pour $X \in \Ob(\textit{LC}_{qc})$ et un faisceau abélien
$\mathcal{F}$ sur $X$, nous avons
$\epsilon_{X, *}a_X^{-1}\mathcal{F} = \pi_X^{-1}\mathcal{F}$ et
$R^i\epsilon_{X, *}(a_X^{-1}\mathcal{F}) = 0$ pour $i > 0$.
\item Pour un morphisme propre $f : X \to Y$ de $\textit{LC}_{qc}$ et un
faisceau abélien $\mathcal{F}$ sur $X$, nous avons
$a_Y^{-1}(R^if_*\mathcal{F}) = R^if_{qc, *}(a_X^{-1}\mathcal{F})$
pour tout $i$.
\item Pour $X \in \Ob(\textit{LC}_{qc})$ et $K$ dans $D^+(X)$, le
morphisme $\pi_X^{-1}K \to R\epsilon_{X, *}(a_X^{-1}K)$ est un isomorphisme.
\item Pour un morphisme propre $f : X \to Y$ de $\textit{LC}_{qc}$ et $K$
dans $D^+(X)$, nous avons
$a_Y^{-1}(Rf_*K) = Rf_{qc, *}(a_X^{-1}K)$.
\end{enumerate}
\end{lemma}

\begin{proof}
Par le lemme \ref{lemma-compare-qc-zar}, tous les lemmes de la
section \ref{section-compare-general} s'appliquent dans notre situation.
Pour traduire les résultats, observons que la catégorie $\mathcal{A}_X$ du
lemme \ref{lemma-A} est l'image essentielle de
$a_X^{-1} : \textit{Ab}(X) \to \textit{Ab}(\textit{LC}_{qc}/X)$.

\medskip\noindent
Le point (1) équivaut à $(V_n)$ pour tout $n$, ce qui est vérifié d'après le
lemme \ref{lemma-V-C-all-n-general}.

\medskip\noindent
Le point (2) résulte de l'application de $\epsilon_Y^{-1}$ à la conclusion
du lemme \ref{lemma-V-implies-C-general}.

\medskip\noindent
Le point (3) résulte du point (1) du
lemme \ref{lemma-V-C-all-n-general}, car $\pi_X^{-1}K$ appartient à
$D^+_{\mathcal{A}'_X}(\textit{LC}_{Zar}/X)$ et
$a_X^{-1} = \epsilon_X^{-1} \circ \pi_X^{-1}$.

\medskip\noindent
Le point (4) résulte du point (2) du
lemme \ref{lemma-V-C-all-n-general} pour la même raison.
\end{proof}

\begin{lemma}
\label{lemma-cohomological-descent-LC}
Soit $X$ un objet de $\textit{LC}_{qc}$. Pour $K \in D^+(X)$, le morphisme
$$
K \longrightarrow Ra_{X, *}a_X^{-1}K
$$
est un isomorphisme, où
$a_X : \Sh(\textit{LC}_{qc}/X) \to \Sh(X)$ est comme ci-dessus.
\end{lemma}

\begin{proof}
Nous nous ramenons d'abord au cas où $K$ est donné par un seul faisceau
abélien. Représentons en effet $K$ par un complexe borné inférieurement
$\mathcal{F}^\bullet$. Le cas d'un faisceau montre que
$\mathcal{F}^n = a_{X, *} a_X^{-1} \mathcal{F}^n$
et que les faisceaux $R^qa_{X, *}a_X^{-1}\mathcal{F}^n$ sont nuls pour
$q > 0$. En appliquant le lemme d'acyclicité de Leray
(Catégories dérivées, lemme \ref{derived-lemma-leray-acyclicity}) à
$a_X^{-1}\mathcal{F}^\bullet$ et au foncteur $a_{X, *}$, nous concluons.
Désormais, supposons $K = \mathcal{F}$.

\medskip\noindent
Par le lemme \ref{lemma-describe-pullback-pi}, nous avons
$a_{X, *}a_X^{-1}\mathcal{F} = \mathcal{F}$. Il suffit donc de montrer que
$R^qa_{X, *}a_X^{-1}\mathcal{F} = 0$ pour $q > 0$.
Pour cela, nous pouvons utiliser $a_X = \pi_X \circ \epsilon_X$ et la suite
spectrale de Leray du lemme \ref{lemma-relative-Leray}.
Par le lemme \ref{lemma-V-C-all-n}, nous avons
$R^i\epsilon_{X, *}(a_X^{-1}\mathcal{F}) = 0$ pour $i > 0$ et
$\epsilon_{X, *}a_X^{-1}\mathcal{F} = \pi_X^{-1}\mathcal{F}$.
Par le lemme \ref{lemma-collect-true-things-Zar}, nous avons
$R^j\pi_{X, *}(\pi_X^{-1}\mathcal{F}) = 0$ pour $j > 0$.
Cela achève la démonstration.
\end{proof}

\begin{lemma}
\label{lemma-compare-cohomology-LC}
Avec $X \in \Ob(\textit{LC}_{qc})$ et
$a_X : \Sh(\textit{LC}_{qc}/X) \to \Sh(X)$ comme ci-dessus :
\begin{enumerate}
\item pour un faisceau abélien $\mathcal{F}$ sur $X$, nous avons
$H^n(X, \mathcal{F}) = H^n_{qc}(X, a_X^{-1}\mathcal{F})$,
\item pour $K \in D^+(X)$, nous avons
$H^n(X, K) = H^n_{qc}(X, a_X^{-1}K)$.
\end{enumerate}
Par exemple, si $A$ est un groupe abélien, alors nous avons
$H^n(X, \underline{A}) = H^n_{qc}(X, \underline{A})$.
\end{lemma}

\begin{proof}
Cela résulte du lemme \ref{lemma-cohomological-descent-LC} et de la
remarque \ref{remark-before-Leray}.
\end{proof}











\section{Suites spectrales pour Ext}
\label{section-spectral-sequence-ext}

\noindent
Dans cette section, nous rassemblons diverses suites spectrales qui interviennent
lorsque l'on considère les foncteurs Ext. Pour toute paire de complexes
$\mathcal{G}^\bullet, \mathcal{F}^\bullet$ de modules sur un site annelé
$(\mathcal{C}, \mathcal{O})$, nous notons
$$
\Ext^n_\mathcal{O}(\mathcal{G}^\bullet, \mathcal{F}^\bullet)
=
\Hom_{D(\mathcal{O})}(\mathcal{G}^\bullet, \mathcal{F}^\bullet[n])
$$
conformément à nos conventions générales de
Catégories dérivées, section \ref{derived-section-ext}.

\begin{example}
\label{example-hom-complex-into-sheaf}
Soit $(\mathcal{C}, \mathcal{O})$ un site annelé.
Soit $\mathcal{K}^\bullet$ un complexe de $\mathcal{O}$-modules borné supérieurement.
Soit $\mathcal{F}$ un $\mathcal{O}$-module. Il existe alors une
suite spectrale dont la page $E_2$ est
$$
E_2^{i, j} =
\Ext_\mathcal{O}^i(H^{-j}(\mathcal{K}^\bullet), \mathcal{F})
\Rightarrow
\Ext_\mathcal{O}^{i + j}(\mathcal{K}^\bullet, \mathcal{F})
$$
et une autre suite spectrale dont la page $E_1$ est
$$
E_1^{i, j} =
\Ext_\mathcal{O}^j(\mathcal{K}^{-i}, \mathcal{F})
\Rightarrow
\Ext_\mathcal{O}^{i + j}(\mathcal{K}^\bullet, \mathcal{F}).
$$
Pour construire ces suites spectrales, choisissons une résolution injective
$\mathcal{F} \to \mathcal{I}^\bullet$ et considérons les deux suites
spectrales provenant du complexe double
$\Hom_\mathcal{O}(\mathcal{K}^\bullet, \mathcal{I}^\bullet)$ ; voir
Homologie, section \ref{homology-section-double-complex}.
\end{example}








\section{Cup-produit}
\label{section-cup-product}

\noindent
Soit $(\mathcal{C}, \mathcal{O})$ un site annelé. Soient
$K, M$ des objets de $D(\mathcal{O})$. Posons
$A = \Gamma(\mathcal{C}, \mathcal{O})$. Dans ce cadre, le cup-produit (global)
est un morphisme
$$
R\Gamma(\mathcal{C}, K) \otimes_A^\mathbf{L}
R\Gamma(\mathcal{C}, M)
\longrightarrow
R\Gamma(\mathcal{C}, K \otimes_\mathcal{O}^\mathbf{L} M)
$$
dans $D(A)$. Nous le définissons comme le cup-produit relatif associé
au morphisme de topos annelés
$(\Sh(\mathcal{C}), \mathcal{O}) \to (\Sh(pt), A)$
comme dans la remarque \ref{remark-cup-product}.

\medskip\noindent
Formulons et démontrons une compatibilité naturelle du
cup-produit relatif. Supposons donc que l'on dispose d'un morphisme
$f : (\Sh(\mathcal{C}), \mathcal{O}_\mathcal{C}) \to
(\Sh(\mathcal{D}), \mathcal{O}_\mathcal{D})$ de topos annelés.
Soient $\mathcal{K}^\bullet$ et $\mathcal{M}^\bullet$
des complexes de $\mathcal{O}_\mathcal{C}$-modules.
Il existe un cup-produit naïf
$$
\text{Tot}(
f_*\mathcal{K}^\bullet
\otimes_{\mathcal{O}_\mathcal{D}}
f_*\mathcal{M}^\bullet)
\longrightarrow
f_*\text{Tot}(\mathcal{K}^\bullet
\otimes_{\mathcal{O}_\mathcal{C}}
\mathcal{M}^\bullet)
$$
Nous affirmons que celui-ci est lié au cup-produit relatif.

\begin{lemma}
\label{lemma-cup-compatible-with-naive}
Dans la situation ci-dessus, le diagramme suivant est commutatif
$$
\xymatrix{
f_*\mathcal{K}^\bullet
\otimes_{\mathcal{O}_\mathcal{D}}^\mathbf{L}
f_*\mathcal{M}^\bullet \ar[r] \ar[d]
&
Rf_*\mathcal{K}^\bullet
\otimes_{\mathcal{O}_\mathcal{D}}^\mathbf{L}
Rf_*\mathcal{M}^\bullet \ar[d]^{\text{Remarque \ref{remark-cup-product}}} \\
\text{Tot}(
f_*\mathcal{K}^\bullet
\otimes_{\mathcal{O}_\mathcal{D}}
f_*\mathcal{M}^\bullet) \ar[d]_{\text{cup-produit naïf}} &
Rf_*(\mathcal{K}^\bullet
\otimes_{\mathcal{O}_\mathcal{C}}^\mathbf{L}
\mathcal{M}^\bullet) \ar[d] \\
f_*\text{Tot}(\mathcal{K}^\bullet
\otimes_{\mathcal{O}_\mathcal{C}}
\mathcal{M}^\bullet) \ar[r] &
Rf_*\text{Tot}(\mathcal{K}^\bullet
\otimes_{\mathcal{O}_\mathcal{C}}
\mathcal{M}^\bullet)
}
$$
\end{lemma}

\begin{proof}
D'après la construction de la remarque \ref{remark-cup-product}, on voit qu'en
parcourant le diagramme dans le sens des aiguilles d'une montre, le morphisme
$$
f_*\mathcal{K}^\bullet
\otimes_{\mathcal{O}_\mathcal{D}}^\mathbf{L}
f_*\mathcal{M}^\bullet 
\longrightarrow
Rf_*\text{Tot}(\mathcal{K}^\bullet
\otimes_{\mathcal{O}_\mathcal{C}}
\mathcal{M}^\bullet)
$$
est adjoint au morphisme
\begin{align*}
Lf^*(f_*\mathcal{K}^\bullet
\otimes_{\mathcal{O}_\mathcal{D}}^\mathbf{L}
f_*\mathcal{M}^\bullet)
& =
Lf^*f_*\mathcal{K}^\bullet
\otimes_{\mathcal{O}_\mathcal{C}}^\mathbf{L}
Lf^*f_*\mathcal{M}^\bullet \\
& \to
Lf^*Rf_*\mathcal{K}^\bullet
\otimes_{\mathcal{O}_\mathcal{C}}^\mathbf{L}
Lf^*Rf_*\mathcal{M}^\bullet \\
& \to
\mathcal{K}^\bullet
\otimes_{\mathcal{O}_\mathcal{C}}^\mathbf{L}
\mathcal{M}^\bullet \\
& \to
\text{Tot}(\mathcal{K}^\bullet
\otimes_{\mathcal{O}_\mathcal{C}}
\mathcal{M}^\bullet)
\end{align*}
D'après le lemme \ref{lemma-adjoints-push-pull-compatibility}, ce morphisme est aussi égal à
\begin{align*}
Lf^*(f_*\mathcal{K}^\bullet
\otimes_{\mathcal{O}_\mathcal{D}}^\mathbf{L}
f_*\mathcal{M}^\bullet)
& =
Lf^*f_*\mathcal{K}^\bullet
\otimes_{\mathcal{O}_\mathcal{C}}^\mathbf{L}
Lf^*f_*\mathcal{M}^\bullet \\
& \to
f^*f_*\mathcal{K}^\bullet
\otimes_{\mathcal{O}_\mathcal{C}}^\mathbf{L}
f^*f_*\mathcal{M}^\bullet \\
& \to
\mathcal{K}^\bullet
\otimes_{\mathcal{O}_\mathcal{C}}^\mathbf{L}
\mathcal{M}^\bullet \\
& \to
\text{Tot}(\mathcal{K}^\bullet
\otimes_{\mathcal{O}_\mathcal{C}}
\mathcal{M}^\bullet)
\end{align*}
En parcourant le diagramme dans le sens inverse des aiguilles d'une montre,
on obtient le morphisme adjoint au morphisme
\begin{align*}
Lf^*(f_*\mathcal{K}^\bullet
\otimes_{\mathcal{O}_\mathcal{D}}^\mathbf{L}
f_*\mathcal{M}^\bullet)
& \to
Lf^*\text{Tot}(
f_*\mathcal{K}^\bullet
\otimes_{\mathcal{O}_\mathcal{D}}
f_*\mathcal{M}^\bullet) \\
& \to
Lf^*f_*\text{Tot}(\mathcal{K}^\bullet
\otimes_{\mathcal{O}_\mathcal{C}}
\mathcal{M}^\bullet) \\
& \to
Lf^*Rf_*\text{Tot}(\mathcal{K}^\bullet
\otimes_{\mathcal{O}_\mathcal{C}}
\mathcal{M}^\bullet) \\
& \to
\text{Tot}(\mathcal{K}^\bullet
\otimes_{\mathcal{O}_\mathcal{C}}
\mathcal{M}^\bullet)
\end{align*}
D'après le lemme \ref{lemma-adjoints-push-pull-compatibility}, ce morphisme est aussi égal à
\begin{align*}
Lf^*(f_*\mathcal{K}^\bullet
\otimes_{\mathcal{O}_\mathcal{D}}^\mathbf{L}
f_*\mathcal{M}^\bullet)
& \to
Lf^*\text{Tot}(
f_*\mathcal{K}^\bullet
\otimes_{\mathcal{O}_\mathcal{D}}
f_*\mathcal{M}^\bullet) \\
& \to
Lf^*f_*\text{Tot}(\mathcal{K}^\bullet
\otimes_{\mathcal{O}_\mathcal{C}}
\mathcal{M}^\bullet) \\
& \to
f^*f_*\text{Tot}(\mathcal{K}^\bullet
\otimes_{\mathcal{O}_\mathcal{C}}
\mathcal{M}^\bullet) \\
& \to
\text{Tot}(\mathcal{K}^\bullet
\otimes_{\mathcal{O}_\mathcal{C}}
\mathcal{M}^\bullet)
\end{align*}
Pour achever la preuve, il suffit maintenant de considérer le diagramme
$$
\xymatrix{
Lf^*(f_*\mathcal{K}^\bullet
\otimes_{\mathcal{O}_\mathcal{D}}^\mathbf{L}
f_*\mathcal{M}^\bullet) \ar[d] \ar[rr] & &
Lf^*f_*\mathcal{K}^\bullet \otimes_{\mathcal{O}_\mathcal{C}}^\mathbf{L}
Lf^*f_*\mathcal{M}^\bullet \ar[d] \\
Lf^*\text{Tot}(
f_*\mathcal{K}^\bullet
\otimes_{\mathcal{O}_\mathcal{D}}
f_*\mathcal{M}^\bullet) \ar[d]_{naive} \ar[r] &
f^*\text{Tot}(
f_*\mathcal{K}^\bullet
\otimes_{\mathcal{O}_\mathcal{D}}
f_*\mathcal{M}^\bullet) \ar[ldd]^{naive} \ar[dd] &
f^*f_*\mathcal{K}^\bullet \otimes_{\mathcal{O}_\mathcal{C}}^\mathbf{L}
f^*f_*\mathcal{M}^\bullet \ar[dd] \ar[ldd] \\
Lf^*f_*\text{Tot}(\mathcal{K}^\bullet
\otimes_{\mathcal{O}_\mathcal{C}}
\mathcal{M}^\bullet) \ar[d] \\
f^*f_*\text{Tot}(\mathcal{K}^\bullet \otimes_{\mathcal{O}_\mathcal{C}}
\mathcal{M}^\bullet) \ar[rd] &
\text{Tot}(f^*f_*\mathcal{K}^\bullet \otimes_{\mathcal{O}_\mathcal{C}}
f^*f_*\mathcal{M}^\bullet) \ar[d] &
\mathcal{K}^\bullet \otimes_{\mathcal{O}_\mathcal{C}}^\mathbf{L}
\mathcal{M}^\bullet \ar[ld] \\
& \text{Tot}(\mathcal{K}^\bullet
\otimes_{\mathcal{O}_\mathcal{C}}
\mathcal{M}^\bullet)
}
$$
Tous les polygones de ce diagramme sont commutatifs. Celui du haut l'est
d'après le lemme \ref{lemma-tensor-pull-compatibility}.
Le carré comportant les deux cup-produits naïfs est commutatif parce que
$Lf^* \to f^*$ est fonctoriel en le complexe de modules.
Il en va de même du carré faisant intervenir les deux morphismes
$\mathcal{A}^\bullet \otimes^\mathbf{L} \mathcal{B}^\bullet \to
\text{Tot}(\mathcal{A}^\bullet \otimes \mathcal{B}^\bullet)$.
Enfin, le dernier carré est commutatif au niveau des complexes et cela peut
être considéré comme la définition du cup-produit naïf (par l'adjonction
entre $f^*$ et $f_*$). La preuve est achevée, car les deux parcours
du bord extérieur du diagramme donnent les deux morphismes
décrits ci-dessus.
\end{proof}

\begin{lemma}
\label{lemma-cup-product-associative}
Soit $f : (\Sh(\mathcal{C}), \mathcal{O}) \to (\Sh(\mathcal{C}'), \mathcal{O}')$
un morphisme de topos annelés. Le cup-produit relatif de la
remarque \ref{remark-cup-product} est associatif au sens où
le diagramme
$$
\xymatrix{
Rf_*K \otimes_{\mathcal{O}'}^\mathbf{L}
Rf_*L \otimes_{\mathcal{O}'}^\mathbf{L}
Rf_*M \ar[r] \ar[d] &
Rf_*(K \otimes_\mathcal{O}^\mathbf{L} L)
\otimes_{\mathcal{O}'}^\mathbf{L} Rf_*M \ar[d] \\
Rf_*K \otimes_{\mathcal{O}'}^\mathbf{L}
Rf_*(L \otimes_\mathcal{O}^\mathbf{L} M) \ar[r] &
Rf_*(K \otimes_\mathcal{O}^\mathbf{L} 
L \otimes_\mathcal{O}^\mathbf{L} M)
}
$$
est commutatif dans $D(\mathcal{O}')$ pour tous $K, L, M$ dans $D(\mathcal{O})$.
\end{lemma}

\begin{proof}
En suivant l'un ou l'autre chemin, on obtient le morphisme adjoint au morphisme évident
\begin{align*}
Lf^*(Rf_*K \otimes_{\mathcal{O}'}^\mathbf{L}
Rf_*L \otimes_{\mathcal{O}'}^\mathbf{L}
Rf_*M) & =
Lf^*(Rf_*K) \otimes_\mathcal{O}^\mathbf{L}
Lf^*(Rf_*L) \otimes_\mathcal{O}^\mathbf{L}
Lf^*(Rf_*M) \\
& \to
K \otimes_\mathcal{O}^\mathbf{L} 
L \otimes_\mathcal{O}^\mathbf{L} M
\end{align*}
dans $D(\mathcal{O})$.
\end{proof}

\begin{lemma}
\label{lemma-cup-product-commutative}
Soit $f : (\Sh(\mathcal{C}), \mathcal{O}) \to (\Sh(\mathcal{C}'), \mathcal{O}')$
un morphisme de topos annelés. Le cup-produit relatif de la
remarque \ref{remark-cup-product} est commutatif au sens où
le diagramme
$$
\xymatrix{
Rf_*K \otimes_{\mathcal{O}'}^\mathbf{L} Rf_*L \ar[r] \ar[d]_\psi &
Rf_*(K \otimes_\mathcal{O}^\mathbf{L} L) \ar[d]^{Rf_*\psi} \\
Rf_*L \otimes_{\mathcal{O}'}^\mathbf{L} Rf_*K \ar[r] &
Rf_*(L \otimes_\mathcal{O}^\mathbf{L} K)
}
$$
est commutatif dans $D(\mathcal{O}')$ pour tous $K, L$ dans $D(\mathcal{O})$.
Ici, $\psi$ est la contrainte de commutativité de la catégorie dérivée
(lemme \ref{lemma-symmetric-monoidal-derived}).
\end{lemma}

\begin{proof}
La démonstration est omise.
\end{proof}

\begin{lemma}
\label{lemma-compose-cup-product}
Soient $f : (\Sh(\mathcal{C}), \mathcal{O}) \to (\Sh(\mathcal{C}'), \mathcal{O}')$
et $f' : (\Sh(\mathcal{C}'), \mathcal{O}') \to
(\Sh(\mathcal{C}''), \mathcal{O}'')$
des morphismes de topos annelés. Le cup-produit relatif de la
remarque \ref{remark-cup-product} est compatible aux compositions
au sens où le diagramme
$$
\xymatrix{
R(f' \circ f)_*K \otimes_{\mathcal{O}''}^\mathbf{L} R(f' \circ f)_*L
\ar@{=}[rr] \ar[d] & &
Rf'_*Rf_*K \otimes_{\mathcal{O}''}^\mathbf{L} Rf'_*Rf_*L \ar[d] \\
R(f' \circ f)_*(K \otimes_\mathcal{O}^\mathbf{L} L) \ar@{=}[r] &
Rf'_*Rf_*(K \otimes_\mathcal{O}^\mathbf{L} L) &
Rf'_*(Rf_*K \otimes_{\mathcal{O}'}^\mathbf{L}  Rf_*L) \ar[l]
}
$$
est commutatif dans $D(\mathcal{O}'')$ pour tous $K, L$ dans $D(\mathcal{O})$.
\end{lemma}

\begin{proof}
Cela résulte du fait qu'en suivant l'un ou l'autre chemin dans le diagramme,
on obtient le morphisme adjoint au morphisme
\begin{align*}
& L(f' \circ f)^*\left(R(f' \circ f)_*K
\otimes_{\mathcal{O}''}^\mathbf{L}
R(f' \circ f)_*L\right) \\
& =
L(f' \circ f)^*R(f' \circ f)_*K
\otimes_\mathcal{O}^\mathbf{L}
L(f' \circ f)^*R(f' \circ f)_*L \\
& \to
K \otimes_\mathcal{O}^\mathbf{L} L
\end{align*}
dans $D(\mathcal{O})$. Pour le voir, on utilise le fait que la composée
des co-unités
$$
L(f' \circ f)^*R(f' \circ f)_* =
Lf^* L(f')^* Rf'_* Rf_*  \to
Lf^* Rf_* \to \text{id}
$$
est la co-unité de l'adjonction entre $L(f' \circ f)^*$ et
$R(f' \circ f)_*$. Voir Catégories, lemme \ref{categories-lemma-compose-counits}.
\end{proof}

\begin{lemma}
\label{lemma-base-change-cup-product}
Considérons un carré commutatif
$$
\xymatrix{
(\Sh(\mathcal{C}'), \mathcal{O}_{\mathcal{C}'}) \ar[r]_{g'} \ar[d]_{f'} &
(\Sh(\mathcal{C}), \mathcal{O}_{\mathcal{C}}) \ar[d]^f \\
(\Sh(\mathcal{D}'), \mathcal{O}_{\mathcal{D}'}) \ar[r]^g &
(\Sh(\mathcal{D}), \mathcal{O}_{\mathcal{D}}) 
}
$$
de topos annelés. Soient $K, L$ dans $D(\mathcal{O}_{\mathcal{C}})$.
Le cup-produit relatif est compatible avec ce carré
au sens où le diagramme
$$
\xymatrix{
Lg^*(Rf_*K \otimes_{\mathcal{O}_{\mathcal{D}}}^\mathbf{L} Rf_*L)
\ar[r] \ar@{=}[d] &
Lg^*(Rf_*(K \otimes_{\mathcal{O}_{\mathcal{C}}}^\mathbf{L} L)) \ar[d] \\
Lg^*Rf_*K \otimes_{\mathcal{O}_{\mathcal{D}'}}^\mathbf{L} Lg^*Rf_*L \ar[d] &
R(f')_*L(g')^*(K \otimes_{\mathcal{O}_{\mathcal{C}}}^\mathbf{L} L) \ar@{=}[d] \\
R(f')_*L(g')^*K \otimes_{\mathcal{O}_{\mathcal{D}'}}^\mathbf{L} R(f')_*L(g')^*L \ar[r] &
R(f')_*(L(g')^*K \otimes_{\mathcal{O}_{\mathcal{C}'}}^\mathbf{L} L(g')^*L)
}
$$
est commutatif dans $D(\mathcal{O}_{\mathcal{D}'})$.
Les flèches horizontales sont données
par le cup-produit relatif (remarque \ref{remark-cup-product})
et les flèches verticales sont données
par le morphisme de changement de base (remarque \ref{remark-base-change})
et le lemme \ref{lemma-pullback-tensor-product}.
\end{lemma}

\begin{proof}
La démonstration est omise.
\end{proof}






\section{Complexes de Hom}
\label{section-hom-complexes}

\noindent
Soit $(\mathcal{C}, \mathcal{O})$ un site annelé. Soient
$\mathcal{L}^\bullet$ et $\mathcal{M}^\bullet$ deux complexes
de $\mathcal{O}$-modules. Nous construisons un complexe
de $\mathcal{O}$-modules
$\SheafHom^\bullet(\mathcal{L}^\bullet, \mathcal{M}^\bullet)$.
Plus précisément, pour tout $n$, nous posons
$$
\SheafHom^n(\mathcal{L}^\bullet, \mathcal{M}^\bullet) =
\prod\nolimits_{n = p + q}
\SheafHom_\mathcal{O}(\mathcal{L}^{-q}, \mathcal{M}^p)
$$
On peut voir $\SheafHom^n$ comme le faisceau de
$\mathcal{O}$-modules formé de toutes les applications $\mathcal{O}$-linéaires
de $\mathcal{L}^\bullet$ dans $\mathcal{M}^\bullet$
(considérés comme des $\mathcal{O}$-modules gradués) qui sont homogènes
de degré $n$. Dans cette terminologie, nous définissons la différentielle par la formule
$$
\text{d}(f) =
\text{d}_\mathcal{M} \circ f - (-1)^n f \circ \text{d}_\mathcal{L}
$$
pour
$f \in \SheafHom^n_\mathcal{O}(\mathcal{L}^\bullet, \mathcal{M}^\bullet)$.
Nous omettons la vérification que $\text{d}^2 = 0$.
Cette construction est un cas particulier d'
Algèbre différentielle graduée, exemple \ref{dga-example-category-complexes}.
Il résulte immédiatement de la construction que l'on a
\begin{equation}
\label{equation-cohomology-hom-complex}
H^n(\Gamma(U, \SheafHom^\bullet(\mathcal{L}^\bullet, \mathcal{M}^\bullet))) =
\Hom_{K(\mathcal{O}_U)}(\mathcal{L}^\bullet|_U, \mathcal{M}^\bullet[n]|_U)
\end{equation}
pour tout $n \in \mathbf{Z}$ et tout $U \in \Ob(\mathcal{C})$. De même,
on a
\begin{equation}
\label{equation-global-cohomology-hom-complex}
H^n(\Gamma(\mathcal{C},
\SheafHom^\bullet(\mathcal{L}^\bullet, \mathcal{M}^\bullet))) =
\Hom_{K(\mathcal{O})}(\mathcal{L}^\bullet, \mathcal{M}^\bullet[n])
\end{equation}
pour le complexe des sections globales.

\begin{lemma}
\label{lemma-compose}
Soit $(\mathcal{C}, \mathcal{O})$ un site annelé.
Étant donnés des complexes $\mathcal{K}^\bullet, \mathcal{L}^\bullet, \mathcal{M}^\bullet$
de $\mathcal{O}$-modules, il existe un isomorphisme
$$
\SheafHom^\bullet(\mathcal{K}^\bullet,
\SheafHom^\bullet(\mathcal{L}^\bullet, \mathcal{M}^\bullet))
=
\SheafHom^\bullet(\text{Tot}(\mathcal{K}^\bullet \otimes_\mathcal{O}
\mathcal{L}^\bullet), \mathcal{M}^\bullet)
$$
de complexes de $\mathcal{O}$-modules, fonctoriel en
$\mathcal{K}^\bullet, \mathcal{L}^\bullet, \mathcal{M}^\bullet$.
\end{lemma}

\begin{proof}
Démonstration omise. Indication : cela se démontre exactement comme dans
Compléments d'algèbre, lemme \ref{more-algebra-lemma-compose}.
\end{proof}

\begin{lemma}
\label{lemma-composition}
Soit $(\mathcal{C}, \mathcal{O})$ un site annelé. Étant donnés des complexes
$\mathcal{K}^\bullet, \mathcal{L}^\bullet, \mathcal{M}^\bullet$
de $\mathcal{O}$-modules, il existe un morphisme canonique
$$
\text{Tot}\left(
\SheafHom^\bullet(\mathcal{L}^\bullet, \mathcal{M}^\bullet)
\otimes_\mathcal{O}
\SheafHom^\bullet(\mathcal{K}^\bullet, \mathcal{L}^\bullet)
\right)
\longrightarrow
\SheafHom^\bullet(\mathcal{K}^\bullet, \mathcal{M}^\bullet)
$$
de complexes de $\mathcal{O}$-modules.
\end{lemma}

\begin{proof}
Démonstration omise. Indication : cela se démontre exactement comme dans
Compléments d'algèbre, lemme \ref{more-algebra-lemma-composition}.
\end{proof}

\begin{lemma}
\label{lemma-diagonal-better}
Soit $(\mathcal{C}, \mathcal{O})$ un site annelé. Étant donnés des complexes
$\mathcal{K}^\bullet, \mathcal{L}^\bullet, \mathcal{M}^\bullet$
de $\mathcal{O}$-modules, il existe un morphisme canonique
$$
\text{Tot}\left(
\mathcal{K}^\bullet \otimes_\mathcal{O}
\SheafHom^\bullet(\mathcal{M}^\bullet, \mathcal{L}^\bullet)
\right)
\longrightarrow
\SheafHom^\bullet(\mathcal{M}^\bullet,
\text{Tot}(\mathcal{K}^\bullet \otimes_\mathcal{O} \mathcal{L}^\bullet))
$$
de complexes de $\mathcal{O}$-modules, fonctoriel par rapport aux trois complexes.
\end{lemma}

\begin{proof}
Démonstration omise. Indication : cela se démontre exactement comme dans
Compléments d'algèbre, lemme \ref{more-algebra-lemma-diagonal-better}.
\end{proof}

\begin{lemma}
\label{lemma-diagonal}
Soit $(\mathcal{C}, \mathcal{O})$ un site annelé. Étant donnés des complexes
$\mathcal{K}^\bullet, \mathcal{L}^\bullet$
de $\mathcal{O}$-modules, il existe un morphisme canonique
$$
\mathcal{K}^\bullet
\longrightarrow
\SheafHom^\bullet(\mathcal{L}^\bullet,
\text{Tot}(\mathcal{K}^\bullet \otimes_\mathcal{O} \mathcal{L}^\bullet))
$$
de complexes de $\mathcal{O}$-modules, fonctoriel par rapport aux deux complexes.
\end{lemma}

\begin{proof}
Démonstration omise. Indication : cela se démontre exactement comme dans
Compléments d'algèbre, lemme \ref{more-algebra-lemma-diagonal}.
\end{proof}

\begin{lemma}
\label{lemma-evaluate-and-more}
Soit $(\mathcal{C}, \mathcal{O})$ un site annelé. Étant donnés des complexes
$\mathcal{K}^\bullet, \mathcal{L}^\bullet, \mathcal{M}^\bullet$
de $\mathcal{O}$-modules, il existe un morphisme canonique
$$
\text{Tot}(\SheafHom^\bullet(\mathcal{L}^\bullet,
\mathcal{M}^\bullet) \otimes_\mathcal{O} \mathcal{K}^\bullet)
\longrightarrow
\SheafHom^\bullet(\SheafHom^\bullet(\mathcal{K}^\bullet,
\mathcal{L}^\bullet), \mathcal{M}^\bullet)
$$
de complexes de $\mathcal{O}$-modules, fonctoriel par rapport aux trois complexes.
\end{lemma}

\begin{proof}
Démonstration omise. Indication : cela se démontre exactement comme dans
Compléments d'algèbre, lemme \ref{more-algebra-lemma-evaluate-and-more}.
\end{proof}

\begin{lemma}
\label{lemma-RHom-into-K-injective}
Soit $(\mathcal{C}, \mathcal{O})$ un site annelé. Soient $L$ et $M$
des objets de $D(\mathcal{O})$. Soit $\mathcal{I}^\bullet$
un complexe K-injectif de $\mathcal{O}$-modules représentant $M$. Soit
$\mathcal{L}^\bullet$ un complexe de $\mathcal{O}$-modules
représentant $L$.
Alors
$$
H^0(\Gamma(U, \SheafHom^\bullet(\mathcal{L}^\bullet, \mathcal{I}^\bullet))) =
\Hom_{D(\mathcal{O}_U)}(L|_U, M|_U)
$$
pour tout $U \in \Ob(\mathcal{C})$. De même,
$H^0(\Gamma(\mathcal{C},
\SheafHom^\bullet(\mathcal{L}^\bullet, \mathcal{I}^\bullet))) =
\Hom_{D(\mathcal{O})}(L, M)$.
\end{lemma}

\begin{proof}
On a
\begin{align*}
H^0(\Gamma(U, \SheafHom^\bullet(\mathcal{L}^\bullet, \mathcal{I}^\bullet)))
& =
\Hom_{K(\mathcal{O}_U)}(\mathcal{L}^\bullet|_U, \mathcal{I}^\bullet|_U) \\
& =
\Hom_{D(\mathcal{O}_U)}(L|_U, M|_U)
\end{align*}
La première égalité est (\ref{equation-cohomology-hom-complex}).
La seconde est vraie parce que $\mathcal{I}^\bullet|_U$
est K-injectif d'après le lemme \ref{lemma-restrict-K-injective-to-open}.
La démonstration de l'égalité portant sur les sections globales est analogue, mais utilise
(\ref{equation-global-cohomology-hom-complex}).
\end{proof}

\begin{lemma}
\label{lemma-RHom-well-defined}
Soit $(\mathcal{C}, \mathcal{O})$ un site annelé. Soit
$(\mathcal{I}')^\bullet \to \mathcal{I}^\bullet$
un quasi-isomorphisme de complexes K-injectifs de $\mathcal{O}$-modules.
Soit $(\mathcal{L}')^\bullet \to \mathcal{L}^\bullet$
un quasi-isomorphisme de complexes de $\mathcal{O}$-modules.
Alors
$$
\SheafHom^\bullet(\mathcal{L}^\bullet, (\mathcal{I}')^\bullet)
\longrightarrow
\SheafHom^\bullet((\mathcal{L}')^\bullet, \mathcal{I}^\bullet)
$$
est un quasi-isomorphisme.
\end{lemma}

\begin{proof}
Soit $M$ l'objet de $D(\mathcal{O})$ représenté par
$\mathcal{I}^\bullet$ et $(\mathcal{I}')^\bullet$.
Soit $L$ l'objet de $D(\mathcal{O})$ représenté par
$\mathcal{L}^\bullet$ et $(\mathcal{L}')^\bullet$.
D'après le lemme \ref{lemma-RHom-into-K-injective},
les faisceaux
$$
H^0(\SheafHom^\bullet(\mathcal{L}^\bullet, (\mathcal{I}')^\bullet))
\quad\text{et}\quad
H^0(\SheafHom^\bullet((\mathcal{L}')^\bullet, \mathcal{I}^\bullet))
$$
sont tous deux égaux au faisceau associé au préfaisceau
$$
U \longmapsto \Hom_{D(\mathcal{O}_U)}(L|_U, M|_U)
$$
Par fonctorialité, le morphisme induit l'identité de ce faisceau.
Le même argument appliqué après décalage montre qu'il induit un
isomorphisme sur tous les faisceaux de cohomologie. C'est donc un quasi-isomorphisme.
\end{proof}

\begin{lemma}
\label{lemma-RHom-from-K-flat-into-K-injective}
Soit $(\mathcal{C}, \mathcal{O})$ un site annelé. Soit $\mathcal{I}^\bullet$
un complexe K-injectif de $\mathcal{O}$-modules. Soit
$\mathcal{L}^\bullet$ un complexe K-plat de $\mathcal{O}$-modules.
Alors $\SheafHom^\bullet(\mathcal{L}^\bullet, \mathcal{I}^\bullet)$
est un complexe K-injectif de $\mathcal{O}$-modules.
\end{lemma}

\begin{proof}
En effet, si $\mathcal{K}^\bullet$ est un complexe acyclique de
$\mathcal{O}$-modules, alors
\begin{align*}
\Hom_{K(\mathcal{O})}(\mathcal{K}^\bullet,
\SheafHom^\bullet(\mathcal{L}^\bullet, \mathcal{I}^\bullet))
& =
H^0(\Gamma(\mathcal{C},
\SheafHom^\bullet(\mathcal{K}^\bullet,
\SheafHom^\bullet(\mathcal{L}^\bullet, \mathcal{I}^\bullet)))) \\
& =
H^0(\Gamma(\mathcal{C}, \SheafHom^\bullet(\text{Tot}(
\mathcal{K}^\bullet \otimes_\mathcal{O} \mathcal{L}^\bullet),
\mathcal{I}^\bullet))) \\
& =
\Hom_{K(\mathcal{O})}(
\text{Tot}(\mathcal{K}^\bullet \otimes_\mathcal{O} \mathcal{L}^\bullet),
\mathcal{I}^\bullet) \\
& =
0
\end{align*}
La première égalité résulte de (\ref{equation-global-cohomology-hom-complex}).
La deuxième résulte du lemme \ref{lemma-compose}.
La troisième résulte de (\ref{equation-global-cohomology-hom-complex}).
Enfin,
$\text{Tot}(\mathcal{K}^\bullet \otimes_\mathcal{O} \mathcal{L}^\bullet)$
est acyclique puisque $\mathcal{L}^\bullet$ est K-plat
(définition \ref{definition-K-flat}) ; la dernière égalité résulte donc du fait que
$\mathcal{I}^\bullet$ est K-injectif.
\end{proof}








\section{Hom interne dans la catégorie dérivée}
\label{section-internal-hom}

\noindent
Soit $(\mathcal{C}, \mathcal{O})$ un site annelé. Soient $L, M$ des objets
de $D(\mathcal{O})$. Nous souhaitons construire un objet
$R\SheafHom(L, M)$ de $D(\mathcal{O})$ tel que, pour tout troisième
objet $K$ de $D(\mathcal{O})$, il existe une bijection canonique
\begin{equation}
\label{equation-internal-hom}
\Hom_{D(\mathcal{O})}(K, R\SheafHom(L, M))
=
\Hom_{D(\mathcal{O})}(K \otimes_\mathcal{O}^\mathbf{L} L, M)
\end{equation}
Observons que cette formule définit $R\SheafHom(L, M)$ à isomorphisme
unique près, par le lemme de Yoneda
(Catégories, lemme \ref{categories-lemma-yoneda}).

\medskip\noindent
Pour construire un tel objet, choisissons un complexe K-injectif de
$\mathcal{O}$-modules $\mathcal{I}^\bullet$ représentant $M$ et un
complexe quelconque de $\mathcal{O}$-modules $\mathcal{L}^\bullet$ représentant $L$.
Posons alors
$$
R\SheafHom(L, M) = \SheafHom^\bullet(\mathcal{L}^\bullet, \mathcal{I}^\bullet)
$$
où le membre de droite est le complexe de $\mathcal{O}$-modules
construit à la section \ref{section-hom-complexes}.
Cette définition est bien fondée par le lemme \ref{lemma-RHom-well-defined}.
On obtient un foncteur
$$
D(\mathcal{O})^{opp} \times D(\mathcal{O}) \longrightarrow D(\mathcal{O}),
\quad
(K, L) \longmapsto R\SheafHom(K, L)
$$
Avant de démontrer (\ref{equation-internal-hom}),
calculons les groupes de cohomologie de $R\SheafHom(K, L)$.

\begin{lemma}
\label{lemma-section-RHom-over-U}
Soit $(\mathcal{C}, \mathcal{O})$ un site annelé. Soient $L, M$ des objets
de $D(\mathcal{O})$. Pour tout objet $U$ de $\mathcal{C}$, on a
$$
H^0(U, R\SheafHom(L, M)) =
\Hom_{D(\mathcal{O}_U)}(L|_U, M|_U)
$$
et l'on a $H^0(\mathcal{C}, R\SheafHom(L, M)) =
\Hom_{D(\mathcal{O})}(L, M)$.
\end{lemma}

\begin{proof}
Choisissons un complexe K-injectif $\mathcal{I}^\bullet$ de
$\mathcal{O}$-modules représentant $M$ et un complexe K-plat
$\mathcal{L}^\bullet$ représentant $L$. Alors
$\SheafHom^\bullet(\mathcal{L}^\bullet, \mathcal{I}^\bullet)$
est K-injectif par le lemme \ref{lemma-RHom-from-K-flat-into-K-injective}.
On peut donc calculer la cohomologie sur $U$ en prenant simplement les sections sur $U$,
et le résultat découle du lemme \ref{lemma-RHom-into-K-injective}.
\end{proof}

\begin{lemma}
\label{lemma-internal-hom}
Soit $(\mathcal{C}, \mathcal{O})$ un site annelé. Soient $K, L, M$ des objets
de $D(\mathcal{O})$. Pour la construction décrite ci-dessus,
il existe un isomorphisme canonique
$$
R\SheafHom(K, R\SheafHom(L, M)) =
R\SheafHom(K \otimes_\mathcal{O}^\mathbf{L} L, M)
$$
dans $D(\mathcal{O})$, fonctoriel en $K, L, M$,
qui redonne (\ref{equation-internal-hom}) après application de $H^0(\mathcal{C}, -)$.
\end{lemma}

\begin{proof}
Choisissons un complexe K-injectif $\mathcal{I}^\bullet$ représentant
$M$ et un complexe K-plat de $\mathcal{O}$-modules $\mathcal{L}^\bullet$
représentant $L$.
Pour tout complexe de $\mathcal{O}$-modules $\mathcal{K}^\bullet$,
on a
$$
\SheafHom^\bullet(\mathcal{K}^\bullet,
\SheafHom^\bullet(\mathcal{L}^\bullet, \mathcal{I}^\bullet))
=
\SheafHom^\bullet(
\text{Tot}(\mathcal{K}^\bullet \otimes_\mathcal{O} \mathcal{L}^\bullet),
\mathcal{I}^\bullet)
$$
par le lemme \ref{lemma-compose}.
Remarquons que le membre de gauche représente
$R\SheafHom(K, R\SheafHom(L, M))$ (utiliser le
lemme \ref{lemma-RHom-from-K-flat-into-K-injective})
et que le membre de droite représente
$R\SheafHom(K \otimes_\mathcal{O}^\mathbf{L} L, M)$.
Cela démontre la formule affichée du lemme.
En prenant les sections globales et en utilisant le lemme \ref{lemma-section-RHom-over-U},
on obtient (\ref{equation-internal-hom}).
\end{proof}

\begin{lemma}
\label{lemma-restriction-RHom-to-U}
Soit $(\mathcal{C}, \mathcal{O})$ un site annelé. Soient $K, L$ des objets
de $D(\mathcal{O})$. La construction de $R\SheafHom(K, L)$
commute aux restrictions, c'est-à-dire que,
pour tout objet $U$ de $\mathcal{C}$, on a
$R\SheafHom(K|_U, L|_U) = R\SheafHom(K, L)|_U$.
\end{lemma}

\begin{proof}
Cela résulte immédiatement de la construction et du
lemme \ref{lemma-restrict-K-injective-to-open}.
\end{proof}

\begin{lemma}
\label{lemma-RHom-triangulated}
Soit $(\mathcal{C}, \mathcal{O})$ un site annelé. Le bifoncteur
$R\SheafHom(- , -)$ transforme les triangles distingués en
triangles distingués dans chacune des deux variables.
\end{lemma}

\begin{proof}
Cela découle du fait que l'application
$$
(\mathcal{L}^\bullet, \mathcal{M}^\bullet) \longmapsto
\SheafHom^\bullet(\mathcal{L}^\bullet, \mathcal{M}^\bullet)
$$
transforme, dans chacune des variables, toute suite exacte courte de complexes
scindée terme à terme en une suite exacte courte scindée terme à terme. Détails omis.
\end{proof}

\begin{lemma}
\label{lemma-internal-hom-evaluate}
Soit $(\mathcal{C}, \mathcal{O})$ un site annelé. Soient $K, L, M$ des objets
de $D(\mathcal{O})$. Il existe un morphisme canonique
$$
R\SheafHom(L, M) \otimes_\mathcal{O}^\mathbf{L} K
\longrightarrow
R\SheafHom(R\SheafHom(K, L), M)
$$
dans $D(\mathcal{O})$, fonctoriel en $K, L, M$.
\end{lemma}

\begin{proof}
Choisissons un complexe K-injectif $\mathcal{I}^\bullet$ représentant $M$,
un complexe K-injectif $\mathcal{J}^\bullet$ représentant $L$, et un complexe
K-plat $\mathcal{K}^\bullet$ représentant $K$. Le morphisme est défini à
l'aide du morphisme
$$
\text{Tot}(\SheafHom^\bullet(\mathcal{J}^\bullet,
\mathcal{I}^\bullet) \otimes_\mathcal{O} \mathcal{K}^\bullet)
\longrightarrow
\SheafHom^\bullet(\SheafHom^\bullet(\mathcal{K}^\bullet,
\mathcal{J}^\bullet), \mathcal{I}^\bullet)
$$
du lemme \ref{lemma-evaluate-and-more}. Avec notre choix particulier de
complexes, le membre de gauche représente
$R\SheafHom(L, M) \otimes_\mathcal{O}^\mathbf{L} K$
et le membre de droite représente
$R\SheafHom(R\SheafHom(K, L), M)$. Nous omettons la démonstration de la
fonctorialité par rapport aux trois objets de $D(\mathcal{O})$.
\end{proof}

\begin{lemma}
\label{lemma-internal-hom-composition}
\begin{slogan}
Composition sur RSheafHom.
\end{slogan}
Soit $(\mathcal{C}, \mathcal{O})$ un site annelé. Pour $K, L, M$ dans
$D(\mathcal{O})$, il existe un morphisme canonique
$$
R\SheafHom(L, M) \otimes_\mathcal{O}^\mathbf{L} R\SheafHom(K, L)
\longrightarrow R\SheafHom(K, M)
$$
dans $D(\mathcal{O})$.
\end{lemma}

\begin{proof}
Choisissons un complexe K-injectif $\mathcal{I}^\bullet$ représentant $M$,
un complexe K-injectif $\mathcal{J}^\bullet$ représentant $L$, et un complexe
quelconque de $\mathcal{O}$-modules $\mathcal{K}^\bullet$ représentant $K$.
Par le lemme \ref{lemma-composition}, il existe un morphisme de complexes
$$
\text{Tot}\left(
\SheafHom^\bullet(\mathcal{J}^\bullet, \mathcal{I}^\bullet)
\otimes_\mathcal{O}
\SheafHom^\bullet(\mathcal{K}^\bullet, \mathcal{J}^\bullet)
\right)
\longrightarrow
\SheafHom^\bullet(\mathcal{K}^\bullet, \mathcal{I}^\bullet)
$$
Les complexes de $\mathcal{O}$-modules
$\SheafHom^\bullet(\mathcal{J}^\bullet, \mathcal{I}^\bullet)$,
$\SheafHom^\bullet(\mathcal{K}^\bullet, \mathcal{J}^\bullet)$ et
$\SheafHom^\bullet(\mathcal{K}^\bullet, \mathcal{I}^\bullet)$
représentent $R\SheafHom(L, M)$, $R\SheafHom(K, L)$ et $R\SheafHom(K, M)$.
Si nous choisissons un complexe K-plat $\mathcal{H}^\bullet$ et un
quasi-isomorphisme
$\mathcal{H}^\bullet \to
\SheafHom^\bullet(\mathcal{K}^\bullet, \mathcal{J}^\bullet)$,
alors il existe un morphisme
$$
\text{Tot}\left(
\SheafHom^\bullet(\mathcal{J}^\bullet, \mathcal{I}^\bullet)
\otimes_\mathcal{O} \mathcal{H}^\bullet
\right)
\longrightarrow
\text{Tot}\left(
\SheafHom^\bullet(\mathcal{J}^\bullet, \mathcal{I}^\bullet)
\otimes_\mathcal{O}
\SheafHom^\bullet(\mathcal{K}^\bullet, \mathcal{J}^\bullet)
\right)
$$
dont la source représente
$R\SheafHom(L, M) \otimes_\mathcal{O}^\mathbf{L} R\SheafHom(K, L)$.
En composant les deux flèches affichées, on obtient le morphisme recherché.
Nous omettons la démonstration de la fonctorialité de la construction.
\end{proof}

\begin{lemma}
\label{lemma-internal-hom-diagonal-better}
Soit $(\mathcal{C}, \mathcal{O})$ un site annelé. Pour $K, L, M$ dans
$D(\mathcal{O})$, il existe un morphisme canonique
$$
K \otimes_\mathcal{O}^\mathbf{L} R\SheafHom(M, L)
\longrightarrow
R\SheafHom(M, K \otimes_\mathcal{O}^\mathbf{L} L)
$$
dans $D(\mathcal{O})$, fonctoriel en $K, L, M$.
\end{lemma}

\begin{proof}
Choisissons un complexe K-plat $\mathcal{K}^\bullet$ représentant $K$, un
complexe K-injectif $\mathcal{I}^\bullet$ représentant $L$, et un complexe
quelconque de $\mathcal{O}$-modules $\mathcal{M}^\bullet$ représentant $M$.
Choisissons un quasi-isomorphisme
$\text{Tot}(\mathcal{K}^\bullet \otimes_{\mathcal{O}_X} \mathcal{I}^\bullet)
\to \mathcal{J}^\bullet$
où $\mathcal{J}^\bullet$ est K-injectif. Nous utilisons alors le morphisme
$$
\text{Tot}\left(
\mathcal{K}^\bullet \otimes_\mathcal{O}
\SheafHom^\bullet(\mathcal{M}^\bullet, \mathcal{I}^\bullet)
\right)
\to
\SheafHom^\bullet(\mathcal{M}^\bullet,
\text{Tot}(\mathcal{K}^\bullet \otimes_\mathcal{O} \mathcal{I}^\bullet))
\to
\SheafHom^\bullet(\mathcal{M}^\bullet, \mathcal{J}^\bullet)
$$
où le premier morphisme est celui du lemme \ref{lemma-diagonal-better}.
\end{proof}

\begin{lemma}
\label{lemma-internal-hom-diagonal}
Soit $(\mathcal{C}, \mathcal{O})$ un site annelé. Pour $K, L$ dans
$D(\mathcal{O})$, il existe un morphisme canonique
$$
K \longrightarrow R\SheafHom(L, K \otimes_\mathcal{O}^\mathbf{L} L)
$$
dans $D(\mathcal{O})$, fonctoriel en $K$ et $L$.
\end{lemma}

\begin{proof}
Choisissons un complexe K-plat $\mathcal{K}^\bullet$ représentant $K$ et un
complexe quelconque de $\mathcal{O}$-modules $\mathcal{L}^\bullet$
représentant $L$. Choisissons un complexe K-injectif
$\mathcal{J}^\bullet$ et un quasi-isomorphisme
$\text{Tot}(\mathcal{K}^\bullet \otimes_\mathcal{O} \mathcal{L}^\bullet)
\to \mathcal{J}^\bullet$. Nous utilisons alors
$$
\mathcal{K}^\bullet \to
\SheafHom^\bullet(\mathcal{L}^\bullet,
\text{Tot}(\mathcal{K}^\bullet \otimes_\mathcal{O} \mathcal{L}^\bullet))
\to
\SheafHom^\bullet(\mathcal{L}^\bullet, \mathcal{J}^\bullet)
$$
où le premier morphisme provient du lemme \ref{lemma-diagonal}.
\end{proof}

\begin{lemma}
\label{lemma-dual}
Soit $(\mathcal{C}, \mathcal{O})$ un site annelé et soit $L$ un objet de
$D(\mathcal{O})$. Posons $L^\vee = R\SheafHom(L, \mathcal{O})$.
Pour $M$ dans $D(\mathcal{O})$, il existe un morphisme canonique
\begin{equation}
\label{equation-eval}
M \otimes^\mathbf{L}_\mathcal{O} L^\vee \longrightarrow R\SheafHom(L, M)
\end{equation}
qui induit un morphisme canonique
$$
H^0(\mathcal{C}, M \otimes_\mathcal{O}^\mathbf{L} L^\vee)
\longrightarrow
\Hom_{D(\mathcal{O})}(L, M)
$$
fonctoriel en $M$ dans $D(\mathcal{O})$.
\end{lemma}

\begin{proof}
Le morphisme (\ref{equation-eval}) est un cas particulier du
lemme \ref{lemma-internal-hom-composition}, en utilisant l'identification
$M = R\SheafHom(\mathcal{O}, M)$.
\end{proof}

\begin{remark}
\label{remark-projection-formula-for-internal-hom}
Soit $f : (\Sh(\mathcal{C}), \mathcal{O}_\mathcal{C}) \to
(\Sh(\mathcal{D}), \mathcal{O}_\mathcal{D})$ un morphisme de topos annelés.
Soient $K, L$ des objets de $D(\mathcal{O}_\mathcal{C})$. Nous affirmons
qu'il existe un morphisme canonique
$$
Rf_*R\SheafHom(L, K) \longrightarrow R\SheafHom(Rf_*L, Rf_*K)
$$
En effet, d'après (\ref{equation-internal-hom}), cela revient à donner un
morphisme
$Rf_*R\SheafHom(L, K) \otimes_{\mathcal{O}_\mathcal{D}}^\mathbf{L} Rf_*L
\to Rf_*K$.
Pour cela, nous pouvons utiliser la composition
$$
Rf_*R\SheafHom(L, K) \otimes_{\mathcal{O}_\mathcal{D}}^\mathbf{L} Rf_*L \to
Rf_*(R\SheafHom(L, K) \otimes_{\mathcal{O}_\mathcal{C}}^\mathbf{L} L) \to
Rf_*K
$$
où la première flèche est le cup-produit relatif
(remarque \ref{remark-cup-product}) et la deuxième est $Rf_*$ appliqué au
morphisme canonique
$R\SheafHom(L, K) \otimes_{\mathcal{O}_\mathcal{C}}^\mathbf{L} L \to K$
provenant du lemme \ref{lemma-internal-hom-composition}
(avec $\mathcal{O}_\mathcal{C}$ à l'une des places).
\end{remark}

\begin{remark}
\label{remark-prepare-fancy-base-change}
Soit $h : (\Sh(\mathcal{C}), \mathcal{O}) \to
(\Sh(\mathcal{C}'), \mathcal{O}')$ un morphisme de topos annelés.
Soient $K, L$ des objets de $D(\mathcal{O}')$. Nous affirmons qu'il existe
un morphisme canonique
$$
Lh^*R\SheafHom(K, L) \longrightarrow R\SheafHom(Lh^*K, Lh^*L)
$$
dans $D(\mathcal{O})$. En effet, d'après (\ref{equation-internal-hom}),
démontré dans le lemme \ref{lemma-internal-hom}, donner un tel morphisme
revient à donner un morphisme
$$
Lh^*R\SheafHom(K, L) \otimes^\mathbf{L} Lh^*K \longrightarrow Lh^*L
$$
La source de cette flèche est
$Lh^*(\SheafHom(K, L) \otimes^\mathbf{L} K)$ d'après le
lemme \ref{lemma-pullback-tensor-product}. Il suffit donc de construire un
morphisme canonique
$$
R\SheafHom(K, L) \otimes^\mathbf{L} K \longrightarrow L.
$$
Pour cela, nous prenons la flèche correspondant à
$$
\text{id} :
R\SheafHom(K, L)
\longrightarrow
R\SheafHom(K, L)
$$
par (\ref{equation-internal-hom}).
\end{remark}

\begin{remark}
\label{remark-fancy-base-change}
Supposons que
$$
\xymatrix{
(\Sh(\mathcal{C}'), \mathcal{O}_{\mathcal{C}'})
\ar[r]_h \ar[d]_{f'} &
(\Sh(\mathcal{C}), \mathcal{O}_\mathcal{C}) \ar[d]^f \\
(\Sh(\mathcal{D}'), \mathcal{O}_{\mathcal{D}'})
\ar[r]^g &
(\Sh(\mathcal{D}), \mathcal{O}_\mathcal{D})
}
$$
soit un diagramme commutatif de topos annelés. Soient $K, L$ des objets de
$D(\mathcal{O}_\mathcal{C})$. Nous affirmons qu'il existe un morphisme
canonique de changement de base
$$
Lg^*Rf_*R\SheafHom(K, L)
\longrightarrow
R(f')_*R\SheafHom(Lh^*K, Lh^*L)
$$
dans $D(\mathcal{O}_{\mathcal{D}'})$. En effet, nous prenons le morphisme
adjoint à la composition
\begin{align*}
L(f')^*Lg^*Rf_*R\SheafHom(K, L)
& =
Lh^*Lf^*Rf_*R\SheafHom(K, L) \\
& \to
Lh^*R\SheafHom(K, L) \\
& \to
R\SheafHom(Lh^*K, Lh^*L)
\end{align*}
où la première flèche utilise le morphisme d'adjonction
$Lf^*Rf_* \to \text{id}$, et la deuxième est le morphisme canonique construit
dans la remarque \ref{remark-prepare-fancy-base-change}.
\end{remark}




\section{Hom dérivé global}
\label{section-global-RHom}

\noindent
Soit $(\Sh(\mathcal{C}), \mathcal{O})$ un topos annelé.
Soient $K, L \in D(\mathcal{O})$.
En utilisant la construction du Hom interne dans la catégorie dérivée, nous
obtenons un objet bien défini
$$
R\Hom_\mathcal{O}(K, L) = R\Gamma(\mathcal{C}, R\SheafHom(K, L))
$$
de $D(\Gamma(\mathcal{C}, \mathcal{O}))$. D'après le
lemme \ref{lemma-section-RHom-over-U}, nous avons
$$
H^0(R\Hom_\mathcal{O}(K, L)) = \Hom_{D(\mathcal{O})}(K, L)
$$
et
$$
H^p(R\Hom_\mathcal{O}(K, L)) = \Ext_{D(\mathcal{O})}^p(K, L)
$$
Si $f : (\mathcal{C}', \mathcal{O}') \to (\mathcal{C}, \mathcal{O})$ est un
morphisme de topos annelés, alors il existe un morphisme canonique
$$
R\Hom_\mathcal{O}(K, L) \longrightarrow R\Hom_{\mathcal{O}'}(Lf^*K, Lf^*L)
$$
dans $D(\Gamma(\mathcal{O}))$, obtenu en prenant les sections globales du
morphisme défini dans la remarque \ref{remark-prepare-fancy-base-change}.









\section{Image directe exceptionnelle dérivée}
\label{section-derived-lower-shriek}

\noindent
Dans cette section, nous étudions des morphismes $g$ de topos annelés pour
lesquels, outre $Lg^*$ et $Rg_*$, il existe aussi un foncteur dérivé $Lg_!$.

\begin{lemma}
\label{lemma-pullback-injective-pre-limp}
Soit $u : \mathcal{C} \to \mathcal{D}$ un foncteur continu et cocontinu de
sites, et soit $g : \Sh(\mathcal{C}) \to \Sh(\mathcal{D})$ le morphisme de
topos correspondant. Soit $\mathcal{O}_\mathcal{D}$ un faisceau d'anneaux et
soit $\mathcal{I}$ un $\mathcal{O}_\mathcal{D}$-module injectif. Alors
$H^p(U, g^{-1}\mathcal{I}) = 0$ pour tout $p > 0$ et tout
$U \in \Ob(\mathcal{C})$.
\end{lemma}

\begin{proof}
L'annulation du lemme résulte du lemme
\ref{lemma-cech-vanish-collection} si nous pouvons démontrer l'annulation de
tous les groupes de cohomologie de {\v C}ech supérieurs
$\check H^p(\mathcal{U}, g^{-1}\mathcal{I})$
pour tout recouvrement $\mathcal{U} = \{U_i \to U\}$ de $\mathcal{C}$.
Comme $u$ est continu, $u(\mathcal{U}) = \{u(U_i) \to u(U)\}$ est un
recouvrement de $\mathcal{D}$, et
$u(U_{i_0} \times_U \ldots \times_U U_{i_n}) =
u(U_{i_0}) \times_{u(U)} \ldots \times_{u(U)} u(U_{i_n})$.
Ainsi, nous avons
$$
\check H^p(\mathcal{U}, g^{-1}\mathcal{I}) =
\check H^p(u(\mathcal{U}), \mathcal{I})
$$
car $g^{-1} = u^p$ d'après Sites, lemme \ref{sites-lemma-when-shriek}.
Comme $\mathcal{I}$ est un $\mathcal{O}_\mathcal{D}$-module injectif, ces
groupes de cohomologie de {\v C}ech sont nuls, voir le
lemme \ref{lemma-injective-module-trivial-cech}.
\end{proof}

\begin{lemma}
\label{lemma-existence-derived-lower-shriek}
Soit $u : \mathcal{C} \to \mathcal{D}$ un foncteur continu et cocontinu de
sites, et soit $g : \Sh(\mathcal{C}) \to \Sh(\mathcal{D})$ le morphisme de
topos correspondant. Soit $\mathcal{O}_\mathcal{D}$ un faisceau d'anneaux et
posons
$\mathcal{O}_\mathcal{C} = g^{-1}\mathcal{O}_\mathcal{D}$.
Le foncteur $g_! : \textit{Mod}(\mathcal{O}_\mathcal{C}) \to
\textit{Mod}(\mathcal{O}_\mathcal{D})$
(voir Modules sur les sites, lemme
\ref{sites-modules-lemma-lower-shriek-modules}) possède un foncteur dérivé à
gauche
$$
Lg_! : D(\mathcal{O}_\mathcal{C}) \longrightarrow D(\mathcal{O}_\mathcal{D})
$$
qui est adjoint à gauche de $g^*$. De plus, pour
$U \in \Ob(\mathcal{C})$, nous avons
$$
Lg_!(j_{U!}\mathcal{O}_U) =
g_!j_{U!}\mathcal{O}_U =
j_{u(U)!} \mathcal{O}_{u(U)}.
$$
où $j_{U!}$ et $j_{u(U)!}$ sont les prolongements par zéro associés aux
morphismes de localisation $j_U : \mathcal{C}/U \to \mathcal{C}$ et
$j_{u(U)} : \mathcal{D}/u(U) \to \mathcal{D}$.
\end{lemma}

\begin{proof}
Nous allons utiliser Catégories dérivées, proposition
\ref{derived-proposition-left-derived-exists}, pour construire $Lg_!$.
Il faut pour cela vérifier les hypothèses (1), (2), (3), (4) et (5) de cette
proposition. Tout d'abord, puisque $g_!$ est un adjoint à gauche, il est exact
à droite et commute à toutes les colimites ; la condition (5) est donc
vérifiée. Les conditions (3) et (4) le sont parce que la catégorie des modules
sur un site annelé est une catégorie abélienne de Grothendieck.
Soit $\mathcal{P} \subset \Ob(\textit{Mod}(\mathcal{O}_\mathcal{C}))$ la
collection des $\mathcal{O}_\mathcal{C}$-modules qui sont des sommes directes
de modules de la forme $j_{U!}\mathcal{O}_U$. Notons que
$g_!j_{U!}\mathcal{O}_U = j_{u(U)!} \mathcal{O}_{u(U)}$, voir la
démonstration de Modules sur les sites, lemme
\ref{sites-modules-lemma-lower-shriek-modules}. Tout
$\mathcal{O}_\mathcal{C}$-module est quotient d'un objet de $\mathcal{P}$,
voir Modules sur les sites, lemme
\ref{sites-modules-lemma-module-quotient-flat}. La condition (1) est donc
vérifiée. Il reste à démontrer (2).
Soit $\mathcal{K}^\bullet$ un complexe acyclique borné supérieurement de
$\mathcal{O}_\mathcal{C}$-modules, avec $\mathcal{K}^n \in \mathcal{P}$ pour
tout $n$. Nous devons montrer que $g_!\mathcal{K}^\bullet$ est exact. Pour
cela, il suffit de montrer, pour tout $\mathcal{O}_\mathcal{D}$-module
injectif $\mathcal{I}$, que
$$
\Hom_{D(\mathcal{O}_\mathcal{D})}(
g_!\mathcal{K}^\bullet, \mathcal{I}[n]) = 0
$$
pour tout $n \in \mathbf{Z}$. Comme $\mathcal{I}$ est injectif, nous avons
\begin{align*}
\Hom_{D(\mathcal{O}_\mathcal{D})}(
g_!\mathcal{K}^\bullet, \mathcal{I}[n])
& =
\Hom_{K(\mathcal{O}_\mathcal{D})}(
g_!\mathcal{K}^\bullet, \mathcal{I}[n]) \\
& =
H^n(\Hom_{\mathcal{O}_\mathcal{D}}(
g_!\mathcal{K}^\bullet, \mathcal{I})) \\
& =
H^n(\Hom_{\mathcal{O}_\mathcal{C}}(
\mathcal{K}^\bullet, g^{-1}\mathcal{I}))
\end{align*}
la dernière égalité venant de l'adjonction entre $g_!$ et $g^{-1}$.

\medskip\noindent
L'annulation de ce groupe serait claire si $g^{-1}\mathcal{I}$ était un
$\mathcal{O}_\mathcal{C}$-module injectif. Mais $g^{-1}\mathcal{I}$ n'est
pas nécessairement un $\mathcal{O}_\mathcal{C}$-module injectif, puisque
$g_!$ n'est en général pas exact. Nous
savons néanmoins que
$$
\Ext^p_{\mathcal{O}_\mathcal{C}}(
j_{U!}\mathcal{O}_U, g^{-1}\mathcal{I}) =
H^p(U, g^{-1}\mathcal{I}) = 0 \text{ pour }p \geq 1
$$
La première égalité vient ici de
$\Hom_{\mathcal{O}_\mathcal{C}}(j_{U!}\mathcal{O}_U, \mathcal{H}) =
\mathcal{H}(U)$ par passage aux foncteurs dérivés, et l'annulation de
$H^p(U, g^{-1}\mathcal{I})$ pour $p > 0$ et
$U \in \Ob(\mathcal{C})$ résulte du
lemme \ref{lemma-pullback-injective-pre-limp}.
Comme chaque $\mathcal{K}^{-q}$ est une somme directe de modules de la forme
$j_{U!}\mathcal{O}_U$, nous voyons que
$$
\Ext^p_{\mathcal{O}_\mathcal{C}}(\mathcal{K}^{-q}, g^{-1}\mathcal{I}) = 0
\text{ pour }p \geq 1\text{ et tout }q
$$
Utilisons la suite spectrale (voir l'exemple
\ref{example-hom-complex-into-sheaf})
$$
E_1^{p, q} = \Ext^q_{\mathcal{O}_\mathcal{C}}(
\mathcal{K}^{-p}, g^{-1}\mathcal{I})
\Rightarrow
\Ext^{p + q}_{\mathcal{O}_\mathcal{C}}(
\mathcal{K}^\bullet, g^{-1}\mathcal{I}) = 0.
$$
Notons que la suite spectrale aboutit à zéro puisque
$\mathcal{K}^\bullet$ est acyclique (il est donc nul dans la catégorie
dérivée et produit des groupes Ext nuls). Par l'annulation des Ext supérieurs
démontrée ci-dessus, les seuls termes non nuls de la page $E_1$ sont
$E_1^{p, 0} = \Hom_{\mathcal{O}_\mathcal{C}}(
\mathcal{K}^{-p}, g^{-1}\mathcal{I})$.
Nous concluons que le complexe
$\Hom_{\mathcal{O}_\mathcal{C}}(
\mathcal{K}^\bullet, g^{-1}\mathcal{I})$
est acyclique, comme voulu.

\medskip\noindent
Ainsi, le foncteur dérivé à gauche $Lg_!$ existe. Il est adjoint à gauche de
$g^{-1} = g^* = Rg^* = Lg^*$, c'est-à-dire que nous avons
\begin{equation}
\label{equation-to-prove}
\Hom_{D(\mathcal{O}_\mathcal{C})}(K, g^*L) =
\Hom_{D(\mathcal{O}_\mathcal{D})}(Lg_!K, L)
\end{equation}
d'après Catégories dérivées, lemme
\ref{derived-lemma-derived-adjoint-functors}. Cela achève la démonstration.
\end{proof}

\begin{remark}
\label{remark-when-derived-shriek-equal}
Attention ! Soient $u : \mathcal{C} \to \mathcal{D}$, $g$,
$\mathcal{O}_\mathcal{D}$ et $\mathcal{O}_\mathcal{C}$ comme dans le
lemme \ref{lemma-existence-derived-lower-shriek}.
En général, le diagramme
$$
\xymatrix{
D(\mathcal{O}_\mathcal{C}) \ar[r]_{Lg_!} \ar[d]_{forget} &
D(\mathcal{O}_\mathcal{D}) \ar[d]^{forget} \\
D(\mathcal{C}) \ar[r]^{Lg^{Ab}_!} &
D(\mathcal{D})
}
$$
ne commute {\bf pas}, où $Lg_!^{Ab}$ est le foncteur construit dans le
lemme \ref{lemma-existence-derived-lower-shriek}, mais en utilisant le
faisceau constant $\mathbf{Z}$ comme faisceau structural sur
$\mathcal{C}$ et $\mathcal{D}$. En général, on n'a même pas
$g_! = g_!^{Ab}$ (voir Modules sur les sites, remarque
\ref{sites-modules-remark-when-shriek-equal}), mais ce phénomène
{\bf peut se produire même si $g_! = g_!^{Ab}$} ! En effet, la construction
de $Lg_!$ dans la démonstration du
lemme \ref{lemma-existence-derived-lower-shriek} montre que $Lg_!$ coïncide
avec $Lg_!^{\textit{Ab}}$ si et seulement si les morphismes canoniques
$$
Lg^{Ab}_!j_{U!}\mathcal{O}_U \longrightarrow j_{u(U)!}\mathcal{O}_{u(U)}
$$
sont des isomorphismes dans $D(\mathcal{D})$ pour tout objet $U$ de
$\mathcal{C}$. En général, tout ce que l'on peut dire est qu'il existe une
transformation naturelle
$$
Lg_!^{Ab} \circ forget \longrightarrow forget \circ Lg_!
$$
\end{remark}

\begin{lemma}
\label{lemma-pullback-injective-limp}
Soit $u : \mathcal{C} \to \mathcal{D}$ un foncteur continu et cocontinu de
sites, et soit $g : \Sh(\mathcal{C}) \to \Sh(\mathcal{D})$ le morphisme de
topos correspondant. Soit $\mathcal{O}_\mathcal{D}$ un faisceau d'anneaux et
soit $\mathcal{I}$ un $\mathcal{O}_\mathcal{D}$-module injectif. Si
$g_!^{Sh} : \Sh(\mathcal{C}) \to \Sh(\mathcal{D})$
commute aux produits fibrés\footnote{C'est le cas si $\mathcal{C}$ possède
des limites connexes finies et si $u$ y commute, voir Sites, lemme
\ref{sites-lemma-preserve-equalizers}.}, alors $g^{-1}\mathcal{I}$ est
totalement acyclique.
\end{lemma}

\begin{proof}
Nous utiliserons le critère du lemme \ref{lemma-characterize-limp}.
La condition (1) est vérifiée d'après le
lemme \ref{lemma-pullback-injective-pre-limp}.
Soit $K' \to K$ un morphisme surjectif de faisceaux d'ensembles sur
$\mathcal{C}$. Comme $g_!^{Sh}$ est un adjoint à gauche, le morphisme
$g_!^{Sh}K' \to g_!^{Sh}K$ est surjectif. Observons que
\begin{align*}
H^0(K' \times_K \ldots \times_K K', g^{-1}\mathcal{I}) 
& =
H^0(g_!^{Sh}(K' \times_K \ldots \times_K K'), \mathcal{I}) \\
& =
H^0(g_!^{Sh}K' \times_{g_!^{Sh}K} \ldots \times_{g_!^{Sh}K} g_!^{Sh}K',
\mathcal{I})
\end{align*}
d'après notre hypothèse sur $g_!^{Sh}$. Comme $\mathcal{I}$ est un module
injectif, il est totalement acyclique par le
lemme \ref{lemma-direct-image-injective-sheaf} (appliqué à l'identité).
Nous pouvons donc utiliser la réciproque du lemme
\ref{lemma-characterize-limp} pour voir que le complexe
$$
0 \to H^0(K, g^{-1}\mathcal{I}) \to H^0(K', g^{-1}\mathcal{I}) \to
H^0(K' \times_K K', g^{-1}\mathcal{I}) \to \ldots
$$
est exact, comme voulu.
\end{proof}

\begin{lemma}
\label{lemma-pullback-same-cohomology}
Soit $u : \mathcal{C} \to \mathcal{D}$ un foncteur continu et cocontinu de
sites, et soit $g : \Sh(\mathcal{C}) \to \Sh(\mathcal{D})$ le morphisme de
topos correspondant. Soit $U \in \Ob(\mathcal{C})$.
\begin{enumerate}
\item Pour $M$ dans $D(\mathcal{D})$, nous avons
$R\Gamma(U, g^{-1}M) = R\Gamma(u(U), M)$.
\item Si $\mathcal{O}_\mathcal{D}$ est un faisceau d'anneaux et
$\mathcal{O}_\mathcal{C} = g^{-1}\mathcal{O}_\mathcal{D}$, alors, pour $M$
dans $D(\mathcal{O}_\mathcal{D})$, nous avons
$R\Gamma(U, g^*M) = R\Gamma(u(U), M)$.
\end{enumerate}
\end{lemma}

\begin{proof}
Dans le cas borné inférieurement, on obtient (1) et (2) en représentant $K$
par un complexe borné inférieurement d'injectifs, puis en utilisant le
lemme \ref{lemma-pullback-injective-pre-limp} ainsi que le lemme d'acyclicité
de Leray. Dans le cas non borné, remarquons d'abord que (1) est un cas
particulier de (2). Pour (2), nous pouvons utiliser
$$
R\Gamma(U, g^*M) =
R\Hom_{\mathcal{O}_\mathcal{C}}(j_{U!}\mathcal{O}_U, g^*M) =
R\Hom_{\mathcal{O}_\mathcal{D}}(j_{u(U)!}\mathcal{O}_{u(U)}, M) =
R\Gamma(u(U), M)
$$
où l'égalité du milieu résulte du
lemme \ref{lemma-existence-derived-lower-shriek}.
\end{proof}

\begin{lemma}
\label{lemma-special-square-cocontinuous}
Supposons donné un diagramme commutatif
$$
\xymatrix{
(\Sh(\mathcal{C}'), \mathcal{O}_{\mathcal{C}'})
\ar[r]_{(g', (g')^\sharp)} \ar[d]_{(f', (f')^\sharp)} &
(\Sh(\mathcal{C}), \mathcal{O}_\mathcal{C}) \ar[d]^{(f, f^\sharp)} \\
(\Sh(\mathcal{D}'), \mathcal{O}_{\mathcal{D}'}) \ar[r]^{(g, g^\sharp)} &
(\Sh(\mathcal{D}), \mathcal{O}_\mathcal{D})
}
$$
de topos annelés. Supposons que :
\begin{enumerate}
\item $f$, $f'$, $g$ et $g'$ correspondent aux foncteurs cocontinus
$u$, $u'$, $v$ et $v'$ comme dans Sites, lemme
\ref{sites-lemma-cocontinuous-morphism-topoi},
\item $v \circ u' = u \circ v'$,
\item $v$ et $v'$ sont continus et cocontinus,
\item pour tout objet $V'$ de $\mathcal{D}'$, le foncteur
${}^{u'}_{V'}\mathcal{I} \to {}^{\ \ \ u}_{v(V')}\mathcal{I}$
défini par $v$ est cofinal,
\item $g^{-1}\mathcal{O}_{\mathcal{D}} = \mathcal{O}_{\mathcal{D}'}$
et $(g')^{-1}\mathcal{O}_{\mathcal{C}} = \mathcal{O}_{\mathcal{C}'}$, et
\item $g'_! : \textit{Ab}(\mathcal{C}') \to \textit{Ab}(\mathcal{C})$
est exact\footnote{C'est le cas si les produits fibrés et les égaliseurs
existent dans $\mathcal{C}'$ et si $v'$ y commute, voir Modules sur les
sites, lemme \ref{sites-modules-lemma-exactness-lower-shriek}.}.
\end{enumerate}
Alors $Rf'_* \circ (g')^* = g^* \circ Rf_*$ comme foncteurs
$D(\mathcal{O}_\mathcal{C}) \to D(\mathcal{O}_{\mathcal{D}'})$.
\end{lemma}

\begin{proof}
Nous avons $g^* = Lg^* = g^{-1}$ et
$(g')^* = L(g')^* = (g')^{-1}$ d'après la condition (5).
Par le lemme \ref{lemma-modules-abelian-unbounded}, il suffit de démontrer le
résultat dans la catégorie dérivée $D(\mathcal{C})$ des faisceaux abéliens.
Choisissons un objet $K \in D(\mathcal{C})$ et soit
$\mathcal{I}^\bullet$ un complexe K-injectif de faisceaux abéliens sur
$\mathcal{C}$ représentant $K$. D'après Catégories dérivées, lemme
\ref{derived-lemma-adjoint-preserve-K-injectives}, et l'hypothèse (6),
$(g')^{-1}\mathcal{I}^\bullet$ est un complexe K-injectif de faisceaux
abéliens sur $\mathcal{C}'$. Par Modules sur les sites, lemme
\ref{sites-modules-lemma-special-square-cocontinuous}
nous obtenons
$f'_*(g')^{-1}\mathcal{I}^\bullet = g^{-1}f_*\mathcal{I}^\bullet$.
Comme $f_*\mathcal{I}^\bullet$ représente $Rf_*K$ et que
$f'_*(g')^{-1}\mathcal{I}^\bullet$ représente $Rf'_*(g')^{-1}K$, nous
concluons.
\end{proof}

\begin{lemma}
\label{lemma-special-square-continuous}
Considérons un diagramme commutatif
$$
\xymatrix{
(\Sh(\mathcal{C}'), \mathcal{O}_{\mathcal{C}'})
\ar[r]_{(g', (g')^\sharp)} \ar[d]_{(f', (f')^\sharp)} &
(\Sh(\mathcal{C}), \mathcal{O}_\mathcal{C}) \ar[d]^{(f, f^\sharp)} \\
(\Sh(\mathcal{D}'), \mathcal{O}_{\mathcal{D}'}) \ar[r]^{(g, g^\sharp)} &
(\Sh(\mathcal{D}), \mathcal{O}_\mathcal{D})
}
$$
de topos annelés et supposons donnés des foncteurs
$$
\xymatrix{
\mathcal{C}' \ar[r]_{v'} &
\mathcal{C} \\
\mathcal{D}' \ar[r]^v \ar[u]^{u'} &
\mathcal{D} \ar[u]_u
}
$$
tels que (avec les notations de Sites, sections
\ref{sites-section-morphism-sites} et
\ref{sites-section-cocontinuous-morphism-topoi}) :
\begin{enumerate}
\item $u$ et $u'$ sont continus et donnent naissance aux morphismes
$f$ et $f'$,
\item $v$ et $v'$ sont cocontinus et donnent naissance aux morphismes
$g$ et $g'$,
\item $u \circ v = v' \circ u'$,
\item $v$ et $v'$ sont continus et cocontinus, et
\item $g^{-1}\mathcal{O}_{\mathcal{D}} = \mathcal{O}_{\mathcal{D}'}$
et $(g')^{-1}\mathcal{O}_{\mathcal{C}} = \mathcal{O}_{\mathcal{C}'}$.
\end{enumerate}
Alors $Rf'_* \circ (g')^* = g^* \circ Rf_*$ comme foncteurs
$D^+(\mathcal{O}_\mathcal{C}) \to D^+(\mathcal{O}_{\mathcal{D}'})$.
Si, de plus,
\begin{enumerate}
\item[(6)] $g'_! : \textit{Ab}(\mathcal{C}') \to \textit{Ab}(\mathcal{C})$
est exact\footnote{C'est le cas si les produits fibrés et les égaliseurs
existent dans $\mathcal{C}'$ et si $v'$ y commute, voir Modules sur les
sites, lemme \ref{sites-modules-lemma-exactness-lower-shriek}.},
\end{enumerate}
alors $Rf'_* \circ (g')^* = g^* \circ Rf_*$ comme foncteurs
$D(\mathcal{O}_\mathcal{C}) \to D(\mathcal{O}_{\mathcal{D}'})$.
\end{lemma}

\begin{proof}
Nous avons $g^* = Lg^* = g^{-1}$ et
$(g')^* = L(g')^* = (g')^{-1}$ d'après la condition (5).
Par le lemme \ref{lemma-modules-abelian-unbounded}, il suffit de démontrer le
résultat dans la catégorie dérivée $D^+(\mathcal{C})$ ou $D(\mathcal{C})$
des faisceaux abéliens.

\medskip\noindent
Choisissons un objet $K \in D^+(\mathcal{C})$. Soit
$\mathcal{I}^\bullet$ un complexe borné inférieurement de faisceaux abéliens
injectifs sur $\mathcal{C}$ représentant $K$. D'après le
lemme \ref{lemma-pullback-injective-pre-limp}, nous avons
$H^p(U', (g')^{-1}\mathcal{I}^q) = 0$ pour tout $p > 0$, tout $q$ et tout
$U' \in \Ob(\mathcal{C}')$. Rappelons que
$R^pf'_*(g')^{-1}\mathcal{I}^q$ est le faisceau associé au préfaisceau
$V' \mapsto H^p(u'(V'), (g')^{-1}\mathcal{I}^q)$, voir le
lemme \ref{lemma-higher-direct-images}. Ainsi,
$(g')^{-1}\mathcal{I}^q$ est acyclique à droite pour le foncteur $f'_*$.
Le lemme d'acyclicité de Leray (Catégories dérivées, lemme
\ref{derived-lemma-leray-acyclicity}) montre alors que
$f'_*(g')^*\mathcal{I}^\bullet$ représente $Rf'_*(g')^{-1}K$.
Par Modules sur les sites, lemme
\ref{sites-modules-lemma-special-square-continuous}
nous obtenons
$f'_*(g')^{-1}\mathcal{I}^\bullet = g^{-1}f_*\mathcal{I}^\bullet$.
Comme $g^{-1}f_*\mathcal{I}^\bullet$ représente $g^{-1}Rf_*K$, nous
concluons.

\medskip\noindent
Choisissons un objet $K \in D(\mathcal{C})$. Soit $\mathcal{I}^\bullet$ un
complexe K-injectif de faisceaux abéliens sur $\mathcal{C}$ représentant
$K$. D'après Catégories dérivées, lemme
\ref{derived-lemma-adjoint-preserve-K-injectives}, et l'hypothèse (6),
$(g')^{-1}\mathcal{I}^\bullet$ est un complexe K-injectif de faisceaux
abéliens sur $\mathcal{C}'$. Par Modules sur les sites, lemme
\ref{sites-modules-lemma-special-square-continuous}
nous obtenons
$f'_*(g')^{-1}\mathcal{I}^\bullet = g^{-1}f_*\mathcal{I}^\bullet$.
Comme $f_*\mathcal{I}^\bullet$ représente $Rf_*K$ et que
$f'_*(g')^{-1}\mathcal{I}^\bullet$ représente $Rf'_*(g')^{-1}K$, nous
concluons.
\end{proof}








\section{Image directe exceptionnelle dérivée pour les catégories fibrées}
\label{section-derived-lower-shriek-fibred}

\noindent
Dans cette section, nous étudions certains cas particuliers de la situation
examinée dans la section \ref{section-derived-lower-shriek}.
Nous nous assurons que les foncteurs image directe exceptionnelle sur les modules
et sur les faisceaux de groupes abéliens coïncident. Nous conseillons au lecteur de sauter
cette section en première lecture.

\begin{situation}
\label{situation-fibred-category}
Soient $(\mathcal{D}, \mathcal{O}_\mathcal{D})$ un site annelé
et $p : \mathcal{C} \to \mathcal{D}$ une catégorie fibrée. Nous munissons
$\mathcal{C}$ de la topologie héritée de $\mathcal{D}$
(Champs, section \ref{stacks-section-topology}). Nous notons
$\pi : \Sh(\mathcal{C}) \to \Sh(\mathcal{D})$ le morphisme de
topos associé à $p$
(Champs, lemme \ref{stacks-lemma-topology-inherited-functorial}).
Nous posons $\mathcal{O}_\mathcal{C} = \pi^{-1}\mathcal{O}_\mathcal{D}$,
de sorte que nous obtenons un morphisme de topos annelés
$$
\pi :
(\Sh(\mathcal{C}), \mathcal{O}_\mathcal{C})
\longrightarrow
(\Sh(\mathcal{D}), \mathcal{O}_\mathcal{D})
$$
\end{situation}

\begin{lemma}
\label{lemma-fibred-category-with-object}
On reprend les hypothèses et les notations de la situation \ref{situation-fibred-category}.
Pour $U \in \Ob(\mathcal{C})$, considérons le morphisme induit
de topos
$$
\pi_U : \Sh(\mathcal{C}/U) \longrightarrow \Sh(\mathcal{D}/p(U))
$$
Alors il existe un morphisme de topos
$$
\sigma : \Sh(\mathcal{D}/p(U)) \to \Sh(\mathcal{C}/U)
$$
tel que $\pi_U \circ \sigma = \text{id}$ et $\sigma^{-1} = \pi_{U, *}$.
\end{lemma}

\begin{proof}
Remarquons que $\pi_U$ est la restriction de $\pi$ aux localisations, voir
Sites, lemme \ref{sites-lemma-localize-cocontinuous}.
Pour un objet $V \to p(U)$ de $\mathcal{D}/p(U)$, notons
$V \times_{p(U)} U \to U$ le morphisme fortement cartésien de $\mathcal{C}$
au-dessus de $\mathcal{D}$, qui existe puisque $p$ est une catégorie fibrée.
Le foncteur
$$
v : \mathcal{D}/p(U) \to \mathcal{C}/U,\quad
V/p(U) \mapsto V \times_{p(U)} U/U
$$
est continu par définition de la topologie sur $\mathcal{C}$.
De plus, il est adjoint à droite de $p$ par définition des morphismes fortement
cartésiens. Nous sommes donc dans la situation examinée dans
Sites, section \ref{sites-section-cocontinuous-adjoint},
et nous voyons que le faisceau $\pi_{U, *}\mathcal{F}$
est égal à $V \mapsto \mathcal{F}(V \times_{p(U)} U)$
(voir en particulier Sites, lemme
\ref{sites-lemma-have-functor-other-way-morphism}).

\medskip\noindent
Mais nous disposons ici d'un résultat plus fort. En effet, le foncteur $v$
est aussi cocontinu (car tous les morphismes des recouvrements de $\mathcal{C}$ 
sont fortement cartésiens). Ainsi, $v$ définit un morphisme $\sigma$ comme
indiqué dans le lemme. L'égalité $\sigma^{-1} = \pi_{U, *}$
résulte immédiatement de la définition. Puisque $\pi_U^{-1}\mathcal{G}$
est donné par la règle $U'/U \mapsto \mathcal{G}(p(U')/p(U))$,
il s'ensuit que $\sigma^{-1} \circ \pi_U^{-1} = \text{id}$,
ce qui prouve l'égalité
$\pi_U \circ \sigma = \text{id}$.
\end{proof}

\begin{situation}
\label{situation-morphism-fibred-categories}
Soit $(\mathcal{D}, \mathcal{O}_\mathcal{D})$ un site annelé.
Soit $u : \mathcal{C}' \to \mathcal{C}$ un $1$-morphisme de catégories
fibrées au-dessus de $\mathcal{D}$
(Catégories, définition \ref{categories-definition-fibred-categories-over-C}).
Munissons $\mathcal{C}$ et $\mathcal{C}'$ de leurs topologies héritées
(Champs, définition \ref{stacks-definition-topology-inherited})
et soient
$\pi : \Sh(\mathcal{C}) \to \Sh(\mathcal{D})$,
$\pi' : \Sh(\mathcal{C}') \to \Sh(\mathcal{D})$, et
$g : \Sh(\mathcal{C}') \to \Sh(\mathcal{C})$
les morphismes de topos correspondants
(Champs, lemme \ref{stacks-lemma-topology-inherited-functorial}).
Posons $\mathcal{O}_\mathcal{C} = \pi^{-1}\mathcal{O}_\mathcal{D}$
et $\mathcal{O}_{\mathcal{C}'} = (\pi')^{-1}\mathcal{O}_\mathcal{D}$.
Remarquons que $g^{-1}\mathcal{O}_\mathcal{C} = \mathcal{O}_{\mathcal{C}'}$,
de sorte que
$$
\xymatrix{
(\Sh(\mathcal{C}'), \mathcal{O}_{\mathcal{C}'}) \ar[rd]_{\pi'} \ar[rr]_g & &
(\Sh(\mathcal{C}), \mathcal{O}_\mathcal{C}) \ar[ld]^\pi \\
& (\Sh(\mathcal{D}), \mathcal{O}_\mathcal{D})
}
$$
est un diagramme commutatif de morphismes de topos annelés.
\end{situation}

\begin{lemma}
\label{lemma-morphism-fibred-categories-with-object}
On reprend les hypothèses et les notations de la
situation \ref{situation-morphism-fibred-categories}.
Pour $U' \in \Ob(\mathcal{C}')$, posons $U = u(U')$ et $V = p'(U')$, puis
considérons les morphismes induits de topos annelés
$$
\xymatrix{
(\Sh(\mathcal{C}'/U'), \mathcal{O}_{U'}) \ar[rd]_{\pi'_{U'}} \ar[rr]_{g'} & &
(\Sh(\mathcal{C}/U), \mathcal{O}_U) \ar[ld]^{\pi_U} \\
& (\Sh(\mathcal{D}/V), \mathcal{O}_V)
}
$$
Alors il existe un morphisme de topos
$$
\sigma' : \Sh(\mathcal{D}/V) \to \Sh(\mathcal{C}'/U'),
$$
tel que, en posant $\sigma = g' \circ \sigma'$, nous ayons
$\pi'_{U'} \circ \sigma' = \text{id}$, $\pi_U \circ \sigma = \text{id}$,
$(\sigma')^{-1} = \pi'_{U', *}$, et $\sigma^{-1} = \pi_{U, *}$.
\end{lemma}

\begin{proof}
Soit $v' : \mathcal{D}/V \to \mathcal{C}'/U'$ le foncteur construit
dans la démonstration du lemme \ref{lemma-fibred-category-with-object} à partir
de $p' : \mathcal{C}' \to \mathcal{D}$ et de l'objet $U'$.
Comme $u$ est un $1$-morphisme de catégories fibrées au-dessus de $\mathcal{D}$,
il transforme les morphismes fortement cartésiens en morphismes fortement cartésiens ;
ainsi, le foncteur $v = u \circ v'$ est le foncteur de
la démonstration du lemme \ref{lemma-fibred-category-with-object}
relatif à $p : \mathcal{C} \to \mathcal{D}$ et à $U$. Le présent lemme
résulte donc de ce lemme.
\end{proof}

\begin{lemma}
\label{lemma-properties-lower-shriek-fibred-category}
On reprend les hypothèses et les notations de la
situation \ref{situation-morphism-fibred-categories}.
\begin{enumerate}
\item Il existe des adjoints à gauche
$g_! : \textit{Mod}(\mathcal{O}_{\mathcal{C}'}) \to
\textit{Mod}(\mathcal{O}_\mathcal{C})$ et
$g_!^{\textit{Ab}} : \textit{Ab}(\mathcal{C}') \to \textit{Ab}(\mathcal{C})$
de $g^* = g^{-1}$ sur les modules et sur les faisceaux abéliens.
\item Le diagramme
$$
\xymatrix{
\textit{Mod}(\mathcal{O}_{\mathcal{C}'}) \ar[d] \ar[r]_{g_!} &
\textit{Mod}(\mathcal{O}_\mathcal{C}) \ar[d] \\
\textit{Ab}(\mathcal{C}') \ar[r]^{g_!^{\textit{Ab}}} &
\textit{Ab}(\mathcal{C})
}
$$
est commutatif.
\item Il existe des adjoints à gauche
$Lg_! : D(\mathcal{O}_{\mathcal{C}'}) \to D(\mathcal{O}_\mathcal{C})$
et
$Lg_!^{\textit{Ab}} : D(\mathcal{C}') \to D(\mathcal{C})$
de $g^* = g^{-1}$ sur les catégories dérivées de modules et de faisceaux abéliens.
\item Le diagramme
$$
\xymatrix{
D(\mathcal{O}_{\mathcal{C}'}) \ar[d] \ar[r]_{Lg_!} &
D(\mathcal{O}_\mathcal{C}) \ar[d] \\
D(\mathcal{C}') \ar[r]^{Lg_!^{\textit{Ab}}} &
D(\mathcal{C})
}
$$
est commutatif.
\end{enumerate}
\end{lemma}

\begin{proof}
Le foncteur $u$ est continu et cocontinu
(Champs, lemme \ref{stacks-lemma-topology-inherited-functorial}).
L'existence des foncteurs $g_!$, $g_!^{\textit{Ab}}$,
$Lg_!$ et $Lg_!^{\textit{Ab}}$ résulte donc de
Modules sur les sites, sections
\ref{sites-modules-section-exactness-lower-shriek} et
\ref{sites-modules-section-lower-shriek-modules}
et de la
section \ref{section-derived-lower-shriek}.

\medskip\noindent
Pour démontrer (2), il suffit de montrer que le morphisme canonique
$$
g_!^{\textit{Ab}}j_{U'!}\mathcal{O}_{U'} \to j_{u(U')!}\mathcal{O}_{u(U')}
$$
est un isomorphisme pour tout objet $U'$ de $\mathcal{C}'$, voir
Modules sur les sites, remarque \ref{sites-modules-remark-when-shriek-equal}.
De même, pour démontrer (4), il suffit de montrer que le morphisme canonique
$$
Lg_!^{\textit{Ab}}j_{U'!}\mathcal{O}_{U'} \to j_{u(U')!}\mathcal{O}_{u(U')}
$$
est un isomorphisme dans $D(\mathcal{C})$ pour tout objet $U'$ de
$\mathcal{C}'$, voir la remarque \ref{remark-when-derived-shriek-equal}.
Cela entraînera aussi la formule précédente ; c'est donc ce que nous allons montrer.

\medskip\noindent
Nous utiliserons le fait que, pour un morphisme de localisation $j$, les
foncteurs $j_!$ et $j_!^{\textit{Ab}}$ coïncident (voir
Modules sur les sites, remarque \ref{sites-modules-remark-localize-shriek-equal})
et que $j_!$ est exact
(Modules sur les sites, lemme \ref{sites-modules-lemma-extension-by-zero-exact}).
Adoptons les notations du
lemme \ref{lemma-morphism-fibred-categories-with-object}.
Puisque $Lg_!^{\textit{Ab}} \circ j_{U'!} = j_{U!} \circ L(g')^{\textit{Ab}}_!$
(par la commutativité du diagramme de Sites, lemme \ref{sites-lemma-localize-cocontinuous},
et l'unicité des foncteurs adjoints), il suffit de prouver que
$L(g')^{\textit{Ab}}_!\mathcal{O}_{U'} = \mathcal{O}_U$. D'après les
résultats du
lemme \ref{lemma-morphism-fibred-categories-with-object},
nous avons, pour tout objet $E$ de $D(\mathcal{C}/u(U'))$, la suite
d'égalités suivante
\begin{align*}
\Hom_{D(\mathcal{C}/U)}(L(g')_!^{\textit{Ab}}\mathcal{O}_{U'}, E)
& =
\Hom_{D(\mathcal{C}'/U')}(\mathcal{O}_{U'}, (g')^{-1}E) \\
& =
\Hom_{D(\mathcal{C}'/U')}((\pi'_{U'})^{-1}\mathcal{O}_V, (g')^{-1}E) \\
& =
\Hom_{D(\mathcal{D}/V)}(\mathcal{O}_V, R\pi'_{U', *}(g')^{-1}E) \\
& =
\Hom_{D(\mathcal{D}/V)}(\mathcal{O}_V, (\sigma')^{-1}(g')^{-1}E) \\
& =
\Hom_{D(\mathcal{D}/V)}(\mathcal{O}_V, \sigma^{-1}E) \\
& =
\Hom_{D(\mathcal{D}/V)}(\mathcal{O}_V, \pi_{U, *}E) \\
& =
\Hom_{D(\mathcal{C}/U)}(\pi_U^{-1}\mathcal{O}_V, E) \\
& =
\Hom_{D(\mathcal{C}/U)}(\mathcal{O}_U, E)
\end{align*}
Le lemme de Yoneda permet de conclure.
\end{proof}

\begin{remark}
\label{remark-fibred-category}
On reprend les hypothèses et les notations de la
situation \ref{situation-fibred-category}. Remarquons qu'en posant
$\mathcal{C}' = \mathcal{D}$ et en prenant pour $u$ le foncteur structural
de $\mathcal{C}$, on obtient une situation comme dans la
situation \ref{situation-morphism-fibred-categories}. Le
lemme \ref{lemma-properties-lower-shriek-fibred-category} nous donne donc les
foncteurs $\pi_!$, $\pi_!^{\textit{Ab}}$, $L\pi_!$ et
$L\pi_!^{\textit{Ab}}$ tels que
$forget \circ \pi_! = \pi_!^{\textit{Ab}} \circ forget$ et
$forget \circ L\pi_! = L\pi_!^{\textit{Ab}} \circ forget$.
\end{remark}

\begin{remark}
\label{remark-morphism-fibred-categories}
On reprend les hypothèses et les notations de la
situation \ref{situation-morphism-fibred-categories}.
Soit $\mathcal{F}$ un faisceau abélien sur $\mathcal{C}$, soit
$\mathcal{F}'$ un faisceau abélien sur $\mathcal{C}'$, et soit
$t : \mathcal{F}' \to g^{-1}\mathcal{F}$ un morphisme.
Nous obtenons alors un morphisme canonique
$$
L\pi'_!(\mathcal{F}') \longrightarrow L\pi_!(\mathcal{F})
$$
en utilisant l'adjoint $g_!\mathcal{F}' \to \mathcal{F}$ de $t$, le
morphisme $Lg_!(\mathcal{F}') \to g_!\mathcal{F}'$ et l'égalité
$L\pi'_! = L\pi_! \circ Lg_!$.
\end{remark}

\begin{lemma}
\label{lemma-compute-pi-shriek}
On reprend les hypothèses et les notations de la
situation \ref{situation-fibred-category}.
Pour $\mathcal{F}$ dans $\textit{Ab}(\mathcal{C})$, le faisceau
$\pi_!\mathcal{F}$ est le faisceau associé au préfaisceau
$$
V \longmapsto \colim_{\mathcal{C}_V^{opp}} \mathcal{F}|_{\mathcal{C}_V}
$$
avec les applications de restriction indiquées dans la démonstration.
\end{lemma}

\begin{proof}
Notons $\mathcal{H}$ la règle du lemme. Pour un morphisme
$h : V' \to V$ de $\mathcal{D}$, il existe un foncteur d'image inverse
$h^* : \mathcal{C}_V \to \mathcal{C}_{V'}$ des catégories fibres
(Catégories, définition
\ref{categories-definition-pullback-functor-fibred-category}).
De plus, pour $U \in \Ob(\mathcal{C}_V)$, il existe un morphisme fortement
cartésien $h^*U \to U$ au-dessus de $h$. La restriction le long de ces
morphismes fortement cartésiens définit une transformation de foncteurs
$$
\mathcal{F}|_{\mathcal{C}_V}
\longrightarrow
\mathcal{F}|_{\mathcal{C}_{V'}} \circ h^*.
$$
et donc un morphisme $\mathcal{H}(V) \to \mathcal{H}(V')$ entre les
colimites, voir Catégories, lemme
\ref{categories-lemma-functorial-colimit}.

\medskip\noindent
Pour démontrer le lemme, nous montrons que
$$
\Mor_{\textit{PSh}(\mathcal{D})}(\mathcal{H}, \mathcal{G}) =
\Mor_{\Sh(\mathcal{C})}(\mathcal{F}, \pi^{-1}\mathcal{G})
$$
pour tout faisceau $\mathcal{G}$ sur $\mathcal{C}$. Un élément du membre de
gauche est un système compatible de morphismes
$\mathcal{F}(U) \to \mathcal{G}(p(U))$ pour tout $U$ de $\mathcal{C}$.
Comme $\pi^{-1}\mathcal{G}(U) = \mathcal{G}(p(U))$ par notre choix de la
topologie sur $\mathcal{C}$, il en est de même pour le membre de droite, ce
qui conclut.
\end{proof}





\section{Homologie d'une catégorie}
\label{section-homology}

\noindent
Dans le cas d'une catégorie au-dessus d'un point, nous appellerons les
foncteurs dérivés à gauche d'image directe exceptionnelle les foncteurs
d'homologie.

\begin{example}[Catégorie au-dessus d'un point]
\label{example-category-to-point}
Soit $\mathcal{C}$ une catégorie. Munissons $\mathcal{C}$ de la topologie
chaotique (Sites, exemple \ref{sites-example-indiscrete}). Les préfaisceaux
et les faisceaux coïncident alors sur $\mathcal{C}$.
Le foncteur $p : \mathcal{C} \to *$, où $*$ est la catégorie avec un seul
objet et un seul morphisme, est cocontinu et continu. Soit
$\pi : \Sh(\mathcal{C}) \to \Sh(*)$ le morphisme de topos correspondant.
Soit $B$ un anneau. Nous munissons $*$ du faisceau d'anneaux $B$ et
$\mathcal{C}$ de $\mathcal{O}_\mathcal{C} = \pi^{-1}B$, que nous noterons
$\underline{B}$. De cette manière,
$$
\pi : (\Sh(\mathcal{C}), \underline{B}) \to (\Sh(*), B)
$$
est un exemple de la situation \ref{situation-fibred-category}.
Par la remarque \ref{remark-fibred-category}, il n'est pas nécessaire de
distinguer $\pi_!$ sur les modules et sur les faisceaux abéliens. Le
lemme \ref{lemma-compute-pi-shriek} montre que
$\pi_!\mathcal{F} = \colim_{\mathcal{C}^{opp}} \mathcal{F}$.
Ainsi, $L_n\pi_!$ est le $n$-ième foncteur dérivé à gauche du foncteur de
colimite. Dans la suite, nous écrivons
$$
H_n(\mathcal{C}, \mathcal{F}) = L_n\pi_!(\mathcal{F})
$$
et nous appelons ceci le {\it $n$-ième groupe d'homologie de $\mathcal{F}$}
sur $\mathcal{C}$.
\end{example}

\begin{example}[Calcul de l'homologie]
\label{example-left-derived-colimits}
Dans l'exemple \ref{example-category-to-point}, on peut calculer les
foncteurs $H_n(\mathcal{C}, -)$ comme suit. Soit
$\mathcal{F} \in \Ob(\textit{Ab}(\mathcal{C}))$. Considérons le complexe de
chaînes
$$
K_\bullet(\mathcal{F}) :
\ \ldots \to
\bigoplus\nolimits_{U_2 \to U_1 \to U_0} \mathcal{F}(U_0)
\to
\bigoplus\nolimits_{U_1 \to U_0} \mathcal{F}(U_0)
\to
\bigoplus\nolimits_{U_0} \mathcal{F}(U_0)
$$
où les morphismes de transition sont donnés par
$$
(U_2 \to U_1 \to U_0, s)
\longmapsto
(U_1 \to U_0, s) - (U_2 \to U_0, s) + (U_2 \to U_1, s|_{U_1})
$$
et de même dans les autres degrés. Par construction,
$$
H_0(\mathcal{C}, \mathcal{F}) =
\colim_{\mathcal{C}^{opp}} \mathcal{F} =
H_0(K_\bullet(\mathcal{F})),
$$
voir Catégories, lemme
\ref{categories-lemma-colimits-coproducts-coequalizers}.
La construction de $K_\bullet(\mathcal{F})$ est fonctorielle en
$\mathcal{F}$ et transforme les suites exactes courtes de
$\textit{Ab}(\mathcal{C})$ en suites exactes courtes de complexes. La suite
de foncteurs $\mathcal{F} \mapsto H_n(K_\bullet(\mathcal{F}))$ forme donc un
$\delta$-foncteur, voir Homologie, définition
\ref{homology-definition-cohomological-delta-functor} et lemme
\ref{homology-lemma-long-exact-sequence-cochain}.
Pour $\mathcal{F} = j_{U!}\mathbf{Z}_U$, le complexe
$K_\bullet(\mathcal{F})$ est celui associé au $\mathbf{Z}$-module libre sur
l'ensemble simplicial $X_\bullet$ dont les termes sont
$$
X_n = \coprod\nolimits_{U_n \to \ldots \to U_1 \to U_0}
\Mor_\mathcal{C}(U_0, U)
$$
Cet ensemble simplicial est homotopiquement équivalent à l'ensemble
simplicial constant sur un singleton $\{*\}$. En effet, le morphisme
$X_\bullet \to \{*\}$ est évident, le morphisme $\{*\} \to X_n$ envoie $*$
sur $(U \to \ldots \to U, \text{id}_U)$, et les morphismes
$$
h_{n, i} : X_n \longrightarrow X_n
$$
(Simplicial, lemme \ref{simplicial-lemma-relations-homotopy})
qui définissent l'homotopie entre les deux morphismes
$X_\bullet \to X_\bullet$ sont donnés par la règle
$$
h_{n, i} :
(U_n \to \ldots \to U_0, f)
\longmapsto
(U_n \to \ldots \to U_i \to U \to \ldots \to U, \text{id})
$$
pour $i > 0$, avec $h_{n, 0} = \text{id}$. Vérifications omises.
Il en résulte que $K_\bullet(j_{U!}\mathbf{Z}_U)$ a une homologie nulle en
degrés négatifs (par la fonctorialité de Simplicial, remarque
\ref{simplicial-remark-homotopy-better}, et le résultat de Simplicial, lemme
\ref{simplicial-lemma-homotopy-s-N}).
Ainsi, $K_\bullet(\mathcal{F})$ calcule les foncteurs dérivés à gauche
$H_n(\mathcal{C}, -)$ de $H_0(\mathcal{C}, -)$, par exemple d'après les
énoncés duaux de Homologie, lemme
\ref{homology-lemma-efface-implies-universal}, et Catégories dérivées, lemme
\ref{derived-lemma-right-derived-delta-functor}.
\end{example}

\begin{example}
\label{example-morphism-categories}
Soit $u : \mathcal{C}' \to \mathcal{C}$ un foncteur. Munissons
$\mathcal{C}'$ et $\mathcal{C}$ de la topologie chaotique comme dans
l'exemple \ref{example-category-to-point}. Les foncteurs $u$,
$\mathcal{C}' \to *$ et $\mathcal{C} \to *$, où $*$ est la catégorie ayant
un seul objet et un seul morphisme, sont cocontinus et continus. Soient
$g : \Sh(\mathcal{C}') \to \Sh(\mathcal{C})$,
$\pi' : \Sh(\mathcal{C}') \to \Sh(*)$, et
$\pi : \Sh(\mathcal{C}) \to \Sh(*)$,
les morphismes de topos correspondants. Soit $B$ un anneau. Nous munissons
$*$ du faisceau d'anneaux $B$, et $\mathcal{C}'$ et $\mathcal{C}$ du
faisceau constant $\underline{B}$. De cette manière,
$$
\xymatrix{
(\Sh(\mathcal{C}'), \underline{B}) \ar[rd]_{\pi'} \ar[rr]_g & &
(\Sh(\mathcal{C}), \underline{B}) \ar[ld]^\pi \\
& (\Sh(*), B)
}
$$
est un exemple de la situation
\ref{situation-morphism-fibred-categories}. Ainsi, le
lemme \ref{lemma-properties-lower-shriek-fibred-category} s'applique à $g$,
et il n'est pas nécessaire de distinguer $g_!$ sur les modules et sur les
faisceaux abéliens. En particulier, la
remarque \ref{remark-morphism-fibred-categories} produit des morphismes
canoniques
$$
H_n(\mathcal{C}', \mathcal{F}')
\longrightarrow
H_n(\mathcal{C}, \mathcal{F})
$$
lorsque $\mathcal{F}$ appartient à $\textit{Ab}(\mathcal{C})$,
$\mathcal{F}'$ à $\textit{Ab}(\mathcal{C}')$, et qu'un morphisme
$t : \mathcal{F}' \to g^{-1}\mathcal{F}$ est donné. En termes du calcul de
l'homologie de l'exemple \ref{example-left-derived-colimits}, ces morphismes
proviennent d'un morphisme de complexes
$$
K_\bullet(\mathcal{F}') \longrightarrow K_\bullet(\mathcal{F})
$$
donné par la règle
$$
(U'_n \to \ldots \to U'_0, s') \longmapsto
(u(U'_n) \to \ldots \to u(U'_0), t(s'))
$$
avec les notations évidentes.
\end{example}

\begin{remark}
\label{remark-map-evaluation-to-derived}
On reprend les notations et les hypothèses de
l'exemple \ref{example-category-to-point}. Soit $\mathcal{F}^\bullet$ un
complexe borné de faisceaux abéliens sur $\mathcal{C}$. Pour tout objet $U$
de $\mathcal{C}$, il existe un morphisme canonique
$$
\mathcal{F}^\bullet(U) \longrightarrow L\pi_!(\mathcal{F}^\bullet)
$$
dans $D(\textit{Ab})$. Si $\mathcal{F}^\bullet$ est un complexe de
$\underline{B}$-modules, ce morphisme appartient à $D(B)$. Pour le montrer,
remarquons que nous calculons $L\pi_!(\mathcal{F}^\bullet)$ en prenant un
quasi-isomorphisme $\mathcal{P}^\bullet \to \mathcal{F}^\bullet$, où
$\mathcal{P}^\bullet$ est un complexe de projectifs. Comme la topologie est
chaotique, $\mathcal{P}^\bullet(U) \to \mathcal{F}^\bullet(U)$ est un
quasi-isomorphisme, donc peut être inversé dans $D(\textit{Ab})$, resp.\
$D(B)$. En le composant avec le morphisme canonique
$\mathcal{P}^\bullet(U) \to \pi_!(\mathcal{P}^\bullet)$ provenant du calcul
de $\pi_!$ comme colimite, nous obtenons la flèche recherchée.
\end{remark}

\begin{lemma}
\label{lemma-initial-final}
On reprend les notations et les hypothèses de
l'exemple \ref{example-category-to-point}. Si $\mathcal{C}$ possède un objet
initial ou un objet final, alors
$L\pi_! \circ \pi^{-1} = \text{id}$ sur $D(\textit{Ab})$, resp.\ $D(B)$.
\end{lemma}

\begin{proof}
Si $\mathcal{C}$ possède un objet initial, alors $\pi_!$ se calcule par
évaluation sur cet objet et l'énoncé est clair. Si $\mathcal{C}$ possède un
objet final, alors $R\pi_*$ se calcule par évaluation sur cet objet, donc
$R\pi_* \circ \pi^{-1} \cong \text{id}$ sur $D(\textit{Ab})$, resp.\
$D(B)$. Il en résulte que
$\pi^{-1} : D(\textit{Ab}) \to D(\mathcal{C})$,
resp.\ $\pi^{-1} : D(B) \to D(\underline{B})$, est pleinement fidèle, voir
Catégories, lemme \ref{categories-lemma-adjoint-fully-faithful}.
Le même lemme implique alors que $L\pi_! \circ \pi^{-1} = \text{id}$,
comme voulu.
\end{proof}

\begin{lemma}
\label{lemma-change-of-rings}
On reprend les notations et les hypothèses de
l'exemple \ref{example-category-to-point}. Soit $B \to B'$ un morphisme
d'anneaux. Considérons le diagramme commutatif de topos annelés
$$
\xymatrix{
(\Sh(\mathcal{C}), \underline{B}) \ar[d]_\pi &
(\Sh(\mathcal{C}), \underline{B'}) \ar[d]^{\pi'} \ar[l]^h \\
(*, B) & (*, B') \ar[l]_f
}
$$
Alors $L\pi'_! \circ Lh^* = Lf^* \circ L\pi_!$.
\end{lemma}

\begin{proof}
Les deux foncteurs sont adjoints à gauche du foncteur évident
$D(B') \to D(\underline{B})$.
\end{proof}

\begin{lemma}
\label{lemma-compute-by-cosimplicial-resolution}
On reprend les notations et les hypothèses de
l'exemple \ref{example-category-to-point}. Soit $U_\bullet$ un objet
cosimplicial de $\mathcal{C}$ tel que, pour tout
$U \in \Ob(\mathcal{C})$, l'ensemble simplicial
$\Mor_\mathcal{C}(U_\bullet, U)$
soit homotopiquement équivalent à l'ensemble simplicial constant sur un
singleton. Alors
$$
L\pi_!(\mathcal{F}) = \mathcal{F}(U_\bullet)
$$
dans $D(\textit{Ab})$, resp.\ $D(B)$, fonctoriellement en $\mathcal{F}$ dans
$\textit{Ab}(\mathcal{C})$, resp.\ $\textit{Mod}(\underline{B})$.
\end{lemma}

\begin{proof}
Comme $L\pi_!$ coïncide pour les modules et les faisceaux abéliens d'après
le lemme \ref{lemma-properties-lower-shriek-fibred-category}, il suffit de
démontrer l'énoncé lorsque $\mathcal{F}$ est un faisceau abélien.
Pour $U \in \Ob(\mathcal{C})$, le faisceau abélien
$j_{U!}\mathbf{Z}_U$ est un objet projectif de
$\textit{Ab}(\mathcal{C})$, car
$\Hom(j_{U!}\mathbf{Z}_U, \mathcal{F}) = \mathcal{F}(U)$
et le foncteur des sections est exact puisque la topologie est chaotique.
Tout faisceau abélien est quotient d'une somme directe de
$j_{U!}\mathbf{Z}_U$ d'après Modules sur les sites, lemme
\ref{sites-modules-lemma-module-quotient-flat}. Nous pouvons donc calculer
$L\pi_!(\mathcal{F})$ en choisissant une résolution
$$
\ldots \to \mathcal{G}^{-1} \to \mathcal{G}^0 \to \mathcal{F} \to 0
$$
dont les termes sont des sommes directes de faisceaux de la forme ci-dessus,
puis en prenant $L\pi_!(\mathcal{F}) = \pi_!(\mathcal{G}^\bullet)$.
Considérons le bicomplexe
$A^{\bullet, \bullet} = \mathcal{G}^\bullet(U_\bullet)$.
Le morphisme $\mathcal{G}^0 \to \mathcal{F}$ donne un morphisme de complexes
$A^{0, \bullet} \to \mathcal{F}(U_\bullet)$.
Comme $\pi_!$ se calcule en prenant la colimite sur
$\mathcal{C}^{opp}$ (lemme \ref{lemma-compute-pi-shriek}), nous voyons que les
deux compositions
$\mathcal{G}^m(U_1) \to \mathcal{G}^m(U_0) \to \pi_!\mathcal{G}^m$
sont égales. Nous obtenons donc un morphisme canonique de complexes
$$
\text{Tot}(A^{\bullet, \bullet})
\longrightarrow
\pi_!(\mathcal{G}^\bullet) = L\pi_!(\mathcal{F})
$$
Pour démontrer le lemme, il suffit de montrer que les complexes
$$
\ldots \to \mathcal{G}^m(U_1) \to \mathcal{G}^m(U_0) \to
\pi_!\mathcal{G}^m \to 0
$$
sont exacts, voir Homologie, lemme
\ref{homology-lemma-double-complex-gives-resolution}.
Comme les faisceaux $\mathcal{G}^m$ sont des sommes directes des faisceaux
$j_{U!}\mathbf{Z}_U$, nous nous ramenons à
$\mathcal{G} = j_{U!}\mathbf{Z}_U$. Le complexe
$j_{U!}\mathbf{Z}_U(U_\bullet)$ est le complexe de groupes abéliens associé
au $\mathbf{Z}$-module libre sur l'ensemble simplicial
$\Mor_\mathcal{C}(U_\bullet, U)$, que nous avons supposé homotopiquement
équivalent à un singleton. Nous en concluons que
$$
j_{U!}\mathbf{Z}_U(U_\bullet) \to \mathbf{Z}
$$
est une équivalence d'homotopie de groupes abéliens, donc un
quasi-isomorphisme (Simplicial, remarque
\ref{simplicial-remark-homotopy-better} et lemme
\ref{simplicial-lemma-homotopy-s-N}). Cela achève la démonstration, puisque
$\pi_!j_{U!}\mathbf{Z}_U = \mathbf{Z}$, comme cela a été montré dans la
démonstration du lemme
\ref{lemma-properties-lower-shriek-fibred-category}.
\end{proof}

\begin{lemma}
\label{lemma-get-it-now}
On reprend les notations et les hypothèses de
l'exemple \ref{example-morphism-categories}. S'il existe un objet
cosimplicial $U'_\bullet$ de $\mathcal{C}'$ tel que le
lemme \ref{lemma-compute-by-cosimplicial-resolution} s'applique à la fois à
$U'_\bullet$ dans $\mathcal{C}'$ et à $u(U'_\bullet)$ dans $\mathcal{C}$,
alors $L\pi'_! \circ g^{-1} = L\pi_!$ comme foncteurs
$D(\mathcal{C}) \to D(\textit{Ab})$,
resp.\ $D(\mathcal{C}, \underline{B}) \to D(B)$.
\end{lemma}

\begin{proof}
Cela résulte immédiatement du
lemme \ref{lemma-compute-by-cosimplicial-resolution} et du fait que
$g^{-1}$ est donné par précomposition avec $u$.
\end{proof}

\begin{lemma}
\label{lemma-product-categories}
Soient $\mathcal{C}_i$, $i = 1, 2$, des catégories. Soient
$u_i : \mathcal{C}_1 \times \mathcal{C}_2 \to \mathcal{C}_i$ les foncteurs
de projection. Soit $B$ un anneau. Soient
$g_i : (\Sh(\mathcal{C}_1 \times \mathcal{C}_2), \underline{B}) \to
(\Sh(\mathcal{C}_i), \underline{B})$ les morphismes de topos annelés
correspondants, voir l'exemple \ref{example-morphism-categories}. Pour
$K_i \in D(\mathcal{C}_i, B)$, nous avons
$$
L(\pi_1 \times \pi_2)_!(
g_1^{-1}K_1 \otimes_{\underline{B}}^\mathbf{L} g_2^{-1}K_2)
=
L\pi_{1, !}(K_1) \otimes_B^\mathbf{L} L\pi_{2, !}(K_2)
$$
dans $D(B)$, avec les notations évidentes.
\end{lemma}

\begin{proof}
Comme les deux membres commutent aux colimites, il suffit de le démontrer pour
$K_1 = j_{U!}\underline{B}_U$ et $K_2 = j_{V!}\underline{B}_V$
avec $U \in \Ob(\mathcal{C}_1)$ et $V \in \Ob(\mathcal{C}_2)$.
Voir la construction de $L\pi_!$ dans le
lemme \ref{lemma-existence-derived-lower-shriek}. Dans ce cas,
$$
g_1^{-1}K_1 \otimes_{\underline{B}}^\mathbf{L} g_2^{-1}K_2 =
g_1^{-1}K_1 \otimes_{\underline{B}} g_2^{-1}K_2 =
j_{(U, V)!}\underline{B}_{(U, V)}
$$
Vérification omise. Le résultat s'ensuit, car les membres de gauche et de
droite de la formule du lemme s'évaluent tous deux en $B$, voir la
construction de $L\pi_!$ dans le
lemme \ref{lemma-existence-derived-lower-shriek}.
\end{proof}

\begin{lemma}
\label{lemma-eilenberg-zilber}
On reprend les notations et les hypothèses de
l'exemple \ref{example-category-to-point}. S'il existe un objet cosimplicial
$U_\bullet$ de $\mathcal{C}$ auquel le
lemme \ref{lemma-compute-by-cosimplicial-resolution} s'applique, alors
$$
L\pi_!(K_1 \otimes^\mathbf{L}_{\underline{B}} K_2) =
L\pi_!(K_1) \otimes^\mathbf{L}_B L\pi_!(K_2)
$$
pour tous $K_i \in D(\underline{B})$.
\end{lemma}

\begin{proof}
Considérons le diagramme de catégories et de foncteurs
$$
\xymatrix{
& & \mathcal{C} \\
\mathcal{C} \ar[r]^-u &
\mathcal{C} \times \mathcal{C} \ar[rd]^{u_2} \ar[ru]_{u_1} \\
& & \mathcal{C}
}
$$
où $u$ est le foncteur diagonal et les $u_i$ sont les foncteurs de projection.
Cela donne des morphismes de topos annelés $g$, $g_1$, $g_2$.
Pour tout objet $(U_1, U_2)$ de $\mathcal{C}$, nous avons
$$
\Mor_{\mathcal{C} \times \mathcal{C}}(u(U_\bullet), (U_1, U_2)) =
\Mor_\mathcal{C}(U_\bullet, U_1) \times \Mor_\mathcal{C}(U_\bullet, U_2)
$$
qui est homotopiquement équivalent à un point d'après Simplicial, lemme
\ref{simplicial-lemma-products-homotopy}. Ainsi, le
lemme \ref{lemma-get-it-now} donne
$L\pi_!(g^{-1}K) = L(\pi \times \pi)_!(K)$ pour tout $K$ dans
$D(\mathcal{C} \times \mathcal{C}, B)$. Prenons
$K = g_1^{-1}K_1 \otimes_B^\mathbf{L} g_2^{-1}K_2$. Alors
$g^{-1}K = K_1 \otimes^\mathbf{L}_{\underline{B}} K_2$, car
$g^{-1} = g^* = Lg^*$ commute au produit tensoriel dérivé
(lemme \ref{lemma-pullback-tensor-product}). Pour finir, nous appliquons le
lemme \ref{lemma-product-categories}.
\end{proof}

\begin{remark}[Modules simpliciaux]
\label{remark-simplicial-modules}
Soit $\mathcal{C} = \Delta$ et soit $B$ un anneau quelconque. C'est un cas
particulier de l'exemple \ref{example-category-to-point} où les hypothèses du
lemme \ref{lemma-compute-by-cosimplicial-resolution} sont vérifiées.
En effet, soit $U_\bullet$ l'objet cosimplicial de $\Delta$ donné par le
foncteur identité. Pour vérifier la condition, il faut montrer que, pour
$[m] \in \Ob(\Delta)$, l'ensemble simplicial
$\Delta[m] : n \mapsto \Mor_\Delta([n], [m])$ est homotopiquement équivalent
à un point. Cela est expliqué dans Simplicial, exemple
\ref{simplicial-example-simplex-contractible}.

\medskip\noindent
Dans cette situation, la catégorie $\textit{Mod}(\underline{B})$ est
simplement celle des $B$-modules simpliciaux, et le foncteur $L\pi_!$ envoie
un $B$-module simplicial $M_\bullet$ sur son complexe associé
$s(M_\bullet)$ de $B$-modules. Les résultats ci-dessus peuvent donc être
réinterprétés comme des résultats sur les modules simpliciaux. Par exemple,
un cas particulier du lemme \ref{lemma-eilenberg-zilber} est le suivant : si
$M_\bullet$, $M'_\bullet$ sont des $B$-modules simpliciaux plats, alors le
complexe $s(M_\bullet \otimes_B M'_\bullet)$ est quasi-isomorphe au complexe
total associé au bicomplexe $s(M_\bullet) \otimes_B s(M'_\bullet)$.
(Indication : utiliser la platitude pour passer des produits tensoriels
dérivés aux produits tensoriels usuels.) C'est un cas particulier du théorème
d'Eilenberg--Zilber que l'on peut trouver dans \cite{Eilenberg-Zilber}.
\end{remark}

\begin{lemma}
\label{lemma-O-homology-qis}
Soit $\mathcal{C}$ une catégorie (munie de la topologie chaotique).
Soit $\mathcal{O} \to \mathcal{O}'$ un morphisme de faisceaux d'anneaux sur
$\mathcal{C}$. Supposons que :
\begin{enumerate}
\item il existe un objet cosimplicial $U_\bullet$ de $\mathcal{C}$ comme
dans le lemme \ref{lemma-compute-by-cosimplicial-resolution}, et
\item $L\pi_!\mathcal{O} \to L\pi_!\mathcal{O}'$ est un isomorphisme.
\end{enumerate}
Pour $K$ dans $D(\mathcal{O})$, nous avons
$$
L\pi_!(K) = L\pi_!(K \otimes_\mathcal{O}^\mathbf{L} \mathcal{O}')
$$
dans $D(\textit{Ab})$.
\end{lemma}

\begin{proof}
Remarque : dans cette démonstration, $L\pi_!$ désigne le foncteur dérivé à
gauche de $\pi_!$ sur les faisceaux abéliens.
Comme $L\pi_!$ commute aux colimites, il suffit de démontrer l'énoncé pour
les complexes bornés supérieurement de $\mathcal{O}$-modules (comparer avec
l'argument de Catégories dérivées, proposition
\ref{derived-proposition-left-derived-exists}, ou se limiter simplement aux
complexes bornés supérieurement). Tout tel complexe est quasi-isomorphe à un
complexe borné supérieurement dont les termes sont des sommes directes de
$j_{U!}\mathcal{O}_U$, avec $U \in \Ob(\mathcal{C})$, voir Modules sur les
sites, lemme \ref{sites-modules-lemma-module-quotient-flat}.
Il suffit donc de démontrer le lemme pour $j_{U!}\mathcal{O}_U$. Par
hypothèse,
$$
S_\bullet = \Mor_\mathcal{C}(U_\bullet, U)
$$
est un ensemble simplicial homotopiquement équivalent à l'ensemble simplicial
constant sur un singleton. Posons $P_n = \mathcal{O}(U_n)$ et
$P'_n = \mathcal{O}'(U_n)$. Observons que le complexe associé au groupe
abélien simplicial
$$
X_\bullet : n \longmapsto \bigoplus\nolimits_{s \in S_n} P_n
$$
calcule $L\pi_!(j_{U!}\mathcal{O}_U)$ d'après le
lemme \ref{lemma-compute-by-cosimplicial-resolution}.
Comme $j_{U!}\mathcal{O}_U$ est un $\mathcal{O}$-module plat, nous avons
$j_{U!}\mathcal{O}_U \otimes^\mathbf{L}_\mathcal{O} \mathcal{O}' =
j_{U!}\mathcal{O}'_U$, et son image par $L\pi_!$ est calculée par le complexe
associé au groupe abélien simplicial
$$
X'_\bullet : n \longmapsto \bigoplus\nolimits_{s \in S_n} P'_n
$$
Comme la règle qui associe à un ensemble simplicial $T_\bullet$ le groupe
abélien simplicial de termes $\bigoplus_{t \in T_n} P_n$ est un foncteur,
$X_\bullet \to P_\bullet$ est une équivalence d'homotopie de groupes
abéliens simpliciaux. De même, la règle qui associe à $T_\bullet$ le groupe
abélien simplicial de termes $\bigoplus_{t \in T_n} P'_n$ est un foncteur.
Ainsi, $X'_\bullet \to P'_\bullet$ est une équivalence d'homotopie de
groupes abéliens simpliciaux.
Par hypothèse, $P_\bullet \to P'_\bullet$ est un quasi-isomorphisme
(puisque $P_\bullet$, resp.\ $P'_\bullet$, calcule $L\pi_!\mathcal{O}$,
resp.\ $L\pi_!\mathcal{O}'$, d'après le
lemme \ref{lemma-compute-by-cosimplicial-resolution}).
Nous concluons que $X_\bullet$ et $X'_\bullet$ sont quasi-isomorphes, comme
voulu.
\end{proof}

\begin{remark}
\label{remark-O-homology-B-homology-general}
Soient $\mathcal{C}$ et $B$ comme dans
l'exemple \ref{example-category-to-point}. Supposons qu'il existe un objet
cosimplicial comme dans le
lemme \ref{lemma-compute-by-cosimplicial-resolution}.
Soit $\mathcal{O} \to \underline{B}$ un morphisme de faisceaux d'anneaux sur
$\mathcal{C}$ qui induit un isomorphisme
$L\pi_!\mathcal{O} \to L\pi_!\underline{B}$.
Dans ce cas, nous obtenons un foncteur exact de catégories triangulées
$$
L\pi_! : D(\mathcal{O}) \longrightarrow D(B)
$$
En effet, pour tout objet $K$ de $D(\mathcal{O})$, nous avons
$L\pi^{\textit{Ab}}_!(K) =
L\pi^{\textit{Ab}}_!(K \otimes_{\mathcal{O}}^\mathbf{L} \underline{B})$
par le lemme \ref{lemma-O-homology-qis}.
Nous pouvons donc définir le foncteur affiché comme la composition de
$- \otimes^\mathbf{L}_\mathcal{O} \underline{B}$ avec le foncteur
$L\pi_! : D(\underline{B}) \to D(B)$.
Autrement dit, nous obtenons une structure de $B$-module sur $L\pi_!(K)$
provenant de l'identification (canonique et fonctorielle) de $L\pi_!(K)$ à
$L\pi_!(K \otimes_\mathcal{O}^\mathbf{L} \underline{B})$ donnée par le
lemme.
\end{remark}







\section{Calcul de l'image directe exceptionnelle dérivée}
\label{section-calculate}

\noindent
Dans cette section, nous appliquons les résultats de la
section \ref{section-homology}
pour calculer
$L\pi_!$ dans la situation \ref{situation-fibred-category} et
$Lg_!$ dans la situation \ref{situation-morphism-fibred-categories}.

\begin{lemma}
\label{lemma-compute-left-derived-pi-shriek-pre}
On reprend les hypothèses et les notations de la situation \ref{situation-fibred-category}.
Pour $\mathcal{F}$ dans $\textit{PAb}(\mathcal{C})$ et $n \geq 0$,
considérons le faisceau abélien $L_n(\mathcal{F})$ sur $\mathcal{D}$
qui est le faisceau associé au préfaisceau
$$
V \longmapsto H_n(\mathcal{C}_V, \mathcal{F}|_{\mathcal{C}_V})
$$
muni des applications de restriction indiquées dans la démonstration. Alors
$L_n(\mathcal{F}) = L_n(\mathcal{F}^\#)$.
\end{lemma}

\begin{proof}
Pour tout morphisme $h : V' \to V$ de $\mathcal{D}$, il existe un
foncteur de changement de base $h^* : \mathcal{C}_V \to \mathcal{C}_{V'}$ entre les
catégories fibres (Catégories, définition
\ref{categories-definition-pullback-functor-fibred-category}).
De plus, pour $U \in \Ob(\mathcal{C}_V)$, il existe un
morphisme fortement cartésien $h^*U \to U$ au-dessus de $h$.
La restriction le long de ces morphismes fortement cartésiens définit une
transformation de foncteurs
$$
\mathcal{F}|_{\mathcal{C}_V}
\longrightarrow
\mathcal{F}|_{\mathcal{C}_{V'}} \circ h^*.
$$
D'après l'exemple \ref{example-morphism-categories},
nous obtenons l'application de restriction voulue
$$
H_n(\mathcal{C}_V, \mathcal{F}|_{\mathcal{C}_V})
\longrightarrow
H_n(\mathcal{C}_{V'}, \mathcal{F}|_{\mathcal{C}_{V'}})
$$
Notons $L_{n, p}(\mathcal{F})$ ce préfaisceau, de sorte que
$L_n(\mathcal{F}) = L_{n, p}(\mathcal{F})^\#$.
Le morphisme canonique $\gamma : \mathcal{F} \to \mathcal{F}^+$
(Sites, théorème \ref{sites-theorem-plus})
définit un morphisme
canonique $L_{n, p}(\mathcal{F}) \to L_{n, p}(\mathcal{F}^+)$.
Il faut montrer que ce morphisme devient un isomorphisme après faisceautisation.

\medskip\noindent
Utilisons le calcul de l'homologie donné dans
l'exemple \ref{example-left-derived-colimits}. Notons
$K_\bullet(\mathcal{F}|_{\mathcal{C}_V})$ le complexe associé à
la restriction de $\mathcal{F}$ à la catégorie fibre $\mathcal{C}_V$.
D'après les remarques ci-dessus, nous obtenons un préfaisceau $K_\bullet(\mathcal{F})$
de complexes
$$
V \longmapsto K_\bullet(\mathcal{F}|_{\mathcal{C}_V})
$$
dont les préfaisceaux d'homologie sont les préfaisceaux $L_{n, p}(\mathcal{F})$.
Il suffit donc de montrer que
$$
K_\bullet(\mathcal{F}) \longrightarrow K_\bullet(\mathcal{F}^+)
$$
devient un isomorphisme après faisceautisation.

\medskip\noindent
Injectivité. Soit $V$ un objet de $\mathcal{D}$ et soit
$\xi \in K_n(\mathcal{F})(V)$ un élément dont l'image
est nulle dans $K_n(\mathcal{F}^+)(V)$. Il faut montrer qu'il existe un
recouvrement $\{V_j \to V\}$ tel que $\xi|_{V_j}$ soit nul dans
$K_n(\mathcal{F})(V_j)$. Écrivons
$$
\xi = \sum (U_{i, n} \to \ldots \to U_{i, 0}, \sigma_i)
$$
avec $\sigma_i \in \mathcal{F}(U_{i, 0})$. Nous pouvons faire en sorte que
chaque suite de morphismes $U_n \to \ldots \to U_0$ de $\mathcal{C}_V$
apparaisse au plus une fois. Comme les sommes intervenant dans la définition
du complexe $K_\bullet$ sont des sommes directes, l'image de cet élément ne peut
être nulle dans $K_n(\mathcal{F}^+)(V)$ que si tous les $\sigma_i$ ont une image
nulle dans $\mathcal{F}^+(U_{i, 0})$. Par construction de
$\mathcal{F}^+$, il existe des recouvrements $\{U_{i, 0, j} \to U_{i, 0}\}$
tels que $\sigma_i|_{U_{i, 0, j}}$ soit nul. Par notre construction de
la topologie sur $\mathcal{C}$, on peut écrire $U_{i, 0, j} \to U_{i, 0}$
comme le changement de base (Catégories, définition
\ref{categories-definition-pullback-functor-fibred-category})
de certains morphismes $V_{i, j} \to V$ et, de plus, chaque
$\{V_{i, j} \to V\}$ est un recouvrement. Choisissons un recouvrement
$\{V_j \to V\}$ qui raffine chacun des recouvrements $\{V_{i, j} \to V\}$.
Il est alors clair que $\xi|_{V_j} = 0$.

\medskip\noindent
Surjectivité. Démonstration omise. Indication : raisonner comme dans la preuve de
l'injectivité.
\end{proof}

\begin{lemma}
\label{lemma-compute-left-derived-pi-shriek}
On reprend les hypothèses et les notations de la situation \ref{situation-fibred-category}.
Pour $\mathcal{F}$ dans $\textit{Ab}(\mathcal{C})$ et $n \geq 0$,
le faisceau $L_n\pi_!(\mathcal{F})$ est égal au faisceau
$L_n(\mathcal{F})$ construit dans le
lemme \ref{lemma-compute-left-derived-pi-shriek-pre}.
\end{lemma}

\begin{proof}
Considérons la suite de foncteurs $\mathcal{F} \mapsto L_n(\mathcal{F})$
de $\textit{PAb}(\mathcal{C}) \to \textit{Ab}(\mathcal{D})$.
Puisque, pour tout $V \in \Ob(\mathcal{D})$, la suite de foncteurs
$H_n(\mathcal{C}_V, - )$ forme un $\delta$-foncteur,
il en va de même des foncteurs $\mathcal{F} \mapsto L_n(\mathcal{F})$.
Notre but est de montrer que ceux-ci forment un $\delta$-foncteur universel.
Pour cela, nous construisons des préfaisceaux abéliens
sur lesquels ces foncteurs s'annulent.

\medskip\noindent
Pour $U' \in \Ob(\mathcal{C})$, considérons le préfaisceau abélien
$\mathcal{F}_{U'} = j_{U'!}^{\textit{PAb}}\mathbf{Z}_{U'}$
(Modules sur les sites, remarque \ref{sites-modules-remark-localize-presheaves}).
Rappelons que
$$
\mathcal{F}_{U'}(U) =
\bigoplus\nolimits_{\Mor_\mathcal{C}(U, U')} \mathbf{Z}
$$
Si $U$ est au-dessus de $V = p(U)$ dans $\mathcal{D}$ et si $U'$ est au-dessus de $V' = p(U')$,
alors tout morphisme $a : U \to U'$ se factorise de manière unique sous la forme $U \to h^*U' \to U'$,
où $h = p(a) : V \to V'$ (voir
Catégories, définition
\ref{categories-definition-pullback-functor-fibred-category}).
Nous voyons donc que
$$
\mathcal{F}_{U'}|_{\mathcal{C}_V}
=
\bigoplus\nolimits_{h \in \Mor_\mathcal{D}(V, V')}
j_{h^*U'!}\mathbf{Z}_{h^*U'}
$$
où $j_{h^*U'} : \Sh(\mathcal{C}_V/h^*U') \to \Sh(\mathcal{C}_V)$
est le morphisme de localisation. Les faisceaux $j_{h^*U'!}\mathbf{Z}_{h^*U'}$
ont leurs groupes d'homologie supérieure nuls (voir
l'exemple \ref{example-left-derived-colimits}).
Nous en concluons que $L_n(\mathcal{F}_{U'}) = 0$ pour tout $n > 0$ et tout $U'$.
Il s'ensuit que tout préfaisceau abélien $\mathcal{F}$ est un quotient
d'un préfaisceau abélien $\mathcal{G}$ tel que $L_n(\mathcal{G}) = 0$ pour
tout $n > 0$ (Modules sur les sites, lemme
\ref{sites-modules-lemma-module-quotient-flat}).
Comme $L_n(\mathcal{F}) = L_n(\mathcal{F}^\#)$, nous voyons
qu'il en va de même pour les faisceaux abéliens. Ainsi,
la suite de foncteurs $L_n(-)$ est un delta-foncteur universel
sur $\textit{Ab}(\mathcal{C})$
(Homologie, lemme \ref{homology-lemma-efface-implies-universal}).
Comme elle coïncide avec
$H^{-n}(L\pi_!(-))$ pour $n = 0$ d'après le
lemme \ref{lemma-compute-pi-shriek},
nous concluons par l'unicité des $\delta$-foncteurs universels
(Homologie, lemme \ref{homology-lemma-uniqueness-universal-delta-functor})
et par
Catégories dérivées, lemme \ref{derived-lemma-right-derived-delta-functor}.
\end{proof}

\begin{lemma}
\label{lemma-compute-left-derived-g-shriek}
On reprend les hypothèses et les notations de la
situation \ref{situation-morphism-fibred-categories}.
Pour un faisceau abélien $\mathcal{F}'$ sur $\mathcal{C}'$, le faisceau
$L_ng_!(\mathcal{F}')$ est le faisceau associé au préfaisceau
$$
U \longmapsto H_n(\mathcal{I}_U, \mathcal{F}'_U)
$$
Pour les notations et les applications de restriction, voir la démonstration.
\end{lemma}

\begin{proof}
Posons $p(U) = V$. La catégorie $\mathcal{I}_U$ est la catégorie des couples
$(U', \varphi)$ où $\varphi : U \to u(U')$ est un morphisme de $\mathcal{C}$
tel que $p(\varphi) = \text{id}_V$, c'est-à-dire que $\varphi$ est un morphisme de la
catégorie fibre $\mathcal{C}_V$. Les morphismes
$(U'_1, \varphi_1) \to (U'_2, \varphi_2)$ sont donnés par les morphismes
$a : U'_1 \to U'_2$ de la catégorie fibre $\mathcal{C}'_V$ tels que
$\varphi_2 = u(a) \circ \varphi_1$. Le préfaisceau $\mathcal{F}'_U$ associe
à $(U', \varphi)$ le groupe $\mathcal{F}'(U')$.
Nous allons construire ci-dessous les applications de restriction.

\medskip\noindent
Choisissons une factorisation
$$
\xymatrix{
\mathcal{C}' \ar@<1ex>[r]^{u'} &
\mathcal{C}'' \ar[r]^{u''} \ar@<1ex>[l]^w & \mathcal{C}
}
$$
de $u$ comme dans
Catégories, lemme \ref{categories-lemma-ameliorate-morphism-fibred-categories}.
Alors $g_! = g''_! \circ g'_!$, et il en va de même pour les foncteurs dérivés.
D'autre part, le foncteur $g'_!$ est exact, voir
Modules sur les sites, lemme \ref{sites-modules-lemma-have-left-adjoint}.
Nous obtenons donc $Lg_!(\mathcal{F}') = Lg''_!(\mathcal{F}'')$, où
$\mathcal{F}'' = g'_!\mathcal{F}'$. Remarquons que
$\mathcal{F}'' = h^{-1}\mathcal{F}'$, où
$h : \Sh(\mathcal{C}'') \to \Sh(\mathcal{C}')$ est le morphisme de topos
associé à $w$, voir
Sites, lemme \ref{sites-lemma-have-left-adjoint}.
Le foncteur $u''$ fait de $\mathcal{C}''$ une catégorie fibrée
sur $\mathcal{C}$ ; le
lemme \ref{lemma-compute-left-derived-pi-shriek}
s'applique donc au calcul de $L_ng''_!$. Le résultat s'ensuit, car la
construction de $\mathcal{C}''$ dans la preuve de
Catégories, lemme \ref{categories-lemma-ameliorate-morphism-fibred-categories}
montre que la catégorie fibre $\mathcal{C}''_U$ est égale à
$\mathcal{I}_U$. De plus, $h^{-1}\mathcal{F}'|_{\mathcal{C}''_U}$
est donné par la règle décrite ci-dessus
(comme $w$ est continu et cocontinu d'après
Champs, lemme \ref{stacks-lemma-topology-inherited-functorial},
nous pouvons appliquer
Sites, lemme \ref{sites-lemma-when-shriek}).
\end{proof}
\section{Modules simpliciaux}
\label{section-simplicial-modules}

\noindent
Soit $A_\bullet$ un anneau simplicial. Rappelons que l'on peut voir
$A_\bullet$ comme un faisceau sur $\Delta$ (muni de la topologie chaotique),
voir Simplicial, section
\ref{simplicial-section-simplicial-presheaves}. Un module simplicial
$M_\bullet$ sur $A_\bullet$ n'est alors autre qu'un faisceau de
$A_\bullet$-modules sur $\Delta$. Autrement dit, pour tout $n \geq 0$, on a
un $A_n$-module $M_n$ et, pour toute application
$\varphi : [n] \to [m]$, une application correspondante
$$
M_\bullet(\varphi) : M_m \longrightarrow M_n
$$
qui est $A_\bullet(\varphi)$-linéaire, ces applications se composant de la
manière usuelle.

\medskip\noindent
Soit $\mathcal{C}$ un site. Un {\it faisceau simplicial d'anneaux}
$\mathcal{A}_\bullet$ sur $\mathcal{C}$ est un objet simplicial de la
catégorie des faisceaux d'anneaux sur $\mathcal{C}$. Dans ce cas,
l'affectation $U \mapsto \mathcal{A}_\bullet(U)$ est un faisceau d'anneaux
simpliciaux et, en fait, les deux notions sont équivalentes. Il en va de même
pour les faisceaux abéliens simpliciaux, les faisceaux simpliciaux d'algèbres
de Lie, et ainsi de suite.

\medskip\noindent
Cependant, comme dans le cas des anneaux simpliciaux ci-dessus, on peut
envisager autrement les faisceaux simpliciaux. Considérons en effet la
projection
$$
p : \Delta \times \mathcal{C} \longrightarrow \mathcal{C}
$$
Elle définit une catégorie fibrée dont les morphismes fortement cartésiens
sont exactement ceux de la forme $([n], U) \to ([n], V)$. On munit la
catégorie $\Delta \times \mathcal{C}$ de la topologie héritée de
$\mathcal{C}$ (voir Stacks, section \ref{stacks-section-topology}). La
description simple des recouvrements de $\Delta \times \mathcal{C}$
(Stacks, lemme \ref{stacks-lemma-topology-inherited}) montre immédiatement
qu'un faisceau simplicial d'anneaux sur $\mathcal{C}$ revient au même qu'un
faisceau d'anneaux sur $\Delta \times \mathcal{C}$.

\medskip\noindent
Par analogie avec les modules simpliciaux sur un anneau simplicial, on
définit comme suit les modules simpliciaux sur un faisceau simplicial
d'anneaux.

\begin{definition}
\label{definition-simplicial-module}
Soit $\mathcal{C}$ un site et soit $\mathcal{A}_\bullet$ un faisceau
simplicial d'anneaux sur $\mathcal{C}$. Un
{\it $\mathcal{A}_\bullet$-module simplicial} $\mathcal{F}_\bullet$
(parfois appelé {\it faisceau simplicial de
$\mathcal{A}_\bullet$-modules}) est un faisceau de modules sur le faisceau
d'anneaux sur $\Delta \times \mathcal{C}$ associé à
$\mathcal{A}_\bullet$.
\end{definition}

\noindent
On obtient une catégorie $\textit{Mod}(\mathcal{A}_\bullet)$ de modules
simpliciaux et la catégorie dérivée correspondante
$D(\mathcal{A}_\bullet)$. Un morphisme
$\mathcal{A}_\bullet \to \mathcal{B}_\bullet$ de faisceaux simpliciaux
d'anneaux donne un foncteur
$$
- \otimes^\mathbf{L}_{\mathcal{A}_\bullet} \mathcal{B}_\bullet :
D(\mathcal{A}_\bullet)
\longrightarrow
D(\mathcal{B}_\bullet)
$$
De plus, les résultats des sections précédentes déterminent un foncteur
$$
L\pi_! : D(\mathcal{A}_\bullet) \longrightarrow D(\mathcal{C})
$$
Pour un module simplicial $\mathcal{F}_\bullet$, l'objet
$L\pi_!(\mathcal{F}_\bullet)$ est représenté par le complexe de chaînes
associé $s(\mathcal{F}_\bullet)$ (Simplicial, section
\ref{simplicial-section-complexes}). Cela résulte des lemmes
\ref{lemma-compute-left-derived-pi-shriek} et
\ref{lemma-compute-by-cosimplicial-resolution}.

\begin{lemma}
\label{lemma-base-change-by-qis}
Soit $\mathcal{C}$ un site et soit
$\mathcal{A}_\bullet \to \mathcal{B}_\bullet$ un homomorphisme de faisceaux
simpliciaux d'anneaux sur $\mathcal{C}$. Si
$L\pi_!\mathcal{A}_\bullet \to L\pi_!\mathcal{B}_\bullet$ est un
isomorphisme dans $D(\mathcal{C})$, alors
$$
L\pi_!(K) =
L\pi_!(K \otimes^\mathbf{L}_{\mathcal{A}_\bullet} \mathcal{B}_\bullet)
$$
pour tout $K$ dans $D(\mathcal{A}_\bullet)$.
\end{lemma}

\begin{proof}
Soit $([n], U)$ un objet de $\Delta \times \mathcal{C}$. Comme $L\pi_!$
commute aux colimites, il suffit d'établir l'assertion pour les complexes
bornés supérieurement de $\mathcal{O}$-modules (comparer à l'argument de
Catégories dérivées, proposition
\ref{derived-proposition-left-derived-exists}, ou se limiter simplement aux
complexes bornés supérieurement). Tout tel complexe est quasi-isomorphe à un
complexe borné supérieurement dont les termes sont des modules plats, voir
Modules on Sites, lemme \ref{sites-modules-lemma-module-quotient-flat}. Il
suffit donc de démontrer le lemme pour un
$\mathcal{A}_\bullet$-module plat $\mathcal{F}$. Dans ce cas, le produit
tensoriel dérivé est le produit tensoriel usuel, qui est encore un faisceau.
Le lemme \ref{lemma-compute-left-derived-pi-shriek} permet donc de calculer
les faisceaux de cohomologie des deux membres de l'égalité par le procédé du
lemme \ref{lemma-compute-left-derived-pi-shriek-pre}. Il suffit ainsi
d'établir le résultat pour la restriction de $\mathcal{F}$ aux catégories
fibres (c'est-à-dire à $\Delta \times U$). Dans ce cas, le résultat résulte du
lemme \ref{lemma-O-homology-qis}.
\end{proof}

\begin{remark}
\label{remark-homology-augmentation}
Soit $\mathcal{C}$ un site et soit
$\epsilon : \mathcal{A}_\bullet \to \mathcal{O}$ une augmentation
(Simplicial, définition \ref{simplicial-definition-augmentation}) dans la
catégorie des faisceaux d'anneaux. Supposons que $\epsilon$ induise un
quasi-isomorphisme $s(\mathcal{A}_\bullet) \to \mathcal{O}$. On obtient alors
un foncteur exact de catégories triangulées
$$
L\pi_! : D(\mathcal{A}_\bullet) \longrightarrow D(\mathcal{O})
$$
En effet, pour tout objet $K$ de $D(\mathcal{A}_\bullet)$, on a
$L\pi_!(K) = L\pi_!(K \otimes_{\mathcal{A}_\bullet}^\mathbf{L} \mathcal{O})$
d'après le lemme \ref{lemma-base-change-by-qis}. On peut donc définir le
foncteur affiché comme la composée de
$- \otimes^\mathbf{L}_{\mathcal{A}_\bullet} \mathcal{O}$ avec le foncteur
$L\pi_! : D(\Delta \times \mathcal{C}, \pi^{-1}\mathcal{O}) \to
D(\mathcal{O})$ de la remarque \ref{remark-fibred-category}. Autrement dit,
on obtient une structure de $\mathcal{O}$-module sur $L\pi_!(K)$ provenant de
l'identification (canonique et fonctorielle) de $L\pi_!(K)$ avec
$L\pi_!(K \otimes_{\mathcal{A}_\bullet}^\mathbf{L} \mathcal{O})$ donnée par
le lemme.
\end{remark}






\section{Cohomologie sur une catégorie}
\label{section-cohomology}

\noindent
Dans la situation de l'exemple \ref{example-category-to-point}, outre le
foncteur dérivé $L\pi_!$, on dispose également du foncteur $R\pi_*$. Pour un
faisceau abélien $\mathcal{F}$ sur $\mathcal{C}$, on a
$H_n(\mathcal{C}, \mathcal{F}) = H^{-n}(L\pi_!\mathcal{F})$ et
$H^n(\mathcal{C}, \mathcal{F}) = H^n(R\pi_*\mathcal{F})$.

\begin{example}[Calcul de la cohomologie]
\label{example-right-derived-limits}
Dans l'exemple \ref{example-category-to-point}, on peut calculer les
foncteurs $H^n(\mathcal{C}, -)$ comme suit. Soit
$\mathcal{F} \in \Ob(\textit{Ab}(\mathcal{C}))$. Considérons le complexe de
cochaînes
$$
K^\bullet(\mathcal{F}) :
\prod\nolimits_{U_0} \mathcal{F}(U_0)
\to
\prod\nolimits_{U_0 \to U_1} \mathcal{F}(U_0)
\to
\prod\nolimits_{U_0 \to U_1 \to U_2} \mathcal{F}(U_0)
\to \ldots
$$
dont les applications de transition sont données par
$$
(s_{U_0 \to U_1})
\longmapsto
((U_0 \to U_1 \to U_2) \mapsto s_{U_0 \to U_1} - s_{U_0 \to U_2}
+ s_{U_1 \to U_2}|_{U_0})
$$
et de même aux autres degrés. Par construction,
$$
H^0(\mathcal{C}, \mathcal{F}) =
\lim_{\mathcal{C}^{opp}} \mathcal{F} =
H^0(K^\bullet(\mathcal{F})),
$$
voir Categories, lemme
\ref{categories-lemma-limits-products-equalizers}. La construction de
$K^\bullet(\mathcal{F})$ est fonctorielle en $\mathcal{F}$ et transforme les
suites exactes courtes de $\textit{Ab}(\mathcal{C})$ en suites exactes
courtes de complexes. La famille de foncteurs
$\mathcal{F} \mapsto H^n(K^\bullet(\mathcal{F}))$ forme donc un
$\delta$-foncteur, voir Homology, définition
\ref{homology-definition-cohomological-delta-functor} et lemme
\ref{homology-lemma-long-exact-sequence-cochain}. Pour un objet $U$ de
$\mathcal{C}$, notons $p_U : \Sh(*) \to \Sh(\mathcal{C})$ le point
correspondant, pour lequel $p_U^{-1}$ est l'évaluation en $U$, voir Sites,
exemple \ref{sites-example-indiscrete-points}. Soit $A$ un groupe abélien et
posons $\mathcal{F} = p_{U, *}A$. Le complexe $K^\bullet(\mathcal{F})$ a alors
pour termes $\text{Map}(X_n, A)$, où
$$
X_n = \coprod\nolimits_{U_0 \to \ldots \to U_{n - 1} \to U_n}
\Mor_\mathcal{C}(U, U_0)
$$
Cet ensemble simplicial est homotopiquement équivalent à l'ensemble
simplicial constant réduit à un point $\{*\}$. En effet, l'application
$X_\bullet \to \{*\}$ est évidente, l'application $\{*\} \to X_n$ envoie
$*$ sur $(U \to \ldots \to U, \text{id}_U)$, et les applications
$$
h_{n, i} : X_n \longrightarrow X_n
$$
(Simplicial, lemme \ref{simplicial-lemma-relations-homotopy}) qui définissent
l'homotopie entre les deux applications
$X_\bullet \to X_\bullet$ sont données par la règle
$$
h_{n, i} :
(U_0 \to \ldots \to U_n, f)
\longmapsto
(U \to \ldots \to U \to U_i \to \ldots \to U_n, \text{id})
$$
pour $i > 0$, et $h_{n, 0} = \text{id}$. Nous omettons les vérifications.
Comme $\text{Map}(-, A)$ est un foncteur contravariant, il s'ensuit que
$K^\bullet(p_{U, *}A)$ a une cohomologie nulle en degrés strictement positifs
(par la fonctorialité de Simplicial, remarque
\ref{simplicial-remark-homotopy-better}, et le résultat de Simplicial, lemme
\ref{simplicial-lemma-homotopy-s-Q}). Il en résulte que
$K^\bullet(\mathcal{F})$ est encore acyclique en degrés strictement positifs
lorsque $\mathcal{F}$ est un produit de faisceaux de la forme $p_{U, *}A$.
Comme tout faisceau abélien sur $\mathcal{C}$ se plonge dans un tel produit,
on conclut que $K^\bullet(\mathcal{F})$ calcule les foncteurs dérivés à droite
$H^n(\mathcal{C}, -)$ de $H^0(\mathcal{C}, -)$, par exemple d'après
Homology, lemme \ref{homology-lemma-efface-implies-universal}, et Derived
Categories, lemme \ref{derived-lemma-right-derived-delta-functor}.
\end{example}

\begin{example}[Calcul des Ext]
\label{example-computing-exts}
Dans l'exemple \ref{example-category-to-point}, supposons en outre donné un
faisceau d'anneaux $\mathcal{O}$ sur $\mathcal{C}$. Soient $\mathcal{F}$ et
$\mathcal{G}$ des $\mathcal{O}$-modules. Considérons le complexe
$K^\bullet(\mathcal{G}, \mathcal{F})$ dont le terme de degré $n$ est
$$
\prod\nolimits_{U_0 \to U_1 \to \ldots \to U_n}
\Hom_{\mathcal{O}(U_n)}(\mathcal{G}(U_n), \mathcal{F}(U_0))
$$
et dont l'application de transition est donnée par
$$
(\varphi_{U_0 \to U_1})
\longmapsto
((U_0 \to U_1 \to U_2) \mapsto
\varphi_{U_0 \to U_1} \circ \rho^{U_2}_{U_1}
- \varphi_{U_0 \to U_2}
+ \rho^{U_1}_{U_0} \circ \varphi_{U_1 \to U_2}
$$
et de même aux autres degrés. Les $\rho$ désignent ici les applications de
restriction. Par construction,
$$
\Hom_\mathcal{O}(\mathcal{G}, \mathcal{F}) =
H^0(K^\bullet(\mathcal{G}, \mathcal{F}))
$$
pour toute paire de $\mathcal{O}$-modules $\mathcal{F}, \mathcal{G}$.
L'affectation
$(\mathcal{G}, \mathcal{F}) \mapsto K^\bullet(\mathcal{G}, \mathcal{F})$
est un bifoncteur qui transforme les sommes directes en la première variable
en produits et qui commute aux produits en la seconde variable. Nous
affirmons que
$$
\Ext^i_\mathcal{O}(\mathcal{G}, \mathcal{F}) =
H^i(K^\bullet(\mathcal{G}, \mathcal{F}))
$$
pour $i \geq 0$, pourvu que l'une des conditions suivantes soit satisfaite :
\begin{enumerate}
\item $\mathcal{G}(U)$ est un $\mathcal{O}(U)$-module projectif pour tout
$U \in \Ob(\mathcal{C})$ ;
\item $\mathcal{F}(U)$ est un $\mathcal{O}(U)$-module injectif pour tout
$U \in \Ob(\mathcal{C})$.
\end{enumerate}
Dans le cas (1), le foncteur $K^\bullet(\mathcal{G}, -)$ est un foncteur exact
de la catégorie des $\mathcal{O}$-modules vers celle des complexes de
cochaînes de groupes abéliens. En raisonnant comme dans l'exemple
\ref{example-right-derived-limits}, il suffit donc de montrer que
$K^\bullet(\mathcal{G}, \mathcal{F})$ est acyclique en degrés strictement
positifs lorsque $\mathcal{F}$ est de la forme $p_{U, *}A$ pour un
$\mathcal{O}(U)$-module $A$.
Choisissons une suite exacte courte
\begin{equation}
\label{equation-split}
0 \to \mathcal{G}' \to \bigoplus j_{U_i!}\mathcal{O}_{U_i} \to
\mathcal{G} \to 0
\end{equation}
voir Modules on Sites, lemme
\ref{sites-modules-lemma-module-quotient-flat}. Comme (1) vaut pour les
faisceaux du milieu et de droite, il vaut aussi pour $\mathcal{G}'$, et
l'évaluation de (\ref{equation-split}) en un objet de $\mathcal{C}$ donne une
suite exacte scindée de modules. On obtient une suite exacte courte de
complexes
$$
0 \to
K^\bullet(\mathcal{G}, \mathcal{F}) \to
\prod K^\bullet(j_{U_i!}\mathcal{O}_{U_i}, \mathcal{F}) \to
K^\bullet(\mathcal{G}', \mathcal{F}) \to 0
$$
pour tout $\mathcal{F}$, en particulier pour
$\mathcal{F} = p_{U, *}A$. En degré $H^0$, on obtient
$$
0 \to \Hom(\mathcal{G}, p_{U, *}A) \to
\Hom(\bigoplus j_{U_i!}\mathcal{O}_{U_i}, p_{U, *}A) \to
\Hom(\mathcal{G}', p_{U, *}A) \to 0
$$
qui est exacte, car
$\Hom(\mathcal{H}, p_{U, *}A) = \Hom_{\mathcal{O}(U)}(\mathcal{H}(U), A)$
et la suite des sections de (\ref{equation-split}) sur $U$ est exacte
scindée. Par décalage dimensionnel, il suffit donc de montrer que
$K^\bullet(j_{U'!}\mathcal{O}_{U'}, p_{U, *}A)$ est acyclique en degrés
strictement positifs pour tous $U, U' \in \Ob(\mathcal{C})$. Dans ce cas,
$K^n(j_{U'!}\mathcal{O}_{U'}, p_{U, *}A)$ est égal à
$$
\prod\nolimits_{U \to U_0 \to U_1 \to \ldots \to U_n \to U'} A
$$
Autrement dit, $K^\bullet(j_{U'!}\mathcal{O}_{U'}, p_{U, *}A)$ est le
complexe de termes $\text{Map}(X_\bullet, A)$, où
$$
X_n = \coprod\nolimits_{U_0 \to \ldots \to U_{n - 1} \to U_n}
\Mor_\mathcal{C}(U, U_0) \times \Mor_\mathcal{C}(U_n, U')
$$
Cet ensemble simplicial est homotopiquement équivalent à l'ensemble
simplicial constant $\Mor_\mathcal{C}(U, U')$ : cela se démontre exactement
comme l'assertion correspondante de l'exemple
\ref{example-right-derived-limits}. Ceci achève la démonstration de
l'affirmation.

\medskip\noindent
L'argument dans le cas (2) est analogue, par dualité.
\end{example}






\section{Modules sur une catégorie}
\label{section-modules-cohomology}

\noindent
Les résultats de cette section serviront à définir une variante de la
catégorie dérivée des modules quasi-cohérents sur un champ en groupoïdes
au-dessus de la catégorie des schémas. Voir Sheaves on Stacks, section
\ref{stacks-sheaves-section-QC}.

\medskip\noindent
Soit $\mathcal{C}$ une catégorie. Considérons $\mathcal{C}$ comme un site muni
de la topologie chaotique. Comme dans l'exemple
\ref{example-computing-exts}, soit
$\mathcal{O}$ un faisceau d'anneaux sur $\mathcal{C}$. Autrement dit,
$\mathcal{O}$ est un préfaisceau d'anneaux sur la catégorie $\mathcal{C}$,
voir Categories, définition \ref{categories-definition-presheaf}.

\begin{definition}
\label{definition-cartesian}
Dans la situation ci-dessus, on note
$\mathit{QC}(\mathcal{C}, \mathcal{O})$, ou simplement
{\it $\mathit{QC}(\mathcal{O})$}, la sous-catégorie pleine de
$D(\mathcal{O}) = D(\mathcal{C}, \mathcal{O})$ formée des objets $K$ tels que,
pour tout $U \to V$ dans $\mathcal{C}$, l'application canonique
$$
R\Gamma(V, K) \otimes_{\mathcal{O}(V)}^\mathbf{L} \mathcal{O}(U)
\longrightarrow
R\Gamma(U, K)
$$
soit un isomorphisme dans $D(\mathcal{O}(U))$.
\end{definition}

\begin{lemma}
\label{lemma-cartesian-good-sub}
Dans la situation ci-dessus, $\mathit{QC}(\mathcal{O})$ est une
sous-catégorie strictement pleine, saturée et triangulée de
$D(\mathcal{O})$, stable par sommes directes arbitraires.
\end{lemma}

\begin{proof}
Soit $U$ un objet de $\mathcal{C}$. Comme la topologie de $\mathcal{C}$ est
chaotique, le foncteur $\mathcal{F} \mapsto \mathcal{F}(U)$ est exact et
commute aux sommes directes. Le foncteur exact
$K \mapsto R\Gamma(U, K)$ se calcule donc en représentant $K$ par un complexe
quelconque $\mathcal{F}^\bullet$ de $\mathcal{O}$-modules et en prenant
$\mathcal{F}^\bullet(U)$. Ainsi, $R\Gamma(U, -)$ commute aux sommes directes,
voir Injectives, lemme \ref{injectives-lemma-derived-products}. De même, pour
un morphisme $U \to V$ de $\mathcal{C}$, le foncteur produit tensoriel dérivé
$- \otimes_{\mathcal{O}(V)}^\mathbf{L} \mathcal{O}(U) :
D(\mathcal{O}(V)) \to D(\mathcal{O}(U))$
est exact et commute aux sommes directes. Le lemme résulte immédiatement de
ces observations ; nous omettons les détails.
\end{proof}

\begin{lemma}
\label{lemma-QC-bounded-above}
Dans la situation ci-dessus, supposons que $M$ soit un objet de
$\mathit{QC}(\mathcal{O})$ et que $b \in \mathbf{Z}$ vérifie $H^i(M) = 0$ pour
tout $i > b$. Alors $H^b(M)$ est un module quasi-cohérent sur
$(\mathcal{C}, \mathcal{O})$ au sens de Modules on Sites, définition
\ref{sites-modules-definition-site-local}.
\end{lemma}

\begin{proof}
D'après Modules on Sites, lemme
\ref{sites-modules-lemma-chaotic-quasi-coherent}, il suffit de montrer que,
pour tout morphisme $U \to V$ de $\mathcal{C}$, l'application
$$
H^b(M)(V) \otimes_{\mathcal{O}(V)} \mathcal{O}(U) \to H^b(M)(U)
$$
est un isomorphisme. Par hypothèse, l'application
$$
R\Gamma(V, M) \otimes_{\mathcal{O}(V)}^\mathbf{L} \mathcal{O}(U)
\to R\Gamma(U, M)
$$
est un isomorphisme. Le résultat découle alors, par exemple, de la suite
spectrale de Tor. Nous omettons les détails.
\end{proof}

\begin{lemma}
\label{lemma-QC-final-object}
Dans la situation ci-dessus, supposons que $\mathcal{C}$ possède un objet
final $X$. Posons $R = \mathcal{O}(X)$ et notons
$f : (\mathcal{C}, \mathcal{O}) \to (pt, R)$ le morphisme évident de sites.
Alors $\mathit{QC}(\mathcal{O}) = D(R)$ via $Lf^*$ et $Rf_*$.
\end{lemma}

\begin{proof}
Omis.
\end{proof}

\begin{lemma}
\label{lemma-QC-hom-out-of}
Dans la situation ci-dessus, supposons que $K$ soit un objet de
$\mathit{QC}(\mathcal{O})$ et que $M$ soit un objet quelconque de
$D(\mathcal{O})$. Pour tout objet $U$ de $\mathcal{C}$, on a
$$
\Hom_{D(\mathcal{O}_U)}(K|_U, M|_U) =
R\Hom_{\mathcal{O}(U)}(R\Gamma(U, K), R\Gamma(U, M))
$$
\end{lemma}

\begin{proof}
On peut remplacer $\mathcal{C}$ par $\mathcal{C}/U$. On peut donc supposer
que $U = X$ est un objet final de $\mathcal{C}$. D'après le lemme
\ref{lemma-QC-final-object}, on a $K = Lf^*P$, où
$P = R\Gamma(U, K) = R\Gamma(X, K) = Rf_*K$. Le résultat s'ensuit, puisque
$Lf^*$ est adjoint à gauche de $Rf_*(-) = R\Gamma(U, -)$.
\end{proof}

\noindent
Soit $(\mathcal{C}, \mathcal{O})$ comme ci-dessus. Pour un complexe
$\mathcal{F}^\bullet$ de $\mathcal{O}$-modules, on définit la
{\it taille} $|\mathcal{F}^\bullet|$ de $\mathcal{F}^\bullet$ par
$$
|\mathcal{F}^\bullet| =
\left|
\coprod\nolimits_{i \in \mathbf{Z},\ U \in \Ob(\mathcal{C})} \mathcal{F}^i(U)
\right|
$$
Pour un objet $K$ de $D(\mathcal{O})$, on appelle {\it taille} $|K|$ de $K$
le cardinal
$$
|K| = \min
\left\{
\left|
\mathcal{F}^\bullet
\right|
\text{ où }\mathcal{F}^\bullet\text{ représente }K
\right\}
$$
Les propriétés des cardinaux assurent l'existence de ce minimum.

\begin{lemma}
\label{lemma-build-from}
Dans la situation ci-dessus, il existe un cardinal $\kappa$ ayant la
propriété suivante : étant donnés un complexe $\mathcal{F}^\bullet$ de
$\mathcal{O}$-modules et des sous-ensembles
$\Omega^i_U \subset \mathcal{F}^i(U)$, il existe un sous-complexe
$\mathcal{H}^\bullet \subset \mathcal{F}^\bullet$ tel que
$\Omega^i_U \subset \mathcal{H}^i(U)$ et
$|\mathcal{H}^\bullet| \leq \max(\kappa, |\bigcup \Omega^i_U|)$.
\end{lemma}

\begin{proof}
Définissons $\mathcal{H}^i(U)$ comme le sous-$\mathcal{O}(U)$-module de
$\mathcal{F}^i(U)$ engendré par les images des $\Omega^i_V$ et de
$\text{d}(\Omega^{i - 1}_U)$ par restriction le long de tout morphisme
$f : U \to V$. Le cardinal de $\mathcal{H}^i(U)$ est majoré par
le maximum de $\aleph_0$, du cardinal de $\mathcal{O}(U)$, du cardinal de
$\text{Arrows}(\mathcal{C})$ et de $|\bigcup \Omega^i_U|$. Nous omettons les
détails.
\end{proof}

\begin{lemma}
\label{lemma-qis-small}
Dans la situation ci-dessus, il existe un cardinal $\kappa$ ayant la
propriété suivante : pour tout complexe $\mathcal{F}^\bullet$ de
$\mathcal{O}$-modules représentant un objet $K$ de $D(\mathcal{O})$, il existe
un sous-complexe $\mathcal{H}^\bullet \subset \mathcal{F}^\bullet$ tel que
$\mathcal{H}^\bullet$ représente $K$ et tel que
$|\mathcal{H}^\bullet| \leq \max(\kappa, |K|)$.
\end{lemma}

\begin{proof}
Pour commencer, choisissons, pour tout $i$ et tout $U$, un sous-ensemble
$\Omega^i_U \subset \Ker(\text{d} :
\mathcal{F}^i(U) \to \mathcal{F}^{i + 1}(U))$
qui s'applique bijectivement sur
$H^i(K)(U) = H^i(\mathcal{F}^\bullet(U))$.
On a donc $|\Omega^i_U| \leq |K|$, puisque l'on peut représenter $K$ par un
complexe de taille $|K|$. En appliquant le lemme \ref{lemma-build-from}, on
obtient un sous-complexe
$\mathcal{S}^\bullet \subset \mathcal{F}^\bullet$ de taille au plus
$\max(\kappa, |K|)$ qui contient les $\Omega^i_U$ ; par conséquent,
$H^i(\mathcal{S}^\bullet) \to H^i(\mathcal{F}^\bullet)$ est une surjection de
faisceaux.

\medskip\noindent
Nous allons construire par récurrence des sous-complexes
$$
\mathcal{S}^\bullet =
\mathcal{S}_0^\bullet \subset
\mathcal{S}_1^\bullet \subset
\mathcal{S}_2^\bullet \subset \ldots \subset \mathcal{F}^\bullet
$$
de taille $\leq \max(\kappa, |K|)$, de sorte que le noyau de
$H^i(\mathcal{S}_n^\bullet) \to H^i(\mathcal{F}^\bullet)$ coïncide avec le
noyau de
$H^i(\mathcal{S}_n^\bullet) \to H^i(\mathcal{S}_{n + 1}^\bullet)$.
Une fois ceci fait, on pourra prendre
$\mathcal{H}^\bullet = \bigcup \mathcal{S}_n^\bullet$.

\medskip\noindent
Construisons $\mathcal{S}_{n + 1}^\bullet$ à partir de
$\mathcal{S}_n^\bullet$. Pour tout $U$ et tout $i$, soit
$\Omega^{i - 1}_U \subset \mathcal{F}^{i - 1}(U)$ un sous-ensemble que
$\text{d} : \mathcal{F}^{i - 1}(U) \to \mathcal{F}^i(U)$ applique
$\Omega^{i - 1}_U$ bijectivement sur
$$
\mathcal{S}_n^i(U) \cap
\text{Im}(\text{d} : \mathcal{F}^{i - 1}(U)\to \mathcal{F}^i(U))
$$
Remarquons que $|\Omega^i_U| \leq \max(\kappa,|K|)$, car
$\mathcal{S}_n^i(U)$ est ainsi majoré. On obtient alors
$\mathcal{S}_{n + 1}^\bullet$ en appliquant le lemme
\ref{lemma-build-from} aux sous-ensembles
$$
\mathcal{S}_n^i(U) \cup \Omega^i_U \subset \mathcal{F}^i(U)
$$
ce qui achève la démonstration.
\end{proof}

\begin{lemma}
\label{lemma-cartesian-colimit-kappa}
Dans la situation ci-dessus, il existe un cardinal $\kappa$ ayant les
propriétés suivantes :
\begin{enumerate}
\item pour tout objet non nul $K$ de $\mathit{QC}(\mathcal{O})$, il existe un
morphisme non nul $E \to K$ dans $\mathit{QC}(\mathcal{O})$ tel que
$|E| \leq \kappa$ ;
\item pour tout morphisme $\alpha : E \to \bigoplus_n K_n$ dans
$\mathit{QC}(\mathcal{O})$ tel que $|E| \leq \kappa$, il existe des
morphismes $E_n \to K_n$ dans $\mathit{QC}(\mathcal{O})$, avec
$|E_n| \leq \kappa$, tels que $\alpha$ se factorise par
$\bigoplus E_n \to \bigoplus K_n$.
\end{enumerate}
\end{lemma}

\begin{proof}
Soit $\kappa$ un majorant de l'ensemble de cardinaux suivant :
\begin{enumerate}
\item $|\coprod_V j_{U!}\mathcal{O}_U(V)|$ pour tout
$U \in \Ob(\mathcal{C})$ ;
\item les cardinaux $\kappa(\mathcal{O}(V) \to \mathcal{O}(U))$ obtenus dans
More on Algebra, lemme \ref{more-algebra-lemma-enlarge}, pour tout morphisme
$U \to V$ de $\mathcal{C}$ ;
\item le cardinal obtenu dans le lemme \ref{lemma-qis-small}.
\end{enumerate}
Nous affirmons que, pour tout complexe $\mathcal{F}^\bullet$ représentant un
objet de $\mathit{QC}(\mathcal{O})$ et tout sous-complexe
$\mathcal{S}^\bullet \subset \mathcal{F}^\bullet$ vérifiant
$|\mathcal{S}^\bullet| \leq \kappa$, il existe un sous-complexe
$\mathcal{H}^\bullet$ de $\mathcal{F}^\bullet$, contenant
$\mathcal{S}^\bullet$, tel que $\mathcal{H}^\bullet$ représente un objet de
$\mathit{QC}(\mathcal{O})$ et satisfait
$|\mathcal{H}^\bullet| \leq \kappa$. Dans les deux paragraphes suivants,
nous montrons que cette affirmation entraîne le lemme.

\medskip\noindent
Pour établir (1), soit $K$ un objet non nul de
$\mathit{QC}(\mathcal{O})$. Supposons que $K$ soit représenté par le complexe
de $\mathcal{O}$-modules $\mathcal{F}^\bullet$. Alors
$H^i(\mathcal{F}^\bullet)$ est non nul pour un certain $i$. Il existe donc un
objet $U$ de $\mathcal{C}$ et une section
$s \in \mathcal{F}^i(U)$ telle que $d(s) = 0$ et qui détermine une section non
nulle de $H^i(\mathcal{F}^\bullet)$ sur $U$. L'image de
$s : j_{U!}\mathcal{O}_U[-i] \to \mathcal{F}^\bullet$ est alors un
sous-complexe $\mathcal{S}^\bullet \subset \mathcal{F}^\bullet$ tel que
$|\mathcal{S}^\bullet| \leq \kappa$. En appliquant l'affirmation, on obtient
un morphisme non nul $\mathcal{H}^\bullet \to \mathcal{F}^\bullet$ dans
$\mathit{QC}(\mathcal{O})$, avec
$|\mathcal{H}^\bullet| \leq \kappa$. Ceci établit (1).

\medskip\noindent
Soit $\alpha : E \to \bigoplus K_n$ comme dans (2). Choisissons des complexes
$\mathcal{K}_n^\bullet$ représentant les $K_n$. Alors
$\bigoplus \mathcal{K}_n^\bullet$ représente $\bigoplus K_n$. Par la
construction de la catégorie dérivée, on peut représenter $E$ par un complexe
$\mathcal{E}^\bullet$ de sorte que $\alpha$ soit représenté par un morphisme
de complexes
$a : \mathcal{E}^\bullet \to \bigoplus \mathcal{K}_n^\bullet$.
D'après le lemme \ref{lemma-qis-small} et le choix de $\kappa$ ci-dessus, on
peut supposer $|\mathcal{E}^\bullet| \leq \kappa$. L'affirmation fournit des
sous-complexes $\mathcal{E}_n^\bullet \subset \mathcal{K}_n^\bullet$
représentant des objets $E_n$ de $\mathit{QC}(\mathcal{O})$ tels que
$|E_n| \leq \kappa$ et contenant l'image de
$a_n : \mathcal{E}^\bullet \to \mathcal{K}_n^\bullet$, comme voulu.

\medskip\noindent
Démonstration de l'affirmation. Soit $\mathcal{F}^\bullet$ un complexe
représentant un objet de $\mathit{QC}(\mathcal{O})$ et soit
$\mathcal{S}^\bullet \subset \mathcal{F}^\bullet$ un sous-complexe de taille
$\leq \kappa$. Nous allons construire par récurrence des sous-complexes
$$
\mathcal{S}^\bullet = \mathcal{S}_0^\bullet \subset
\mathcal{S}_1^\bullet \subset
\mathcal{S}_2^\bullet \subset \ldots \subset \mathcal{F}^\bullet
$$
de taille $\leq \kappa$ tels que, pour tout morphisme $f : U \to V$ de
$\mathcal{C}$ et tout $i \in \mathbf{Z}$,
\begin{enumerate}
\item le noyau de la flèche
$H^i(\mathcal{S}_n^\bullet(V)
\otimes_{\mathcal{O}(V)}^\mathbf{L} \mathcal{O}(U)) \to
H^i(\mathcal{S}_n^\bullet(U))$
s'envoie sur zéro dans
$H^i(\mathcal{S}_{n + 1}^\bullet(V)
\otimes_{\mathcal{O}(V)}^\mathbf{L} \mathcal{O}(U))$,
\item l'image de la flèche
$H^i(\mathcal{S}_n^\bullet(U)) \to H^i(\mathcal{S}_{n + 1}^\bullet(U))$
est contenue dans l'image de
$H^i(\mathcal{S}_{n + 1}^\bullet(V)
\otimes_{\mathcal{O}(V)}^\mathbf{L} \mathcal{O}(U)) \to
H^i(\mathcal{S}_{n + 1}^\bullet(U))$,
\end{enumerate}
Une fois ceci fait, on pourra poser
$\mathcal{H}^\bullet = \bigcup \mathcal{S}_n^\bullet$. En effet, le produit
tensoriel dérivé et la prise de cohomologie des complexes de modules sur des
anneaux commutent aux colimites filtrantes ; les conditions (1) et (2)
garantiront donc ensemble que
$$
\mathcal{H}^\bullet(V) \otimes_{\mathcal{O}(V)}^\mathbf{L} \mathcal{O}(U)
\longrightarrow
\mathcal{H}^\bullet(U)
$$
est un isomorphisme en cohomologie à tout degré, donc un isomorphisme dans
$D(\mathcal{O}(U))$, pour tout $f : U \to V$ dans $\mathcal{C}$. Ainsi,
$\mathcal{H}^\bullet$ représente un objet de $\mathit{QC}(\mathcal{O})$,
comme voulu.

\medskip\noindent
Construisons $\mathcal{S}_{n + 1}$ à partir de $\mathcal{S}_n$. Pour tout
morphisme $f : U \to V$ de $\mathcal{C}$, considérons le diagramme commutatif
$$
\xymatrix{
\mathcal{S}_n^\bullet(V)
\ar[r] \ar[d] &
\mathcal{S}_n^\bullet(U) \ar[d] \\
\mathcal{F}^\bullet(V)
\ar[r] &
\mathcal{F}^\bullet(U)
}
$$
C'est un diagramme du type considéré dans More on Algebra, lemme
\ref{more-algebra-lemma-enlarge}, pour l'homomorphisme d'anneaux
$\mathcal{O}(V) \to \mathcal{O}(U)$ ; autrement dit, la ligne inférieure
induit un isomorphisme
$$
\mathcal{F}^\bullet(V) \otimes_{\mathcal{O}(V)}^\mathbf{L} \mathcal{O}(U)
\longrightarrow
\mathcal{F}^\bullet(U)
$$
dans $D(\mathcal{O}(U))$. On peut donc choisir des sous-complexes
$$
\mathcal{S}_n^\bullet(V) \subset M^\bullet_f \subset \mathcal{F}^\bullet(V)
\quad\text{et}\quad
\mathcal{S}_n^\bullet(U) \subset N^\bullet_f \subset \mathcal{F}^\bullet(U)
$$
comme dans More on Algebra, lemme \ref{more-algebra-lemma-enlarge} ; en
particulier, $|N^i_f|, |M^i_f| \leq \kappa$. Appliquons ensuite le lemme
\ref{lemma-build-from} aux sous-ensembles
$$
\mathcal{S}_n^i(U) \amalg \coprod\nolimits_{f : U \to V} N^i_f
\amalg \coprod\nolimits_{g : W \to U} M^i_g
\subset
\mathcal{F}^i(U)
$$
afin d'obtenir un sous-complexe
$$
\mathcal{S}_n^\bullet \subset
\mathcal{S}_{n + 1}^\bullet \subset
\mathcal{F}^\bullet
$$
qui contient ces sous-ensembles et vérifie
$|\mathcal{S}_{n + 1}^\bullet| \leq \kappa$. Les conditions (1) et (2) sont
satisfaites, car les assertions correspondantes valent pour
$\mathcal{S}_n^\bullet(V) \subset M^\bullet_f$ et
$\mathcal{S}_n^\bullet(U) \subset N^\bullet_f$ par la construction de More
on Algebra, lemme \ref{more-algebra-lemma-enlarge}. La démonstration est donc
achevée.
\end{proof}

\begin{proposition}
\label{proposition-cartesian-brown}
Soit $\mathcal{C}$ une catégorie considérée comme un site muni de la
topologie chaotique et soit $\mathcal{O}$ un faisceau d'anneaux sur
$\mathcal{C}$. Pour $\mathit{QC}(\mathcal{O})$ défini comme dans la définition
\ref{definition-cartesian}, on a :
\begin{enumerate}
\item $\mathit{QC}(\mathcal{O})$ est une sous-catégorie strictement pleine,
saturée et triangulée de $D(\mathcal{O})$, stable par sommes directes
arbitraires ;
\item tout foncteur cohomologique contravariant
$H : \mathit{QC}(\mathcal{O}) \to \textit{Ab}$
qui transforme les sommes directes en produits est représentable ;
\item tout foncteur exact $F : \mathit{QC}(\mathcal{O}) \to \mathcal{D}$ de
catégories triangulées qui transforme les sommes directes en sommes directes
admet un adjoint à droite exact ;
\item le foncteur d'inclusion
$\mathit{QC}(\mathcal{O}) \to D(\mathcal{O})$ admet un adjoint à droite exact.
\end{enumerate}
\end{proposition}

\begin{proof}
L'assertion (1) est le lemme \ref{lemma-cartesian-good-sub}. L'assertion (2)
résulte du lemme \ref{lemma-cartesian-colimit-kappa} et de Catégories dérivées,
lemme \ref{derived-lemma-brown-bis}. L'assertion (3) résulte du lemme
\ref{lemma-cartesian-colimit-kappa} et de Catégories dérivées, proposition
\ref{derived-proposition-brown-bis}. Enfin, (4) est un cas particulier de
(3).
\end{proof}

\noindent
Soit $u : \mathcal{C}' \to \mathcal{C}$ un foncteur entre catégories. Si
l'on considère $\mathcal{C}$ et $\mathcal{C}'$ comme des sites munis de la
topologie chaotique, alors $u$ est un foncteur continu et cocontinu. On obtient
donc un morphisme de topos
$g : \Sh(\mathcal{C}') \to \Sh(\mathcal{C})$, voir Sites, lemme
\ref{sites-lemma-cocontinuous-morphism-topoi}. Supposons en outre donnés des
faisceaux d'anneaux $\mathcal{O}$ sur $\mathcal{C}$ et $\mathcal{O}'$ sur
$\mathcal{C}'$, ainsi qu'un morphisme
$g^\sharp : g^{-1}\mathcal{O} \to \mathcal{O}'$. On note simplement le
morphisme correspondant de topos annelés par
$g : (\Sh(\mathcal{C}'), \mathcal{O}') \to (\Sh(\mathcal{C}), \mathcal{O})$,
voir Modules on Sites, section \ref{sites-modules-section-ringed-topoi}.

\begin{lemma}
\label{lemma-cartesian-functorial}
Soit
$g : (\Sh(\mathcal{C}'), \mathcal{O}') \to (\Sh(\mathcal{C}), \mathcal{O})$
comme ci-dessus. Alors le foncteur
$Lg^* : D(\mathcal{O}) \to D(\mathcal{O}')$ envoie
$\mathit{QC}(\mathcal{O})$ dans $\mathit{QC}(\mathcal{O}')$.
\end{lemma}

\begin{proof}
Soit $U' \in \Ob(\mathcal{C}')$, d'image $U = u(U')$ dans $\mathcal{C}$. Notons
$pt$ la catégorie ayant un seul objet et un seul morphisme, et notons
$(\Sh(pt), \mathcal{O}'(U'))$ et $(\Sh(pt), \mathcal{O}(U))$ les topos annelés
indiqués. On identifie bien entendu les catégories dérivées de modules sur
ces topos annelés à $D(\mathcal{O}'(U'))$ et $D(\mathcal{O}(U))$. On dispose
alors d'un diagramme commutatif de topos annelés
$$
\xymatrix{
(\Sh(pt), \mathcal{O}'(U')) \ar[rr]_{U'} \ar[d] & &
(\Sh(\mathcal{C}'), \mathcal{O}') \ar[d]^g \\
(\Sh(pt), \mathcal{O}(U)) \ar[rr]^U & &
(\Sh(\mathcal{C}), \mathcal{O})
}
$$
L'image inverse le long du morphisme horizontal inférieur envoie $K$ dans
$D(\mathcal{O})$ sur $R\Gamma(U, K)$. Celle le long de la flèche verticale de
gauche envoie $M$ sur
$M \otimes_{\mathcal{O}(U)}^\mathbf{L} \mathcal{O}'(U')$. Les deux parcours du
diagramme donnent le même résultat (lemme
\ref{lemma-derived-pullback-composition}) ; par conséquent,
$$
R\Gamma(U', Lg^*K) =
R\Gamma(U, K) \otimes_{\mathcal{O}(U)}^\mathbf{L} \mathcal{O}'(U')
$$
Enfin, soit $f' : U' \to V'$ un morphisme de $\mathcal{C}'$ et notons
$f = u(f') : U = u(U') \to V = u(V')$ son image dans $\mathcal{C}$. Si $K$ est
dans $\mathit{QC}(\mathcal{O})$, alors
\begin{align*}
R\Gamma(V', Lg^*K) \otimes_{\mathcal{O}'(V')}^\mathbf{L} \mathcal{O}'(U')
& =
R\Gamma(V, K) \otimes_{\mathcal{O}(V)}^\mathbf{L} \mathcal{O}'(V')
\otimes_{\mathcal{O}'(V')}^\mathbf{L} \mathcal{O}'(U') \\
& =
R\Gamma(V, K) \otimes_{\mathcal{O}(V)}^\mathbf{L} \mathcal{O}'(U') \\
& =
R\Gamma(V, K) \otimes_{\mathcal{O}(V)}^\mathbf{L} \mathcal{O}(U)
\otimes_{\mathcal{O}(U)}^\mathbf{L} \mathcal{O}'(U') \\
& =
R\Gamma(U, K) \otimes_{\mathcal{O}(U)}^\mathbf{L} \mathcal{O}'(U') \\
& =
R\Gamma(U', Lg^*K)
\end{align*}
comme voulu. Nous avons utilisé l'observation précédente à la fois pour $U'$
et pour $V'$.
\end{proof}

\begin{lemma}
\label{lemma-cartesian-flat-quasi-coherent}
Soit $\mathcal{C}$ une catégorie considérée comme un site muni de la
topologie chaotique et soit $\mathcal{O}$ un faisceau d'anneaux sur
$\mathcal{C}$. Supposons que, pour tout $U \to V$ dans $\mathcal{C}$,
l'homomorphisme de restriction
$\mathcal{O}(V) \to \mathcal{O}(U)$ soit plat. Alors
$\mathit{QC}(\mathcal{O})$ coïncide avec la sous-catégorie
$D_\QCoh(\mathcal{O}) \subset D(\mathcal{O})$ des complexes dont les
faisceaux de cohomologie sont quasi-cohérents.
\end{lemma}

\begin{proof}
Rappelons que, sous nos hypothèses,
$\QCoh(\mathcal{O}) \subset \textit{Mod}(\mathcal{O})$ est une sous-catégorie
de Serre faible, voir Modules on Sites, lemme
\ref{sites-modules-lemma-chaotic-flat-quasi-coherent}. Il est donc légitime de
considérer la sous-catégorie pleine
$$
D_\QCoh(\mathcal{O}) = D_{\QCoh(\mathcal{O})}(\textit{Mod}(\mathcal{O}))
$$
de $D(\mathcal{O})$, voir Catégories dérivées, section
\ref{derived-section-triangulated-sub}. (À strictement parler, ceci n'est pas
nécessaire dans la démonstration du lemme.)

\medskip\noindent
Soit $M$ un objet de $\mathit{QC}(\mathcal{O})$. Pour tout morphisme
$U \to V$ de $\mathcal{C}$, l'homomorphisme de restriction
$\mathcal{O}(V) \to \mathcal{O}(U)$ est plat ; on a donc
\begin{align*}
H^i(M)(U)
& =
H^i(R\Gamma(U, M)) \\
& =
H^i(R\Gamma(V, M) \otimes_{\mathcal{O}(V)}^\mathbf{L} \mathcal{O}(U)) \\
& =
H^i(R\Gamma(V, M)) \otimes_{\mathcal{O}(V)} \mathcal{O}(U) \\
& =
H^i(M)(V) \otimes_{\mathcal{O}(V)} \mathcal{O}(U)
\end{align*}
et $H^i(M)$ est donc quasi-cohérent d'après Modules on Sites, lemme
\ref{sites-modules-lemma-chaotic-quasi-coherent}. La première et la dernière
égalité ci-dessus résultent du fait que la prise des sections sur un objet de
$\mathcal{C}$ est un foncteur exact, puisque la topologie sur $\mathcal{C}$
est chaotique.

\medskip\noindent
Réciproquement, si $M$ est un objet de $D_\QCoh(\mathcal{O})$, alors Modules
on Sites, lemme \ref{sites-modules-lemma-chaotic-quasi-coherent}, montre que
l'application $R\Gamma(V, M) \to R\Gamma(U, M)$ induit des isomorphismes
$H^i(M)(V) \otimes_{\mathcal{O}(V)} \mathcal{O}(U) \to H^i(M)(U)$.
Par conséquent,
$R\Gamma(V, M) \otimes_{\mathcal{O}(V)}^\mathbf{L} \mathcal{O}(U)
\to R\Gamma(U, M)$ est un isomorphisme dans $D(\mathcal{O}(U))$, par platitude
de $\mathcal{O}(V) \to \mathcal{O}(U)$, et l'on conclut que $M$ appartient à
$\mathit{QC}(\mathcal{O})$.
\end{proof}

\begin{lemma}
\label{lemma-cartesion-plus-topology}
Soit $\epsilon : (\mathcal{C}_\tau, \mathcal{O}_\tau) \to
(\mathcal{C}_{\tau'}, \mathcal{O}_{\tau'})$ comme dans la section
\ref{section-compare}. Supposons que
\begin{enumerate}
\item $\tau'$ soit la topologie chaotique sur la catégorie $\mathcal{C}$ ;
\item pour tout $U \in \Ob(\mathcal{C})$ et tout complexe K-plat de
$\mathcal{O}(U)$-modules $M^\bullet$, l'application
$$
M^\bullet \longrightarrow
R\Gamma((\mathcal{C}/U)_\tau,
(M^\bullet \otimes_{\mathcal{O}(U)} \mathcal{O}_U)^\#)
$$
soit un quasi-isomorphisme (voir la démonstration pour une explication).
\end{enumerate}
Alors $\epsilon^*$ et $R\epsilon_*$ définissent des équivalences quasi-inverses
l'une de l'autre entre $\mathit{QC}(\mathcal{O})$ et la sous-catégorie pleine
de $D(\mathcal{C}_\tau, \mathcal{O}_\tau)$ formée des objets $K$ tels que
$R\epsilon_*K$ appartienne à $\mathit{QC}(\mathcal{O})$\footnote{Cela signifie
que
$R\Gamma(V, K) \otimes_{\mathcal{O}(V)}^\mathbf{L} \mathcal{O}(U)
\to R\Gamma(U, K)$ est un isomorphisme pour tout $U \to V$ dans
$\mathcal{C}$.}.
\end{lemma}

\begin{proof}
Nous utiliserons sans autre mention les observations de la section
\ref{section-compare}. Comme $R\epsilon_*$ est pleinement fidèle et
$\epsilon^* \circ R\epsilon_* = \text{id}$, il suffit, pour démontrer le
lemme, d'établir que, pour $M$ dans $\mathit{QC}(\mathcal{O})$, on a
$R\epsilon_*(\epsilon^*M) = M$. La
condition (2) est précisément celle qui permet de le voir. Choisissons en
effet un complexe K-plat $\mathcal{M}^\bullet$ de $\mathcal{O}$-modules, à
termes plats, qui représente $M$. Alors $\epsilon^*M$ est représenté par la
$\tau$-faisceautisation $(\mathcal{M}^\bullet)^\#$ de
$\mathcal{M}^\bullet$. Soit $U \in \Ob(\mathcal{C})$. Par Leray, on obtient
$$
R\Gamma(U, R\epsilon_*(\epsilon^*M)) =
R\Gamma((\mathcal{C}/U)_\tau, (\mathcal{M}^\bullet)^\#|_{\mathcal{C}/U}) =
R\Gamma((\mathcal{C}/U)_\tau, (\mathcal{M}^\bullet|_{\mathcal{C}/U})^\#)
$$
La dernière égalité résulte du fait que la faisceautisation commute à la
restriction à $\mathcal{C}/U$. Comme d'habitude, notons $\mathcal{O}_U$ la
restriction de $\mathcal{O}$ à $\mathcal{C}/U$. Considérons l'application
$$
\mathcal{M}^\bullet(U) \otimes_{\mathcal{O}(U)} \mathcal{O}_U
\longrightarrow
\mathcal{M}^\bullet|_{\mathcal{C}/U}
$$
de complexes de $\mathcal{O}_U$-modules (pour la topologie $\tau'$). Par le
choix de $\mathcal{M}^\bullet$, le complexe $\mathcal{M}^\bullet(U)$ est un
complexe K-plat de $\mathcal{O}(U)$-modules ; voir le lemme
\ref{lemma-pullback-K-flat} et utiliser que l'inclusion de $U$ dans
$\mathcal{C}$ définit un morphisme de topos annelés
$(\Sh(pt), \mathcal{O}(U)) \to (\Sh(\mathcal{C}_{\tau'}), \mathcal{O})$.
Comme $M$ appartient à $\mathit{QC}(\mathcal{O})$, la flèche affichée est un
quasi-isomorphisme. La faisceautisation étant exacte, elle reste un
quasi-isomorphisme après faisceautisation. Par conséquent,
$$
R\Gamma(U, R\epsilon_*(\epsilon^*M)) =
R\Gamma((\mathcal{C}/U)_\tau,
(\mathcal{M}^\bullet(U) \otimes_{\mathcal{O}(U)} \mathcal{O}_U)^\#)
$$
et l'hypothèse (2) montre que ceci est égal à
$R\Gamma(U, M) = \mathcal{M}^\bullet(U)$, comme voulu.
\end{proof}

\begin{lemma}
\label{lemma-QC-hom-out-of-plus-topology}
Conservons les notations et les hypothèses du lemme
\ref{lemma-cartesion-plus-topology}. Supposons que $K$ soit un objet de
$\mathit{QC}(\mathcal{O})$ et que $M$ soit un objet quelconque de
$D(\mathcal{O}_\tau)$. Pour tout objet $U$ de $\mathcal{C}$, on a
$$
\Hom_{D((\mathcal{O}_U)_\tau)}(\epsilon^*K|_U, M|_U) =
R\Hom_{\mathcal{O}(U)}(R\Gamma(U, K), R\Gamma(U, M))
$$
\end{lemma}

\begin{proof}
On a
$$
\Hom_{D((\mathcal{O}_U)_\tau)}(\epsilon^*K|_U, M|_U) =
\Hom_{D((\mathcal{O}_U)_{\tau'})}(K|_U, R\epsilon_*M|_U)
$$
par adjonction. Le résultat découle donc du lemme
\ref{lemma-QC-hom-out-of} et du fait que
$$
R\Gamma(U, M) = R\Gamma(U, R\epsilon_*M)
$$
par Leray.
\end{proof}







\section{Complexes strictement parfaits}
\label{section-strictly-perfect}

\noindent
Cette section est l'analogue de Cohomology, section
\ref{cohomology-section-strictly-perfect}.

\begin{definition}
\label{definition-strictly-perfect}
Soit $(\mathcal{C}, \mathcal{O})$ un site annelé et soit
$\mathcal{E}^\bullet$ un complexe de $\mathcal{O}$-modules. On dit que
$\mathcal{E}^\bullet$ est {\it strictement parfait} si
$\mathcal{E}^i$ est nul sauf pour un nombre fini de $i$ et si
$\mathcal{E}^i$ est facteur direct d'un $\mathcal{O}$-module libre de type fini
pour tout $i$.
\end{definition}

\noindent
Soit $U$ un objet de $\mathcal{C}$. Nous dirons souvent « Soit
$\mathcal{E}^\bullet$ un complexe strictement parfait de
$\mathcal{O}_U$-modules » pour signifier que $\mathcal{E}^\bullet$ est un
complexe strictement parfait de modules sur le site annelé
$(\mathcal{C}/U, \mathcal{O}_U)$, voir Modules on Sites, définition
\ref{sites-modules-definition-localize-ringed-site}.

\begin{lemma}
\label{lemma-cone}
Le cône d'un morphisme de complexes strictement parfaits est strictement
parfait.
\end{lemma}

\begin{proof}
C'est immédiat d'après les définitions.
\end{proof}

\begin{lemma}
\label{lemma-tensor}
Le complexe total associé au produit tensoriel de deux complexes strictement
parfaits est strictement parfait.
\end{lemma}

\begin{proof}
Omis.
\end{proof}

\begin{lemma}
\label{lemma-strictly-perfect-pullback}
Soit $(f, f^\sharp) : (\mathcal{C}, \mathcal{O}_\mathcal{C}) \to
(\mathcal{D}, \mathcal{O}_\mathcal{D})$
un morphisme de topos annelés. Si $\mathcal{F}^\bullet$ est un complexe
strictement parfait de $\mathcal{O}_\mathcal{D}$-modules, alors
$f^*\mathcal{F}^\bullet$ est un complexe strictement parfait de
$\mathcal{O}_\mathcal{C}$-modules.
\end{lemma}

\begin{proof}
Nous avons vu dans Modules on Sites, lemme
\ref{sites-modules-lemma-global-pullback}, que l'image inverse d'un module
libre de type fini est libre de type fini. Le foncteur $f^*$ est additif et
préserve donc les facteurs directs. Le lemme en résulte.
\end{proof}

\begin{lemma}
\label{lemma-local-lift-map}
Soit $(\mathcal{C}, \mathcal{O})$ un site annelé et soit $U$ un objet de
$\mathcal{C}$. Étant donné un diagramme de $\mathcal{O}_U$-modules dont les
flèches pleines sont
$$
\xymatrix{
\mathcal{E} \ar@{..>}[dr] \ar[r] & \mathcal{F} \\
& \mathcal{G} \ar[u]_p
}
$$
où $\mathcal{E}$ est facteur direct d'un $\mathcal{O}_U$-module libre de type
fini et $p$ est surjectif, il existe un recouvrement $\{U_i \to U\}$ tel que,
sur chaque $U_i$, on puisse tracer une flèche pointillée rendant le diagramme
commutatif.
\end{lemma}

\begin{proof}
On peut supposer $\mathcal{E} = \mathcal{O}_U^{\oplus n}$ pour un certain $n$.
Trouver la flèche pointillée revient alors à relever les images des éléments
de la base dans $\Gamma(U, \mathcal{F})$. Cela est possible localement d'après
la caractérisation des morphismes surjectifs de faisceaux (Sites, section
\ref{sites-section-sheaves-injective}).
\end{proof}

\begin{lemma}
\label{lemma-local-homotopy}
Soit $(\mathcal{C}, \mathcal{O})$ un site annelé et soit $U$ un objet de
$\mathcal{C}$.
\begin{enumerate}
\item Soit $\alpha : \mathcal{E}^\bullet \to \mathcal{F}^\bullet$ un
morphisme de complexes de $\mathcal{O}_U$-modules, avec
$\mathcal{E}^\bullet$ strictement parfait et $\mathcal{F}^\bullet$ acyclique.
Il existe un recouvrement $\{U_i \to U\}$ tel que chaque
$\alpha|_{U_i}$ soit homotope à zéro.
\item Soit $\alpha : \mathcal{E}^\bullet \to \mathcal{F}^\bullet$ un
morphisme de complexes de $\mathcal{O}_U$-modules, avec
$\mathcal{E}^\bullet$ strictement parfait, $\mathcal{E}^i = 0$ pour $i < a$, et
$H^i(\mathcal{F}^\bullet) = 0$ pour $i \geq a$. Il existe un recouvrement
$\{U_i \to U\}$ tel que chaque $\alpha|_{U_i}$ soit homotope à zéro.
\end{enumerate}
\end{lemma}

\begin{proof}
La première assertion résulte de la seconde ; il suffit donc de démontrer
(2). Procédons par récurrence sur la longueur du complexe
$\mathcal{E}^\bullet$. Si
$\mathcal{E}^\bullet \cong \mathcal{E}[-n]$, où $\mathcal{E}$ est facteur
direct d'un $\mathcal{O}$-module libre de type fini et $n \geq a$, alors le
résultat découle du lemme \ref{lemma-local-lift-map} et du fait que
$\mathcal{F}^{n - 1} \to \Ker(\mathcal{F}^n \to \mathcal{F}^{n + 1})$
est surjective par l'annulation supposée de
$H^n(\mathcal{F}^\bullet)$. Si $\mathcal{E}^i$ est nul sauf pour
$i \in [a, b]$, on dispose d'une suite exacte scindée de complexes
$$
0 \to \mathcal{E}^b[-b] \to \mathcal{E}^\bullet \to
\sigma_{\leq b - 1}\mathcal{E}^\bullet \to 0
$$
qui détermine un triangle distingué dans $K(\mathcal{O}_U)$, donc une suite
exacte
$$
\Hom_{K(\mathcal{O}_U)}(
\sigma_{\leq b - 1}\mathcal{E}^\bullet, \mathcal{F}^\bullet)
\to
\Hom_{K(\mathcal{O}_U)}(\mathcal{E}^\bullet, \mathcal{F}^\bullet)
\to
\Hom_{K(\mathcal{O}_U)}(\mathcal{E}^b[-b], \mathcal{F}^\bullet)
$$
d'après les axiomes des catégories triangulées. D'après ce qui précède, la
composée $\mathcal{E}^b[-b] \to \mathcal{F}^\bullet$ est homotope à zéro sur
les membres d'un recouvrement de $U$. On peut donc supposer que notre
morphisme provient d'un élément du membre de gauche de la suite exacte
affichée. Par l'hypothèse de récurrence, cet élément est nul sur les membres
d'un recouvrement de $U$.
\end{proof}

\begin{lemma}
\label{lemma-lift-through-quasi-isomorphism}
Soit $(\mathcal{C}, \mathcal{O})$ un site annelé et soit $U$ un objet de
$\mathcal{C}$. Étant donné un diagramme de complexes de
$\mathcal{O}_U$-modules dont les flèches pleines sont
$$
\xymatrix{
\mathcal{E}^\bullet \ar@{..>}[dr] \ar[r]_\alpha & \mathcal{F}^\bullet \\
& \mathcal{G}^\bullet \ar[u]_f
}
$$
où $\mathcal{E}^\bullet$ est strictement parfait,
$\mathcal{E}^j = 0$ pour $j < a$, et $H^j(f)$ est un isomorphisme pour $j > a$
et une surjection pour $j = a$, il existe un recouvrement $\{U_i \to U\}$ et,
pour
chaque $i$, une flèche pointillée sur $U_i$ qui rend le diagramme commutatif à
homotopie près.
\end{lemma}

\begin{proof}
Les hypothèses sur $f$ impliquent que les faisceaux de cohomologie du cône
$C(f)^\bullet$ s'annulent aux degrés $\geq a$. Le lemme
\ref{lemma-local-homotopy} fournit donc un recouvrement $\{U_i \to U\}$ tel que
la composée
$\mathcal{E}^\bullet \to \mathcal{F}^\bullet \to C(f)^\bullet$
soit homotope à zéro sur $U_i$. Comme
$$
\mathcal{G}^\bullet \to \mathcal{F}^\bullet \to C(f)^\bullet \to
\mathcal{G}^\bullet[1]
$$
se restreint en un triangle distingué dans $K(\mathcal{O}_{U_i})$, on peut
relever $\alpha|_{U_i}$, à homotopie près, en un morphisme
$\alpha_i : \mathcal{E}^\bullet|_{U_i} \to \mathcal{G}^\bullet|_{U_i}$
comme voulu.
\end{proof}

\begin{lemma}
\label{lemma-local-actual}
Soit $(\mathcal{C}, \mathcal{O})$ un site annelé, soit $U$ un objet de
$\mathcal{C}$, et soient $\mathcal{E}^\bullet$, $\mathcal{F}^\bullet$ des
complexes de $\mathcal{O}_U$-modules, avec $\mathcal{E}^\bullet$ strictement
parfait.
\begin{enumerate}
\item Pour tout élément
$\alpha \in \Hom_{D(\mathcal{O}_U)}(\mathcal{E}^\bullet, \mathcal{F}^\bullet)$
il existe un recouvrement $\{U_i \to U\}$ tel que $\alpha|_{U_i}$ soit donné
par un morphisme de complexes
$\alpha_i : \mathcal{E}^\bullet|_{U_i} \to \mathcal{F}^\bullet|_{U_i}$.
\item Soit un morphisme de complexes
$\alpha : \mathcal{E}^\bullet \to \mathcal{F}^\bullet$ dont l'image dans le
groupe
$\Hom_{D(\mathcal{O}_U)}(\mathcal{E}^\bullet, \mathcal{F}^\bullet)$
est nulle. Il existe un recouvrement $\{U_i \to U\}$ tel que
$\alpha|_{U_i}$ soit homotope à zéro.
\end{enumerate}
\end{lemma}

\begin{proof}
Démontrons (1). Par la construction de la catégorie dérivée, on peut trouver
un quasi-isomorphisme
$f : \mathcal{F}^\bullet \to \mathcal{G}^\bullet$ et un morphisme de
complexes $\beta : \mathcal{E}^\bullet \to \mathcal{G}^\bullet$ tels que
$\alpha = f^{-1}\beta$. Le résultat découle alors du lemme
\ref{lemma-lift-through-quasi-isomorphism}. Nous omettons la démonstration de
(2).
\end{proof}

\begin{lemma}
\label{lemma-Rhom-strictly-perfect}
Soit $(\mathcal{C}, \mathcal{O})$ un site annelé. Soient
$\mathcal{E}^\bullet$, $\mathcal{F}^\bullet$ des complexes de
$\mathcal{O}$-modules, avec $\mathcal{E}^\bullet$ strictement parfait. Alors le
Hom interne $R\SheafHom(\mathcal{E}^\bullet, \mathcal{F}^\bullet)$ est
représenté par le complexe $\mathcal{H}^\bullet$ dont les termes sont
$$
\mathcal{H}^n =
\bigoplus\nolimits_{n = p + q}
\SheafHom_\mathcal{O}(\mathcal{E}^{-q}, \mathcal{F}^p)
$$
et dont la différentielle est décrite dans la section
\ref{section-internal-hom}.
\end{lemma}

\begin{proof}
Choisissons un quasi-isomorphisme
$\mathcal{F}^\bullet \to \mathcal{I}^\bullet$ vers un complexe K-injectif.
Soit $(\mathcal{H}')^\bullet$ le complexe dont les termes sont
$$
(\mathcal{H}')^n =
\prod\nolimits_{n = p + q}
\SheafHom_\mathcal{O}(\mathcal{E}^{-q}, \mathcal{I}^p)
$$
Il représente $R\SheafHom(\mathcal{E}^\bullet, \mathcal{F}^\bullet)$ par la
construction de la section \ref{section-internal-hom}. Il suffit de montrer
que l'application
$$
\mathcal{H}^\bullet \longrightarrow (\mathcal{H}')^\bullet
$$
est un quasi-isomorphisme. Pour un objet $U$ de $\mathcal{C}$, on constate que
$$
H^0(\mathcal{H}^\bullet(U)) =
\Hom_{K(\mathcal{O}_U)}(\mathcal{E}^\bullet|_U, \mathcal{I}^\bullet|_U)
\to
H^0((\mathcal{H}')^\bullet(U)) =
\Hom_{D(\mathcal{O}_U)}(\mathcal{E}^\bullet|_U, \mathcal{I}^\bullet|_U)
$$
D'après le lemme \ref{lemma-local-actual}, la faisceautisation de
$U \mapsto H^0(\mathcal{H}^\bullet(U))$ est égale à celle de
$U \mapsto H^0((\mathcal{H}')^\bullet(U))$. Le même argument vaut pour les
autres faisceaux de cohomologie. Ainsi, $\mathcal{H}^\bullet$ est
quasi-isomorphe à $(\mathcal{H}')^\bullet$, ce qui démontre le lemme.
\end{proof}

\begin{lemma}
\label{lemma-Rhom-complex-of-direct-summands-finite-free}
Soit $(\mathcal{C}, \mathcal{O})$ un site annelé. Soient
$\mathcal{E}^\bullet$, $\mathcal{F}^\bullet$ des complexes de
$\mathcal{O}$-modules tels que
\begin{enumerate}
\item $\mathcal{F}^n = 0$ pour $n \ll 0$ ;
\item $\mathcal{E}^n = 0$ pour $n \gg 0$ ;
\item $\mathcal{E}^n$ soit isomorphe à un facteur direct d'un
$\mathcal{O}$-module libre de type fini.
\end{enumerate}
Alors le Hom interne
$R\SheafHom(\mathcal{E}^\bullet, \mathcal{F}^\bullet)$ est représenté par le
complexe $\mathcal{H}^\bullet$ dont les termes sont
$$
\mathcal{H}^n =
\bigoplus\nolimits_{n = p + q}
\SheafHom_\mathcal{O}(\mathcal{E}^{-q}, \mathcal{F}^p)
$$
et dont la différentielle est décrite dans la section
\ref{section-internal-hom}.
\end{lemma}

\begin{proof}
Choisissons un quasi-isomorphisme
$\mathcal{F}^\bullet \to \mathcal{I}^\bullet$, où $\mathcal{I}^\bullet$ est un
complexe d'injectifs borné inférieurement. Remarquons que
$\mathcal{I}^\bullet$ est K-injectif (Catégories dérivées, lemme
\ref{derived-lemma-bounded-below-injectives-K-injective}). La construction de
la section \ref{section-internal-hom} montre donc que
$R\SheafHom(\mathcal{E}^\bullet, \mathcal{F}^\bullet)$ est représenté par le
complexe $(\mathcal{H}')^\bullet$ dont les termes sont
$$
(\mathcal{H}')^n =
\prod\nolimits_{n = p + q}
\SheafHom_\mathcal{O}(\mathcal{E}^{-q}, \mathcal{I}^p) =
\bigoplus\nolimits_{n = p + q}
\SheafHom_\mathcal{O}(\mathcal{E}^{-q}, \mathcal{I}^p)
$$
Cette égalité vaut parce qu'il n'y a qu'un nombre fini de termes non nuls.
Remarquons que $\mathcal{H}^\bullet$ est le complexe total associé au complexe
double de termes
$\SheafHom_\mathcal{O}(\mathcal{E}^{-q}, \mathcal{F}^p)$
et de même pour $(\mathcal{H}')^\bullet$. L'application naturelle
$\mathcal{H}^\bullet \to (\mathcal{H}')^\bullet$ provient d'un morphisme de
complexes doubles. Pour montrer qu'elle est un quasi-isomorphisme, on peut
donc utiliser la suite spectrale d'un complexe double (Homology, lemme
\ref{homology-lemma-first-quadrant-ss})
$$
{}'E_1^{p, q} =
H^p(\SheafHom_\mathcal{O}(\mathcal{E}^{-q}, \mathcal{F}^\bullet))
$$
qui converge vers $H^{p + q}(\mathcal{H}^\bullet)$, et de même pour
$(\mathcal{H}')^\bullet$. Pour achever la démonstration, il suffit de montrer
que $\mathcal{F}^\bullet \to \mathcal{I}^\bullet$ induit un isomorphisme
$$
H^p(\SheafHom_\mathcal{O}(\mathcal{E}, \mathcal{F}^\bullet))
\longrightarrow
H^p(\SheafHom_\mathcal{O}(\mathcal{E}, \mathcal{I}^\bullet))
$$
sur les faisceaux de cohomologie lorsque $\mathcal{E}$ est facteur direct
d'un $\mathcal{O}$-module libre de type fini. C'est clair lorsque
$\mathcal{E}$ est libre de type fini, d'où le résultat.
\end{proof}









\section{Modules pseudo-cohérents}
\label{section-pseudo-coherent}

\noindent
Dans cette section, nous étudions les complexes pseudo-cohérents.

\begin{definition}
\label{definition-pseudo-coherent}
Soit $(\mathcal{C}, \mathcal{O})$ un site annelé, soit
$\mathcal{E}^\bullet$ un complexe de $\mathcal{O}$-modules et soit
$m \in \mathbf{Z}$.
\begin{enumerate}
\item On dit que $\mathcal{E}^\bullet$ est {\it $m$-pseudo-cohérent} si, pour
tout objet $U$ de $\mathcal{C}$, il existe un recouvrement $\{U_i \to U\}$ et,
pour chaque $i$, un morphisme de complexes
$\alpha_i : \mathcal{E}_i^\bullet \to \mathcal{E}^\bullet|_{U_i}$
où $\mathcal{E}_i^\bullet$ est un complexe strictement parfait de
$\mathcal{O}_{U_i}$-modules, $H^j(\alpha_i)$ est un isomorphisme pour $j > m$ et
$H^m(\alpha_i)$ est surjectif.
\item On dit que $\mathcal{E}^\bullet$ est {\it pseudo-cohérent} s'il est
$m$-pseudo-cohérent pour tout $m$.
\item On dit qu'un objet $E$ de $D(\mathcal{O})$ est
{\it $m$-pseudo-cohérent} (resp.\ {\it pseudo-cohérent}) si et seulement s'il
peut être représenté par un complexe $m$-pseudo-cohérent
(resp.\ pseudo-cohérent) de $\mathcal{O}$-modules.
\end{enumerate}
\end{definition}

\noindent
Si $\mathcal{C}$ possède un objet final quasi-compact $X$ (par exemple, si
tout recouvrement de $X$ admet un raffinement fini), alors tout objet
$m$-pseudo-cohérent de $D(\mathcal{O})$ appartient à $D^-(\mathcal{O})$.
Ce n'est toutefois pas nécessairement le cas en général.

\begin{lemma}
\label{lemma-pseudo-coherent-independent-representative}
Soit $(\mathcal{C}, \mathcal{O})$ un site annelé.
Soit $E$ un objet de $D(\mathcal{O})$.
\begin{enumerate}
\item Si $\mathcal{C}$ possède un objet final $X$ et s'il existe un recouvrement
$\{U_i \to X\}$, des complexes strictement parfaits $\mathcal{E}_i^\bullet$
de $\mathcal{O}_{U_i}$-modules et des morphismes
$\alpha_i : \mathcal{E}_i^\bullet \to E|_{U_i}$ dans
$D(\mathcal{O}_{U_i})$ tels que $H^j(\alpha_i)$ soit un isomorphisme pour
$j > m$ et $H^m(\alpha_i)$ soit surjectif, alors $E$ est
$m$-pseudo-cohérent.
\item Si $E$ est $m$-pseudo-cohérent, alors tout complexe de
$\mathcal{O}$-modules représentant $E$ est $m$-pseudo-cohérent.
\item Si, pour tout objet $U$ de $\mathcal{C}$, il existe un recouvrement
$\{U_i \to U\}$ tel que $E|_{U_i}$ soit $m$-pseudo-cohérent, alors
$E$ est $m$-pseudo-cohérent.
\end{enumerate}
\end{lemma}

\begin{proof}
Soit $\mathcal{F}^\bullet$ un complexe quelconque représentant $E$, et
soient $X$, $\{U_i \to X\}$ et
$\alpha_i : \mathcal{E}_i \to E|_{U_i}$ comme en (1). Nous allons montrer
que $\mathcal{F}^\bullet$ est $m$-pseudo-cohérent comme complexe, ce qui
prouvera (1) et (2) lorsque $\mathcal{C}$ possède un objet final. Par le lemme
\ref{lemma-local-actual}, après avoir raffiné le recouvrement
$\{U_i \to X\}$, nous pouvons représenter les morphismes $\alpha_i$ par des
morphismes de complexes
$\alpha_i : \mathcal{E}_i^\bullet \to \mathcal{F}^\bullet|_{U_i}$.
Par hypothèse, les $H^j(\alpha_i)$ sont des isomorphismes pour $j > m$ et
$H^m(\alpha_i)$ est surjectif ; ainsi, $\mathcal{F}^\bullet$ est
$m$-pseudo-cohérent.

\medskip\noindent
Démonstration de (2). D'après ce qui précède,
$\mathcal{F}^\bullet|_U$ est $m$-pseudo-cohérent comme complexe de
$\mathcal{O}_U$-modules pour tout objet $U$ de $\mathcal{C}$. Il résulte
formellement des définitions que $\mathcal{F}^\bullet$ est
$m$-pseudo-cohérent.

\medskip\noindent
Démonstration de (3). Cela résulte des définitions et de Sites,
définition \ref{sites-definition-site}, partie (2).
\end{proof}

\begin{lemma}
\label{lemma-pseudo-coherent-pullback}
Soit $(f, f^\sharp) : (\mathcal{C}, \mathcal{O}_\mathcal{C}) \to
(\mathcal{D}, \mathcal{O}_\mathcal{D})$
un morphisme de sites annelés. Soit $E$ un objet de
$D(\mathcal{O}_\mathcal{C})$. Si $E$ est $m$-pseudo-cohérent,
alors $Lf^*E$ est $m$-pseudo-cohérent.
\end{lemma}

\begin{proof}
Supposons $f$ donné par le foncteur $u : \mathcal{D} \to \mathcal{C}$.
Soit $U$ un objet de $\mathcal{C}$. Par Sites, lemme
\ref{sites-lemma-morphism-of-sites-covering}, nous pouvons trouver un
recouvrement $\{U_i \to U\}$ et, pour chaque $i$, un morphisme
$U_i \to u(V_i)$ pour un certain objet $V_i$ de $\mathcal{D}$.
Par le lemme \ref{lemma-pseudo-coherent-independent-representative},
il suffit de montrer que $Lf^*E|_{U_i}$ est $m$-pseudo-cohérent.
Pour cela, il suffit de montrer que $Lf^*E|_{u(V_i)}$ est
$m$-pseudo-cohérent, puisque $Lf^*E|_{U_i}$ est la restriction de
$Lf^*E|_{u(V_i)}$ à $\mathcal{C}/U_i$ (via Modules on Sites, lemme
\ref{sites-modules-lemma-relocalize}).
Par le diagramme commutatif de Modules on Sites, lemme
\ref{sites-modules-lemma-localize-morphism-ringed-sites}, il suffit de
démontrer le lemme pour le morphisme de sites annelés
$(\mathcal{C}/u(V_i), \mathcal{O}_{u(V_i)}) \to
(\mathcal{D}/V_i, \mathcal{O}_{V_i})$.
Nous pouvons donc supposer que $\mathcal{D}$ possède un objet final $Y$ tel
que $X = u(Y)$ soit un objet final de $\mathcal{C}$.

\medskip\noindent
Soit $\{V_i \to Y\}$ un recouvrement tel que, pour chaque $i$, il existe
un complexe strictement parfait $\mathcal{F}_i^\bullet$ de
$\mathcal{O}_{V_i}$-modules et un morphisme
$\alpha_i : \mathcal{F}_i^\bullet \to E|_{V_i}$ de $D(\mathcal{O}_{V_i})$
tel que $H^j(\alpha_i)$ soit un isomorphisme pour $j > m$ et
$H^m(\alpha_i)$ soit surjectif. En raisonnant comme ci-dessus, il suffit de
démontrer le résultat pour
$(\mathcal{C}/u(V_i), \mathcal{O}_{u(V_i)}) \to
(\mathcal{D}/V_i, \mathcal{O}_{V_i})$. Nous pouvons donc supposer qu'il
existe un complexe strictement parfait $\mathcal{F}^\bullet$ de
$\mathcal{O}_\mathcal{D}$-modules et un morphisme
$\alpha : \mathcal{F}^\bullet \to E$ de $D(\mathcal{O}_\mathcal{D})$ tel
que $H^j(\alpha)$ soit un isomorphisme pour $j > m$ et que
$H^m(\alpha)$ soit surjectif. Dans ce cas, choisissons un triangle distingué
$$
\mathcal{F}^\bullet \to E \to C \to \mathcal{F}^\bullet[1]
$$
L'hypothèse sur $\alpha$ signifie exactement que les faisceaux de cohomologie
$H^j(C)$ sont nuls pour tout $j \geq m$. En appliquant $Lf^*$, nous obtenons
le triangle distingué
$$
Lf^*\mathcal{F}^\bullet \to Lf^*E \to Lf^*C \to Lf^*\mathcal{F}^\bullet[1]
$$
Par la construction de $Lf^*$ comme foncteur dérivé à gauche, on voit que
$H^j(Lf^*C) = 0$ pour $j \geq m$ (par l'énoncé dual de Catégories dérivées,
lemme \ref{derived-lemma-negative-vanishing}). Ainsi $H^j(Lf^*\alpha)$ est
un isomorphisme pour $j > m$ et $H^m(Lf^*\alpha)$ est surjectif.
D'autre part, puisque $\mathcal{F}^\bullet$ est un complexe borné supérieurement
de $\mathcal{O}_\mathcal{D}$-modules plats, on voit que
$Lf^*\mathcal{F}^\bullet = f^*\mathcal{F}^\bullet$.
On conclut en appliquant le lemme \ref{lemma-strictly-perfect-pullback}.
\end{proof}

\begin{lemma}
\label{lemma-cone-pseudo-coherent}
Soit $(\mathcal{C}, \mathcal{O})$ un site annelé et $m \in \mathbf{Z}$.
Soit $(K, L, M, f, g, h)$ un triangle distingué dans $D(\mathcal{O})$.
\begin{enumerate}
\item Si $K$ est $(m + 1)$-pseudo-cohérent et $L$ est $m$-pseudo-cohérent,
alors $M$ est $m$-pseudo-cohérent.
\item Si $K$ et $M$ sont $m$-pseudo-cohérents, alors $L$ est
$m$-pseudo-cohérent.
\item Si $L$ est $(m + 1)$-pseudo-cohérent et $M$ est
$m$-pseudo-cohérent, alors $K$ est $(m + 1)$-pseudo-cohérent.
\end{enumerate}
\end{lemma}

\begin{proof}
Démonstration de (1). Soit $U$ un objet de $\mathcal{C}$. Choisissons un
recouvrement $\{U_i \to U\}$ et des morphismes
$\alpha_i : \mathcal{K}_i^\bullet \to K|_{U_i}$ dans
$D(\mathcal{O}_{U_i})$, où $\mathcal{K}_i^\bullet$ est strictement parfait,
tels que les $H^j(\alpha_i)$ soient des isomorphismes pour $j > m + 1$ et
soient surjectifs pour $j = m + 1$. Nous pouvons remplacer
$\mathcal{K}_i^\bullet$ par
$\sigma_{\geq m + 1}\mathcal{K}_i^\bullet$ et donc supposer que
$\mathcal{K}_i^j = 0$ pour $j < m + 1$. Après avoir raffiné le recouvrement,
nous pouvons choisir des morphismes
$\beta_i : \mathcal{L}_i^\bullet \to L|_{U_i}$ dans
$D(\mathcal{O}_{U_i})$, où $\mathcal{L}_i^\bullet$ est strictement parfait,
tels que $H^j(\beta_i)$ soit un isomorphisme pour $j > m$ et surjectif pour
$j = m$. Par le lemme \ref{lemma-lift-through-quasi-isomorphism}, après avoir
raffiné le recouvrement, nous pouvons trouver des morphismes de complexes
$\gamma_i : \mathcal{K}_i^\bullet \to \mathcal{L}_i^\bullet$ tels que les
diagrammes
$$
\xymatrix{
K|_{U_i} \ar[r] & L|_{U_i} \\
\mathcal{K}_i^\bullet \ar[u]^{\alpha_i} \ar[r]^{\gamma_i} &
\mathcal{L}_i^\bullet \ar[u]_{\beta_i}
}
$$
sont commutatifs dans $D(\mathcal{O}_{U_i})$ (cela exige de représenter les
morphismes $\alpha_i$, $\beta_i$ et $K|_{U_i} \to L|_{U_i}$ par de véritables
morphismes de complexes ; certains détails sont omis). Le cône
$C(\gamma_i)^\bullet$ est strictement parfait (lemme \ref{lemma-cone}).
La commutativité du diagramme implique qu'il existe un morphisme de triangles
distingués
$$
(\mathcal{K}_i^\bullet, \mathcal{L}_i^\bullet, C(\gamma_i)^\bullet)
\longrightarrow
(K|_{U_i}, L|_{U_i}, M|_{U_i}).
$$
Il résulte du morphisme induit sur les suites exactes longues de cohomologie
et de Homology, lemmes \ref{homology-lemma-four-lemma} et
\ref{homology-lemma-five-lemma}, que
$C(\gamma_i)^\bullet \to M|_{U_i}$ induit un isomorphisme en cohomologie aux
degrés $> m$ et une surjection au degré $m$. Ainsi, $M$ est
$m$-pseudo-cohérent par le lemme
\ref{lemma-pseudo-coherent-independent-representative}.

\medskip\noindent
Les assertions (2) et (3) résultent de (1) en faisant tourner le triangle
distingué.
\end{proof}

\begin{lemma}
\label{lemma-tensor-pseudo-coherent}
Soit $(\mathcal{C}, \mathcal{O})$ un site annelé. Soient $K, L$ des objets
de $D(\mathcal{O})$.
\begin{enumerate}
\item Si $K$ est $n$-pseudo-cohérent et $H^i(K) = 0$ pour $i > a$, et si
$L$ est $m$-pseudo-cohérent et $H^j(L) = 0$ pour $j > b$, alors
$K \otimes_\mathcal{O}^\mathbf{L} L$ est $t$-pseudo-cohérent, où
$t = \max(m + a, n + b)$.
\item Si $K$ et $L$ sont pseudo-cohérents, alors
$K \otimes_\mathcal{O}^\mathbf{L} L$ est pseudo-cohérent.
\end{enumerate}
\end{lemma}

\begin{proof}
Démonstration de (1). Soit $U$ un objet de $\mathcal{C}$. En remplaçant $U$
par les membres d'un recouvrement et $\mathcal{C}$ par la localisation
$\mathcal{C}/U$, nous pouvons supposer qu'il existe des complexes strictement
parfaits $\mathcal{K}^\bullet$ et $\mathcal{L}^\bullet$, ainsi que des
morphismes $\alpha : \mathcal{K}^\bullet \to K$ et
$\beta : \mathcal{L}^\bullet \to L$, tels que $H^i(\alpha)$ soit un
isomorphisme pour $i > n$ et surjectif pour $i = n$, et que $H^i(\beta)$ soit
un isomorphisme pour $i > m$ et surjectif pour $i = m$. Alors le morphisme
$$
\alpha \otimes^\mathbf{L} \beta :
\text{Tot}(\mathcal{K}^\bullet \otimes_\mathcal{O} \mathcal{L}^\bullet)
\to K \otimes_\mathcal{O}^\mathbf{L} L
$$
induit des isomorphismes sur les faisceaux de cohomologie au degré $i$ pour
$i > t$ et une surjection pour $i = t$. Cela résulte de la suite spectrale
des Tor (détails omis).

\medskip\noindent
Démonstration de (2). Soit $U$ un objet de $\mathcal{C}$. Nous pouvons d'abord
remplacer $U$ par les membres d'un recouvrement et $\mathcal{C}$ par la
localisation $\mathcal{C}/U$ afin de nous ramener au cas où $K$ et $L$ sont
bornés supérieurement. L'assertion résulte alors immédiatement du cas (1).
\end{proof}

\begin{lemma}
\label{lemma-summands-pseudo-coherent}
Soit $(\mathcal{C}, \mathcal{O})$ un site annelé. Soit $m \in \mathbf{Z}$.
Si $K \oplus L$ est $m$-pseudo-cohérent (resp.\ pseudo-cohérent) dans
$D(\mathcal{O})$, il en est de même de $K$ et de $L$.
\end{lemma}

\begin{proof}
Supposons que $K \oplus L$ soit $m$-pseudo-cohérent. Soit $U$ un objet de
$\mathcal{C}$. Après avoir remplacé $U$ par les membres d'un recouvrement,
nous pouvons supposer que $K \oplus L \in D^-(\mathcal{O}_U)$, donc que
$L \in D^-(\mathcal{O}_U)$. Notons qu'il existe un triangle distingué
$$
(K \oplus L, K \oplus L, L \oplus L[1]) =
(K, K, 0) \oplus (L, L, L \oplus L[1])
$$
voir Catégories dérivées, lemme
\ref{derived-lemma-direct-sum-triangles}. Par le lemme
\ref{lemma-cone-pseudo-coherent}, on voit que $L \oplus L[1]$ est
$m$-pseudo-cohérent. Par conséquent, $L[1] \oplus L[2]$ est également
$m$-pseudo-cohérent. Par récurrence, $L[n] \oplus L[n + 1]$ est
$m$-pseudo-cohérent. Puisque $L$ est borné supérieurement, on voit que
$L[n]$ est $m$-pseudo-cohérent pour $n$ assez grand. En remontant alors et en
utilisant les triangles distingués
$$
(L[n], L[n] \oplus L[n - 1], L[n - 1])
$$
nous concluons que $L[n - 1], L[n - 2], \ldots, L$ sont
$m$-pseudo-cohérents, comme souhaité.
\end{proof}

\begin{lemma}
\label{lemma-finite-cohomology}
Soit $(\mathcal{C}, \mathcal{O})$ un site annelé. Soit $K$ un objet de
$D(\mathcal{O})$. Soit $m \in \mathbf{Z}$.
\begin{enumerate}
\item Si $K$ est $m$-pseudo-cohérent et $H^i(K) = 0$ pour $i > m$, alors
$H^m(K)$ est un $\mathcal{O}$-module de type fini.
\item Si $K$ est $m$-pseudo-cohérent et $H^i(K) = 0$ pour $i > m + 1$, alors
$H^{m + 1}(K)$ est un $\mathcal{O}$-module de présentation finie.
\end{enumerate}
\end{lemma}

\begin{proof}
Démonstration de (1). Soit $U$ un objet de $\mathcal{C}$. Nous devons montrer
que $H^m(K)$ peut être engendré par un nombre fini de sections sur les membres
d'un recouvrement de $U$ (voir Modules on Sites, définition
\ref{sites-modules-definition-site-local}). Ainsi, au cours de la
démonstration, nous pouvons (un nombre fini de fois) choisir un recouvrement
$\{U_i \to U\}$, remplacer $\mathcal{C}$ par $\mathcal{C}/U_i$ et $U$ par
$U_i$. En particulier, d'après nos définitions, nous pouvons supposer qu'il
existe un complexe strictement parfait $\mathcal{E}^\bullet$ et un morphisme
$\alpha : \mathcal{E}^\bullet \to K$ qui induit un isomorphisme en cohomologie
aux degrés $> m$ et une surjection au degré $m$. Il suffit de démontrer le
résultat pour $\mathcal{E}^\bullet$. Soit $n$ le plus grand entier tel que
$\mathcal{E}^n \not = 0$. Si $n = m$, alors
$H^m(\mathcal{E}^\bullet)$ est un quotient de $\mathcal{E}^n$ et le résultat
est clair. Si $n > m$, alors $\mathcal{E}^{n - 1} \to \mathcal{E}^n$ est
surjectif puisque $H^n(\mathcal{E}^\bullet) = 0$. Par le lemme
\ref{lemma-local-lift-map}, nous pouvons (après avoir remplacé $U$ par les
membres d'un recouvrement) trouver une section de cette surjection et écrire
$\mathcal{E}^{n - 1} = \mathcal{E}' \oplus \mathcal{E}^n$. Il suffit donc de
démontrer le résultat pour le complexe $(\mathcal{E}')^\bullet$, qui est le
même que $\mathcal{E}^\bullet$, sauf qu'il a $\mathcal{E}'$ au degré $n - 1$
et $0$ au degré $n$. On conclut par récurrence sur $n$.

\medskip\noindent
Démonstration de (2). Choisissons un objet $U$ de $\mathcal{C}$. Comme dans
la démonstration de (1), nous pouvons travailler localement sur $U$. Nous
pouvons donc supposer qu'il existe un complexe strictement parfait
$\mathcal{E}^\bullet$ et un morphisme
$\alpha : \mathcal{E}^\bullet \to K$ qui induit un isomorphisme en cohomologie
aux degrés $> m$ et une surjection au degré $m$. Comme dans la démonstration
de (1), nous pouvons nous ramener au cas où $\mathcal{E}^i = 0$ pour
$i > m + 1$. Alors on voit que
$H^{m + 1}(K) \cong H^{m + 1}(\mathcal{E}^\bullet) =
\Coker(\mathcal{E}^m \to \mathcal{E}^{m + 1})$
est de présentation finie.
\end{proof}








\section{Dimension de Tor}
\label{section-tor}

\noindent
Dans cette section, nous étudions plus en détail les résolutions par des
modules plats.

\begin{definition}
\label{definition-tor-amplitude}
Soit $(\mathcal{C}, \mathcal{O})$ un site annelé.
Soit $E$ un objet de $D(\mathcal{O})$.
Soient $a, b \in \mathbf{Z}$ avec $a \leq b$.
\begin{enumerate}
\item On dit que $E$ a une {\it amplitude de Tor dans $[a, b]$} si
$H^i(E \otimes_\mathcal{O}^\mathbf{L} \mathcal{F}) = 0$ pour tout
$\mathcal{O}$-module $\mathcal{F}$ et tout $i \not \in [a, b]$.
\item On dit que $E$ a une {\it dimension de Tor finie} s'il a une amplitude
de Tor dans $[a, b]$ pour certains $a, b$.
\item On dit que $E$ {\it est localement de dimension de Tor finie} si, pour
tout objet $U$ de $\mathcal{C}$, il existe un recouvrement $\{U_i \to U\}$
tel que $E|_{U_i}$ soit de dimension de Tor finie pour tout $i$.
\end{enumerate}
Un $\mathcal{O}$-module $\mathcal{F}$ a une {\it dimension de Tor $\leq d$}
si $\mathcal{F}[0]$, considéré comme un objet de $D(\mathcal{O})$, a une
amplitude de Tor dans $[-d, 0]$.
\end{definition}

\noindent
Notons que si $E$, comme dans la définition, est de dimension de Tor finie,
alors $E$ est un objet de $D^b(\mathcal{O})$, comme on le voit en prenant
$\mathcal{F} = \mathcal{O}$ dans la définition ci-dessus.

\begin{lemma}
\label{lemma-last-one-flat}
Soit $(\mathcal{C}, \mathcal{O})$ un site annelé. Soit
$\mathcal{E}^\bullet$ un complexe borné supérieurement de
$\mathcal{O}$-modules plats, d'amplitude de Tor dans $[a, b]$. Alors
$\Coker(d_{\mathcal{E}^\bullet}^{a - 1})$ est un $\mathcal{O}$-module plat.
\end{lemma}

\begin{proof}
Comme $\mathcal{E}^\bullet$ est un complexe borné supérieurement de modules
plats, on voit que
$\mathcal{E}^\bullet \otimes_\mathcal{O} \mathcal{F} =
\mathcal{E}^\bullet \otimes_\mathcal{O}^{\mathbf{L}} \mathcal{F}$
pour tout $\mathcal{O}$-module $\mathcal{F}$. Par conséquent, pour tout
$\mathcal{O}$-module $\mathcal{F}$, la suite
$$
\mathcal{E}^{a - 2} \otimes_\mathcal{O} \mathcal{F} \to
\mathcal{E}^{a - 1} \otimes_\mathcal{O} \mathcal{F} \to
\mathcal{E}^a \otimes_\mathcal{O} \mathcal{F}
$$
est exacte au milieu. Puisque
$\mathcal{E}^{a - 2} \to \mathcal{E}^{a - 1} \to \mathcal{E}^a \to
\Coker(d^{a - 1}) \to 0$
est une résolution plate, cela implique que
$\text{Tor}_1^\mathcal{O}(\Coker(d^{a - 1}), \mathcal{F}) = 0$
pour tout $\mathcal{O}$-module $\mathcal{F}$. Cela signifie que
$\Coker(d^{a - 1})$ est plat, voir le lemme \ref{lemma-flat-tor-zero}.
\end{proof}

\begin{lemma}
\label{lemma-tor-amplitude}
Soit $(\mathcal{C}, \mathcal{O})$ un site annelé. Soit $E$ un objet de
$D(\mathcal{O})$. Soient $a, b \in \mathbf{Z}$ avec $a \leq b$. Les
conditions suivantes sont équivalentes :
\begin{enumerate}
\item $E$ a une amplitude de Tor dans $[a, b]$.
\item $E$ est représenté par un complexe $\mathcal{E}^\bullet$ de
$\mathcal{O}$-modules plats tel que $\mathcal{E}^i = 0$ pour
$i \not \in [a, b]$.
\end{enumerate}
\end{lemma}

\begin{proof}
Si (2) est satisfaite, alors nous pouvons calculer
$E \otimes_\mathcal{O}^\mathbf{L} \mathcal{F} =
\mathcal{E}^\bullet \otimes_\mathcal{O} \mathcal{F}$
et il est clair que (1) est satisfaite.

\medskip\noindent
Supposons (1) satisfaite. Nous pouvons représenter $E$ par un complexe borné
supérieurement de $\mathcal{O}$-modules plats, noté $\mathcal{K}^\bullet$, voir
la section \ref{section-flat}. Soit $n$ le plus grand entier tel que
$\mathcal{K}^n \not = 0$. Si $n > b$, alors
$\mathcal{K}^{n - 1} \to \mathcal{K}^n$ est surjectif puisque
$H^n(\mathcal{K}^\bullet) = 0$. Comme $\mathcal{K}^n$ est plat, on voit que
$\Ker(\mathcal{K}^{n - 1} \to \mathcal{K}^n)$ est plat (Modules on Sites,
lemme \ref{sites-modules-lemma-flat-ses}). Nous pouvons donc remplacer
$\mathcal{K}^\bullet$ par $\tau_{\leq n - 1}\mathcal{K}^\bullet$. Ainsi, par
récurrence sur $n$, nous nous ramenons au cas où $\mathcal{K}^\bullet$ est un complexe
de $\mathcal{O}$-modules plats tel que $\mathcal{K}^i = 0$ pour $i > b$.

\medskip\noindent
Posons $\mathcal{E}^\bullet = \tau_{\geq a}\mathcal{K}^\bullet$.
Tout est clair, sauf le fait que $\mathcal{E}^a$ soit plat, lequel résulte
immédiatement du lemme \ref{lemma-last-one-flat} et des définitions.
\end{proof}

\begin{lemma}
\label{lemma-bounded-below-tor-amplitude}
Soit $(\mathcal{C}, \mathcal{O})$ un site annelé. Soit $E$ un objet de
$D(\mathcal{O})$. Soit $a \in \mathbf{Z}$. Les conditions suivantes sont
équivalentes :
\begin{enumerate}
\item $E$ a une amplitude de Tor dans $[a, \infty]$.
\item $E$ peut être représenté par un complexe K-plat
$\mathcal{E}^\bullet$ de $\mathcal{O}$-modules plats tel que
$\mathcal{E}^i = 0$ pour $i \not \in [a, \infty]$.
\end{enumerate}
De plus, nous pouvons choisir $\mathcal{E}^\bullet$ de sorte que toute image
inverse par un morphisme de sites annelés soit un complexe K-plat à termes
plats.
\end{lemma}

\begin{proof}
L'implication (2) $\Rightarrow$ (1) est immédiate. Supposons (1) satisfaite.
Choisissons d'abord un complexe K-plat $\mathcal{K}^\bullet$ à termes plats
représentant $E$, voir le lemme \ref{lemma-K-flat-resolution}. Pour tout
$\mathcal{O}$-module $\mathcal{M}$, la cohomologie de
$$
\mathcal{K}^{n - 1} \otimes_\mathcal{O} \mathcal{M} \to
\mathcal{K}^n \otimes_\mathcal{O} \mathcal{M} \to
\mathcal{K}^{n + 1} \otimes_\mathcal{O} \mathcal{M}
$$
calcule $H^n(E \otimes_\mathcal{O}^\mathbf{L} \mathcal{M})$. Celui-ci est
toujours nul pour $n < a$. Par conséquent, si nous appliquons le lemme
\ref{lemma-last-one-flat} au complexe
$\ldots \to \mathcal{K}^{a - 1} \to \mathcal{K}^a \to \mathcal{K}^{a + 1}$
nous concluons que
$\mathcal{N} = \Coker(\mathcal{K}^{a - 1} \to \mathcal{K}^a)$ est un
$\mathcal{O}$-module plat. Posons
$$
\mathcal{E}^\bullet = \tau_{\geq a}\mathcal{K}^\bullet =
(\ldots \to 0 \to \mathcal{N} \to \mathcal{K}^{a + 1} \to \ldots )
$$
Le noyau $\mathcal{L}^\bullet$ de
$\mathcal{K}^\bullet \to \mathcal{E}^\bullet$ est le complexe
$$
\mathcal{L}^\bullet = (\ldots \to \mathcal{K}^{a - 1} \to
\mathcal{I} \to 0 \to \ldots)
$$
où $\mathcal{I} \subset \mathcal{K}^a$ est l'image de
$\mathcal{K}^{a - 1} \to \mathcal{K}^a$. Puisque nous avons la suite exacte
courte
$0 \to \mathcal{I} \to \mathcal{K}^a \to \mathcal{N} \to 0$
on voit que $\mathcal{I}$ est un $\mathcal{O}$-module plat. Ainsi,
$\mathcal{L}^\bullet$ est un complexe borné supérieurement de modules plats,
donc K-plat par le lemme \ref{lemma-bounded-flat-K-flat}. Il s'ensuit que
$\mathcal{E}^\bullet$ est K-plat par le lemme
\ref{lemma-K-flat-two-out-of-three-ses}.

\medskip\noindent
Démonstration de l'assertion finale. Soit
$f : (\mathcal{C}', \mathcal{O}') \to (\mathcal{C}, \mathcal{O})$
un morphisme de sites annelés. Par le lemme \ref{lemma-pullback-K-flat}, le
complexe $f^*\mathcal{K}^\bullet$ est K-plat à termes plats. Le complexe
$f^*\mathcal{L}^\bullet$ est K-plat, car c'est un complexe borné
supérieurement de $\mathcal{O}'$-modules plats. Nous avons une suite exacte
courte de complexes de $\mathcal{O}'$-modules
$$
0 \to f^*\mathcal{L}^\bullet \to f^*\mathcal{K}^\bullet \to
f^*\mathcal{E}^\bullet \to 0
$$
car la suite exacte courte
$0 \to \mathcal{I} \to \mathcal{K}^a \to \mathcal{N} \to 0$
de modules plats se tire en arrière en une suite exacte courte. Par le lemme
\ref{lemma-K-flat-two-out-of-three-ses}, le complexe
$f^*\mathcal{E}^\bullet$ est K-plat, ce qui achève la démonstration.
\end{proof}

\begin{lemma}
\label{lemma-tor-amplitude-pullback}
Soit $(f, f^\sharp) : (\mathcal{C}, \mathcal{O}_\mathcal{C}) \to
(\mathcal{D}, \mathcal{O}_\mathcal{D})$
un morphisme de sites annelés. Soit $E$ un objet de
$D(\mathcal{O}_\mathcal{D})$. Si $E$ a une amplitude de Tor dans $[a, b]$,
alors $Lf^*E$ a une amplitude de Tor dans $[a, b]$.
\end{lemma}

\begin{proof}
Supposons que $E$ ait une amplitude de Tor dans $[a, b]$. Par le lemme
\ref{lemma-tor-amplitude}, nous pouvons représenter $E$ par un complexe
$\mathcal{E}^\bullet$ de $\mathcal{O}$-modules plats tel que
$\mathcal{E}^i = 0$ pour $i \not \in [a, b]$. Alors $Lf^*E$ est représenté
par $f^*\mathcal{E}^\bullet$. Par Modules on Sites, lemme
\ref{sites-modules-lemma-pullback-flat}, les modules
$f^*\mathcal{E}^i$ sont plats. Par le lemme \ref{lemma-tor-amplitude}, nous
concluons donc que $Lf^*E$ a une amplitude de Tor dans $[a, b]$.
\end{proof}

\begin{lemma}
\label{lemma-cone-tor-amplitude}
Soit $(\mathcal{C}, \mathcal{O})$ un site annelé. Soit
$(K, L, M, f, g, h)$ un triangle distingué dans $D(\mathcal{O})$.
Soient $a, b \in \mathbf{Z}$.
\begin{enumerate}
\item Si $K$ a une amplitude de Tor dans $[a + 1, b + 1]$ et $L$ une
amplitude de Tor dans $[a, b]$, alors $M$ a une amplitude de Tor dans
$[a, b]$.
\item Si $K$ et $M$ ont une amplitude de Tor dans $[a, b]$, alors $L$ a une
amplitude de Tor dans $[a, b]$.
\item Si $L$ a une amplitude de Tor dans $[a + 1, b + 1]$ et $M$ une
amplitude de Tor dans $[a, b]$, alors $K$ a une amplitude de Tor dans
$[a + 1, b + 1]$.
\end{enumerate}
\end{lemma}

\begin{proof}
Démonstration omise. Indication : cela résulte simplement de la suite exacte
longue de cohomologie associée à un triangle distingué et du fait que
$- \otimes_\mathcal{O}^{\mathbf{L}} \mathcal{F}$
préserve les triangles distingués. L'assertion la plus facile à démontrer est
(2), et les autres s'en déduisent par translation.
\end{proof}

\begin{lemma}
\label{lemma-tensor-tor-amplitude}
Soit $(\mathcal{C}, \mathcal{O})$ un site annelé. Soient $K, L$ des objets
de $D(\mathcal{O})$. Si $K$ a une amplitude de Tor dans $[a, b]$ et $L$ une
amplitude de Tor dans $[c, d]$, alors $K \otimes_\mathcal{O}^\mathbf{L} L$
a une amplitude de Tor dans $[a + c, b + d]$.
\end{lemma}

\begin{proof}
Démonstration omise. Indication : utiliser la suite spectrale des Tor.
\end{proof}

\begin{lemma}
\label{lemma-summands-tor-amplitude}
Soit $(\mathcal{C}, \mathcal{O})$ un site annelé. Soient
$a, b \in \mathbf{Z}$. Si $K$, $L$ sont des objets de $D(\mathcal{O})$ et si
$K \oplus L$ a une amplitude de Tor dans $[a, b]$, il en est de même de $K$
et de $L$.
\end{lemma}

\begin{proof}
C'est clair puisque les foncteurs Tor sont additifs.
\end{proof}

\begin{lemma}
\label{lemma-bounded}
Soit $(\mathcal{C}, \mathcal{O})$ un site annelé. Soit
$\mathcal{I} \subset \mathcal{O}$ un faisceau d'idéaux. Soit $K$ un objet
de $D(\mathcal{O})$.
\begin{enumerate}
\item Si $K \otimes_\mathcal{O}^\mathbf{L} \mathcal{O}/\mathcal{I}$ est
borné supérieurement, alors
$K \otimes_\mathcal{O}^\mathbf{L} \mathcal{O}/\mathcal{I}^n$
est uniformément borné supérieurement pour tout $n$.
\item Si $K \otimes_\mathcal{O}^\mathbf{L} \mathcal{O}/\mathcal{I}$,
considéré comme un objet de $D(\mathcal{O}/\mathcal{I})$, a une amplitude de
Tor dans $[a, b]$, alors
$K \otimes_\mathcal{O}^\mathbf{L} \mathcal{O}/\mathcal{I}^n$, considéré
comme un objet de $D(\mathcal{O}/\mathcal{I}^n)$, a une amplitude de Tor dans
$[a, b]$ pour tout $n$.
\end{enumerate}
\end{lemma}

\begin{proof}
Démonstration de (1). Supposons que
$K \otimes_\mathcal{O}^\mathbf{L} \mathcal{O}/\mathcal{I}$
soit borné supérieurement, disons que
$H^i(K \otimes_\mathcal{O}^\mathbf{L} \mathcal{O}/\mathcal{I}) = 0$
pour $i > b$. Notons que nous avons les triangles distingués
$$
K \otimes_\mathcal{O}^\mathbf{L}
\mathcal{I}^n/\mathcal{I}^{n + 1} \to
K \otimes_\mathcal{O}^\mathbf{L}
\mathcal{O}/\mathcal{I}^{n + 1} \to
K \otimes_\mathcal{O}^\mathbf{L}
\mathcal{O}/\mathcal{I}^n \to
K \otimes_\mathcal{O}^\mathbf{L}
\mathcal{I}^n/\mathcal{I}^{n + 1}[1]
$$
et que
$$
K \otimes_\mathcal{O}^\mathbf{L}
\mathcal{I}^n/\mathcal{I}^{n + 1} =
\left(
K \otimes_\mathcal{O}^\mathbf{L}
\mathcal{O}/\mathcal{I}\right)
\otimes_{\mathcal{O}/\mathcal{I}}^\mathbf{L}
\mathcal{I}^n/\mathcal{I}^{n + 1}
$$
Par récurrence, nous concluons que
$H^i(K \otimes_\mathcal{O}^\mathbf{L} \mathcal{O}/\mathcal{I}^n) = 0$
pour $i > b$ et pour tout $n$.

\medskip\noindent
Démonstration de (2). Supposons que
$K \otimes_\mathcal{O}^\mathbf{L} \mathcal{O}/\mathcal{I}$, considéré comme
un objet de $D(\mathcal{O}/\mathcal{I})$, ait une amplitude de Tor dans
$[a, b]$. Soit $\mathcal{F}$ un faisceau de
$\mathcal{O}/\mathcal{I}^n$-modules. Nous avons alors une filtration finie
$$
0 \subset \mathcal{I}^{n - 1}\mathcal{F} \subset \ldots
\subset \mathcal{I}\mathcal{F} \subset \mathcal{F}
$$
dont les quotients successifs sont des faisceaux de
$\mathcal{O}/\mathcal{I}$-modules. Ainsi, pour démontrer que
$K \otimes_\mathcal{O}^\mathbf{L} \mathcal{O}/\mathcal{I}^n$ a une amplitude
de Tor dans $[a, b]$, il suffit de montrer que
$H^i(K \otimes_\mathcal{O}^\mathbf{L} \mathcal{O}/\mathcal{I}^n
\otimes_{\mathcal{O}/\mathcal{I}^n}^\mathbf{L} \mathcal{G})$
est nul pour $i \not \in [a, b]$ et pour tout
$\mathcal{O}/\mathcal{I}$-module $\mathcal{G}$. Puisque
$$
\left(K \otimes_\mathcal{O}^\mathbf{L} \mathcal{O}/\mathcal{I}^n\right)
\otimes_{\mathcal{O}/\mathcal{I}^n}^\mathbf{L} \mathcal{G}
=
\left(K \otimes_\mathcal{O}^\mathbf{L} \mathcal{O}/\mathcal{I}\right)
\otimes_{\mathcal{O}/\mathcal{I}}^\mathbf{L} \mathcal{G}
$$
pour tout faisceau de $\mathcal{O}/\mathcal{I}$-modules $\mathcal{G}$, le
résultat s'ensuit.
\end{proof}

\begin{lemma}
\label{lemma-tor-amplitude-stalk}
Soit $(\mathcal{C}, \mathcal{O})$ un site annelé. Soit $E$ un objet de
$D(\mathcal{O})$. Soient $a, b \in \mathbf{Z}$.
\begin{enumerate}
\item Si $E$ a une amplitude de Tor dans $[a, b]$, alors, pour tout point
$p$ du site $\mathcal{C}$, l'objet $E_p$ de $D(\mathcal{O}_p)$ a une amplitude
de Tor dans $[a, b]$.
\item Si $\mathcal{C}$ possède suffisamment de points, la réciproque est
vraie.
\end{enumerate}
\end{lemma}

\begin{proof}
Démonstration de (1). Cela résulte du fait que prendre les germes en $p$
revient à tirer en arrière par le morphisme de sites annelés
$(p, \mathcal{O}_p) \to (\mathcal{C}, \mathcal{O})$ ; nous pouvons donc
appliquer le lemme \ref{lemma-tor-amplitude-pullback}.

\medskip\noindent
Démonstration de (2). Si $\mathcal{C}$ possède suffisamment de points, alors
nous pouvons vérifier l'annulation de
$H^i(E \otimes_\mathcal{O}^\mathbf{L} \mathcal{F})$
sur les germes, voir Modules on Sites, lemme
\ref{sites-modules-lemma-check-exactness-stalks}. Puisque
$H^i(E \otimes_\mathcal{O}^\mathbf{L} \mathcal{F})_p =
H^i(E_p \otimes_{\mathcal{O}_p}^\mathbf{L} \mathcal{F}_p)$, nous concluons.
\end{proof}










\section{Complexes parfaits}
\label{section-perfect}

\noindent
Dans cette section, nous étudions les propriétés des complexes parfaits sur
les sites annelés.

\begin{definition}
\label{definition-perfect}
Soit $(\mathcal{C}, \mathcal{O})$ un site annelé. Soit
$\mathcal{E}^\bullet$ un complexe de $\mathcal{O}$-modules. On dit que
$\mathcal{E}^\bullet$ est {\it parfait} si, pour tout objet $U$ de
$\mathcal{C}$, il existe un recouvrement $\{U_i \to U\}$ tel que, pour tout
$i$, il existe un morphisme de complexes
$\mathcal{E}_i^\bullet \to \mathcal{E}^\bullet|_{U_i}$
qui soit un quasi-isomorphisme, où $\mathcal{E}_i^\bullet$ est strictement
parfait. Un objet $E$ de $D(\mathcal{O})$ est {\it parfait} s'il peut être
représenté par un complexe parfait de $\mathcal{O}$-modules.
\end{definition}

\noindent
Si $\Sh(\mathcal{C})$ est quasi-compact (Sites, section
\ref{sites-section-quasi-compact}), alors tout objet parfait de
$D(\mathcal{O})$ appartient à $D^b(\mathcal{O})$. Ce n'est toutefois pas
nécessairement le cas sinon.

\begin{lemma}
\label{lemma-perfect-independent-representative}
Soit $(\mathcal{C}, \mathcal{O})$ un site annelé.
Soit $E$ un objet de $D(\mathcal{O})$.
\begin{enumerate}
\item Si $\mathcal{C}$ possède un objet final $X$ et s'il existe un
recouvrement $\{U_i \to X\}$, des complexes strictement parfaits
$\mathcal{E}_i^\bullet$ de $\mathcal{O}_{U_i}$-modules et des isomorphismes
 $\alpha_i : \mathcal{E}_i^\bullet \to E|_{U_i}$ dans
$D(\mathcal{O}_{U_i})$, alors $E$ est parfait.
\item Si $E$ est parfait, alors tout complexe représentant $E$ est parfait.
\end{enumerate}
\end{lemma}

\begin{proof}
Identique à la démonstration du lemme
\ref{lemma-pseudo-coherent-independent-representative}.
\end{proof}

\begin{lemma}
\label{lemma-perfect-precise}
Soit $(\mathcal{C}, \mathcal{O})$ un site annelé. Soit $E$ un objet de
$D(\mathcal{O})$. Soient $a \leq b$ des entiers. Si $E$ a une amplitude de
Tor dans $[a, b]$ et est $(a - 1)$-pseudo-cohérent, alors $E$ est parfait.
\end{lemma}

\begin{proof}
Soit $U$ un objet de $\mathcal{C}$. Après avoir remplacé $U$ par les membres
d'un recouvrement et $\mathcal{C}$ par la localisation $\mathcal{C}/U$, nous
pouvons supposer qu'il existe un complexe strictement parfait
$\mathcal{E}^\bullet$ et un morphisme
$\alpha : \mathcal{E}^\bullet \to E$ tel que $H^i(\alpha)$ soit un
isomorphisme pour $i \geq a$. Nous pouvons remplacer, et remplaçons,
$\mathcal{E}^\bullet$ par $\sigma_{\geq a - 1}\mathcal{E}^\bullet$.
Choisissons un triangle distingué
$$
\mathcal{E}^\bullet \to E \to C \to \mathcal{E}^\bullet[1]
$$
De l'annulation des faisceaux de cohomologie de $E$ et de
$\mathcal{E}^\bullet$, ainsi que de l'hypothèse sur $\alpha$, nous obtenons
$C \cong \mathcal{K}[2 - a]$, où
$\mathcal{K} = \Ker(\mathcal{E}^{a - 1} \to \mathcal{E}^a)$. Soit
$\mathcal{F}$ un $\mathcal{O}$-module. En appliquant
$- \otimes_\mathcal{O}^\mathbf{L} \mathcal{F}$, l'hypothèse selon laquelle
$E$ a une amplitude de Tor dans $[a, b]$ implique que le morphisme
$\mathcal{K} \otimes_\mathcal{O} \mathcal{F} \to
\mathcal{E}^{a - 1} \otimes_\mathcal{O} \mathcal{F}$ a pour image
$\Ker(\mathcal{E}^{a - 1} \otimes_\mathcal{O} \mathcal{F}
\to \mathcal{E}^a \otimes_\mathcal{O} \mathcal{F})$.
Il s'ensuit que
$\text{Tor}_1^\mathcal{O}(\mathcal{E}', \mathcal{F}) = 0$, où
$\mathcal{E}' = \Coker(\mathcal{E}^{a - 1} \to \mathcal{E}^a)$. Ainsi,
$\mathcal{E}'$ est plat (lemme \ref{lemma-flat-tor-zero}). Par Modules on
Sites, lemme \ref{sites-modules-lemma-flat-locally-finite-presentation}, il
existe donc un recouvrement $\{U_i \to U\}$ tel que
$\mathcal{E}'|_{U_i}$ soit facteur direct d'un module libre de type fini.
Ainsi, le complexe
$$
\mathcal{E}'|_{U_i} \to \mathcal{E}^{a + 1}|_{U_i} \to \ldots \to
\mathcal{E}^b|_{U_i}
$$
est quasi-isomorphe à $E|_{U_i}$ et $E$ est parfait.
\end{proof}

\begin{lemma}
\label{lemma-perfect}
Soit $(\mathcal{C}, \mathcal{O})$ un site annelé. Soit $E$ un objet de
$D(\mathcal{O})$. Les conditions suivantes sont équivalentes :
\begin{enumerate}
\item $E$ est parfait, et
\item $E$ est pseudo-cohérent et est localement de dimension de Tor finie.
\end{enumerate}
\end{lemma}

\begin{proof}
Supposons (1). Soit $U$ un objet de $\mathcal{C}$. Par définition, il existe
un recouvrement $\{U_i \to U\}$ tel que $E|_{U_i}$ soit représenté par un
complexe strictement parfait. Ainsi, $E$ est pseudo-cohérent (c'est-à-dire
$m$-pseudo-cohérent pour tout $m$) par le lemme
\ref{lemma-pseudo-coherent-independent-representative}. De plus, un facteur
direct d'un module libre de type fini est plat ; $E|_{U_i}$ est donc de
dimension de Tor finie par le lemme \ref{lemma-tor-amplitude}. Ainsi, (2) est
satisfaite.

\medskip\noindent
Supposons (2). Soit $U$ un objet de $\mathcal{C}$. Après avoir remplacé $U$
par les membres d'un recouvrement, nous pouvons supposer qu'il existe des
entiers $a \leq b$ tels que $E|_U$ ait une amplitude de Tor dans $[a, b]$.
Puisque $E|_U$ est $m$-pseudo-cohérent pour tout $m$, nous concluons à l'aide
du lemme \ref{lemma-perfect-precise}.
\end{proof}

\begin{lemma}
\label{lemma-perfect-pullback}
Soit $(f, f^\sharp) : (\mathcal{C}, \mathcal{O}_\mathcal{C}) \to
(\mathcal{D}, \mathcal{O}_\mathcal{D})$
un morphisme de sites annelés. Soit $E$ un objet de
$D(\mathcal{O}_\mathcal{D})$. Si $E$ est parfait dans
$D(\mathcal{O}_\mathcal{D})$, alors $Lf^*E$ est parfait dans
$D(\mathcal{O}_\mathcal{C})$.
\end{lemma}

\begin{proof}
Cela résulte des lemmes \ref{lemma-perfect},
\ref{lemma-tor-amplitude-pullback} et
\ref{lemma-pseudo-coherent-pullback}.
\end{proof}

\begin{lemma}
\label{lemma-two-out-of-three-perfect}
Soit $(\mathcal{C}, \mathcal{O})$ un site annelé. Soit $(K, L, M, f, g, h)$
un triangle distingué dans $D(\mathcal{O})$. Si deux des trois objets
$K, L, M$ sont parfaits, alors le troisième l'est également.
\end{lemma}

\begin{proof}
Première démonstration : combiner les lemmes \ref{lemma-perfect},
\ref{lemma-cone-pseudo-coherent} et \ref{lemma-cone-tor-amplitude}.
Deuxième démonstration (esquisse) : supposons $K$ et $L$ parfaits. Soit $U$
un objet de $\mathcal{C}$. Après avoir remplacé $U$ par les membres d'un
recouvrement, nous pouvons supposer que $K|_U$ et $L|_U$ sont représentés
par des complexes strictement parfaits $\mathcal{K}^\bullet$ et
$\mathcal{L}^\bullet$. Après avoir à nouveau remplacé $U$ par les membres
d'un recouvrement, nous pouvons supposer que le morphisme
$K|_U \to L|_U$ est donné par un morphisme de complexes
$\alpha : \mathcal{K}^\bullet \to \mathcal{L}^\bullet$, voir le lemme
\ref{lemma-local-actual}. Alors $M|_U$ est isomorphe au cône de $\alpha$, qui
est strictement parfait par le lemme \ref{lemma-cone}.
\end{proof}

\begin{lemma}
\label{lemma-tensor-perfect}
Soit $(\mathcal{C}, \mathcal{O})$ un site annelé. Si $K, L$ sont des objets
parfaits de $D(\mathcal{O})$, alors
$K \otimes_\mathcal{O}^\mathbf{L} L$ l'est également.
\end{lemma}

\begin{proof}
Cela résulte des lemmes \ref{lemma-perfect},
\ref{lemma-tensor-pseudo-coherent} et
\ref{lemma-tensor-tor-amplitude}.
\end{proof}

\begin{lemma}
\label{lemma-summands-perfect}
Soit $(\mathcal{C}, \mathcal{O})$ un site annelé. Si $K \oplus L$ est un
objet parfait de $D(\mathcal{O})$, alors $K$ et $L$ le sont également.
\end{lemma}

\begin{proof}
Cela résulte des lemmes \ref{lemma-perfect},
\ref{lemma-summands-pseudo-coherent} et
\ref{lemma-summands-tor-amplitude}.
\end{proof}














\section{Duaux}
\label{section-duals}

\noindent
Dans cette section, nous caractérisons les objets dualisables de la catégorie
des complexes et de la catégorie dérivée. En particulier, nous verrons qu'un
objet de $D(\mathcal{O})$ possède un dual si et seulement s'il est parfait
(cela résulte de l'exemple \ref{example-dual-derived} et du lemme
\ref{lemma-left-dual-derived}).

\begin{lemma}
\label{lemma-symmetric-monoidal-cat-complexes}
Soit $(\mathcal{C}, \mathcal{O})$ un site annelé. La catégorie des complexes
de $\mathcal{O}$-modules, munie du produit tensoriel défini par
$\mathcal{F}^\bullet \otimes \mathcal{G}^\bullet =
\text{Tot}(\mathcal{F}^\bullet \otimes_\mathcal{O} \mathcal{G}^\bullet)$
est une catégorie monoïdale symétrique.
\end{lemma}

\begin{proof}
Démonstration omise. Indications : comme unité $\mathbf{1}$, on prend le
complexe ayant $\mathcal{O}$ au degré $0$ et zéro aux autres degrés, avec les
isomorphismes évidents
$\text{Tot}(\mathbf{1} \otimes_\mathcal{O} \mathcal{G}^\bullet) =
\mathcal{G}^\bullet$ et
$\text{Tot}(\mathcal{F}^\bullet \otimes_\mathcal{O} \mathbf{1}) =
\mathcal{F}^\bullet$. Pour démontrer le lemme, il faut vérifier la
commutativité de divers diagrammes, voir Categories, définitions
\ref{categories-definition-monoidal-category} et
\ref{categories-definition-symmetric-monoidal-category}.
Dans chaque cas, les vérifications sont immédiates.
\end{proof}

\begin{example}
\label{example-dual}
Soit $(\mathcal{C}, \mathcal{O})$ un site annelé. Soit
$\mathcal{F}^\bullet$ un complexe de $\mathcal{O}$-modules tel que, pour tout
$U \in \Ob(\mathcal{C})$, il existe un recouvrement $\{U_i \to U\}$ tel que
$\mathcal{F}^\bullet|_{U_i}$ soit strictement parfait. Considérons le complexe
$$
\mathcal{G}^\bullet = \SheafHom^\bullet(\mathcal{F}^\bullet, \mathcal{O})
$$
comme dans la section \ref{section-hom-complexes}. Soient
$$
\eta :
\mathcal{O}
\to
\text{Tot}(\mathcal{F}^\bullet \otimes_\mathcal{O} \mathcal{G}^\bullet)
\quad\text{et}\quad
\epsilon :
\text{Tot}(\mathcal{G}^\bullet \otimes_\mathcal{O} \mathcal{F}^\bullet)
\to
\mathcal{O}
$$
définis par $\eta = \sum \eta_n$ et $\epsilon = \sum \epsilon_n$, où
$\eta_n : \mathcal{O} \to
\mathcal{F}^n \otimes_\mathcal{O} \mathcal{G}^{-n}$
et
$\epsilon_n : \mathcal{G}^{-n} \otimes_\mathcal{O} \mathcal{F}^n
\to \mathcal{O}$ sont comme dans Modules on Sites, exemple
\ref{sites-modules-example-dual}. Alors
$\mathcal{G}^\bullet, \eta, \epsilon$ est un dual à gauche de
$\mathcal{F}^\bullet$ au sens de Categories, définition
\ref{categories-definition-dual}. Nous omettons la vérification de
$(1 \otimes \epsilon) \circ (\eta \otimes 1) = \text{id}_{\mathcal{F}^\bullet}$
et de
$(\epsilon \otimes 1) \circ (1 \otimes \eta) =
\text{id}_{\mathcal{G}^\bullet}$. Comparer avec More on Algebra, lemme
\ref{more-algebra-lemma-left-dual-complex}.
\end{example}

\begin{lemma}
\label{lemma-left-dual-complex}
Soit $(\mathcal{C}, \mathcal{O})$ un site annelé. Soit
$\mathcal{F}^\bullet$ un complexe de $\mathcal{O}$-modules. Si
$\mathcal{F}^\bullet$ possède un dual à gauche dans la catégorie monoïdale
des complexes de $\mathcal{O}$-modules (Categories, définition
\ref{categories-definition-dual}), alors, pour tout objet $U$ de
$\mathcal{C}$, il existe un recouvrement $\{U_i \to U\}$ tel que
$\mathcal{F}^\bullet|_{U_i}$ soit strictement parfait et que le dual à gauche
soit celui construit dans l'exemple \ref{example-dual}.
\end{lemma}

\begin{proof}
Par unicité des duaux à gauche (Categories, remarque
\ref{categories-remark-left-dual-adjoint}), nous obtenons l'assertion finale
dès que nous avons montré que $\mathcal{F}^\bullet$ est comme annoncé. Soit
$\mathcal{G}^\bullet, \eta, \epsilon$ un dual à gauche. Écrivons
$\eta = \sum \eta_n$ et $\epsilon = \sum \epsilon_n$, où
$\eta_n : \mathcal{O} \to
\mathcal{F}^n \otimes_\mathcal{O} \mathcal{G}^{-n}$
et
$\epsilon_n : \mathcal{G}^{-n} \otimes_\mathcal{O} \mathcal{F}^n
\to \mathcal{O}$. Puisque
$(1 \otimes \epsilon) \circ (\eta \otimes 1) = \text{id}_{\mathcal{F}^\bullet}$
et
$(\epsilon \otimes 1) \circ (1 \otimes \eta) = \text{id}_{\mathcal{G}^\bullet}$
par Categories, définition \ref{categories-definition-dual}, on voit
immédiatement que
$(1 \otimes \epsilon_n) \circ (\eta_n \otimes 1) = \text{id}_{\mathcal{F}^n}$
et
$(\epsilon_n \otimes 1) \circ (1 \otimes \eta_n) =
\text{id}_{\mathcal{G}^{-n}}$.
Autrement dit, $\mathcal{G}^{-n}$ est un dual à gauche de
$\mathcal{F}^n$, et Modules on Sites, lemme
\ref{sites-modules-lemma-left-dual-module}, s'applique à chaque
$\mathcal{F}^n$. Soit $U$ un objet de $\mathcal{C}$. Il existe un recouvrement
$\{U_i \to U\}$ tel que, pour tout $i$, seul un nombre fini des
$\eta_n|_{U_i}$ soient non nuls. Ainsi, après avoir remplacé $U$ par $U_i$,
nous pouvons supposer que seul un nombre fini des $\eta_n|_U$ sont non nuls ;
par le lemme cité, cela implique que seul un nombre fini des
$\mathcal{F}^n|_U$ sont non nuls. En utilisant à nouveau le lemme, nous
pouvons alors trouver un recouvrement $\{U_i \to U\}$ tel que chaque
$\mathcal{F}^n|_{U_i}$ soit facteur direct d'un $\mathcal{O}$-module libre de
type fini, ce qui achève la démonstration.
\end{proof}

\begin{lemma}
\label{lemma-dual-perfect-complex}
Soit $(\mathcal{C}, \mathcal{O})$ un site annelé. Soit $K$ un objet parfait
de $D(\mathcal{O})$. Alors $K^\vee = R\SheafHom(K, \mathcal{O})$ est également
un objet parfait et $(K^\vee)^\vee \cong K$. Il existe des isomorphismes
fonctoriels
$$
M \otimes^\mathbf{L}_\mathcal{O} K^\vee = R\SheafHom_\mathcal{O}(K, M)
$$
et
$$
H^0(\mathcal{C}, M \otimes^\mathbf{L}_\mathcal{O} K^\vee) =
\Hom_{D(\mathcal{O})}(K, M)
$$
pour $M$ dans $D(\mathcal{O})$.
\end{lemma}

\begin{proof}
Nous utiliserons sans autre mention le fait que la formation du Hom interne
commute à la restriction (lemme \ref{lemma-restriction-RHom-to-U}). Soit $U$
un objet quelconque de $\mathcal{C}$. Pour vérifier que $K^\vee$ est parfait,
il suffit de montrer qu'il existe un recouvrement $\{U_i \to U\}$ tel que
$K^\vee|_{U_i}$ soit parfait pour tout $i$. Il existe un morphisme canonique
$$
K = R\SheafHom(\mathcal{O}, \mathcal{O})
\otimes_{\mathcal{O}}^\mathbf{L} K \longrightarrow
R\SheafHom(R\SheafHom(K, \mathcal{O}), \mathcal{O}) =
(K^\vee)^\vee
$$
voir le lemme \ref{lemma-internal-hom-evaluate}. Il suffit de démontrer qu'il
existe un recouvrement $\{U_i \to U\}$ tel que la restriction de ce morphisme
à $\mathcal{C}/U_i$ soit un isomorphisme pour tout $i$. Par le lemme
\ref{lemma-dual}, pour obtenir l'assertion finale, il suffit de vérifier que
le morphisme (\ref{equation-eval})
$$
M \otimes^\mathbf{L}_\mathcal{O} K^\vee
\longrightarrow
R\SheafHom(K, M)
$$
est un isomorphisme. C'est également une question locale (au sens ci-dessus).
Il suffit donc de démontrer le lemme lorsque $K$ est représenté par un
complexe strictement parfait.

\medskip\noindent
Supposons $K$ représenté par le complexe strictement parfait
$\mathcal{E}^\bullet$. Il résulte alors du lemme
\ref{lemma-Rhom-strictly-perfect} que $K^\vee$ est représenté par le complexe
dont les termes sont
$(\mathcal{E}^n)^\vee =
\SheafHom_\mathcal{O}(\mathcal{E}^n, \mathcal{O})$
au degré $-n$. Puisque $\mathcal{E}^n$ est facteur direct d'un
$\mathcal{O}$-module libre de type fini, il en est de même de
$(\mathcal{E}^n)^\vee$. Ainsi, $K^\vee$ est lui aussi représenté par un
complexe strictement parfait, et l'on voit que $K^\vee$ est parfait. Le
morphisme $K \to (K^\vee)^\vee$ est un isomorphisme, car il est donné, au
signe près, par les morphismes d'évaluation
$\mathcal{E}^n \to ((\mathcal{E}^n)^\vee)^\vee$, qui sont des isomorphismes.
Pour voir que (\ref{equation-eval}) est un isomorphisme, représentons $M$ par
un complexe K-plat $\mathcal{F}^\bullet$. Par le lemme
\ref{lemma-Rhom-strictly-perfect}, le complexe $R\SheafHom(K, M)$ est
représenté par le complexe de termes
$$
\bigoplus\nolimits_{n = p + q}
\SheafHom_\mathcal{O}(\mathcal{E}^{-q}, \mathcal{F}^p)
$$
D'autre part, l'objet $M \otimes^\mathbf{L}_\mathcal{O} K^\vee$ est
représenté par le complexe de termes
$$
\bigoplus\nolimits_{n = p + q}
\mathcal{F}^p \otimes_\mathcal{O} (\mathcal{E}^{-q})^\vee
$$
Ainsi, l'assertion selon laquelle (\ref{equation-eval}) est un isomorphisme se
ramène à l'assertion selon laquelle le morphisme canonique
$$
\mathcal{F}
\otimes_\mathcal{O}
\SheafHom_\mathcal{O}(\mathcal{E}, \mathcal{O})
\longrightarrow
\SheafHom_\mathcal{O}(\mathcal{E}, \mathcal{F})
$$
est un isomorphisme lorsque $\mathcal{E}$ est facteur direct d'un
$\mathcal{O}$-module libre de type fini et $\mathcal{F}$ est un
$\mathcal{O}$-module quelconque. Cela résulte immédiatement de l'assertion
correspondante lorsque $\mathcal{E}$ est libre de type fini.
\end{proof}

\begin{lemma}
\label{lemma-symmetric-monoidal-derived}
Soit $(\mathcal{C}, \mathcal{O})$ un site annelé. La catégorie dérivée
$D(\mathcal{O})$ est une catégorie monoïdale symétrique, dont le produit
tensoriel est le produit tensoriel dérivé, avec les contraintes usuelles
d'associativité et de commutativité (pour les règles de signes, voir More on
Algebra, section \ref{more-algebra-section-sign-rules}).
\end{lemma}

\begin{proof}
Démonstration omise. Comparer avec le lemme
\ref{lemma-symmetric-monoidal-cat-complexes}.
\end{proof}

\begin{example}
\label{example-dual-derived}
Soit $(\mathcal{C}, \mathcal{O})$ un site annelé. Soit $K$ un objet parfait
de $D(\mathcal{O})$. Posons $K^\vee = R\SheafHom(K, \mathcal{O})$ comme dans
le lemme \ref{lemma-dual-perfect-complex}. Alors le morphisme
$$
K \otimes_\mathcal{O}^\mathbf{L} K^\vee \longrightarrow R\SheafHom(K, K)
$$
est un isomorphisme (par le lemme). Notons
$$
\eta :
\mathcal{O}
\longrightarrow
K \otimes_\mathcal{O}^\mathbf{L} K^\vee
$$
le morphisme qui envoie $1$ sur la section correspondant à $\text{id}_K$
par l'isomorphisme ci-dessus. Notons
$$
\epsilon : 
K^\vee
\otimes_\mathcal{O}^\mathbf{L} K
\longrightarrow
\mathcal{O}
$$
le morphisme d'évaluation (pour le construire, on peut par exemple utiliser
le lemme \ref{lemma-internal-hom-composition}). Alors
$K^\vee, \eta, \epsilon$ est un dual à gauche de $K$ au sens de Categories,
définition \ref{categories-definition-dual}. Nous omettons la vérification de
$(1 \otimes \epsilon) \circ (\eta \otimes 1) = \text{id}_K$
et de
$(\epsilon \otimes 1) \circ (1 \otimes \eta) =
\text{id}_{K^\vee}$.
\end{example}

\begin{lemma}
\label{lemma-left-dual-derived}
Soit $(\mathcal{C}, \mathcal{O})$ un site annelé. Soit $M$ un objet de
$D(\mathcal{O})$. Si $M$ possède un dual à gauche dans la catégorie monoïdale
$D(\mathcal{O})$ (Categories, définition
\ref{categories-definition-dual}), alors $M$ est parfait et le dual à gauche
est celui construit dans l'exemple \ref{example-dual-derived}.
\end{lemma}

\begin{proof}
Soit $N, \eta, \epsilon$ un dual à gauche. Observons que, pour tout objet $U$
de $\mathcal{C}$, la restriction $N|_U, \eta|_U, \epsilon|_U$ est un dual à
gauche de $M|_U$.

\medskip\noindent
Soit $U$ un objet de $\mathcal{C}$. Il suffit de trouver un recouvrement
$\{U_i \to U\}_{i \in I}$ de $\mathcal{C}$ tel que $M|_{U_i}$ soit un objet
parfait de $D(\mathcal{O}_{U_i})$. Nous pouvons donc remplacer
$\mathcal{C}, \mathcal{O}, M, N, \eta, \epsilon$ par
$\mathcal{C}/U, \mathcal{O}_U, M|_U, N|_U, \eta|_U, \epsilon|_U$
et supposer que $\mathcal{C}$ possède un objet final $X$. De plus, au cours
de la démonstration, nous pouvons (un nombre fini de fois) remplacer $X$ par
les membres d'un recouvrement $\{U_i \to X\}$ de $X$.

\medskip\noindent
Nous allons utiliser plusieurs fois l'argument suivant. Choisissons un
complexe quelconque $\mathcal{M}^\bullet$ de $\mathcal{O}$-modules
représentant $M$. Choisissons un complexe K-plat $\mathcal{N}^\bullet$
représentant $N$ et dont les termes sont des $\mathcal{O}$-modules plats,
voir le lemme \ref{lemma-K-flat-resolution}. Considérons le morphisme
$$
\eta : \mathcal{O} \to
\text{Tot}(\mathcal{M}^\bullet \otimes_\mathcal{O} \mathcal{N}^\bullet)
$$
Après avoir remplacé $X$ par les membres d'un recouvrement, nous pouvons
trouver un entier $N$ et, pour $i = 1, \ldots, N$, des entiers
$n_i \in \mathbf{Z}$ ainsi que des sections $f_i$ et $g_i$ de
$\mathcal{M}^{n_i}$ et de $\mathcal{N}^{-n_i}$ telles que
$$
\eta(1) = \sum\nolimits_i f_i \otimes g_i
$$
Soit $\mathcal{K}^\bullet \subset \mathcal{M}^\bullet$ un sous-complexe
quelconque de $\mathcal{O}$-modules contenant les sections $f_i$ pour
$i = 1, \ldots, N$. Puisque
$\text{Tot}(\mathcal{K}^\bullet \otimes_\mathcal{O} \mathcal{N}^\bullet)
\subset
\text{Tot}(\mathcal{M}^\bullet \otimes_\mathcal{O} \mathcal{N}^\bullet)$
par platitude des modules $\mathcal{N}^n$, on voit que $\eta$ se factorise par
$$
\tilde \eta :
\mathcal{O} \to
\text{Tot}(\mathcal{K}^\bullet \otimes_\mathcal{O} \mathcal{N}^\bullet)
$$
En notant $K$ l'objet de $D(\mathcal{O})$ représenté par
$\mathcal{K}^\bullet$, nous obtenons un diagramme commutatif
$$
\xymatrix{
M \ar[rr]_-{\eta \otimes 1} \ar[rrd]_{\tilde \eta \otimes 1} & &
M \otimes^\mathbf{L} N \otimes^\mathbf{L} M
\ar[r]_-{1 \otimes \epsilon} &
M \\
& &
K \otimes^\mathbf{L} N \otimes^\mathbf{L} M
\ar[u] \ar[r]^-{1 \otimes \epsilon} &
K \ar[u]
}
$$
Puisque la composée de la ligne supérieure est l'identité de $M$, nous
concluons que $M$ est facteur direct de $K$ dans $D(\mathcal{O})$.

\medskip\noindent
Comme première application de l'argument ci-dessus, nous pouvons choisir le
sous-complexe
$\mathcal{K}^\bullet = \sigma_{\geq a} \tau_{\leq b}\mathcal{M}^\bullet$
avec $a < n_i < b$ pour $i = 1, \ldots, N$. Ainsi, $M$ est facteur direct
dans $D(\mathcal{O})$ d'un complexe borné, et nous concluons que nous pouvons
supposer $M$ dans $D^b(\mathcal{O})$. (Rappelons que le procédé ci-dessus
comporte le remplacement de $X$ par les membres d'un recouvrement.)

\medskip\noindent
Puisque $M$ appartient à $D^b(\mathcal{O})$, nous pouvons choisir
$\mathcal{M}^\bullet$ comme un complexe borné supérieurement de modules plats
(par Modules, lemme \ref{modules-lemma-module-quotient-flat}, et Derived
Categories, lemme \ref{derived-lemma-subcategory-left-resolution}). Dans
l'argument ci-dessus, nous pouvons alors choisir
$\mathcal{K}^\bullet = \sigma_{\geq a}\mathcal{M}^\bullet$ avec $a < n_i$
pour $i = 1, \ldots, N$. Ainsi, nous pouvons supposer que $M$ est facteur
direct dans $D(\mathcal{O})$ d'un complexe borné de modules plats. En
particulier, $M$ a une amplitude de Tor finie.

\medskip\noindent
Disons que $M$ a une amplitude de Tor dans $[a, b]$. En supposant $M$
$m$-pseudo-cohérent, nous allons montrer qu'après avoir remplacé $X$ par les
membres d'un recouvrement, nous pouvons supposer $M$
$(m - 1)$-pseudo-cohérent. Cela achèvera la démonstration par le lemme
\ref{lemma-perfect-precise} et le fait que $M$ est de toute façon
$(b + 1)$-pseudo-cohérent. Après avoir remplacé $X$ par les membres d'un
recouvrement, nous pouvons supposer qu'il existe un complexe strictement
parfait $\mathcal{E}^\bullet$ et un morphisme
$\alpha : \mathcal{E}^\bullet \to M$ dans $D(\mathcal{O})$ tel que
$H^i(\alpha)$ soit un isomorphisme pour $i > m$ et surjectif pour $i = m$.
Nous pouvons supposer, et supposons, que $\mathcal{E}^i = 0$ pour $i < m$.
Choisissons un triangle distingué
$$
\mathcal{E}^\bullet \to M \to L \to \mathcal{E}^\bullet[1]
$$
Observons que $H^i(L) = 0$ pour $i \geq m$. Nous pouvons donc représenter $L$
par un complexe $\mathcal{L}^\bullet$ tel que $\mathcal{L}^i = 0$ pour
$i \geq m$. Le morphisme $L \to \mathcal{E}^\bullet[1]$ est donné par un
morphisme de complexes $\mathcal{L}^\bullet \to \mathcal{E}^\bullet[1]$ qui
est nul à tous les degrés sauf au degré $m - 1$, où nous obtenons un morphisme
$\mathcal{L}^{m - 1} \to \mathcal{E}^m$, voir Catégories dérivées, lemme
\ref{derived-lemma-negative-exts}. Alors $M$ est représenté par le complexe
$$
\mathcal{M}^\bullet :
\ldots \to
\mathcal{L}^{m - 2} \to
\mathcal{L}^{m - 1} \to
\mathcal{E}^m \to
\mathcal{E}^{m + 1} \to \ldots
$$
Appliquons à ce complexe la discussion du deuxième paragraphe pour obtenir
des sections $f_i$ de $\mathcal{M}^{n_i}$ pour $i = 1, \ldots, N$. Pour
$n < m$, soit $\mathcal{K}^n \subset \mathcal{L}^n$ le
$\mathcal{O}$-sous-module engendré par les sections $f_i$ pour $n_i = n$ et
par les $d(f_i)$ pour $n_i = n - 1$. Pour $n \geq m$, posons
$\mathcal{K}^n = \mathcal{E}^n$. Nous avons clairement un morphisme de
triangles distingués
$$
\xymatrix{
\mathcal{E}^\bullet \ar[r] &
\mathcal{M}^\bullet \ar[r] &
\mathcal{L}^\bullet \ar[r] &
\mathcal{E}^\bullet[1] \\
\mathcal{E}^\bullet \ar[r] \ar[u] &
\mathcal{K}^\bullet \ar[r] \ar[u] &
\sigma_{\leq m - 1}\mathcal{K}^\bullet \ar[r] \ar[u] &
\mathcal{E}^\bullet[1] \ar[u]
}
$$
où tous les morphismes sont ceux indiqués ci-dessus. Notons $K$ l'objet de
$D(\mathcal{O})$ correspondant au complexe $\mathcal{K}^\bullet$. Par les
arguments du deuxième paragraphe de la démonstration, nous obtenons un
morphisme $s : M \to K$ dans $D(\mathcal{O})$ tel que la composée
$M \to K \to M$ soit l'identité de $M$. Nous ne savons pas si le diagramme
$$
\xymatrix{
\mathcal{E}^\bullet \ar[r] &
\mathcal{K}^\bullet \ar@{=}[r] &
K \\
\mathcal{E}^\bullet \ar[u]^{\text{id}} \ar[r]^i &
\mathcal{M}^\bullet \ar@{=}[r] &
M \ar[u]_s
}
$$
est commutatif, mais nous savons qu'il le devient après composition avec le
morphisme $K \to M$. Par le lemme \ref{lemma-local-actual}, après avoir
remplacé $X$ par les membres d'un recouvrement, nous pouvons supposer que
$s \circ i$ est donné par un morphisme de complexes
$\sigma : \mathcal{E}^\bullet \to \mathcal{K}^\bullet$. Par le même lemme,
nous pouvons supposer que la composée de $\sigma$ avec l'inclusion
$\mathcal{K}^\bullet \subset \mathcal{M}^\bullet$ est homotope à zéro par une
certaine homotopie
$\{h^i : \mathcal{E}^i \to \mathcal{M}^{i - 1}\}$. Ainsi, après avoir
remplacé $\mathcal{K}^{m - 1}$ par
$\mathcal{K}^{m - 1} + \Im(h^m)$ (notons qu'après cette opération,
$\mathcal{K}^{m - 1}$ est encore engendré par un nombre fini de sections
globales), on voit que $\sigma$ lui-même est homotope à zéro ! Cela signifie
que nous avons un diagramme commutatif en traits pleins
$$
\xymatrix{
\mathcal{E}^\bullet \ar[r] &
M \ar[r] &
\mathcal{L}^\bullet \ar[r] &
\mathcal{E}^\bullet[1] \\
\mathcal{E}^\bullet \ar[r] \ar[u] &
K \ar[r] \ar[u] &
\sigma_{\leq m - 1}\mathcal{K}^\bullet \ar[r] \ar[u] &
\mathcal{E}^\bullet[1] \ar[u] \\
\mathcal{E}^\bullet \ar[r] \ar[u] &
M \ar[r] \ar[u]^s &
\mathcal{L}^\bullet \ar[r] \ar@{..>}[u] &
\mathcal{E}^\bullet[1] \ar[u]
}
$$
Par les axiomes des catégories triangulées, nous obtenons une flèche en
pointillés complétant le diagramme. En considérant les faisceaux de
cohomologie au degré $m - 1$, nous obtenons
$$
\xymatrix{
H^{m - 1}(M) \ar[r] &
H^{m - 1}(\mathcal{L}^\bullet) \ar[r] &
H^m(\mathcal{E}^\bullet) \\
H^{m - 1}(K) \ar[r] \ar[u] &
H^{m - 1}(\sigma_{\leq m - 1}\mathcal{K}^\bullet) \ar[r] \ar[u] &
H^m(\mathcal{E}^\bullet) \ar[u] \\
H^{m - 1}(M) \ar[r] \ar[u] &
H^{m - 1}(\mathcal{L}^\bullet) \ar[r] \ar[u] &
H^m(\mathcal{E}^\bullet) \ar[u]
}
$$
Puisque les composées verticales sont l'identité dans les colonnes de gauche
et de droite, nous concluons que la composée verticale
$H^{m - 1}(\mathcal{L}^\bullet) \to
H^{m - 1}(\sigma_{\leq m - 1}\mathcal{K}^\bullet) \to
H^{m - 1}(\mathcal{L}^\bullet)$ au milieu est surjective ! En particulier,
$H^{m - 1}(\sigma_{\leq m - 1}\mathcal{K}^\bullet) \to
H^{m - 1}(\mathcal{L}^\bullet)$ est surjectif. En utilisant le morphisme
induit entre les suites exactes longues de faisceaux de cohomologie par le
morphisme de triangles ci-dessus, une chasse au diagramme montre que
$H^i(K) \to H^i(M)$ est un isomorphisme pour $i \geq m$ et est surjectif pour
$i = m - 1$. Par construction, nous pouvons choisir un $r \geq 0$ et une
surjection $\mathcal{O}^{\oplus r} \to \mathcal{K}^{m - 1}$. Alors la composée
$$
(\mathcal{O}^{\oplus r} \to \mathcal{E}^m \to
\mathcal{E}^{m + 1} \to \ldots ) \longrightarrow
K \longrightarrow M
$$
induit un isomorphisme sur les faisceaux de cohomologie aux degrés $\geq m$
et une surjection au degré $m - 1$, ce qui achève la démonstration.
\end{proof}

\begin{lemma}
\label{lemma-colim-and-lim-of-duals}
\begin{slogan}
Dualité triviale pour les systèmes d'objets parfaits.
\end{slogan}
Soit $(\mathcal{C}, \mathcal{O})$ un site annelé. Soit
$(K_n)_{n \in \mathbf{N}}$ un système d'objets parfaits de $D(\mathcal{O})$.
Soit $K = \text{hocolim} K_n$ la colimite dérivée (Catégories dérivées,
définition \ref{derived-definition-derived-colimit}). Alors, pour tout objet
$E$ de $D(\mathcal{O})$, nous avons
$$
R\SheafHom(K, E) = R\lim E \otimes^\mathbf{L}_\mathcal{O} K_n^\vee
$$
où $(K_n^\vee)$ est le système inverse de complexes parfaits duaux.
\end{lemma}

\begin{proof}
Par le lemme \ref{lemma-dual-perfect-complex}, nous avons
$R\lim E \otimes^\mathbf{L}_\mathcal{O} K_n^\vee =
R\lim R\SheafHom(K_n, E)$
qui s'insère dans le triangle distingué
$$
R\lim R\SheafHom(K_n, E) \to
\prod R\SheafHom(K_n, E) \to
\prod R\SheafHom(K_n, E)
$$
Puisque $K$ s'insère de même dans le triangle distingué
$\bigoplus K_n \to \bigoplus K_n \to K$, il suffit de montrer que
$\prod R\SheafHom(K_n, E) = R\SheafHom(\bigoplus K_n, E)$. C'est une
conséquence formelle de (\ref{equation-internal-hom}) et du fait que le
produit tensoriel dérivé commute aux sommes directes.
\end{proof}






\section{Objets inversibles dans la catégorie dérivée}
\label{section-invertible-D-or-R}

\noindent
Nous caractérisons les objets inversibles dans la catégorie dérivée d'un site
annelé (à la fois dans le cas d'un topos localement annelé et dans le cas
général).

\begin{lemma}
\label{lemma-category-summands-finite-free}
Soit $(\mathcal{C}, \mathcal{O})$ un site annelé. Posons
$R = \Gamma(\mathcal{C}, \mathcal{O})$. La catégorie des
$\mathcal{O}$-modules qui sont facteurs directs de $\mathcal{O}$-modules
libres de type fini est équivalente à la catégorie des $R$-modules projectifs
de type fini.
\end{lemma}

\begin{proof}
Observons qu'un $R$-module projectif de type fini est la même chose qu'un
facteur direct d'un $R$-module libre de type fini. L'équivalence est donnée
par le foncteur
$\mathcal{E} \mapsto \Gamma(\mathcal{C}, \mathcal{E})$. Le foncteur inverse
est donné par la construction suivante. Considérons le morphisme de topos
$f : \Sh(\mathcal{C}) \to \Sh(\text{pt})$, où $f_*$ est donné par la prise
des sections globales et $f^{-1}$ envoie un ensemble $S$, c'est-à-dire un
objet de $\Sh(\text{pt})$, sur le faisceau constant de valeur $S$. Nous
obtenons un morphisme
$(f, f^\sharp) : (\Sh(\mathcal{C}), \mathcal{O}) \to (\Sh(\text{pt}), R)$
de topos annelés en utilisant le morphisme identité $R \to f_*\mathcal{O}$.
Le foncteur inverse est alors donné par $f^*$.
\end{proof}

\begin{lemma}
\label{lemma-invertible-derived}
Soit $(\mathcal{C}, \mathcal{O})$ un site annelé. Soit $M$ un objet de
$D(\mathcal{O})$. Les conditions suivantes sont équivalentes :
\begin{enumerate}
\item $M$ est inversible dans $D(\mathcal{O})$, voir Categories, définition
\ref{categories-definition-invertible}, et
\item il existe une décomposition en produit direct localement
finie\footnote{Cela signifie que, pour tout objet $U$ de $\mathcal{C}$, il
existe un recouvrement $\{U_i \to U\}$ tel que, pour tout $i$, le faisceau
$\mathcal{O}_n|_{U_i}$ ne soit non nul que pour un nombre fini de $n$.}
$$
\mathcal{O} = \prod\nolimits_{n \in \mathbf{Z}} \mathcal{O}_n
$$
et, pour chaque $n$, il existe un $\mathcal{O}_n$-module inversible
$\mathcal{H}^n$ (Modules on Sites, définition
\ref{sites-modules-definition-invertible-sheaf}), et
$M = \bigoplus \mathcal{H}^n[-n]$ dans $D(\mathcal{O})$.
\end{enumerate}
Si (1) et (2) sont satisfaites, alors $M$ est un objet parfait de
$D(\mathcal{O})$. Si $(\mathcal{C}, \mathcal{O})$ est un site localement
annelé, ces conditions sont également équivalentes à :
\begin{enumerate}
\item[(3)] pour tout objet $U$ de $\mathcal{C}$, il existe un recouvrement
$\{U_i \to U\}$ et, pour chaque $i$, un entier $n_i$ tel que $M|_{U_i}$ soit
représenté par un $\mathcal{O}_{U_i}$-module inversible placé au degré $n_i$.
\end{enumerate}
\end{lemma}

\begin{proof}
Supposons (2). Considérons l'objet $R\SheafHom(M, \mathcal{O})$ et le
morphisme de composition
$$
R\SheafHom(M, \mathcal{O}) \otimes_\mathcal{O}^\mathbf{L} M \to \mathcal{O}
$$
Pour montrer que c'est un isomorphisme, nous pouvons travailler localement.
Nous pouvons donc supposer que
$\mathcal{O} = \prod_{a \leq n \leq b} \mathcal{O}_n$ et
$M = \bigoplus_{a \leq n \leq b} \mathcal{H}^n[-n]$. Il suffit alors de
montrer que
$$
R\SheafHom(\mathcal{H}^m, \mathcal{O})
\otimes_\mathcal{O}^\mathbf{L} \mathcal{H}^n
$$
est nul si $n \not = m$ et égal à $\mathcal{O}_n$ si $n = m$. Le cas
$n \not = m$ résulte du fait que $\mathcal{O}_n$ et $\mathcal{O}_m$ sont des
$\mathcal{O}$-algèbres plates telles que
$\mathcal{O}_n \otimes_\mathcal{O} \mathcal{O}_m = 0$. En utilisant la
structure locale des $\mathcal{O}$-modules inversibles (Modules on Sites,
lemme \ref{sites-modules-lemma-invertible}) et en travaillant localement,
l'isomorphisme dans le cas $n = m$ s'obtient immédiatement ; nous omettons
les détails. Puisque $D(\mathcal{O})$ est monoïdale symétrique, nous concluons
que $M$ est inversible.

\medskip\noindent
Supposons (1). La description de (2) montre que nous avons un candidat pour
$\mathcal{O}_n$, à savoir
$\SheafHom_\mathcal{O}(H^n(M), H^n(M))$. Si cette famille de faisceaux
d'anneaux est localement finie et si
$\mathcal{O} = \prod \mathcal{O}_n$, alors nous obtenons immédiatement la
décomposition en somme directe $M = \bigoplus H^n(M)[-n]$ en utilisant les
idempotents de $\mathcal{O}$ provenant de la décomposition en produit. Cela
montre que, pour démontrer (2), nous pouvons travailler localement au sens
suivant. Soit $U$ un objet de $\mathcal{C}$. Nous devons montrer qu'il existe
un recouvrement $\{U_i \to U\}$ de $U$ tel qu'avec les $\mathcal{O}_n$
ci-dessus, les assertions précédentes et celles de (2) soient satisfaites
après restriction à $\mathcal{C}/U_i$. Nous pouvons donc supposer que
$\mathcal{C}$ possède un objet final $X$ et, au cours de la démonstration de
(2), remplacer un nombre fini de fois $X$ par les membres d'un recouvrement
de $X$.

\medskip\noindent
Choisissons un objet $N$ de $D(\mathcal{O})$ et un isomorphisme
$M \otimes_\mathcal{O}^\mathbf{L} N \cong \mathcal{O}$. Alors $N$ est un dual
à gauche de $M$ dans la catégorie monoïdale $D(\mathcal{O})$, et nous
concluons que $M$ est parfait par le lemme \ref{lemma-left-dual-derived}. Par
symétrie, $N$ est parfait. Après avoir remplacé $X$ par les membres d'un
recouvrement, nous pouvons supposer que $M$ et $N$ sont représentés par des
complexes strictement parfaits $\mathcal{E}^\bullet$ et
$\mathcal{F}^\bullet$. Alors $M \otimes_\mathcal{O}^\mathbf{L} N$ est
représenté par
$\text{Tot}(\mathcal{E}^\bullet \otimes_\mathcal{O} \mathcal{F}^\bullet)$.
Après avoir remplacé $X$ par les membres d'un recouvrement de $X$, nous
pouvons supposer que les isomorphismes réciproques
$\mathcal{O} \to M \otimes_\mathcal{O}^\mathbf{L} N$ et
$M \otimes_\mathcal{O}^\mathbf{L} N \to \mathcal{O}$ sont donnés par des
morphismes de complexes
$$
\alpha : \mathcal{O} \to
\text{Tot}(\mathcal{E}^\bullet \otimes_\mathcal{O} \mathcal{F}^\bullet)
\quad\text{et}\quad
\beta :
\text{Tot}(\mathcal{E}^\bullet \otimes_\mathcal{O} \mathcal{F}^\bullet)
\to \mathcal{O}
$$
voir le lemme \ref{lemma-local-actual}. Alors $\beta \circ \alpha = 1$ comme
morphismes de complexes et $\alpha \circ \beta = 1$ comme morphisme dans
$D(\mathcal{O})$. Après avoir remplacé $X$ par les membres d'un recouvrement
de $X$, nous pouvons supposer que la composée $\alpha \circ \beta$ est
homotope à $1$ par une certaine homotopie $\theta$ de composantes
$$
\theta^n :
\text{Tot}^n(\mathcal{E}^\bullet \otimes_\mathcal{O} \mathcal{F}^\bullet)
\to
\text{Tot}^{n - 1}(
\mathcal{E}^\bullet \otimes_\mathcal{O} \mathcal{F}^\bullet)
$$
par le même lemme que précédemment. Posons
$R = \Gamma(\mathcal{C}, \mathcal{O})$. Par le lemme
\ref{lemma-category-summands-finite-free}, nous obtenons :
\begin{enumerate}
\item $M^\bullet = \Gamma(X, \mathcal{E}^\bullet)$
est un complexe borné de $R$-modules projectifs de type fini,
\item $N^\bullet = \Gamma(X, \mathcal{F}^\bullet)$
est un complexe borné de $R$-modules projectifs de type fini,
\item $\alpha$ et $\beta$ correspondent aux morphismes de complexes
$a : R \to \text{Tot}(M^\bullet \otimes_R N^\bullet)$ et
$b : \text{Tot}(M^\bullet \otimes_R N^\bullet) \to R$,
\item $\theta^n$ correspond au morphisme
$h^n : \text{Tot}^n(M^\bullet \otimes_R N^\bullet) \to
\text{Tot}^{n - 1}(M^\bullet \otimes_R N^\bullet)$, et
\item $b \circ a = 1$ et $a \circ b - 1 = dh + hd$,
\end{enumerate}
Il s'ensuit que $M^\bullet$ et $N^\bullet$ définissent des objets inverses
l'un de l'autre dans $D(R)$. Par More on Algebra, lemme
\ref{more-algebra-lemma-invertible-derived}, nous obtenons une décomposition
en produit $R = \prod_{a \leq n \leq b} R_n$ et des $R_n$-modules inversibles
$H^n$ tels que
$M^\bullet \cong \bigoplus_{a \leq n \leq b} H^n[-n]$. Cet isomorphisme dans
$D(R)$ peut être relevé en un morphisme
$$
\bigoplus H^n[-n] \longrightarrow M^\bullet
$$
de complexes, car chaque $H^n$ est projectif comme $R$-module. De manière
correspondante, en utilisant à nouveau le lemme
\ref{lemma-category-summands-finite-free}, nous obtenons un morphisme
$$
\bigoplus H^n \otimes_R \mathcal{O}[-n] \to \mathcal{E}^\bullet
$$
qui est un isomorphisme dans $D(\mathcal{O})$. Ici,
$M \otimes_R \mathcal{O}$ désigne le foncteur des $R$-modules projectifs de
type fini vers les $\mathcal{O}$-modules construit dans la démonstration du
lemme \ref{lemma-category-summands-finite-free}. En posant
$\mathcal{O}_n = R_n \otimes_R \mathcal{O}$, nous concluons que (2) est vraie.

\medskip\noindent
Si $(\mathcal{C}, \mathcal{O})$ est un site localement annelé, étant donnés
un objet $U$ et une décomposition finie en produit
$\mathcal{O}|_U = \prod_{a \leq n \leq b} \mathcal{O}_n|_U$, nous pouvons
trouver un recouvrement $\{U_i \to U\}$ tel que, pour tout $i$, il existe au
plus un $n$ pour lequel $\mathcal{O}_n|_{U_i}$ est non nul. Cela résulte
aisément de la partie (2) de Modules on Sites, lemme
\ref{sites-modules-lemma-locally-ringed}, et de la définition des sites
localement annelés donnée dans Modules on Sites, définition
\ref{sites-modules-definition-locally-ringed}. L'implication
(2) $\Rightarrow$ (3) s'en déduit facilement. L'implication
(3) $\Rightarrow$ (2) est vraie sans aucune hypothèse sur le site annelé.
Nous omettons les détails.
\end{proof}










\section{Formule de projection}
\label{section-projection-formula}

\noindent
Soit $f : (\Sh(\mathcal{C}), \mathcal{O}_\mathcal{C}) \to
(\Sh(\mathcal{D}), \mathcal{O}_\mathcal{D})$ un morphisme de topos annelés.
Soient $E \in D(\mathcal{O}_\mathcal{C})$ et
$K \in D(\mathcal{O}_\mathcal{D})$. Sans aucune hypothèse supplémentaire, il
existe un morphisme
\begin{equation}
\label{equation-projection-formula-map}
Rf_*E \otimes^\mathbf{L}_{\mathcal{O}_\mathcal{D}} K
\longrightarrow
Rf_*(E \otimes^\mathbf{L}_{\mathcal{O}_\mathcal{C}} Lf^*K)
\end{equation}
À savoir, c'est l'adjoint du morphisme canonique
$$
Lf^*(Rf_*E \otimes^\mathbf{L}_{\mathcal{O}_\mathcal{D}} K) =
Lf^*Rf_*E \otimes^\mathbf{L}_{\mathcal{O}_\mathcal{C}} Lf^*K
\longrightarrow
E \otimes^\mathbf{L}_{\mathcal{O}_\mathcal{C}} Lf^*K
$$
provenant du morphisme $Lf^*Rf_*E \to E$ et des lemmes
\ref{lemma-pullback-tensor-product} et \ref{lemma-adjoint}. Une version assez
générale de la formule de projection est la suivante.

\begin{lemma}
\label{lemma-projection-formula}
Soit $f : (\Sh(\mathcal{C}), \mathcal{O}_\mathcal{C}) \to
(\Sh(\mathcal{D}), \mathcal{O}_\mathcal{D})$ un morphisme de topos annelés.
Soient $E \in D(\mathcal{O}_\mathcal{C})$ et
$K \in D(\mathcal{O}_\mathcal{D})$. Si $K$ est parfait, alors
$$
Rf_*E \otimes^\mathbf{L}_{\mathcal{O}_\mathcal{D}} K =
Rf_*(E \otimes^\mathbf{L}_{\mathcal{O}_\mathcal{C}} Lf^*K)
$$
dans $D(\mathcal{O}_\mathcal{D})$.
\end{lemma}

\begin{proof}
Pour vérifier que (\ref{equation-projection-formula-map}) est un isomorphisme,
nous pouvons travailler localement sur $\mathcal{D}$, c'est-à-dire que, pour
tout objet $V$ de $\mathcal{D}$, nous devons trouver un recouvrement
$\{V_j \to V\}$ tel que la restriction du morphisme à $V_j$ soit un
isomorphisme. Par définition des objets parfaits, cela signifie que nous
pouvons supposer $K$ représenté par un complexe strictement parfait de
$\mathcal{O}_\mathcal{D}$-modules. Notons que, tout à fait généralement,
l'assertion est vraie pour $K = K_1 \oplus K_2$ si et seulement si elle est
vraie pour $K_1$ et $K_2$. Nous pouvons donc supposer que $K$ est un complexe
fini de $\mathcal{O}_\mathcal{D}$-modules libres de type fini. Dans ce cas, un
argument simple utilisant les troncations bêtes ramène l'assertion au cas où
$K$ est représenté par un $\mathcal{O}_\mathcal{D}$-module libre de type fini.
Puisque l'assertion est invariante par passage aux facteurs directs finis
dans la variable $K$, nous concluons qu'il suffit de la démontrer pour
$K = \mathcal{O}_\mathcal{D}[n]$, cas où elle est triviale.
\end{proof}

\begin{remark}
\label{remark-compatible-with-diagram}
Le morphisme (\ref{equation-projection-formula-map}) est compatible avec le
morphisme de changement de base de la remarque \ref{remark-base-change} au
sens suivant. Supposons en effet que
$$
\xymatrix{
(\Sh(\mathcal{C}'), \mathcal{O}_{\mathcal{C}'})
\ar[r]_{g'} \ar[d]_{f'} &
(\Sh(\mathcal{C}), \mathcal{O}_\mathcal{C}) \ar[d]^f \\
(\Sh(\mathcal{D}'), \mathcal{O}_{\mathcal{D}'})
\ar[r]^g &
(\Sh(\mathcal{D}), \mathcal{O}_\mathcal{D})
}
$$
soit un diagramme commutatif de topos annelés. Soient
$E \in D(\mathcal{O}_\mathcal{C})$ et $K \in D(\mathcal{O}_\mathcal{D})$.
Alors le diagramme
$$
\xymatrix{
Lg^*(Rf_*E \otimes^\mathbf{L}_{\mathcal{O}_\mathcal{D}} K)
\ar[r]_p \ar[d]_t &
Lg^*Rf_*(E \otimes^\mathbf{L}_{\mathcal{O}_\mathcal{C}} Lf^*K)
\ar[d]_b \\
Lg^*Rf_*E \otimes^\mathbf{L}_{\mathcal{O}_{\mathcal{D}'}}
Lg^*K \ar[d]_b &
Rf'_*L(g')^*(E \otimes^\mathbf{L}_{\mathcal{O}_\mathcal{C}}
Lf^*K) \ar[d]_t \\
Rf'_*L(g')^*E \otimes^\mathbf{L}_{\mathcal{O}_{\mathcal{D}'}}
Lg^*K \ar[rd]_p &
Rf'_*(L(g')^*E \otimes^\mathbf{L}_{\mathcal{O}_{\mathcal{D}'}}
L(g')^*Lf^*K) \ar[d]_c \\
& Rf'_*(L(g')^*E \otimes^\mathbf{L}_{\mathcal{O}_{\mathcal{D}'}}
L(f')^*Lg^*K)
}
$$
est commutatif. Ici, les flèches étiquetées $t$ sont obtenues par une
application du lemme \ref{lemma-pullback-tensor-product}, celles étiquetées
$b$ par une application de la remarque \ref{remark-base-change}, celles
étiquetées $p$ par une application de
(\ref{equation-projection-formula-map}), et $c$ provient de
$L(g')^* \circ Lf^* = L(f')^* \circ Lg^*$. Nous omettons la vérification.
\end{remark}






\section{Objets faiblement contractiles}
\label{section-w-contractible}

\noindent
Un objet $U$ d'un site est {\it faiblement contractile} si toute surjection
$\mathcal{F} \to \mathcal{G}$ de faisceaux d'ensembles donne lieu à une
surjection $\mathcal{F}(U) \to \mathcal{G}(U)$, voir Sites, définition
\ref{sites-definition-w-contractible}.

\begin{lemma}
\label{lemma-w-contractible}
Soit $\mathcal{C}$ un site. Soit $U$ un objet faiblement contractile de
$\mathcal{C}$. Alors :
\begin{enumerate}
\item le foncteur $\mathcal{F} \mapsto \mathcal{F}(U)$ est un foncteur exact
$\textit{Ab}(\mathcal{C}) \to \textit{Ab}$,
\item $H^p(U, \mathcal{F}) = 0$ pour tout faisceau abélien $\mathcal{F}$ et
tout $p \geq 1$, et
\item pour tout faisceau de groupes $\mathcal{G}$, tout
$\mathcal{G}$-torseur possède une section sur $U$.
\end{enumerate}
\end{lemma}

\begin{proof}
La première assertion résulte immédiatement de la définition (voir aussi
Homology, section \ref{homology-section-functors}). Les foncteurs dérivés
supérieurs s'annulent par Catégories dérivées, lemme
\ref{derived-lemma-right-derived-exact-functor}. Soit $\mathcal{F}$ un
$\mathcal{G}$-torseur. Alors $\mathcal{F} \to *$ est un morphisme surjectif de
faisceaux. Ainsi, (3) résulte également de la définition.
\end{proof}

\noindent
Il est utile d'énumérer ici quelques conséquences de l'existence de
suffisamment d'objets faiblement contractiles.

\begin{proposition}
\label{proposition-enough-weakly-contractibles}
Soit $\mathcal{C}$ un site. Soit $\mathcal{B} \subset \Ob(\mathcal{C})$ tel
que tout $U \in \mathcal{B}$ soit faiblement contractile et que tout objet de
$\mathcal{C}$ possède un recouvrement par des éléments de $\mathcal{B}$. Soit
$\mathcal{O}$ un faisceau d'anneaux sur $\mathcal{C}$. Alors :
\begin{enumerate}
\item Un complexe $\mathcal{F}_1 \to \mathcal{F}_2 \to \mathcal{F}_3$ de
$\mathcal{O}$-modules est exact si et seulement si
$\mathcal{F}_1(U) \to \mathcal{F}_2(U) \to \mathcal{F}_3(U)$
est exact pour tout $U \in \mathcal{B}$.
\item Tout objet $K$ de $D(\mathcal{O})$ est la limite dérivée de ses
troncations canoniques : $K = R\lim \tau_{\geq -n} K$.
\item Étant donné un système inverse
$\ldots \to \mathcal{F}_3 \to \mathcal{F}_2 \to \mathcal{F}_1$
dont les morphismes de transition sont surjectifs, la projection
$\lim \mathcal{F}_n \to \mathcal{F}_1$ est surjective.
\item Les produits sont exacts dans $\textit{Mod}(\mathcal{O})$.
\item Les produits dans $D(\mathcal{O})$ peuvent être calculés en prenant les
produits de complexes représentants quelconques.
\item Si $(\mathcal{F}_n)$ est un système inverse de $\mathcal{O}$-modules,
alors $R^p\lim \mathcal{F}_n = 0$ pour tout $p > 1$ et
$$
R^1\lim \mathcal{F}_n  =
\Coker(\prod \mathcal{F}_n \to \prod \mathcal{F}_n)
$$
où le morphisme est $(x_n) \mapsto (x_n - f(x_{n + 1}))$.
\item Si $(K_n)$ est un système inverse d'objets de $D(\mathcal{O})$, alors il
existe des suites exactes courtes
$$
0 \to R^1\lim H^{p - 1}(K_n) \to H^p(R\lim K_n) \to
\lim H^p(K_n) \to 0
$$
\end{enumerate}
\end{proposition}

\begin{proof}
Démonstration de (1). Si la suite est exacte, son évaluation en tout élément
faiblement contractile de $\mathcal{C}$ donne une suite exacte par le lemme
\ref{lemma-w-contractible}. Réciproquement, supposons que
$\mathcal{F}_1(U) \to \mathcal{F}_2(U) \to \mathcal{F}_3(U)$
soit exact pour tout $U \in \mathcal{B}$. Soit $V$ un objet de
$\mathcal{C}$ et soit $s \in \mathcal{F}_2(V)$ un élément du noyau de
$\mathcal{F}_2 \to \mathcal{F}_3$. Par hypothèse, il existe un recouvrement
$\{U_i \to V\}$ avec $U_i \in \mathcal{B}$. Alors $s|_{U_i}$ se relève en
une section $s_i \in \mathcal{F}_1(U_i)$. Ainsi, $s$ est une section du
faisceau image $\Im(\mathcal{F}_1 \to \mathcal{F}_2)$. Autrement dit, la
suite $\mathcal{F}_1 \to \mathcal{F}_2 \to \mathcal{F}_3$ est exacte.

\medskip\noindent
Démonstration de (2). Cela résulte du lemme
\ref{lemma-is-limit-dimension} avec $d = 0$.

\medskip\noindent
Démonstration de (3). Soit $(\mathcal{F}_n)$ un système comme en (3), et
posons $\mathcal{F} = \lim \mathcal{F}_n$. Si $U \in \mathcal{B}$, alors
$\mathcal{F}(U) = \lim \mathcal{F}_n(U)$
se surjecte sur $\mathcal{F}_1(U)$ puisque tous les morphismes de transition
$\mathcal{F}_{n + 1}(U) \to \mathcal{F}_n(U)$ sont surjectifs. Ainsi,
$\mathcal{F} \to \mathcal{F}_1$ est surjectif par Sites, définition
\ref{sites-definition-sheaves-injective-surjective}, et par l'hypothèse que
tout objet possède un recouvrement par des éléments de $\mathcal{B}$.

\medskip\noindent
Démonstration de (4). Soit
$\mathcal{F}_{i, 1} \to \mathcal{F}_{i, 2} \to \mathcal{F}_{i, 3}$
une famille de suites exactes de $\mathcal{O}$-modules. Nous voulons montrer
que
$\prod \mathcal{F}_{i, 1} \to \prod \mathcal{F}_{i, 2} \to
\prod \mathcal{F}_{i, 3}$ est exacte. Nous utilisons le critère de (1).
Soit $U \in \mathcal{B}$. Alors
$$
(\prod \mathcal{F}_{i, 1})(U) \to
(\prod \mathcal{F}_{i, 2})(U) \to
(\prod \mathcal{F}_{i, 3})(U)
$$
est la même suite que
$$
\prod \mathcal{F}_{i, 1}(U) \to
\prod \mathcal{F}_{i, 2}(U) \to
\prod \mathcal{F}_{i, 3}(U)
$$
Chacune des suites
$\mathcal{F}_{i, 1}(U) \to \mathcal{F}_{i, 2}(U) \to \mathcal{F}_{i, 3}(U)$
est exacte par (1). Ainsi, les suites affichées sont exactes par Homology,
lemme \ref{homology-lemma-product-abelian-groups-exact}. Nous concluons en
appliquant à nouveau (1).

\medskip\noindent
Démonstration de (5). Cela résulte de (4) et de Catégories dérivées, lemme
\ref{derived-lemma-products} (légèrement généralisé).

\medskip\noindent
Démonstration de (6) et (7). Nous renvoyons à la section
\ref{section-derived-limits} pour une discussion des limites dérivées, des
limites homotopiques et de leur relation. Par Catégories dérivées, définition
\ref{derived-definition-derived-limit}, nous avons un triangle distingué
$$
R\lim K_n \to \prod K_n \to \prod K_n \to R\lim K_n[1]
$$
En prenant la suite exacte longue des faisceaux de cohomologie, nous obtenons
$$
H^{p - 1}(\prod K_n) \to H^{p - 1}(\prod K_n) \to
H^p(R\lim K_n) \to H^p(\prod K_n) \to H^p(\prod K_n)
$$
Puisque les produits sont exacts par (4), cela devient
$$
\prod H^{p - 1}(K_n) \to \prod H^{p - 1}(K_n) \to
H^p(R\lim K_n) \to \prod H^p(K_n) \to \prod H^p(K_n)
$$
Appliquons d'abord ceci au cas $K_n = \mathcal{F}_n[0]$, où
$(\mathcal{F}_n)$ est comme en (6). Nous concluons que (6) est satisfaite.
Appliquons-le ensuite à $(K_n)$ comme en (7) et nous concluons que (7) est
satisfaite.
\end{proof}






\section{Objets compacts}
\label{section-compact}

\noindent
Dans cette section, nous étudions les objets compacts de la catégorie dérivée
des modules sur un site annelé. Rappelons que les objets compacts sont définis
dans Catégories dérivées, définition \ref{derived-definition-compact-object}.

\begin{lemma}
\label{lemma-compact-in-terms-of-generators}
Soit $\mathcal{A}$ une catégorie abélienne de Grothendieck. Soit
$S \subset \Ob(\mathcal{A})$ un ensemble d'objets tel que :
\begin{enumerate}
\item tout objet de $\mathcal{A}$ est un quotient d'une somme directe
d'éléments de $S$, et
\item pour tout $E \in S$, le foncteur $\Hom_\mathcal{A}(E, -)$ commute aux
sommes directes.
\end{enumerate}
Alors tout objet compact de $D(\mathcal{A})$ est facteur direct dans
$D(\mathcal{A})$ d'un complexe fini de sommes directes finies d'éléments de
$S$.
\end{lemma}

\begin{proof}
Supposons que $K \in D(\mathcal{A})$ soit un objet compact. Représentons $K$
par un complexe $K^\bullet$ et considérons le morphisme
$$
K^\bullet
\longrightarrow
\bigoplus\nolimits_{n \geq 0} \tau_{\geq n} K^\bullet
$$
où nous avons utilisé les troncations canoniques, voir Homology, section
\ref{homology-section-truncations}. Cela a un sens, car, à chaque degré, la
somme directe de droite est finie. Par hypothèse, ce morphisme se factorise
par une somme directe finie. Nous concluons que
$K \to \tau_{\geq n} K$ est nul pour au moins un $n$, c'est-à-dire que $K$
appartient à $D^{-}(\mathcal{A})$.

\medskip\noindent
Nous pouvons représenter $K$ par un complexe borné supérieurement
$K^\bullet$, dont chaque terme est une somme directe d'objets de $S$, voir
Catégories dérivées, lemme \ref{derived-lemma-subcategory-left-resolution}.
Notons que
$$
K^\bullet = \bigcup\nolimits_{n \leq 0} \sigma_{\geq n}K^\bullet
$$
où nous avons utilisé les troncations bêtes, voir Homology, section
\ref{homology-section-truncations}. Par Catégories dérivées, lemmes
\ref{derived-lemma-colim-hocolim} et
\ref{derived-lemma-commutes-with-countable-sums}, nous voyons donc que
$1 : K^\bullet \to K^\bullet$ se factorise par
$\sigma_{\geq n}K^\bullet \to K^\bullet$ dans $D(\mathcal{A})$. Ainsi,
$1 : K^\bullet \to K^\bullet$ se factorise sous la forme
$$
K^\bullet \xrightarrow{\varphi} L^\bullet \xrightarrow{\psi} K^\bullet
$$
dans $D(\mathcal{A})$, pour un certain complexe borné $L^\bullet$ dont les
termes sont des sommes directes d'éléments de $S$. Disons que $L^i$ est nul
pour $i \not \in [a, b]$. Soit $c$ le plus grand entier $\leq b + 1$ tel que
$L^i$ soit une somme directe finie d'éléments de $S$ pour $i < c$.
Affirmation : si $c < b + 1$, alors nous pouvons modifier $L^\bullet$ pour
augmenter $c$. Par récurrence, cette affirmation montrera que nous avons une
factorisation de $1_K$ sous la forme
$$
K \xrightarrow{\varphi} L \xrightarrow{\psi} K
$$
dans $D(\mathcal{A})$, où $L$ peut être représenté par un complexe fini de
sommes directes finies d'éléments de $S$. Notons que
$e = \varphi \circ \psi \in \text{End}_{D(\mathcal{A})}(L)$ est un
idempotent. Par Catégories dérivées, lemme
\ref{derived-lemma-projectors-have-images-triangulated}, on voit que
$L = \Ker(e) \oplus \Ker(1 - e)$. Le morphisme $\varphi : K \to L$ induit un
isomorphisme avec $\Ker(1 - e)$ dans $D(\mathcal{A})$, et nous concluons.

\medskip\noindent
Démonstration de l'affirmation. Écrivons
$L^c = \bigoplus_{\lambda \in \Lambda} E_\lambda$. Puisque $L^{c - 1}$ est
une somme directe finie d'éléments de $S$, l'hypothèse (2) nous permet de
trouver un sous-ensemble fini $\Lambda' \subset \Lambda$ tel que
$L^{c - 1} \to L^c$ se factorise par
$\bigoplus_{\lambda \in \Lambda'} E_\lambda \subset L^c$. Considérons le
morphisme de complexes
$$
\pi :
L^\bullet
\longrightarrow
(\bigoplus\nolimits_{\lambda \in \Lambda \setminus \Lambda'} E_\lambda)[-c]
$$
donné par la projection sur les facteurs correspondant à
$\Lambda \setminus \Lambda'$ au degré $c$. Par notre hypothèse sur $K$, on
voit qu'après avoir éventuellement remplacé $\Lambda'$ par un sous-ensemble
fini plus grand, nous pouvons supposer que $\pi \circ \varphi = 0$ dans
$D(\mathcal{A})$. Soit $(L')^\bullet \subset L^\bullet$ le noyau de $\pi$.
Comme $\pi$ est surjectif, nous obtenons une suite exacte courte de complexes,
qui donne un triangle distingué dans $D(\mathcal{A})$ (voir Derived
Categories, lemme \ref{derived-lemma-derived-canonical-delta-functor}).
Puisque $\Hom_{D(\mathcal{A})}(K, -)$ est homologique (voir Derived
Categories, lemme \ref{derived-lemma-representable-homological}) et que
$\pi \circ \varphi = 0$, nous pouvons trouver un morphisme
$\varphi' : K^\bullet \to (L')^\bullet$ dans $D(\mathcal{A})$ dont la
composée avec $(L')^\bullet \to L^\bullet$ est $\varphi$. En prenant pour
$\psi'$ la composée de $\psi$ avec $(L')^\bullet \to L^\bullet$, nous
obtenons une nouvelle factorisation. Puisque $(L')^\bullet$ coïncide avec
$L^\bullet$ sauf au degré $c$ et que
$(L')^c = \bigoplus_{\lambda \in \Lambda'} E_\lambda$, l'affirmation est
démontrée.
\end{proof}

\begin{lemma}
\label{lemma-compact-objects-if-enough-qc}
Soit $(\mathcal{C}, \mathcal{O})$ un site annelé. Supposons que tout objet de
$\mathcal{C}$ possède un recouvrement par des objets quasi-compacts. Alors
tout objet compact de $D(\mathcal{O})$ est facteur direct dans
$D(\mathcal{O})$ d'un complexe fini dont les termes sont des sommes directes
finies de $\mathcal{O}$-modules de la forme $j_!\mathcal{O}_U$, où $U$ est un
objet quasi-compact de $\mathcal{C}$.
\end{lemma}

\begin{proof}
Appliquons le lemme \ref{lemma-compact-in-terms-of-generators}, où
$S \subset \Ob(\textit{Mod}(\mathcal{O}))$ est l'ensemble des modules de la
forme $j_!\mathcal{O}_U$ avec $U \in \Ob(\mathcal{C})$ quasi-compact.
L'hypothèse (1) est satisfaite par Modules on Sites, lemme
\ref{sites-modules-lemma-module-quotient-flat}, et par l'hypothèse que tout
$U$ peut être recouvert par des objets quasi-compacts. L'hypothèse (2) résulte
de
$$
\Hom_\mathcal{O}(j_!\mathcal{O}_U, \mathcal{F}) = \mathcal{F}(U)
$$
qui commute aux sommes directes par Sites, lemme
\ref{sites-lemma-directed-colimits-sections}.
\end{proof}

\noindent
Dans la situation du lemme ci-dessus, il n'est pas toujours vrai que les
modules $j_!\mathcal{O}_U$ soient des objets compacts de $D(\mathcal{O})$
(même si $U$ est un objet quasi-compact de $\mathcal{C}$). Voici deux lemmes
qui traitent cette question.


\begin{lemma}
\label{lemma-when-jshriek-lower-compact}
Soit $(\mathcal{C}, \mathcal{O})$ un site annelé. Soit $U$ un objet de
$\mathcal{C}$. Supposons que les foncteurs
$\mathcal{F} \mapsto H^p(U, \mathcal{F})$ commutent aux sommes directes. Alors
le $\mathcal{O}$-module $j_!\mathcal{O}_U$ est un objet compact de
$D^+(\mathcal{O})$ au sens suivant : si
$M = \bigoplus_{i \in I} M_i$ dans $D(\mathcal{O})$ est borné inférieurement,
alors $\Hom(j_{U!}\mathcal{O}_U, M) =
\bigoplus_{i \in I} \Hom(j_{U!}\mathcal{O}_U, M_i)$.
\end{lemma}

\begin{proof}
Puisque $\Hom(j_{U!}\mathcal{O}_U, -)$ est le même foncteur que
$\mathcal{F} \mapsto \mathcal{F}(U)$ par Modules on Sites, équation
(\ref{sites-modules-equation-map-lower-shriek-OU-into-module}), il suffit de
montrer que $H^p(U, M) = \bigoplus H^p(U, M_i)$. Soit $\mathcal{I}_i$,
$i \in I$, une famille de $\mathcal{O}$-modules injectifs. Par hypothèse,
nous avons
$$
H^p(U, \bigoplus\nolimits_{i \in I} \mathcal{I}_i) =
\bigoplus\nolimits_{i \in I} H^p(U, \mathcal{I}_i) = 0
$$
pour tout $p > 0$. Puisque $M = \bigoplus M_i$ est borné inférieurement, il
existe un $a \in \mathbf{Z}$ tel que $H^n(M_i) = 0$ pour $n < a$. Nous
pouvons donc choisir des complexes de $\mathcal{O}$-modules injectifs
$\mathcal{I}_i^\bullet$ représentant $M_i$ avec $\mathcal{I}_i^n = 0$ pour
$n < a$, voir Catégories dérivées, lemme
\ref{derived-lemma-injective-resolutions-exist}. Par Injectives, lemme
\ref{injectives-lemma-derived-products}, on voit que le complexe somme directe
$\bigoplus \mathcal{I}_i^\bullet$ représente $M$. Par acyclicité de Leray
(Catégories dérivées, lemme \ref{derived-lemma-leray-acyclicity}), on voit que
$$
R\Gamma(U, M) = \Gamma(U, \bigoplus \mathcal{I}_i^\bullet) =
\bigoplus \Gamma(U, \mathcal{I}_i^\bullet) =
\bigoplus R\Gamma(U, M_i)
$$
comme souhaité.
\end{proof}

\begin{lemma}
\label{lemma-when-jshriek-lower-compact-worked-out}
Soit $(\mathcal{C}, \mathcal{O})$ un site annelé dont l'ensemble des
recouvrements est $\text{Cov}_\mathcal{C}$. Soient
$\mathcal{B} \subset \Ob(\mathcal{C})$ et
$\text{Cov} \subset \text{Cov}_\mathcal{C}$
des sous-ensembles. Supposons que :
\begin{enumerate}
\item Pour tout $\mathcal{U} \in \text{Cov}$, nous avons
$\mathcal{U} = \{U_i \to U\}_{i \in I}$ avec $I$ fini,
$U, U_i \in \mathcal{B}$ et tout
$U_{i_0} \times_U \ldots \times_U U_{i_p} \in \mathcal{B}$.
\item Pour tout $U \in \mathcal{B}$, les recouvrements de $U$ appartenant à
$\text{Cov}$ forment un système cofinal de recouvrements de $U$.
\end{enumerate}
Alors, pour $U \in \mathcal{B}$, l'objet $j_{U!}\mathcal{O}_U$ est un objet
compact de $D^+(\mathcal{O})$ au sens suivant : si
$M = \bigoplus_{i \in I} M_i$ dans $D(\mathcal{O})$ est borné inférieurement,
alors $\Hom(j_{U!}\mathcal{O}_U, M) =
\bigoplus_{i \in I} \Hom(j_{U!}\mathcal{O}_U, M_i)$.
\end{lemma}

\begin{proof}
Cela résulte des lemmes \ref{lemma-when-jshriek-lower-compact} et
\ref{lemma-colim-works-over-collection}.
\end{proof}

\begin{lemma}
\label{lemma-when-jshriek-compact}
Soit $(\mathcal{C}, \mathcal{O})$ un site annelé. Soit $U$ un objet de
$\mathcal{C}$. Le $\mathcal{O}$-module $j_!\mathcal{O}_U$ est un objet
compact de $D(\mathcal{O})$ s'il existe un entier $d$ tel que :
\begin{enumerate}
\item $H^p(U, \mathcal{F}) = 0$ pour tout $p > d$, et
\item les foncteurs $\mathcal{F} \mapsto H^p(U, \mathcal{F})$ commutent aux
sommes directes.
\end{enumerate}
\end{lemma}

\begin{proof}
Supposons (1) et (2). Rappelons que
$\Hom(j_!\mathcal{O}_U, K) = R\Gamma(U, K)$ pour $K$ dans $D(\mathcal{O})$.
Nous devons donc montrer que $R\Gamma(U, -)$ commute aux sommes directes. La
première hypothèse signifie que le foncteur $F = H^0(U, -)$ est de dimension
cohomologique finie. De plus, la seconde hypothèse implique que toute somme
directe de modules injectifs est acyclique pour $F$. Soit $K_i$ une famille
d'objets de $D(\mathcal{O})$. Choisissons des représentants K-injectifs
$I_i^\bullet$ à termes injectifs représentant $K_i$, voir Injectives,
théorème \ref{injectives-theorem-K-injective-embedding-grothendieck}. Puisque
nous pouvons calculer $RF$ en appliquant $F$ à tout complexe d'acycliques
(Catégories dérivées, lemme \ref{derived-lemma-unbounded-right-derived}) et
que $\bigoplus K_i$ est représenté par $\bigoplus I_i^\bullet$ (Injectives,
lemme \ref{injectives-lemma-derived-products}), nous concluons que
$R\Gamma(U, \bigoplus K_i)$ est représenté par
$\bigoplus H^0(U, I_i^\bullet)$. Ainsi, $R\Gamma(U, -)$ commute aux sommes
directes, comme souhaité.
\end{proof}

\begin{lemma}
\label{lemma-quasi-compact-weakly-contractible-compact}
Soit $(\mathcal{C}, \mathcal{O})$ un site annelé. Soit $U$ un objet de
$\mathcal{C}$ qui est quasi-compact et faiblement contractile. Alors
$j_!\mathcal{O}_U$ est un objet compact de $D(\mathcal{O})$.
\end{lemma}

\begin{proof}
Combiner les lemmes \ref{lemma-when-jshriek-compact} et
\ref{lemma-w-contractible} avec Modules on Sites, lemme
\ref{sites-modules-lemma-sections-over-quasi-compact}.
\end{proof}

\begin{lemma}
\label{lemma-perfect-is-compact}
Soit $(\mathcal{C}, \mathcal{O})$ un site annelé. Supposons que
$\mathcal{C}$ possède les propriétés suivantes :
\begin{enumerate}
\item $\mathcal{C}$ possède un objet final quasi-compact $X$,
\item tout objet quasi-compact de $\mathcal{C}$ possède un système cofinal de
recouvrements finis constitués d'objets quasi-compacts,
\item pour un recouvrement fini $\{U_i \to U\}_{i \in I}$, où $U$ et les
$U_i$ sont quasi-compacts, les produits fibrés $U_i \times_U U_j$ sont
quasi-compacts.
\end{enumerate}
Soit $K$ un objet parfait de $D(\mathcal{O})$. Alors :
\begin{enumerate}
\item[(a)] $K$ est un objet compact de $D^+(\mathcal{O})$ au sens suivant :
si $M = \bigoplus_{i \in I} M_i$ est borné inférieurement, alors
$\Hom(K, M) = \bigoplus_{i \in I} \Hom(K, M_i)$.
\item[(b)] Si $(\mathcal{C}, \mathcal{O})$ est de dimension cohomologique
finie, c'est-à-dire s'il existe un $d$ tel que
$H^i(X, \mathcal{F}) = 0$ pour $i > d$ et pour tout $\mathcal{O}$-module
$\mathcal{F}$, alors $K$ est un objet compact de $D(\mathcal{O})$.
\end{enumerate}
\end{lemma}

\begin{proof}
Soit $K^\vee$ le dual de $K$, voir le lemme
\ref{lemma-dual-perfect-complex}. Alors nous avons
$$
\Hom_{D(\mathcal{O})}(K, M) =
H^0(X, K^\vee \otimes_\mathcal{O}^\mathbf{L} M)
$$
fonctoriellement en $M$ dans $D(\mathcal{O})$. Puisque
$K^\vee \otimes_\mathcal{O}^\mathbf{L} -$ commute aux sommes directes, il
suffit de montrer que $R\Gamma(X, -)$ commute aux sommes directes concernées.

\medskip\noindent
Démonstration de (a). Après reformulation comme ci-dessus, c'est un cas
particulier du lemme \ref{lemma-when-jshriek-lower-compact-worked-out} avec
$U = X$.

\medskip\noindent
Démonstration de (b). Puisque $R\Gamma(X, K) = R\Hom(\mathcal{O}, K)$ et que
$H^p(X, -)$ commute aux sommes directes par le lemme
\ref{lemma-colim-works-over-collection}, c'est un cas particulier du lemme
\ref{lemma-when-jshriek-compact}.
\end{proof}






\section{Complexes à faisceaux de cohomologie localement constants}
\label{section-locally-constant}

\noindent
Les faisceaux localement constants sont introduits dans Modules on Sites,
section \ref{sites-modules-section-locally-constant}. Soit $\mathcal{C}$ un
site. Soit $\Lambda$ un anneau. Nous notons $D(\mathcal{C}, \Lambda)$ la
catégorie dérivée de la catégorie abélienne des
$\underline{\Lambda}$-modules sur $\mathcal{C}$.

\begin{lemma}
\label{lemma-locally-constant}
Soit $\mathcal{C}$ un site possédant un objet final $X$. Soit $\Lambda$ un
anneau noethérien. Soit $K \in D^b(\mathcal{C}, \Lambda)$ tel que les
$H^i(K)$ soient des faisceaux localement constants de $\Lambda$-modules de
type fini. Alors il existe un recouvrement $\{U_i \to X\}$ tel que chaque
$K|_{U_i}$ soit représenté par un complexe de faisceaux localement constants
de $\Lambda$-modules de type fini.
\end{lemma}

\begin{proof}
Soient $a \leq b$ tels que $H^i(K) = 0$ pour $i \not \in [a, b]$. Par
récurrence sur $b - a$, nous allons démontrer qu'il existe un recouvrement
$\{U_i \to X\}$ tel que $K|_{U_i}$ puisse être représenté par un complexe
$\underline{M^\bullet}_{U_i}$, où $M^p$ est un $\Lambda$-module de type fini
et $M^p = 0$ pour $p \not \in [a, b]$. Si $b = a$, c'est clair. En général,
nous pouvons remplacer $X$ par les membres d'un recouvrement et supposer que
$H^b(K)$ est constant, disons $H^b(K) = \underline{M}$. Par Modules on Sites,
lemme \ref{sites-modules-lemma-locally-constant-finite-type}, le module $M$
est un $\Lambda$-module de type fini. Choisissons une surjection
$\Lambda^{\oplus r} \to M$ donnée par des générateurs $x_1, \ldots, x_r$ de
$M$.

\medskip\noindent
Par une légère généralisation du lemme
\ref{lemma-kill-cohomology-class-on-covering} (détails omis), il existe un
recouvrement $\{U_i \to X\}$ tel que, pour tout $i$ et tout $j = 1, \ldots, r$,
la classe $x_j \in H^0(X, H^b(K))$ se relève en un élément de $H^b(U_i, K)$.
Ainsi, après avoir remplacé $X$ par les $U_i$, nous
parvenons à une situation où il existe un morphisme
$\underline{\Lambda^{\oplus r}}[-b] \to K$ induisant une surjection sur les
faisceaux de cohomologie au degré $b$. Choisissons un triangle distingué
$$
\underline{\Lambda^{\oplus r}}[-b] \to K \to L \to
\underline{\Lambda^{\oplus r}}[-b + 1]
$$
Les faisceaux de cohomologie de $L$ ne sont alors non nuls que dans
l'intervalle $[a, b - 1]$, coïncident avec ceux de $K$ dans l'intervalle
$[a, b - 2]$, et il existe une suite exacte courte
$$
0 \to H^{b - 1}(K) \to H^{b - 1}(L) \to
\underline{\Ker(\Lambda^{\oplus r} \to M)} \to 0
$$
au degré $b - 1$. Par Modules on Sites, lemme
\ref{sites-modules-lemma-kernel-finite-locally-constant}, on voit que
$H^{b - 1}(L)$ est localement constant de type fini. Par l'hypothèse de
récurrence, nous obtenons un isomorphisme
$\underline{M^\bullet} \to L$ dans $D(\mathcal{C}, \Lambda)$,
où $M^p$ est un $\Lambda$-module de type fini et $M^p = 0$ pour
$p \not \in [a, b - 1]$. Le morphisme
$L \to \underline{\Lambda^{\oplus r}}[-b + 1]$ donne un morphisme
$\underline{M^{b - 1}} \to \underline{\Lambda^{\oplus r}}$ qui est
localement constant (Modules on Sites, lemme
\ref{sites-modules-lemma-morphism-locally-constant}). Nous pouvons donc
supposer qu'il est donné par un morphisme
$M^{b - 1} \to \Lambda^{\oplus r}$. Le triangle distingué montre que la
composée $M^{b - 2} \to M^{b - 1} \to \Lambda^{\oplus r}$ est nulle, et les
axiomes des catégories triangulées produisent un isomorphisme
$$
\underline{M^a \to \ldots \to M^{b - 1} \to \Lambda^{\oplus r}}
\longrightarrow K
$$
dans $D(\mathcal{C}, \Lambda)$.
\end{proof}

\noindent
Soit $\mathcal{C}$ un site. Soit $\Lambda$ un anneau. En utilisant le
morphisme $\Sh(\mathcal{C}) \to \Sh(pt)$, on voit qu'il existe un foncteur
$D(\Lambda) \to D(\mathcal{C}, \Lambda)$, $K \mapsto \underline{K}$.

\begin{lemma}
\label{lemma-map-out-of-locally-constant}
Soit $\mathcal{C}$ un site possédant un objet final $X$. Soit $\Lambda$ un
anneau. Soient :
\begin{enumerate}
\item $K$ un objet parfait de $D(\Lambda)$,
\item $K^\bullet$ un complexe fini de $\Lambda$-modules projectifs de type
fini représentant $K$,
\item $\mathcal{L}^\bullet$ un complexe de faisceaux de $\Lambda$-modules, et
\item $\varphi : \underline{K} \to \mathcal{L}^\bullet$ un morphisme dans
$D(\mathcal{C}, \Lambda)$.
\end{enumerate}
Alors il existe un recouvrement $\{U_i \to X\}$ et des morphismes de complexes
$\alpha_i : \underline{K}^\bullet|_{U_i} \to \mathcal{L}^\bullet|_{U_i}$
représentant $\varphi|_{U_i}$.
\end{lemma}

\begin{proof}
Cela résulte immédiatement du lemme \ref{lemma-local-actual}.
\end{proof}

\begin{lemma}
\label{lemma-locally-constant-map}
Soit $\mathcal{C}$ un site possédant un objet final $X$. Soit $\Lambda$ un
anneau. Soient $K, L$ des objets de $D(\Lambda)$, avec $K$ parfait. Soit
$\varphi : \underline{K} \to \underline{L}$ un morphisme dans
$D(\mathcal{C}, \Lambda)$. Il existe un recouvrement $\{U_i \to X\}$ tel que
$\varphi|_{U_i}$ soit égal à $\underline{\alpha_i}$ pour un certain morphisme
$\alpha_i : K \to L$ dans $D(\Lambda)$.
\end{lemma}

\begin{proof}
Cela résulte du lemme \ref{lemma-map-out-of-locally-constant} et de Modules on
Sites, lemme \ref{sites-modules-lemma-morphism-locally-constant}.
\end{proof}

\begin{lemma}
\label{lemma-locally-constant-tensor-product}
Soit $\mathcal{C}$ un site. Soit $\Lambda$ un anneau noethérien. Soient
$K, L \in D^-(\mathcal{C}, \Lambda)$. Si les faisceaux de cohomologie de $K$
et de $L$ sont des faisceaux localement constants de $\Lambda$-modules de
type fini, alors les faisceaux de cohomologie de
$K \otimes_\Lambda^\mathbf{L} L$ sont des faisceaux localement constants de
$\Lambda$-modules de type fini.
\end{lemma}

\begin{proof}
Nous allons le démontrer comme application du lemme
\ref{lemma-locally-constant}. Notons que
$H^i(K \otimes_\Lambda^\mathbf{L} L)$ est le même que
$H^i(\tau_{\geq i - 1}K \otimes_\Lambda^\mathbf{L} \tau_{\geq i - 1}L)$.
Nous pouvons donc supposer que $K$ et $L$ sont bornés. Par le lemme
\ref{lemma-locally-constant}, nous pouvons supposer que $K$ et $L$ sont
représentés par des complexes de faisceaux localement constants de
$\Lambda$-modules de type fini. Nous pouvons alors remplacer ces complexes
par des complexes bornés supérieurement de $\Lambda$-modules libres de type
fini. Dans ce cas, le résultat est clair.
\end{proof}

\begin{lemma}
\label{lemma-locally-constant-bounded}
Soit $\mathcal{C}$ un site. Soit $\Lambda$ un anneau noethérien. Soit
$I \subset \Lambda$ un idéal. Soit $K \in D^-(\mathcal{C}, \Lambda)$. Si les
faisceaux de cohomologie de
$K \otimes_\Lambda^\mathbf{L} \underline{\Lambda/I}$ sont des faisceaux
localement constants de $\Lambda/I$-modules de type fini, alors les faisceaux
de cohomologie de
$K \otimes_\Lambda^\mathbf{L} \underline{\Lambda/I^n}$ sont des faisceaux
localement constants de $\Lambda/I^n$-modules de type fini pour tout
$n \geq 1$.
\end{lemma}

\begin{proof}
Rappelons que les faisceaux localement constants de $\Lambda$-modules de type
fini forment une sous-catégorie de Serre faible de tous les
$\underline{\Lambda}$-modules, voir Modules on Sites, lemme
\ref{sites-modules-lemma-kernel-finite-locally-constant}. Ainsi, la
sous-catégorie de $D(\mathcal{C}, \Lambda)$ constituée des complexes dont les
faisceaux de cohomologie sont des faisceaux localement constants de
$\Lambda$-modules de type fini forme une sous-catégorie triangulée strictement
pleine et saturée de $D(\mathcal{C}, \Lambda)$, voir Catégories dérivées,
lemme \ref{derived-lemma-cohomology-in-serre-subcategory}. Considérons ensuite
les triangles distingués
$$
K \otimes_\Lambda^\mathbf{L} \underline{I^n/I^{n + 1}} \to
K \otimes_\Lambda^\mathbf{L} \underline{\Lambda/I^{n + 1}} \to
K \otimes_\Lambda^\mathbf{L} \underline{\Lambda/I^n} \to
K \otimes_\Lambda^\mathbf{L} \underline{I^n/I^{n + 1}}[1]
$$
et les isomorphismes
$$
K \otimes_\Lambda^\mathbf{L} \underline{I^n/I^{n + 1}}
=
\left(K \otimes_\Lambda^\mathbf{L} \underline{\Lambda/I}\right)
\otimes_{\Lambda/I}^\mathbf{L} \underline{I^n/I^{n + 1}}
$$
En les combinant avec le lemme \ref{lemma-locally-constant-tensor-product},
nous obtenons le résultat.
\end{proof}






\begin{multicols}{2}[\section{Autres chapitres}]
\noindent
Préliminaires
\begin{enumerate}
\item \hyperref[introduction-section-phantom]{Introduction}
\item \hyperref[conventions-section-phantom]{Conventions}
\item \hyperref[sets-section-phantom]{Théorie des ensembles}
\item \hyperref[categories-section-phantom]{Catégories}
\item \hyperref[topology-section-phantom]{Topologie}
\item \hyperref[sheaves-section-phantom]{Faisceaux sur les espaces}
\item \hyperref[sites-section-phantom]{Sites et faisceaux}
\item \hyperref[stacks-section-phantom]{Champs}
\item \hyperref[fields-section-phantom]{Corps}
\item \hyperref[algebra-section-phantom]{Algèbre commutative}
\item \hyperref[brauer-section-phantom]{Groupes de Brauer}
\item \hyperref[homology-section-phantom]{Algèbre homologique}
\item \hyperref[derived-section-phantom]{Catégories dérivées}
\item \hyperref[simplicial-section-phantom]{Méthodes simpliciales}
\item \hyperref[more-algebra-section-phantom]{Compléments d'algèbre}
\item \hyperref[smoothing-section-phantom]{Lissage des morphismes d'anneaux}
\item \hyperref[modules-section-phantom]{Faisceaux de modules}
\item \hyperref[sites-modules-section-phantom]{Modules sur les sites}
\item \hyperref[injectives-section-phantom]{Objets injectifs}
\item \hyperref[cohomology-section-phantom]{Cohomologie des faisceaux}
\item \hyperref[sites-cohomology-section-phantom]{Cohomologie sur les sites}
\item \hyperref[dga-section-phantom]{Algèbre différentielle graduée}
\item \hyperref[dpa-section-phantom]{Algèbre à puissances divisées}
\item \hyperref[sdga-section-phantom]{Faisceaux différentiels gradués}
\item \hyperref[hypercovering-section-phantom]{Hyperrecouvrements}
\end{enumerate}
Schémas
\begin{enumerate}
\setcounter{enumi}{25}
\item \hyperref[schemes-section-phantom]{Schémas}
\item \hyperref[constructions-section-phantom]{Constructions de schémas}
\item \hyperref[properties-section-phantom]{Propriétés des schémas}
\item \hyperref[morphisms-section-phantom]{Morphismes de schémas}
\item \hyperref[coherent-section-phantom]{Cohomologie des schémas}
\item \hyperref[divisors-section-phantom]{Diviseurs}
\item \hyperref[limits-section-phantom]{Limites de schémas}
\item \hyperref[varieties-section-phantom]{Variétés}
\item \hyperref[topologies-section-phantom]{Topologies sur les schémas}
\item \hyperref[descent-section-phantom]{Descente}
\item \hyperref[perfect-section-phantom]{Catégories dérivées de schémas}
\item \hyperref[more-morphisms-section-phantom]{Compléments sur les morphismes}
\item \hyperref[flat-section-phantom]{Compléments sur la platitude}
\item \hyperref[groupoids-section-phantom]{Schémas en groupoïdes}
\item \hyperref[more-groupoids-section-phantom]{Compléments sur les schémas en groupoïdes}
\item \hyperref[etale-section-phantom]{Morphismes étales de schémas}
\end{enumerate}
Thèmes de théorie des schémas
\begin{enumerate}
\setcounter{enumi}{41}
\item \hyperref[chow-section-phantom]{Homologie de Chow}
\item \hyperref[intersection-section-phantom]{Théorie de l'intersection}
\item \hyperref[pic-section-phantom]{Schémas de Picard des courbes}
\item \hyperref[weil-section-phantom]{Théories cohomologiques de Weil}
\item \hyperref[adequate-section-phantom]{Modules adéquats}
\item \hyperref[dualizing-section-phantom]{Complexes dualisants}
\item \hyperref[duality-section-phantom]{Dualité pour les schémas}
\item \hyperref[discriminant-section-phantom]{Discriminants et différentes}
\item \hyperref[derham-section-phantom]{Cohomologie de de Rham}
\item \hyperref[local-cohomology-section-phantom]{Cohomologie locale}
\item \hyperref[algebraization-section-phantom]{Géométrie algébrique et formelle}
\item \hyperref[curves-section-phantom]{Courbes algébriques}
\item \hyperref[resolve-section-phantom]{Résolution des surfaces}
\item \hyperref[models-section-phantom]{Réduction semi-stable}
\item \hyperref[functors-section-phantom]{Foncteurs et morphismes}
\item \hyperref[equiv-section-phantom]{Catégories dérivées de variétés}
\item \hyperref[pione-section-phantom]{Groupes fondamentaux de schémas}
\item \hyperref[etale-cohomology-section-phantom]{Cohomologie étale}
\item \hyperref[crystalline-section-phantom]{Cohomologie cristalline}
\item \hyperref[proetale-section-phantom]{Cohomologie pro-étale}
\item \hyperref[relative-cycles-section-phantom]{Cycles relatifs}
\item \hyperref[more-etale-section-phantom]{Compléments de cohomologie étale}
\item \hyperref[trace-section-phantom]{La formule des traces}
\end{enumerate}
Espaces algébriques
\begin{enumerate}
\setcounter{enumi}{64}
\item \hyperref[spaces-section-phantom]{Espaces algébriques}
\item \hyperref[spaces-properties-section-phantom]{Propriétés des espaces algébriques}
\item \hyperref[spaces-morphisms-section-phantom]{Morphismes d'espaces algébriques}
\item \hyperref[decent-spaces-section-phantom]{Espaces algébriques décents}
\item \hyperref[spaces-cohomology-section-phantom]{Cohomologie des espaces algébriques}
\item \hyperref[spaces-limits-section-phantom]{Limites d'espaces algébriques}
\item \hyperref[spaces-divisors-section-phantom]{Diviseurs sur les espaces algébriques}
\item \hyperref[spaces-over-fields-section-phantom]{Espaces algébriques sur des corps}
\item \hyperref[spaces-topologies-section-phantom]{Topologies sur les espaces algébriques}
\item \hyperref[spaces-descent-section-phantom]{Descente et espaces algébriques}
\item \hyperref[spaces-perfect-section-phantom]{Catégories dérivées d'espaces}
\item \hyperref[spaces-more-morphisms-section-phantom]{Compléments sur les morphismes d'espaces}
\item \hyperref[spaces-flat-section-phantom]{Platitude sur les espaces algébriques}
\item \hyperref[spaces-groupoids-section-phantom]{Groupoïdes dans les espaces algébriques}
\item \hyperref[spaces-more-groupoids-section-phantom]{Compléments sur les groupoïdes dans les espaces}
\item \hyperref[bootstrap-section-phantom]{Amorçage}
\item \hyperref[spaces-pushouts-section-phantom]{Sommes amalgamées d'espaces algébriques}
\end{enumerate}
Thèmes de géométrie
\begin{enumerate}
\setcounter{enumi}{81}
\item \hyperref[spaces-chow-section-phantom]{Groupes de Chow des espaces}
\item \hyperref[groupoids-quotients-section-phantom]{Quotients de groupoïdes}
\item \hyperref[spaces-more-cohomology-section-phantom]{Compléments sur la cohomologie des espaces}
\item \hyperref[spaces-simplicial-section-phantom]{Espaces simpliciaux}
\item \hyperref[spaces-duality-section-phantom]{Dualité pour les espaces}
\item \hyperref[formal-spaces-section-phantom]{Espaces algébriques formels}
\item \hyperref[restricted-section-phantom]{Algébrisation des espaces formels}
\item \hyperref[spaces-resolve-section-phantom]{Résolution des surfaces revisitée}
\end{enumerate}
Théorie des déformations
\begin{enumerate}
\setcounter{enumi}{89}
\item \hyperref[formal-defos-section-phantom]{Théorie formelle des déformations}
\item \hyperref[defos-section-phantom]{Théorie des déformations}
\item \hyperref[cotangent-section-phantom]{Le complexe cotangent}
\item \hyperref[examples-defos-section-phantom]{Problèmes de déformation}
\end{enumerate}
Champs algébriques
\begin{enumerate}
\setcounter{enumi}{93}
\item \hyperref[algebraic-section-phantom]{Champs algébriques}
\item \hyperref[examples-stacks-section-phantom]{Exemples de champs}
\item \hyperref[stacks-sheaves-section-phantom]{Faisceaux sur les champs algébriques}
\item \hyperref[criteria-section-phantom]{Critères de représentabilité}
\item \hyperref[artin-section-phantom]{Axiomes d'Artin}
\item \hyperref[quot-section-phantom]{Espaces de Quot et de Hilbert}
\item \hyperref[stacks-properties-section-phantom]{Propriétés des champs algébriques}
\item \hyperref[stacks-morphisms-section-phantom]{Morphismes de champs algébriques}
\item \hyperref[stacks-limits-section-phantom]{Limites de champs algébriques}
\item \hyperref[stacks-cohomology-section-phantom]{Cohomologie des champs algébriques}
\item \hyperref[stacks-perfect-section-phantom]{Catégories dérivées de champs}
\item \hyperref[stacks-introduction-section-phantom]{Introduction aux champs algébriques}
\item \hyperref[stacks-more-morphisms-section-phantom]{Compléments sur les morphismes de champs}
\item \hyperref[stacks-geometry-section-phantom]{La géométrie des champs}
\end{enumerate}
Thèmes de théorie des espaces de modules
\begin{enumerate}
\setcounter{enumi}{107}
\item \hyperref[moduli-section-phantom]{Champs de modules}
\item \hyperref[moduli-curves-section-phantom]{Espaces de modules de courbes}
\end{enumerate}
Divers
\begin{enumerate}
\setcounter{enumi}{109}
\item \hyperref[examples-section-phantom]{Exemples}
\item \hyperref[exercises-section-phantom]{Exercices}
\item \hyperref[guide-section-phantom]{Guide bibliographique}
\item \hyperref[desirables-section-phantom]{Desiderata}
\item \hyperref[coding-section-phantom]{Style de programmation}
\item \hyperref[obsolete-section-phantom]{Obsolète}
\item \hyperref[fdl-section-phantom]{Licence de documentation libre GNU}
\item \hyperref[index-section-phantom]{Index généré automatiquement}
\end{enumerate}
\end{multicols}


\bibliography{my}
\bibliographystyle{amsalpha}

\end{document}
